% =============================================================================
% MOLAMBUDOS — O Diário do Paciente 1.260
% Template Victoria-Regia by @jancarauma
% 84 fragmentos · 223 rotas de leitura · 5 partes
% Atualizado com revisoes R286-R295:
%   Aviso de Conteudo, Protocolo de Leitura, escalas dinamicas,
%   MEM-27, LUC-Contraprova/Entrevista/Escolha, CONT-01 a CONT-13 (nova edição R376)
% =============================================================================

\documentclass[12pt,openany]{extbook}
\newif\ifmolambudoskdpprint
\ifdefined\KDPPRINT
  \molambudoskdpprinttrue
\else
  \molambudoskdpprintfalse
\fi
\newcommand{\goldenratio}{1.618}
\usepackage[T1]{fontenc}

% --- Book metadata ---
\newcommand{\authorname}{Marcelo Dias de Carvalho Filho}
\newcommand{\coauthorname}{Marcelo Claro}
\newcommand{\translatorname}{}
\newcommand{\booktitle}{Molambudos --- O Diário do Paciente 1.260}
\newcommand{\subtitle}{Uma narrativa em fragmentos}
\newcommand{\publisher}{OpenCode Books}
\newcommand{\publishercity}{Crateús}
\newcommand{\editionyear}{2026}
\newcommand{\isbn}{9798189170492}

% --- Template packages + customizations ---

% --- Geometria da edição Grimório (160 x 230 mm, margem externa ampliada) ---
\def\MolambudosPaperWidth{160mm}
\def\MolambudosPaperHeight{230mm}
\def\MolambudosInnerMargin{22mm}
\def\MolambudosOuterMargin{30mm}
\def\MolambudosTopMargin{20mm}
\def\MolambudosBottomMargin{22mm}

\usepackage{misc/options}
\usepackage{float}
\usepackage{etoolbox}
\usepackage{xstring}

% --- Sistema visual R1076: cor por código, capitular e abertura de fragmento ---
\newcommand{\fragcode}[1]{%
  \IfBeginWith{#1}{MEM}{\textcolor{mem}{\textbf{#1}}}{%
  \IfBeginWith{#1}{DOC}{\textcolor{doc}{\textbf{#1}}}{%
  \IfBeginWith{#1}{LUC}{\textcolor{luc}{\textbf{#1}}}{%
  \IfBeginWith{#1}{CONT}{\textcolor{cont}{\textbf{#1}}}{\textbf{#1}}}}}%
}

\newcommand{\fragfamily}[1]{%
  \ifstrequal{#1}{MEM}{Memória}{%
  \ifstrequal{#1}{DOC}{Documento}{%
  \ifstrequal{#1}{LUC}{Lucidez}{%
  \ifstrequal{#1}{CONT}{Contaminação}{#1}}}}%
}

\newcommand{\fragtitle}[5]{%
  % #1 cor, #2 inicial, #3 código, #4 número, #5 título
  \par\addvspace{2.2em}%
  \noindent
  \begin{minipage}[c]{0.14\textwidth}
    \centering
    {\fontsize{40}{40}\selectfont\bfseries\color{#1}#2}
  \end{minipage}%
  \begin{minipage}[c]{0.82\textwidth}
    {\fontfamily{cmss}\selectfont\bfseries\footnotesize\color{#1}%
      \fragfamily{#3}\ifx&#4&\else\,#4\fi}\par
    \vspace{0.2em}
    {\fontsize{18}{22}\selectfont\bfseries #5\par}
    \vspace{0.35em}
    {\color{#1}\rule{0.28\textwidth}{1.2pt}}
  \end{minipage}\par
  \addvspace{1.4em}%
}

\newcommand{\fragcap}[2]{%
  \lettrine[lines=2,findent=2pt,nindent=4pt]{\textcolor{#1}{#2}}{}%
}


% --- Diretório de figuras (fotografias sépia do diário) ---
\graphicspath{{figures/}}

% ============================================================================
% Customizações específicas do Molambudos
% ============================================================================

% Unicode characters used in fragments
% (only for pdflatex; xelatex handles Unicode natively)
\ifdefined\DeclareUnicodeCharacter
\DeclareUnicodeCharacter{21AA}{\ensuremath{\hookrightarrow}}
\DeclareUnicodeCharacter{2610}{\ensuremath{\square}}
\DeclareUnicodeCharacter{2713}{\checkmark}
\DeclareUnicodeCharacter{2717}{\texttimes}
\DeclareUnicodeCharacter{2611}{\ensuremath{\checkmark}}
\DeclareUnicodeCharacter{2639}{\textfrown}
\DeclareUnicodeCharacter{263A}{\textsmile}
\DeclareUnicodeCharacter{2665}{\ensuremath{\heartsuit}}
\DeclareUnicodeCharacter{2714}{\checkmark}
\DeclareUnicodeCharacter{2605}{\ensuremath{\bigstar}}
\DeclareUnicodeCharacter{27E1}{\ensuremath{\blacktriangleright}}
\DeclareUnicodeCharacter{2022}{\textbullet}
\DeclareUnicodeCharacter{25CF}{\textbullet}
\DeclareUnicodeCharacter{25CB}{\ensuremath{\circ}}
\DeclareUnicodeCharacter{2726}{\ensuremath{\star}}
\DeclareUnicodeCharacter{25C6}{\ensuremath{\blacklozenge}}
\DeclareUnicodeCharacter{26A0}{\textbf{!}}
\DeclareUnicodeCharacter{2192}{$\rightarrow$}
\DeclareUnicodeCharacter{2190}{$\leftarrow$}
\DeclareUnicodeCharacter{00B0}{$^\circ$}
\DeclareUnicodeCharacter{2013}{--}
\DeclareUnicodeCharacter{2014}{---}
\DeclareUnicodeCharacter{2018}{`}
\DeclareUnicodeCharacter{2019}{'}
\DeclareUnicodeCharacter{201C}{``}
\DeclareUnicodeCharacter{201D}{''}
\fi

% Hipertexto condicional — reconhece o percurso de leitura
% \ifpath{origem1}{origem2}{texto se veio de 1}{texto se veio de 2}
% Usa o sistema de rótulos do LaTeX para detectar qual fragmento
% foi lido primeiro com base na ordem de aparição no .aux.
% \ifpath{origem1}{origem2}{texto se veio de 1}{texto se veio de 2}
% Em PDF, o caminho real é indeterminável: exibe ambos como possibilidades.
% O leitor reconhece o próprio percurso entre as alternativas.
\newcommand{\ifpath}[4]{%
  \textit{Se você veio de #1:} #3\par
  \vspace{0.3em}
  \textit{Se você veio de #2:} #4\par
  \vspace{0.3em}
  \textit{(Ambos os caminhos levam ao mesmo lugar. A entidade não precisa saber por onde você veio. Ela só precisa que você tenha chegado.)}%
}

% Rota fantasma — linka para um fragmento que não existe.
% O leitor descobre que a rota não leva a lugar nenhum.
% Usado para criar o "fragmento que você não leu".
\newcommand{\rotaphantom}[1]{%
  \textbf{#1}{\footnotesize(p.~---)}%
}

% Molambudos color palette
% Paleta R1076 — legibilidade sobre papel sépia (#F8F0D7)
% Contrastes medidos: MEM 6,76:1; DOC 6,76:1; LUC 9,23:1; CONT 7,18:1.
\definecolor{sangue}{HTML}{8B1A1A}   % alerta/nota; 8,16:1
\definecolor{sangueclaro}{HTML}{FF4D4D} % apenas fundo preto; 6,42:1 sobre #000
\definecolor{olho}{HTML}{FFD700}     % apenas sobre fundo escuro
\definecolor{mem}{HTML}{8A3B1F}      % cobre queimado
\definecolor{doc}{HTML}{1F5C57}      % teal profundo
\definecolor{cont}{HTML}{9B1B1B}     % carmim
\definecolor{luc}{HTML}{552A7A}      % violeta
\definecolor{epi}{HTML}{6E4E10}      % âmbar escuro; 6,68:1
% Paleta sépia imersiva (fundo de página + tinta) — apelo sensorial.
\definecolor{sepiabg}{HTML}{F2E8CF}
\definecolor{sepiaink}{HTML}{2B2117}

% Hyperlink setup (internal hyperlinks between fragments)
\ifmolambudoskdpprint
  \hypersetup{
    colorlinks=false,
    hidelinks
  }
\else
  \hypersetup{
    colorlinks=true,
    linkcolor=linkred,
    citecolor=linkred,
    urlcolor=linkred
  }
\fi

% Metadados editoriais do PDF (Title/Author/Subject/Keywords).
% Sem isso, o arquivo sai com os campos vazios mesmo tendo hyperref.
\hypersetup{
  pdftitle={\booktitle},
  pdfauthor={\authorname},
  pdfsubject={Ficção brasileira contemporânea; horror; romance em fragmentos},
  pdfkeywords={Molambudos, Paciente 1.260, Marcelo Dias de Carvalho Filho,
    ficção brasileira, horror, Ceará, romance fragmentado},
  pdfcreator={LaTeX with hyperref},
  pdfproducer={pdfTeX}
}

% Abertura de parte — título + medalhão circular temático (icone/N.png)
% numa única página limpa. #1 título exibido, #2 nome curto (TOC), #3 ícone.
% Mapeamento: 2=Sertão, 4=Colônia, 1=Diário/Laudos,
% 3=Investigação Lúcia, 5=Contaminação.
\newcommand{\partopener}[3]{%
  \cleardoublepage
  \thispagestyle{empty}%
  \phantomsection\addcontentsline{toc}{part}{#2}%
  \markboth{}{}%
  \begin{tikzpicture}[remember picture,overlay]
    \fill[black] (current page.south west) rectangle (current page.north east);
    \node[anchor=center] at (current page.center){%
      \begin{tabular}{c}
        {\bfseries\LARGE\color{customgold} #1}\\[2.4em]
        \includegraphics[width=0.58\textwidth]{icone/#3}\\
      \end{tabular}};
  \end{tikzpicture}%
  \null\vfill\clearpage}

% Rota de leitura — link para um fragmento + seu número de página, para que
% o leitor navegue a obra fora da ordem linear. \rota{ID} usa \pageref sobre
% os \label{frag:ID} colocados em cada abertura de fragmento (ver \mainmatter).
\ifmolambudoskdpprint
  \newcommand{\rota}[1]{\textbf{#1}\,{\footnotesize(p.~\pageref{frag:#1})}}
\else
  \newcommand{\rota}[1]{\hyperlink{#1}{\textbf{#1}}\,{\footnotesize(p.~\pageref{frag:#1})}}
\fi
\newcommand{\rotasep}{\quad\textbullet\quad}

% Marginália dos fragmentos de Contaminação — anotações que parecem escritas
% à mão na cópia do leitor (vermelho-sangue), para intensificar a quebra da
% quarta parede. Margem estreita (livro 6x9), então notas curtas.
\setlength{\marginparwidth}{0.75in}
\setlength{\marginparsep}{6pt}
\ifmolambudoskdpprint
  % Modo impressão KDP: as notas de margem são suprimidas para evitar
  % reprovação por "texto fora das margens" no previewer da Amazon.
  \newcommand{\margnota}[1]{}
\else
  \newcommand{\margnota}[1]{\marginpar{\raggedright\footnotesize\itshape\color{sangue}#1}}
\fi

% Nota do Editor — anotação paratextual para moldura editorial crítica.
% Agora implementada como nota de rodapé (fonte ~10pt) para discrição.
\newcommand{\NE}[1]{\footnote{\footnotesize\textbf{[N.E.]} #1}}

% Contador de fragmento — garante que cada \label{frag:ID} tenha âncora única
% em vez de herdar o \phantomsection do \partopener (que faz o \pageref
% apontar para a página do ato em vez da página do fragmento).
\newcounter{fragctr}
\newcommand{\fragdef}[1]{%
  \refstepcounter{fragctr}%
  \hypertarget{#1}{}%
  \label{frag:#1}%
}


% ============================================================================
% CAMADA RPG / GRIMÓRIO — tipografia, caixas, selos e dados
% ============================================================================
\usepackage{yfonts}
\usepackage{tcolorbox}
\tcbuselibrary{skins,breakable}
\usepackage{marginnote}
\usepackage{fontawesome5}
\usepackage{pifont}
\usepackage{tabularx}
\usetikzlibrary{shapes.geometric}

% Paleta de grimório
\definecolor{rpgink}{HTML}{1A120C}
\definecolor{rpgpaper}{HTML}{EFE0BF}
\definecolor{rpggold}{HTML}{B08D57}
\definecolor{rpgblood}{HTML}{7A1E1E}
\definecolor{rpgteal}{HTML}{1F5C57}
\definecolor{rpgviolet}{HTML}{4B2A5E}
\definecolor{rpgsepia}{HTML}{6B4A2E}

% Cores das famílias, redefinidas para o tom de grimório
\definecolor{mem}{HTML}{6E3B24}
\definecolor{doc}{HTML}{1F5C57}
\definecolor{luc}{HTML}{4B2A5E}
\definecolor{cont}{HTML}{7A1E1E}
\definecolor{epi}{HTML}{8A6A2F}
\definecolor{sangue}{HTML}{7A1E1E}
\definecolor{sepiabg}{HTML}{EFE0BF}
\definecolor{sepiaink}{HTML}{241A12}

% Textura sutil de papel (pontos a 3,5% de opacidade)


% Cabeçalho e folio de grimório
\pagestyle{fancy}
\setlength{\headheight}{25pt}
\fancyhf{}
\renewcommand{\headrulewidth}{0.4pt}
\renewcommand{\headrule}{\hbox to\headwidth{\color{rpggold}\leaders\hrule height \headrulewidth\hfill}}
\fancyhead[LE]{\textsc{\footnotesize\color{rpgink}Molambudos · Arquivo 1.260}}
\fancyhead[RO]{\textsc{\footnotesize\color{rpgink}Edição Grimório}}
\fancyfoot[C]{\textcolor{rpggold}{\small$\blacklozenge$}~\thepage~\textcolor{rpggold}{\small$\blacklozenge$}}
\hypersetup{hidelinks}

% Caixas
\newtcolorbox{rpgbox}[1][]{enhanced, breakable, colback=rpgpaper!88!white,
  colframe=rpggold, boxrule=0.7pt, arc=1.2mm, left=3.5mm, right=3.5mm,
  top=2.5mm, bottom=2.5mm, fonttitle=\scshape\bfseries\small,
  coltitle=rpgink, colbacktitle=rpggold!32!rpgpaper, title={#1}}
\newtcolorbox{rpgdark}[1][]{enhanced, breakable, colback=rpgink,
  coltext=rpggold, colframe=rpggold, boxrule=0.9pt, arc=1.2mm,
  left=3.5mm, right=3.5mm, top=2.5mm, bottom=2.5mm,
  fonttitle=\scshape\bfseries\small, coltitle=rpggold,
  colbacktitle=rpgink, title={#1}}
\newtcolorbox{rpgcard}[1][]{enhanced, colback=rpgpaper!92!white,
  colframe=rpggold!70!black, boxrule=0.6pt, arc=1.5mm,
  left=2.5mm, right=2.5mm, top=2mm, bottom=2mm,
  fonttitle=\scshape\bfseries\footnotesize, coltitle=rpgink,
  colbacktitle=rpggold!25!rpgpaper, title={#1}}

% Linhas, selos e dados
\newcommand{\rpgrule}{\par\noindent\centerline{%
  \textcolor{rpggold}{\rule{0.22\textwidth}{0.5pt}}\hspace{0.6em}%
  \textcolor{rpgblood}{\faDiceD20}\hspace{0.6em}%
  \textcolor{rpggold}{\rule{0.22\textwidth}{0.5pt}}}\par}
\newcommand{\rpglabel}[1]{{\scshape\bfseries\footnotesize\color{rpgblood}#1}}
\newcommand{\rpgnota}[1]{\marginnote{\raggedright\footnotesize\itshape\textcolor{rpgblood}{#1}}}
\newcommand{\rpgdot}{\raisebox{0.15ex}{\tiny$\square$}}

% Capitular gótica: o inicial é sempre uma letra sem acento (M, D, L, C)
\renewcommand{\fragcap}[2]{%
  \lettrine[lines=3,findent=2pt,nindent=4pt]{\textcolor{#1}{\textgoth{#2}}}{}%
}

% Abertura de fragmento: selo gótico + título em Garamond + filete dourado
\renewcommand{\fragtitle}[5]{%
  \par\addvspace{2.2em}%
  \noindent
  \begin{tikzpicture}[baseline=(ini.center)]
    \node[circle, draw=rpggold, line width=0.8pt, fill=#1, text=white,
      inner sep=3pt, font=\gothfamily\fontsize{20}{20}\selectfont] (ini) {#2};
  \end{tikzpicture}\hspace{0.55em}%
  \begin{minipage}[b]{0.80\textwidth}
    {\scshape\bfseries\footnotesize\color{#1}\fragfamily{#3}\ifx&#4&\else\enspace #4\fi}\par
    \vspace{0.15em}
    {\fontsize{17}{20}\selectfont\bfseries\color{rpgink}#5\par}
    \vspace{0.22em}
    {\color{rpggold}\rule{0.38\textwidth}{0.8pt}}\par
  \end{minipage}\par
  \addvspace{1.1em}%
}

% Abertura de parte: página preta, numeral gótico, título em versaletes
\newcounter{partopenerctr}
\renewcommand{\partopener}[3]{%
  \cleardoublepage
  \thispagestyle{empty}%
  \stepcounter{partopenerctr}%
  \phantomsection\addcontentsline{toc}{part}{#2}%
  \markboth{}{}%
  \begin{tikzpicture}[remember picture,overlay]
    \fill[rpgink] (current page.south west) rectangle (current page.north east);
    \draw[rpggold, line width=1pt]
      ([xshift=12mm,yshift=12mm]current page.south west) rectangle
      ([xshift=-12mm,yshift=-12mm]current page.north east);
    \draw[rpggold, line width=0.35pt]
      ([xshift=15mm,yshift=15mm]current page.south west) rectangle
      ([xshift=-15mm,yshift=-15mm]current page.north east);
    \node[anchor=center, text=rpggold, font=\gothfamily\fontsize{64}{64}\selectfont]
      at ([yshift=38mm]current page.center) {\Roman{partopenerctr}};
    \node[anchor=center, text=rpggold, font=\scshape\bfseries\fontsize{22}{26}\selectfont,
      text width=0.78\paperwidth, align=center]
      at ([yshift=2mm]current page.center) {#1};
    \node[anchor=center] at ([yshift=-28mm]current page.center)
      {\includegraphics[width=0.24\textwidth]{icone/#3}};
    \node[anchor=center, text=rpggold!80, font=\footnotesize\scshape]
      at ([yshift=-52mm]current page.center) {Arquivo 1.260 · Acesso restrito};
  \end{tikzpicture}%
  \null\vfill\clearpage}

\renewcommand{\NE}[1]{\footnote{\footnotesize\textcolor{rpgblood}{\textbf{[N.E.]}} #1}}

\begin{document}

% Uma única sequência romana cobre capa e todos os pré-textuais. As âncoras
% automáticas de página ficam desativadas até o início arábico do corpus.
\hypersetup{pageanchor=false}
\frontmatter

\thispagestyle{empty}
\begin{tikzpicture}[remember picture,overlay]
  \fill[rpgink] (current page.south west) rectangle (current page.north east);
  \draw[rpggold, line width=1pt]
    ([xshift=12mm,yshift=12mm]current page.south west) rectangle
    ([xshift=-12mm,yshift=-12mm]current page.north east);
  \draw[rpggold, line width=0.35pt]
    ([xshift=15mm,yshift=15mm]current page.south west) rectangle
    ([xshift=-15mm,yshift=-15mm]current page.north east);
  \node[anchor=north, text=rpggold, font=\gothfamily\fontsize{28}{30}\selectfont,
    text width=132mm, align=center] at ([yshift=-28mm]current page.north) {MOLAMBUDOS};
  \node[anchor=north, text=rpggold, font=\scshape\fontsize{14}{18}\selectfont,
    text width=122mm, align=center] at ([yshift=-43mm]current page.north) {O DIÁRIO DO PACIENTE 1.260};
  \node[anchor=north, text=rpggold!80, font=\itshape\small,
    text width=110mm, align=center] at ([yshift=-55mm]current page.north) {Uma narrativa em fragmentos};
  \node[anchor=center] at ([yshift=12mm]current page.center) {%
    \tikz{
      \draw[rpggold, line width=1pt] (0,0) circle (14mm);
      \draw[rpggold, line width=0.6pt] (0,12mm) -- (-10.4mm,-6mm) -- (10.4mm,-6mm) -- cycle;
      \node[text=rpggold, font=\bfseries\large] at (0,0) {1.260};
    }};
  \node[anchor=center, text=rpggold, font=\scshape\bfseries\large]
    at ([yshift=-28mm]current page.center) {EDIÇÃO GRIMÓRIO};
  \node[anchor=south, text=rpggold, font=\small\scshape]
    at ([yshift=34mm]current page.south) {MARCELO DIAS DE CARVALHO FILHO};
  \node[anchor=south, text=rpggold!85, font=\footnotesize\scshape]
    at ([yshift=27mm]current page.south) {COAUTOR E ILUSTRADOR: MARCELO CLARO};
  \node[anchor=south, text=rpggold!70, font=\footnotesize\scshape]
    at ([yshift=20mm]current page.south) {OPENCODE BOOKS · CRATEÚS · 2026};
\end{tikzpicture}
\null\vfill\clearpage


% =================
% AVISO AO LEITOR — trigger warning formal (não diegético)
% Página oficial antes do início da imersão ficcional
% =================
\cleardoublepage
\thispagestyle{empty}
\begingroup
\setlength{\fboxsep}{8mm}
\begin{center}
{\color{sangue}\scshape\Large AVISO AO LEITOR}\\[1.2em]
{\color{sangue}\rule{0.4\paperwidth}{0.6pt}}\\[2.2em]
\end{center}

\noindent Este livro contém, em sua parte ficcional, representações diretas de:

\begin{itemize}
\itemsep0.4em
\item violência institucional psiquiátrica --- incluindo eletrochoque, lobotomia e cárcere em hospital-colônia (Hospital Colônia de Barbacena, 1903--1994);
\item fome, morte e enterro coletivo de crianças e adultos (secas do Ceará, 1915--1932);
\item descrições médicas e necroscópicas detalhadas, com linguagem clínica explícita;
\item morte, luto, abandono e isolamento de longa duração;
\item narrativa em segunda pessoa que convoca diretamente o leitor a ocupar o lugar do \textit{Paciente 1.263}, com instruções simuladas de leitura.
\end{itemize}

\vspace{1em}

\noindent \textbf{Sobre as instruções simuladas:} nenhuma instrução presente no livro produz efeito físico, químico ou biológico real. A ``contaminação'' é um dispositivo narrativo. As páginas foram impressas com papel e tinta convencionais, sem qualquer substância aplicada. O leitor pode interromper a leitura a qualquer momento, em qualquer fragmento, sem nenhuma consequência além daquela que a própria história sugerir.

\vspace{1em}

\noindent \textbf{Sobre sonhos, sono e sintomas:} não há evidência de que a leitura deste livro induza sonhos específicos, altere o ciclo sono-vigília ou produza qualquer sintoma físico. As descrições de reação corporal são recurso formal, não afirmação empírica.

\vspace{1em}

\noindent \textbf{Sobre os códigos e escalas clínicas:} os protocolos diagnósticos, códigos CID e escalas de avaliação que aparecem na parte ficcional são construção literária. Os transtornos citados são condições clínicas reais, mas sua aplicação a um agente narrativo não tem qualquer respaldo médico. \textbf{Este livro não diagnostica o leitor}: ele simula um diagnóstico, como dispositivo de ficção.

\vspace{1em}

\noindent \textbf{Sobre a distinção ficção/realidade:} os personagens Joaquim Antônio Correia, Dr. Heitor Oliveira, Dra. Lúcia Mendes e a criatura são ficcionais. Os currais do governo (1915--1932) e o Hospital Colônia de Barbacena (1903--1994) são fatos históricos documentados. A distinção completa encontra-se na \textit{Nota Histórica}, ao final do volume (p. \pageref{nota-historica}).

\vspace{1em}

\noindent \textbf{Se você é sensível a algum desses temas:} esta obra pode não ser adequada neste momento. Não há vergonha em adiar a leitura ou abandoná-la. Livros não devem ser provas de resistência.

\vspace{1em}

\noindent \textbf{Se você precisar de apoio:} Centro de Valorização da Vida --- CVV, 188 (ligação gratuita, 24 horas) ou \texttt{cvv.org.br}; Unidades de Saúde Mental (CAPS) em todo o território nacional; ou sua rede de confiança.

\vspace{2.2em}
\begin{center}
{\color{sangue}\rule{0.4\paperwidth}{0.6pt}}\\[1em]
{\small\itshape O Arquivista}
\end{center}
\endgroup
\null\vfill
\clearpage


\ifdefined\MolambudosKDPNoFolios
  \makeatletter
  \let\ps@plain\ps@empty
  \makeatother
  \pagestyle{empty}
\fi
\ifdefined\MolambudosKDPSafeFolios
  % Modo impressão KDP com navegação física preservada:
  % sem cabeçalhos corridos, mas com folio central seguro no rodapé.
  \setlength{\footskip}{6mm}
  \fancypagestyle{molambudos-kdp-folios}{%
    \fancyhf{}%
    \renewcommand{\headrulewidth}{0pt}%
    \renewcommand{\footrulewidth}{0pt}%
    \fancyfoot[C]{\normalfont\footnotesize\thepage}%
  }%
  \fancypagestyle{plain}{%
    \fancyhf{}%
    \renewcommand{\headrulewidth}{0pt}%
    \renewcommand{\footrulewidth}{0pt}%
    \fancyfoot[C]{\normalfont\footnotesize\thepage}%
  }%
  \pagestyle{molambudos-kdp-folios}
\fi

% Fundo sépia (pergaminho) + tinta sépia escura — apelo sensorial imersivo.
\pagecolor{sepiabg}
\color{sepiaink}

% =================
% Meia-folha + epígrafe — primeira inquietação
% =================
\cleardoublepage
\thispagestyle{empty}
\vspace*{\fill}
\begin{center}
{\scshape\Large Molambudos}\\[0.4em]
{\small\itshape O Diário do Paciente 1.260}\\[6em]

{\large\itshape
``Ela não é um monstro.\\[0.5em]
Ela é uma consequência.''}
\end{center}
\vspace*{\fill}
\clearpage

% =================
% CUIDADO — página de choque (preta), quebra sensorial no folhear
% =================
\cleardoublepage
\thispagestyle{empty}
\begin{tikzpicture}[remember picture,overlay]
  \fill[black] (current page.south west) rectangle (current page.north east);
  \node[anchor=center, text width=0.74\paperwidth, align=center]
    at (current page.center){%
    {\color{sangueclaro}\scshape\Huge CUIDADO}\\[0.4em]
    {\color{sangueclaro}\rule{0.35\paperwidth}{0.6pt}}\\[2.2em]
    {\color{sepiabg}\large Este livro contém um ciclo de\\
     contaminação narrativa.}\\[1.8em]
    {\color{sepiabg}\large Ao ler, você se torna\\ o
     {\color{olho}\bfseries Paciente 1.263}.}\\[3em]
    {\color{sepiabg!55}\footnotesize\itshape ao final, não haverá como devolver o livro à estante\\
     como se você não o tivesse aberto.}\\[3em]
    {\color{sangueclaro}\footnotesize\itshape Se um zumbido surgir durante a leitura, não pare.\\
     Já é tarde demais para parar.}\\[1.2em]
    {\color{sepiabg}\footnotesize\itshape Esta página é parte da ficção. Se precisar parar, pare.}%
  };
\end{tikzpicture}%
\null\vfill\clearpage

% =================
% Copyright + Aviso ao Leitor — legível (raggedright), uma página
% =================
\cleardoublepage
\thispagestyle{empty}
{\small
\setlength{\parindent}{0em}
\setlength{\parskip}{0.7em}
~
\vspace{0.3em}
\vspace{1em}

2.ª edição revista, \editionyear{}\\
\publishercity{}: \publisher{}

Copyright \copyright{} \editionyear{} \authorname{} e \coauthorname{}

Todos os direitos reservados. Nenhuma parte deste livro pode ser reproduzida
por qualquer meio sem a autorização prévia e por escrito do autor, exceto
para citações breves em resenhas e artigos.

\textbf{ISBN:} \isbn{}. \quad Impresso no Brasil.

\textbf{Ficha catalográfica (CIP):} ver página seguinte.

\vspace{0.4em}
{\footnotesize
\textbf{Aviso ao leitor.} Este romance é uma obra de ficção baseada em
eventos históricos documentados. Os campos de concentração do Ceará (1915),
o Hospital Colônia de Barbacena e as cerca de 60.000 mortes em suas
dependências são fatos reais. Joaquim Antônio Correia, o Dr. Heitor Oliveira
e a Dra. Lúcia Mendes são personagens ficcionais. O ciclo de contaminação
descrito neste livro é uma construção narrativa.

\textit{Qualquer semelhança com a sua experiência de leitura é inteiramente
intencional.}

O autor se isenta de qualquer responsabilidade por noites mal dormidas;
elas fazem parte do pacto ficcional.}

\vfill
\begin{center}
\includegraphics[width=0.16\linewidth]{frontmatter/logo-black.png}\\[0.4em]
{\scshape Publicado por \publisher{} \quad·\quad \editionyear{}}
\end{center}
}
\clearpage

% =================
% Ficha catalográfica (CIP) — oficial
% =================
\cleardoublepage
\thispagestyle{empty}
{\small
\setlength{\parindent}{0pt}\setlength{\parskip}{0.7em}
~\vfill

\begin{center}
\textbf{Dados Internacionais de Catalogação na Publicação (CIP) --- ficha provisória}
\end{center}

\vspace{1em}

\noindent\fbox{%
\begin{minipage}{0.92\textwidth}
Carvalho Filho, Marcelo Dias de.\\
\hspace*{1.5em}\textit{Molambudos — O Diário do Paciente 1.260} / Marcelo Dias
de Carvalho Filho, Marcelo Claro. — 2. ed. rev. — \publishercity: \publisher,
\editionyear{}.\\
\hspace*{1.5em}ISBN \isbn{}.\\[0.5em]
\hspace*{1.5em}1. Ficção brasileira. 2. Horror psicológico. 3. Narrativa
fragmentária. 4. Literatura contemporânea. 5. Memória e trauma.
I. Carvalho Filho, Marcelo Dias de. II. Claro, Marcelo. III. Título.\\[0.5em]
\hspace*{9em}CDD: 869.93 \quad CDU: 821.134.3(81)-3
\end{minipage}}

\vspace{1em}

\noindent\textbf{Índice para catálogo sistemático:}\\
1. Ficção brasileira: 869.93

\vfill
}
\clearpage

\cleardoublepage
% =============================================================================
% NOTA DO ARQUIVISTA — Moldura Diegética do Arquivo 1.260
% =============================================================================
\cleardoublepage
\thispagestyle{empty}
\chapter*{Nota do Arquivista}
\addcontentsline{toc}{chapter}{Nota do Arquivista}
\markboth{Nota do Arquivista}{}

\begin{center}
\itshape\small
Este envelope contém a transcrição integral do diário do Paciente 1.260.\\
A leitura do material original causou alterações psicossomáticas nos examinadores.\\[0.5em]
\upshape\scriptsize --- Anotação em caneta-tinteiro no envelope pericial nº 1.260/1979
\end{center}

\vspace{1.2em}

\begingroup

O volume que você tem em mãos não foi escrito para ser publicado. É a reconstituição pericial do diário de \textbf{Joaquim Antônio Correia}, o Paciente 1.260 do Hospital Colônia de Barbacena.

Sobrevivente da Seca de 1915 e dos campos de concentração de retirantes no Ceará, Joaquim foi internado no Colônia em 1917 e lá permaneceu por 62 anos, até sua morte em 1979. Ao longo desse cárcere, produziu 847 páginas de anotações fragmentárias. O caderno original foi encontrado amarrado ao seu peito por tiras de pano quando seu corpo foi exumado da vala comum --- 62 dias após o sepultamento.

\vspace{0.4em}
\noindent Segundo o relato pericial do coveiro, o papel emanava um odor adocicado de umidade e terra, e a tinta parecia não secar. O laudo do IML registrou o achado como objeto de custódia e o anexou ao prontuário.

\vspace{0.6em}
Esta edição reúne 84 fragmentos organizados em cinco dimensões: as memórias da infância no sertão; o isolamento e as torturas no hospital psiquiátrico; a devassa dos laudos clínicos; a investigação forense conduzida pela Dra. Lúcia Mendes em 2026; e o registro da contaminação do próprio leitor.

\vspace{0.6em}
Não se trata de uma simples história de horror. A dor histórica que inspira estas páginas não é ficção. O prontuário de Joaquim, porém, é uma reconstituição ficcional; os números abaixo são estimativas atribuídas à bibliografia, não contagens deste arquivo.

\vspace{1.2em}
\noindent{\itshape
Recomenda-se não ler este livro no escuro.\\[0.35em]
Recomenda-se não ler este livro em solidão absoluta.\\[0.35em]
Se o cheiro de terra molhada surgir durante a leitura, feche o volume imediatamente.}

\vspace{1.2em}
\noindent A ferida histórica sobre a qual esta narrativa se assenta é real. O Hospital Colônia de Barbacena (1903--1994) internou à revelia cerca de 60 mil pessoas; 1.853 corpos foram exumados de suas valas comuns, segundo Daniela Arbex. Joaquim é um personagem, mas a dor que ele carrega --- fome, deslocamento, confinamento, tortura psiquiátrica --- pertenceu a milhares de brasileiros anônimos.

\vspace{0.6em}
\noindent O horror aqui evocado não deve ser consumido como entretenimento. A manipulação do leitor é uma escolha ética e formal --- não para enganar, mas para fazer o corpo sentir o que a mente sabe: que a vala não é uma metáfora, e que o sofrimento documentado existe para ser testemunhado, não estetizado.

\vspace{0.6em}
\noindent{\itshape Você será manipulado. Você será contaminado. Mas, ao final, saberá por quê.}

\vspace{1.4em}
\hfill{\scshape --- O Arquivista}\hspace{1em}

\endgroup
\clearpage

\cleardoublepage
\thispagestyle{empty}
\begin{rpgdark}[PROTOCOLO DE MESA]
\centering
{\textgoth{\fontsize{24}{26}\selectfont PROTOCOLO}}\\[0.25em]
{\scshape\bfseries\small QUATRO MODOS DE ENTRAR NO ARQUIVO}\\[0.35em]
{\color{rpggold}\rule{0.55\textwidth}{0.7pt}}
\end{rpgdark}

\vspace{0.5em}
\noindent{\small\itshape Este é um aparato de leitura, não um jogo completo. A obra pode ser interrompida a qualquer momento; a rota não muda o texto, só decide por onde você entra.}

\vspace{0.6em}
\begin{rpgbox}[1 · LEITURA LINEAR]
{\color{rpgblood}\faDiceD20}\enspace \textbf{Siga a ordem das páginas.} Do Sertão (1914) ao Epílogo. É o caminho recomendado para a primeira leitura: acompanha Joaquim e a propagação da entidade em ordem temporal.
\end{rpgbox}

\begin{rpgbox}[2 · ROTA INVESTIGATIVA]
{\color{rpgblood}\faDiceD20}\enspace \textbf{Comece pela Parte 4 (Investigação Lúcia, 2026).} Descubra as anomalias junto com a psiquiatra forense e volte às Partes 1, 2 e 3 para remontar as provas como uma autópsia.
\end{rpgbox}

\begin{rpgbox}[3 · NAVEGAÇÃO HIPERTEXTUAL]
{\color{rpgblood}\faDiceD20}\enspace \textbf{Siga as rotas ao final de cada fragmento.} Cada salto reflete a fragmentação do arquivo. Nenhum caminho é um beco: de qualquer fragmento existe percurso até o Epílogo.
\end{rpgbox}

\begin{rpgbox}[4 · ROTAS CURADAS]
{\color{rpgblood}\faDiceD20}\enspace \textbf{Sequências pré-selecionadas.} Rota do Terror: CONT-03 $\to$ CONT-04 $\to$ MEM-12 $\to$ CONT-05. Rota da Compaixão: MEM-01 $\to$ MEM-05 $\to$ MEM-10 $\to$ MEM-15. Rota Forense: DOC-01 $\to$ LUC-01 $\to$ DOC-16 $\to$ CONT-05.
\end{rpgbox}

\vspace{0.5em}
\begin{rpgbox}[ROTA ALEATÓRIA · role 1d20]
\centering
\small
\begin{tabular}{@{}cl@{}}
\toprule
\textbf{d20} & \textbf{Entrada} \\
\midrule
1--5 & Sertão \\
6--10 & Colônia \\
11--15 & Diário de Oliveira e Laudos \\
16--20 & Investigação Lúcia \\
\bottomrule
\end{tabular}

\vspace{0.4em}
{\itshape A Parte 5 (Contaminação) deve ser lida por último em qualquer percurso. Ela não é sobre a história. Ela é sobre você.}
\end{rpgbox}

\vspace{0.5em}
\begin{rpgbox}[CÓDIGOS DO ARQUIVO]
\small
\textcolor{mem}{\textbf{MEM}} Memória · \textcolor{doc}{\textbf{DOC}} Documento · \textcolor{luc}{\textbf{LUC}} Lucidez · \textcolor{cont}{\textbf{CONT}} Contaminação
\end{rpgbox}
\clearpage

\cleardoublepage
\thispagestyle{empty}
\begin{rpgdark}[FICHA DO PACIENTE 1.263]
\centering
{\textgoth{\fontsize{26}{28}\selectfont FICHA}}\\[0.2em]
{\scshape\bfseries\small ARQUIVO 1.260 · FORMULÁRIO DE ACOMPANHAMENTO}\\[0.3em]
{\color{rpggold}\rule{0.6\textwidth}{0.7pt}}
\end{rpgdark}

\vspace{0.5em}
\begin{rpgbox}[IDENTIFICAÇÃO]
\small
\renewcommand{\arraystretch}{1.35}
\begin{tabularx}{\textwidth}{@{}lX@{}}
\textbf{Nome:} & \hrulefill \\
\textbf{Idade:} & \hrulefill \\
\textbf{Ocupação:} & \hrulefill \\
\textbf{Procedência:} & \hrulefill \\
\end{tabularx}
\end{rpgbox}

\vspace{0.4em}
\begin{rpgbox}[ATRIBUTOS]
\small
\renewcommand{\arraystretch}{1.6}
\begin{tabularx}{\textwidth}{@{}XX@{}}
\textbf{\textcolor{mem}{MEMÓRIA}} & \textbf{\textcolor{doc}{LUCIDEZ}} \\
\rpgdot\;\rpgdot\;\rpgdot\;\rpgdot\;\rpgdot & \rpgdot\;\rpgdot\;\rpgdot\;\rpgdot\;\rpgdot \\
\textbf{\textcolor{cont}{VONTADE}} & \textbf{\textcolor{luc}{CONTAMINAÇÃO}} \\
\rpgdot\;\rpgdot\;\rpgdot\;\rpgdot\;\rpgdot & \rpgdot\;\rpgdot\;\rpgdot\;\rpgdot\;\rpgdot \\
\end{tabularx}
\end{rpgbox}

\vspace{0.4em}
\begin{rpgbox}[TRILHA DE CONTAMINAÇÃO · 0--20]
\centering
\begin{tikzpicture}[baseline]
  \foreach \i in {1,...,20} {
    \node[draw=rpggold, fill=rpgpaper, minimum size=3.3mm, inner sep=0pt]
      at (\i*4.7mm,0) {\tiny\color{rpgsepia}\i};
  }
\end{tikzpicture}

\vspace{0.35em}
{\footnotesize\itshape Marque um quadrado por contato direto com o diário, por noite sem sono ou por leitura da Parte 5.}
\end{rpgbox}

\vspace{0.4em}
\begin{rpgbox}[SINTOMAS E SINAIS]
\small
\begin{itemize}[leftmargin=1.4em,itemsep=0.15em]
  \item[$\square$] Cheiro adocicado de terra molhada.
  \item[$\square$] Olho direito amarelo no espelho.
  \item[$\square$] Despertar às 3h14 sem motivo.
  \item[$\square$] Sonho com a vala e com a voz que canta.
  \item[$\square$] Vontade de passar o diário adiante.
\end{itemize}
\end{rpgbox}

\vspace{0.4em}
\begin{rpgbox}[REGISTRO]
\small
\renewcommand{\arraystretch}{1.5}
\begin{tabularx}{\textwidth}{@{}lX@{}}
\textbf{Assinatura:} & \hrulefill \\
\textbf{Data:} & \hrulefill \\
\end{tabularx}
\end{rpgbox}

\vspace{0.4em}
\begin{rpgbox}[NOTAS DO JOGADOR]
\small
\foreach \i in {1,...,10} {\noindent\rule{\textwidth}{0.25pt}\par\vspace{0.55em}}
\end{rpgbox}

\vfill
\noindent{\footnotesize\itshape Formulário opcional, sem valor clínico. Uso ficcional. Não deve ser enviado a ninguém.}
\clearpage

% =================
% LEGENDA DOS CÓDIGOS — sistema visual R1076
% =================
\chapter*{Legenda dos Códigos}
\addcontentsline{toc}{chapter}{Legenda dos Códigos}
\markboth{Legenda dos Códigos}{}

\noindent Cada fragmento pertence a uma família. A cor e a capitular de abertura
indicam essa família antes de o leitor começar a ler.

\vspace{1.2em}
\begin{description}
  \item[\textcolor{mem}{\large\bfseries MEM}] \textbf{Memória.} A voz de Joaquim: sertão, fome, colônia, o diário. Cor: cobre queimado.
  \item[\textcolor{doc}{\large\bfseries DOC}] \textbf{Documento.} Laudos, prontuários, cartas, registros e provas. Cor: teal profundo.
  \item[\textcolor{luc}{\large\bfseries LUC}] \textbf{Lucidez.} A investigação de Lúcia Mendes em 2026. Cor: violeta.
  \item[\textcolor{cont}{\large\bfseries CONT}] \textbf{Contaminação.} A segunda pessoa, o leitor, o ciclo. Cor: carmim.
\end{description}

\vspace{1.2em}
\noindent\textit{O Epílogo usa o âmbar do arquivo; a capitular de cada fragmento é a inicial da
família (\textbf{M}emória, \textbf{D}ocumento, \textbf{L}ucidez,
\textbf{C}ontaminação), e o filete abaixo do título usa a mesma cor.}

\vspace{1em}
\noindent\textit{Contraste medido sobre o papel sépia: MEM 6,76:1; DOC 6,76:1;
LUC 9,23:1; CONT 7,18:1. Todas as cores de leitura passam o mínimo de 4,5:1.}

% =================
% Table of contents (uma página)
% =================
\begingroup
\tableofcontents
\thispagestyle{empty}
\endgroup
\cleardoublepage

\clearpage
\thispagestyle{empty}
\begin{rpgdark}[TABELA DE ROTAS · d20]
\centering
{\textgoth{\fontsize{22}{24}\selectfont TABELA DE ROTAS}}\\[0.25em]
{\scshape\bfseries\small VINTE PORTAS PARA O ARQUIVO}\\[0.3em]
{\color{rpggold}\rule{0.5\textwidth}{0.6pt}}
\end{rpgdark}

\vspace{0.4em}
\noindent{\small\itshape Role 1d20 e escolha uma entrada. O texto não muda; a rota só decide por onde você entra. As páginas são as desta edição.}

\vspace{0.25em}
\begin{rpgbox}[ENTRADAS]
\footnotesize
\renewcommand{\arraystretch}{0.98}
\begin{tabularx}{\textwidth}{@{}c l >{\raggedright\arraybackslash}X r@{}}
\toprule
\textbf{d20} & \textbf{Fragmento} & \textbf{Título} & \textbf{Página} \\
\midrule
1 & MEM-01 & O Sertão Antes do Fim & \pageref{frag:MEM-01} \\
2 & MEM-04 & A Vala & \pageref{frag:MEM-04} \\
3 & MEM-10 & O Primeiro Homicídio & \pageref{frag:MEM-10} \\
4 & MEM-14 & O Diário Emerge & \pageref{frag:MEM-14} \\
5 & MEM-20 & A Morte do Vizinho de Cela & \pageref{frag:MEM-20} \\
6 & MEM-25 & A Morte no Espelho & \pageref{frag:MEM-25} \\
7 & MEM-27 & A Página Seguinte & \pageref{frag:MEM-27} \\
8 & DOC-01 & A Exumação & \pageref{frag:DOC-01} \\
9 & DOC-03 & Laudo de Necropsia & \pageref{frag:DOC-03} \\
10 & DOC-06 & Relatório de Sessão III & \pageref{frag:DOC-06} \\
11 & DOC-13 & Diário de Oliveira III & \pageref{frag:DOC-13} \\
12 & DOC-19 & Holocausto Brasileiro & \pageref{frag:DOC-19} \\
13 & DOC-20 & Termo de Cadeia de Custódia & \pageref{frag:DOC-20} \\
14 & LUC-01 & O Convite & \pageref{frag:LUC-01} \\
15 & LUC-02 & O Arquivo & \pageref{frag:LUC-02} \\
16 & LUC-09 & A Exumação de Oliveira & \pageref{frag:LUC-09} \\
17 & LUC-17 & O Medo & \pageref{frag:LUC-17} \\
18 & LUC-20 & O Arquivo Aberto & \pageref{frag:LUC-20} \\
19 & CONT-01 & Você Abriu o Diário & \pageref{frag:CONT-01} \\
20 & CONT-13 & A Fila & \pageref{frag:CONT-13} \\
\bottomrule
\end{tabularx}
\end{rpgbox}

\vspace{0.4em}
\noindent{\footnotesize\itshape A tabela é um aparato de leitura; não altera a ordem, o texto ou o desfecho. O índice completo está nas páginas de navegação.}
\clearpage


% ===== Navegação pré-textual =====
\pagestyle{fancy}
\ifdefined\MolambudosKDPNoFolios
  \pagestyle{empty}
\fi
\ifdefined\MolambudosKDPSafeFolios
  \pagestyle{molambudos-kdp-folios}
\fi

% ============================================
% SEÇÃO DE NAVEGAÇÃO — Leitura Não-Linear
% ============================================
\chapter*{Navegação --- Rotas de Leitura}
\addcontentsline{toc}{chapter}{Navegação --- Rotas de Leitura}

\noindent\textit{Este arquivo contém 84 fragmentos organizados em 5 partes.}
\textit{Cada fragmento é independente e pode ser lido em qualquer ordem.}
\textit{As setas ↪ ao final de cada fragmento indicam rotas de leitura possíveis.}
\textit{A leitura linear recomendada segue a ordem das partes abaixo.}

\vspace{1em}

\section*{Índice de Fragmentos}

% Compact fraglink for the index (reduces paragraph spacing locally)
\ifmolambudoskdpprint
  \newcommand{\fraglink}[3]{%
    \par\noindent%
    \fragcode{#1} --- #2\,{\footnotesize(p.~\pageref{frag:#1})}%
    \ifx&#3&\else\hfill {\itshape\textcolor{textgray}{#3}}\fi%
    \ifdefined\MolambudosCompactIndex\par\vspace{0pt}\else\par\vspace{0.1em}\fi}
\else
  \newcommand{\fraglink}[3]{%
    \par\noindent%
    \hyperlink{#1}{\fragcode{#1}} --- #2\,{\footnotesize(p.~\pageref{frag:#1})}%
    \ifx&#3&\else\hfill {\itshape\textcolor{textgray}{#3}}\fi%
    \par\vspace{0.1em}}
\fi

\subsection*{Parte 1 --- Sertão (9 fragmentos)}
{\setlength{\parskip}{0.15em}
\fraglink{MEM-01}{O Sertão Antes do Fim}{1914}
\fraglink{MEM-02}{A Caminhada dos Retirantes}{1915}
\fraglink{MEM-03}{A Morte do Pai}{1915}
\fraglink{MEM-04}{A Vala}{1915}
\fraglink{MEM-05}{O Olho Amarelo}{1915}
\fraglink{MEM-06}{O Curral do Governo}{1915}
\fraglink{MEM-07}{A Mãe Levada}{1915}
\fraglink{DOC-17}{Carta de um Soldado}{1915}
\fraglink{DOC-18}{Relatório Oficial do Curral}{1916}
}

\subsection*{Parte 2 --- Colônia (16 fragmentos)}
{\setlength{\parskip}{0.15em}
\fraglink{MEM-08}{A Chegada ao Colônia}{1917}
\fraglink{DOC-02}{Prontuário do Paciente 1.260}{1917}
\fraglink{MEM-09}{Os Primeiros Anos}{1920--1930}
\fraglink{MEM-10}{O Primeiro Homicídio}{1922}
\fraglink{MEM-11}{O Isolamento}{1923--1945}
\fraglink{MEM-17}{O Eletrochoque}{1937}
\fraglink{MEM-18}{A Tentativa de Fuga}{1943}
\fraglink{MEM-12}{A Rasga Mortalha}{1945}
\fraglink{MEM-21}{O Jornal}{1948}
\fraglink{MEM-19}{O Visitante}{1950}
\fraglink{DOC-10}{Registro de Eletroconvulsoterapia}{1937}
\fraglink{MEM-22}{As Vozes}{1952}
\fraglink{MEM-23}{A Neve de Cinzas}{1960}
\fraglink{MEM-20}{A Morte do Vizinho de Cela}{1963}
\fraglink{MEM-24}{A Enfermeira que Chorava}{1964}
\fraglink{MEM-25}{A Morte no Espelho}{1970}
}

\subsection*{Parte 3 --- Diário de Oliveira\\\hspace*{2.5em}e Laudos (25 fragmentos)}
{\setlength{\parskip}{0.15em}
\fraglink{DOC-01}{A Exumação}{1979}
\fraglink{DOC-03}{Laudo de Necropsia}{1979}
\fraglink{DOC-04}{Relatório de Sessão I}{1979}
\fraglink{DOC-05}{Relatório de Sessão II}{1979}
\fraglink{DOC-06}{Relatório de Sessão III}{1979}
\fraglink{DOC-07}{Diagnóstico Psiquiátrico (1979)}{1979}
\fraglink{MEM-13}{O Médico Novo}{1979}
\fraglink{MEM-14}{O Diário Emerge}{1979}
\fraglink{MEM-15}{A Entidade Fala}{1979}
\fraglink{MEM-16}{A Morte no Colônia}{1979}
\fraglink{MEM-26}{A Última Página}{1979}
\fraglink{DOC-11}{Diário de Oliveira I}{1979}
\fraglink{DOC-12}{Diário de Oliveira II}{1979}
\fraglink{DOC-13}{Diário de Oliveira III}{1979}
\fraglink{DOC-14}{Diário de Oliveira IV}{1979}
\fraglink{DOC-15}{Diário de Oliveira V}{1979}
\fraglink{DOC-16}{Diário de Oliveira VI}{1979}
\fraglink{DOC-19}{Holocausto Brasileiro}{Daniela Arbex}
\fraglink{DOC-20}{Termo de Cadeia de Custódia}{2026}
\fraglink{DOC-21}{Carta de Heitor Oliveira (não enviada)}{1979}
\fraglink{DOC-22}{Memorando Interno nº~47/2026}{2026}
\fraglink{DOC-23}{Registro de Ocorrência do Arquivo}{2026}
\fraglink{DOC-24}{Parecer Técnico de Conservação}{2026}
\fraglink{DOC-25}{O Homem do Pano Preto}{2026}
\fraglink{DOC-26}{Anotações de Oliveira — O Mecanismo}{1979}
}

\subsection*{Parte 4 --- Investigação Lúcia (26 fragmentos)}
{\setlength{\parskip}{0.15em}
\fraglink{DOC-08}{Re-Diagnóstico Psiquiátrico (2026)}{Lúcia Mendes}
\fraglink{DOC-09}{Escala de Avaliação do Paciente 1.263}{Instrumento clínico}
\fraglink{DOC-27}{A Origem (1853)}{Padre José Inácio}
\fraglink{LUC-01}{O Convite}{2026}
\fraglink{LUC-15}{A Mulher que Não Voltou para Casa}{2026}
\fraglink{LUC-02}{O Arquivo}{2026}
\fraglink{LUC-03}{As Ruínas}{2026}
\fraglink{LUC-04}{A Fotografia}{2026}
\fraglink{LUC-05}{O Sonho de Lúcia}{2026}
\fraglink{LUC-Contraprova}{A Contraprova}{2026}
\fraglink{LUC-Entrevista}{A Entrevista}{2026}
\fraglink{LUC-16}{A Mãe no Telefone}{2026}
\fraglink{LUC-Escolha}{A Escolha}{2026}
\fraglink{LUC-06}{O Desaparecimento}{2026}
\fraglink{LUC-07}{O Rastro de Oliveira}{2026}
\fraglink{LUC-08}{A Última Sessão de Lúcia}{2026}
\fraglink{LUC-09}{A Exumação de Oliveira}{2026}
\fraglink{LUC-10}{O Laudo de Lúcia}{2026}
\fraglink{LUC-11}{O Diário Digital}{2026}
\fraglink{LUC-12}{A Última Foto}{2026}
\fraglink{LUC-17}{O Medo}{2026}
\fraglink{LUC-13}{A Noite do Diagnóstico}{2026}
\fraglink{LUC-18}{O Que Ronaldo Não Disse}{2026}
\fraglink{LUC-14}{Autoavaliação Clínica (DSM-5-TR)}{2026}
\fraglink{LUC-19}{A Última Noite em Casa}{2026}
\fraglink{LUC-20}{O Arquivo Aberto}{2026}
}

\subsection*{Parte 5 --- Contaminação (7 fragmentos + Epílogo)}
{\setlength{\parskip}{0.15em}
\fraglink{CONT-01}{Você Abriu o Diário}{}
\fraglink{CONT-03}{Você Sente o Cheiro}{}
\fraglink{CONT-04}{Você Começa a Sonhar}{}
\fraglink{CONT-05}{LAUDO DE AVALIAÇÃO CLÍNICA}{}
\fraglink{CONT-07}{O Diagnóstico Final}{}
\fraglink{CONT-09}{O Espelho}{}
\fraglink{CONT-13}{A Fila}{}
\fraglink{MEM-27}{A Página Seguinte}{Joaquim, 1907--1979}
}

\newpage

% ============================================
% GRAFOS NARRATIVOS — Navegação Visual da Obra
% ============================================
\chapter*{Navegação Visual --- Os Mapas do Arquivo}
\addcontentsline{toc}{chapter}{Navegação Visual --- Mapas}

\noindent\textit{Nas páginas seguintes, três mapas ajudam a orientar sua leitura.}
\textit{Cada um serve a um propósito diferente: explorar, compreender ou seguir.}
\textit{Nenhum é obrigatório. Todos são rotas possíveis.}

\vspace{0.6em}
\noindent\textbf{Mapa 1} --- O \textbf{grafo de rotas} mostra todos os fragmentos e suas conexões. Use-o para \textit{navegação não linear}: escolha um fragmento e siga as setas.

\noindent\textbf{Mapa 2} --- O \textbf{grafo do enredo} revela a arquitetura secreta da obra em três atos. Nele estão as causas, as transmissões e o ciclo que conecta os quatro hospedeiros --- inclusive você.

\noindent\textbf{Mapa 3} --- O \textbf{grafo linear} é o percurso tradicional, do Sertão (1914) ao Epílogo. Recomendado para quem prefere começar pelo começo.

\vspace{0.8em}
\noindent{\small\textit{A escolha do caminho não altera o destino: todas as rotas convergem para o mesmo ciclo.}}

%% --- GRAFO 1: MAPA DE ROTAS (rotacionado 90°) ---
\newpage
\thispagestyle{empty}
\begingroup
\vbox to \textheight{%
  \vfill
  \centering
  \setlength{\fboxsep}{3pt}%
  \setlength{\fboxrule}{0.6pt}%
  \fbox{\includegraphics[angle=90,width=0.98\textwidth,height=0.90\textheight,keepaspectratio]{misc/grafo_narrativo.png}}%
  \vfill
}%
\endgroup
\newpage

\section*{Mapa 1 --- Grafo de Leitura: Rotas entre Fragmentos}

\noindent Este mapa representa os \textbf{84 fragmentos} (cada um é um nodo) e suas \textbf{223 conexões} (as arestas que os ligam). Cada fragmento aponta para outros ao final do texto através dos links \texttt{↪} acompanhados do número da página.

\noindent Quanto mais conexões um fragmento tem, mais central ele é na teia da narrativa. Os mais conectados são:

\begin{itemize}
\itemsep0em
\item \textbf{DOC-01} (A Exumação) --- 19 rotas chegam até ele; o diário emerge do passado
\item \textbf{LUC-02} (O Arquivo) --- 13 rotas; a investigação de Lúcia é o segundo polo
\item \textbf{MEM-12} (A Rasga Mortalha) e \textbf{CONT-07} (O Diagnóstico Final) --- 7 rotas cada; memória e diagnóstico disputam o centro
\end{itemize}

\vspace{0.5em}
\noindent\textbf{Legenda:}
\textcolor{mem}{MEM (marrom)} --- Memórias do Paciente 1.260 \quad
\textcolor{doc}{DOC (cinza)} --- Documentos e laudos \quad
\textcolor{cont}{CONT (vermelho)} --- Contaminação do leitor \quad
\textcolor{luc}{LUC (índigo)} --- Investigação forense

%% --- GRAFO 2: ESTRUTURA DO ENREDO ---
\newpage
\thispagestyle{empty}
\begingroup
\vbox to \textheight{%
  \vfill
  \centering
  \setlength{\fboxsep}{3pt}%
  \setlength{\fboxrule}{0.6pt}%
  \fbox{\includegraphics[angle=90,width=0.98\textwidth,height=0.90\textheight,keepaspectratio]{misc/grafo_enredo.png}}%
  \vfill
}%
\endgroup
\newpage

\section*{Mapa 2 --- Grafo do Enredo: Arquitetura em Três Atos}

\begingroup\raggedright

\noindent Diferente do mapa anterior (que mostra \textit{conexões}), este grafo revela a \textit{estrutura causal} da obra --- o esqueleto narrativo que sustenta os 84 fragmentos.

\vspace{1em}
\noindent\textbf{Os três atos:}

\vspace{0.4em}
\noindent\textcolor{mem}{\textbf{ATO 1 --- O Trauma}} \textit{(marrom)}:
a seca de 1853 que pariu a entidade, a seca de 1915 que levou Joaquim à vala,
o pacto de um menino de oito anos com a fome que nunca acaba.\\[0.35em]
\noindent\textcolor{doc}{\textbf{ATO 2 --- Institucionalização}} \textit{(cinza)}:
62 anos no Colônia, o diário que testemunha, a exumação que o devolve ao
mundo, e o médico que ao ler se torna o próximo elo.\\[0.35em]
\noindent\textcolor{cont}{\textbf{ATO 3 --- O Ciclo}} \textit{(vermelho)}:
a investigação de Lúcia, a descoberta de que o horror é um ciclo,
a contaminação do leitor --- e a constatação final de que
\textit{o ciclo recomeça a cada leitor}.

\vspace{0.9em}
\noindent\textbf{A fila dos hospedeiros:} as arestas tracejadas
(\textcolor{sangue}{-\,-\,-}) ligam quem carrega a criatura, um após o outro
--- Joaquim (1.260), Oliveira (1.261), Lúcia (1.262) e você (1.263) ---,
numa corrente que começou em 1853 e não tem fim anunciado.

\vspace{0.5em}
\noindent\textit{Nota: os fragmentos DOC-27 (A Origem, 1853) e LUC-13 (A Noite do
Diagnóstico) foram acrescentados ao arquivo durante a pesquisa de Lúcia.
DOC-27 é o documento mais antigo do acervo (1853) e revela a origem da
entidade; LUC-13 registra o momento exato em que Lúcia deixa de ser
investigadora para tornar-se hospedeira.}

\endgroup

%% --- GRAFO 3: PERCURSO LINEAR (rotacionado 90°) ---
\newpage
\thispagestyle{empty}
\begingroup
\vbox to \textheight{%
  \vfill
  \centering
  \setlength{\fboxsep}{3pt}%
  \setlength{\fboxrule}{0.6pt}%
  \fbox{\includegraphics[angle=90,width=0.98\textwidth,height=0.90\textheight,keepaspectratio]{misc/grafo_leitura_linear.png}}%
  \vfill
}%
\endgroup
\newpage

\section*{Mapa 3 --- Grafo de Leitura Linear: Percurso Tradicional}

\noindent Para quem prefere uma leitura sequencial, este grafo organiza os
\textbf{84 fragmentos} na ordem cronológico-narrativa recomendada:
do Sertão (1914) ao Epílogo --- passando pelo Colônia, pelos diários de
Oliveira e pela investigação de Lúcia, até a contaminação final.

\vspace{0.8em}

\noindent\textbf{Percurso linear:}
\textcolor{mem}{Sertão (9)} $\to$ \textcolor{doc}{Colônia (16)} $\to$
\textcolor{doc}{Diário de Oliveira e Laudos (25)} $\to$
\textcolor{luc}{Investigação Lúcia (26)} $\to$
\textcolor{cont}{Contaminação (7)} $\to$
\textcolor{olho}{MEM-27 (1)}

\vspace{0.3em}
\noindent{\small $\bigstar$ = Fragmentos marcados com estrela são obrigatórios na rota essencial mínima.}

\vspace{0.7em}
\noindent\textit{Dica de leitura:} se esta é sua primeira vez no arquivo,
comece pelo Sertão. As memórias de Joaquim preparam o terreno para tudo
o que vem depois. Se preferir o mistério, comece pela Investigação de Lúcia
(Parte 4) e descubra os fragmentos anteriores como quem escava um passado.

\vspace{0.3em}
\noindent{\small \textcolor{sangue}{\textbf{Atenção:}} a Parte 5 (Contaminação) deve ser lida por último em qualquer percurso.}
\textit{Ela não é sobre a história. Ela é sobre você.}

\newpage

% ============================================
% LEITURA LINEAR — 84 FRAGMENTOS EM 5 PARTES + EPÍLOGO
% ============================================

%% ============================================
%% PARTE 1 — SERTÃO (1915--1928)
%% ============================================
\partopener{Sertão (1914--1917)}{Sertão}{2.png}

% O primeiro fragmento é a página arábica 1; não inserir \newpage antes dele.
\mainmatter
\hypersetup{pageanchor=true}
\fragdef{MEM-01}% ============================================
% MEM-01
% Gerado de molambudos.md — build_miolo.py
% ============================================
\fragtitle{mem}{M}{MEM}{01}{O Sertão Antes do Fim}
\markboth{MEM-01 --- O Sertão Antes do Fim}{}

\noindent \textbf{1914 --- Quixeramobim, Ceará}

\vspace{0.5em}\noindent\rule{\textwidth}{0.5pt}\vspace{0.5em}

\noindent Meu pai dizia que a terra tinha um jeito de avisar.

\noindent Primeiro o gado morre. Eles sentem antes da gente. O gado começa a cair --- não de doença, de fraqueza. O pasto some, o gado emagrece, os ossos aparecem por baixo do couro, e um dia o boi deita e não levanta mais. Meu pai matava os que ainda estavam em pé para não perder a carne. Mas não tinha água pra salgar. A carne apodrecia antes de a gente comer.

\noindent O calor não dava trégua nem de noite. A poeira entrava pela porta e cobria tudo com uma camada fina, cor de osso. Eu sentia o cheiro de carne estragada chegar antes do vento --- um cheiro doce e podre que, anos depois, eu reconheceria de novo, vindo de dentro de mim mesmo.

\noindent Depois do gado, o rio seca. O Quixeramobim era um rio que cantava. Eu ouvia ele da porta de casa --- um barulho de água nas pedras, baixinho, como quem conversa sozinho. Em 1914, ele parou de cantar. Primeiro ficou raso. Depois ficou barrento. Depois ficou seco. O leito parecia uma estrada de osso branco.

\noindent Depois do rio, as pessoas mudam.

\noindent Meu pai mudou primeiro. Ele começou a ficar horas sentado na soleira da porta, olhando pro nada. Minha mãe perguntava o que ele tinha. Ele dizia: \textit{\textquotedblleft{}Tô esperando.\textquotedblright{}} Não dizia o quê. Acho que ele mesmo não sabia.

\noindent Depois mudaram os vizinhos. Seu Lázaro, que era o homem mais alegre da região, parou de falar. Dona Sinhá, que fazia pão pra todo mundo, começou a guardar comida debaixo da cama. As crianças pararam de brincar.

\noindent Quando a comida acabou de vez, ninguém se despediu.

\noindent A gente simplesmente saiu.

\noindent Minha mãe arrumou um saco com três peças de roupa, uma cuia, e uma imagem de Nossa Senhora das Dores que era da avó dela. Meu pai não arrumou nada. Ele olhou a casa uma última vez. A porta aberta. A mesa vazia.

\noindent \textit{\textquotedblleft{}Joaquim\textquotedblright{}\textit{, ele disse. }\textquotedblleft{}Vai olhar pra frente. Não olha pra trás.\textquotedblright{}}

\noindent Eu olhei pra trás.

\noindent A casa estava queimando.

\noindent O fogo era a luz do sol poente que pegava no telhado de palha e fazia a casa parecer em chamas. Parecia que o sertão inteiro estava pegando fogo.

\noindent Eu pensei: \textit{a gente não vai voltar.}

\noindent A gente não voltou.

\vspace{0.5em}
\noindent\textit{\textbf{↪ Rotas:} \rota{MEM-02} \rotasep \rota{DOC-02}}
\newpage\fragdef{MEM-02}% ============================================
% MEM-02
% Gerado de molambudos.md — build_miolo.py
% ============================================
\fragtitle{mem}{M}{MEM}{02}{A Caminhada dos Retirantes}
\markboth{MEM-02 --- A Caminhada dos Retirantes}{}

\noindent \textit{1915 --- A estrada de Senador Pompeu a Fortaleza}

\noindent \textbf{1915 --- A estrada de Senador Pompeu a Fortaleza}

\vspace{0.5em}\noindent\rule{\textwidth}{0.5pt}\vspace{0.5em}

\noindent Depois de meses de retirada desde Quixeramobim, chegamos a Senador Pompeu. Em 1915, partimos dali na direção de Fortaleza, seguindo o rumor de que no litoral ainda havia comida.

\noindent No primeiro dia, a gente andou até o sol se pôr.

\noindent Meu pai na frente. Minha mãe atrás, com o saco de roupa na cabeça. Eu no meio, segurando a barra da saia dela. O chão era seco e rachado, e cada passo levantava uma nuvenzinha de poeira que entrava no nariz e deixava a boca com gosto de terra.

\noindent De vez em quando a gente passava por uma casa. As portas estavam abertas. As janelas também. Não tinha ninguém dentro. As casas estavam vazias como cascas de cigarra.

\noindent No segundo dia, a gente começou a encontrar os outros.

\noindent Eles vinham de todo lado. Famílias inteiras. Gente sozinha. Velhos carregando crianças. Crianças carregando velhos. Todo mundo andando na mesma direção.

\noindent Ninguém falava.

\noindent Falar gasta saliva. Saliva precisa de água. Água não tinha.

\noindent A gente andava.

\noindent No terceiro dia, a fila era tão grande que eu não via o fim. A estrada estava cheia de gente que se movia devagar, como um rio de lama. As pessoas tropeçavam umas nas outras. Algumas caíam e não levantavam. A gente passava por cima.

\noindent Minha mãe disse: \textit{\textquotedblleft{}Não olha, Joaquim.\textquotedblright{}}

\noindent Eu olhava.

\noindent No quarto dia, meu pai começou a andar diferente. Os passos ficaram mais curtos. A respiração ficou mais alta. Ele parou de olhar pra frente. Ele olhava para o chão, como se o chão tivesse uma resposta.

\noindent \textit{\textquotedblleft{}Pai\textquotedblright{}\textit{, eu disse. }\textquotedblleft{}Pai, eu quero água.\textquotedblright{}}

\noindent Ele não respondeu.

\noindent \textit{\textquotedblleft{}Pai.\textquotedblright{}}

\noindent Ele parou. Virou o rosto. A boca dele estava aberta. A língua estava branca. Os olhos estavam vermelhos.

\noindent \textit{\textquotedblleft{}Pai?\textquotedblright{}}

\noindent Ele tentou falar. Não saiu som. Só ar. Um ar que cheirava a estômago vazio.

\noindent Ele sentou no chão.

\noindent Minha mãe sentou do lado dele.

\noindent Eu sentei do lado dela.

\noindent A fila de retirantes passou por nós. Ninguém parou. Ninguém perguntou. Todo mundo estava ocupado morrendo.

\noindent No quinto dia, ele parou debaixo de um juazeiro seco. Não foi para descansar.

\noindent Minha mãe ajoelhou na frente dele. Ficaram se olhando. Ela não disse nada. Ele também não. Só balançou a cabeça.

\noindent Ela se levantou. Pegou minha mão.

\noindent \textit{\textquotedblleft{}Anda, Joaquim.\textquotedblright{}}

\noindent Eu olhei para trás. Ele ainda estava sentado. A cabeça caída. O braço estendido no colo.

\noindent Minha mãe puxou minha mão.

\noindent \textit{\textquotedblleft{}Anda.\textquotedblright{}}

\noindent A gente continuou andando.

\noindent A fila na frente. A fila atrás.

\noindent Eu no meio.

\noindent Aprendendo que o mundo é um lugar onde as pessoas ficam para trás e a gente continua andando.

\vspace{0.5em}
\noindent\textit{\textbf{↪ Rotas:} \rota{MEM-03} \rotasep \rota{MEM-01} \rotasep \rota{MEM-06}}
\newpage\fragdef{MEM-03}% ============================================
% MEM-03
% Gerado de molambudos.md — build_miolo.py
% ============================================
\fragtitle{mem}{M}{MEM}{03}{A Morte do Pai}
\markboth{MEM-03 --- A Morte do Pai}{}

\noindent \textit{O segundo dia sem água}

\noindent \textbf{1915 --- O segundo dia sem água}

\vspace{0.5em}\noindent\rule{\textwidth}{0.5pt}\vspace{0.5em}

\noindent Meu pai sentou debaixo de um juazeiro.

\noindent A árvore estava seca --- não tinha folha, não tinha sombra. Mas ele sentou. Encostou as costas no tronco e fechou os olhos.

\noindent \textit{\textquotedblleft{}Pai, vamos.\textquotedblright{}}

\noindent \textit{\textquotedblleft{}Vão vocês.\textquotedblright{}}

\noindent Minha mãe ajoelhou na frente dele.

\noindent \textit{\textquotedblleft{}Antônio, levanta.\textquotedblright{}}

\noindent \textit{\textquotedblleft{}Não dá.\textquotedblright{}}

\noindent \textit{\textquotedblleft{}Dá sim. Tu sempre disse que dava.\textquotedblright{}}

\noindent Ele abriu os olhos. Olhou para ela. O olhar dele estava longe. Como se ele já estivesse noutro lugar.

\noindent \textit{\textquotedblleft{}Maria, eu não sinto minhas pernas.\textquotedblright{}}

\noindent Ela olhou para as pernas dele. Estavam inchadas. Os pés estavam rachados, com sangue seco nos calcanhares.

\noindent \textit{\textquotedblleft{}Isso é cansaço\textquotedblright{}}, ela disse.

\noindent \textit{\textquotedblleft{}Isso é morte.\textquotedblright{}}

\noindent \textit{\textquotedblleft{}Não fala assim.\textquotedblright{}}

\noindent \textit{\textquotedblleft{}Maria.\textquotedblright{}\textit{ Ele colocou a mão no rosto dela. A mão dele era só osso. }\textquotedblleft{}Leva o menino. Vai. Não para.\textquotedblright{}}

\noindent \textit{\textquotedblleft{}Eu não vou deixar tu aqui.\textquotedblright{}}

\noindent \textit{\textquotedblleft{}Já estou morto. Só falta o corpo parar.\textquotedblright{}}

\noindent Minha mãe começou a chorar. Era a primeira vez que eu via ela chorar. O choro dela era seco --- não tinha lágrima. O corpo não tinha água para lágrimas.

\noindent \textit{\textquotedblleft{}Joaquim\textquotedblright{}\textit{, disse meu pai. }\textquotedblleft{}Vem cá.\textquotedblright{}}

\noindent Eu fui.

\noindent Ele colocou a mão na minha cabeça.

\noindent \textit{\textquotedblleft{}Tu é homem agora. Tu cuida da tua mãe.\textquotedblright{}}

\noindent \textit{\textquotedblleft{}Sim, pai.\textquotedblright{}}

\noindent \textit{\textquotedblleft{}Tu não deixa ela passar fome.\textquotedblright{}}

\noindent \textit{\textquotedblleft{}...Sim, pai.\textquotedblright{}}

\noindent \textit{\textquotedblleft{}Tu come o que tiver que comer. Mesmo que...\textquotedblright{}\textit{ Ele parou. Olhou para o horizonte. }\textquotedblleft{}Mesmo que não pareça comida.\textquotedblright{}}

\noindent \textit{\textquotedblleft{}Sim, pai.\textquotedblright{}}

\noindent Ele puxou meu rosto perto do dele. O hálito dele cheirava a estômago vazio.

\noindent \textit{\textquotedblleft{}Tu não morre de fome, Joaquim. Tu me escutou? Tu não morre de fome. Tu faz o que tiver que fazer. Mas tu não morre de fome.\textquotedblright{}}

\noindent \textit{\textquotedblleft{}Sim, pai.\textquotedblright{}}

\noindent Ele soltou meu rosto. Fechou os olhos.

\noindent \textit{\textquotedblleft{}Podem ir.\textquotedblright{}}

\noindent Minha mãe ficou parada.

\noindent \textit{\textquotedblleft{}Podem ir\textquotedblright{}\textit{, ele repetiu. }\textquotedblleft{}Já disse.\textquotedblright{}}

\noindent Ela levantou. Pegou minha mão.

\noindent A gente andou.

\noindent Quando eu olhei para trás, ele ainda estava sentado no juazeiro seco. A cabeça caída. O braço estendido no colo.

\noindent Minha mãe puxou minha mão. \textit{\textquotedblleft{}Vamos.\textquotedblright{}} Eu podia ter virado o rosto. Ela não olhou. Mas eu escolhi guardar aquela imagem. A última.

\noindent A última vez que eu vi meu pai, ele parecia um santo.

\noindent Um santo que morreu de fome.

\vspace{0.5em}
\noindent\textit{\textbf{↪ Rotas:} \rota{MEM-02} \rotasep \rota{MEM-07} \rotasep \rota{MEM-04}}
\newpage\fragdef{MEM-04}% ============================================
% MEM-04
% Gerado de molambudos.md — build_miolo.py
% ============================================
\fragtitle{mem}{M}{MEM}{04}{A Vala}
\markboth{MEM-04 --- A Vala}{}

\noindent \textit{O encontro}

\noindent \textbf{1915 --- Na rota de Senador Pompeu a Fortaleza}

\vspace{0.5em}\noindent\rule{\textwidth}{0.5pt}\vspace{0.5em}

\noindent Meu pai morreu de manhã. O sol nasceu e ele já estava duro, com os olhos abertos para o céu que não tinha nuvem. Minha mãe passou a mão nos olhos dele e disse: \textit{\textquotedblleft{}Anda, Joaquim. Teu pai já foi.\textquotedblright{}}

\noindent Eu tinha oito anos.

\noindent A gente vinha andando há quatro dias. Antes éramos seis. Agora éramos três: eu, minha mãe e o velho Anselmo, que não era parente mas caminhava junto porque sozinho ninguém sobrevive. O velho Anselmo disse que tinha visto a lagoa secar em 1877. Disse que aquilo sim foi seca. Eu olhei para ele e vi que ele não acreditava no que dizia. Os olhos dele estavam iguais aos do meu pai.

\noindent No quinto dia, a gente encontrou o comboio.

\noindent Eram soldados. Magros, sujos, com fardas remendadas. Um deles levava uma espingarda. O outro levava um chicote.

\noindent \textit{\textquotedblleft{}Pra onde vão?\textquotedblright{}}, perguntou o do chicote.

\noindent Minha mãe respondeu que íamos para o litoral. Que ouviu dizer que lá tinha comida.

\noindent O soldado riu. \textit{\textquotedblleft{}Comida? O litoral tá cheio de retirante que nem você. Tá todo mundo morrendo na praia.\textquotedblright{}}

\noindent Minha mãe não respondeu.

\noindent \textit{\textquotedblleft{}Vem com a gente\textquotedblright{}\textit{, disse o soldado. }\textquotedblleft{}Tem o Campo do Alagadiço, perto de Fortaleza. Lá dá comida.\textquotedblright{}}

\noindent Ela hesitou. Eu senti a mão dela apertar a minha.

\noindent \textit{\textquotedblleft{}É ordem\textquotedblright{}}, disse o soldado, e o chicote estalou no ar. Não bateu em ninguém. Mas o som ficou. Como se tivesse marcado a gente.

\vspace{0.5em}\noindent\rule{\textwidth}{0.5pt}\vspace{0.5em}

\noindent O Curral do Governo era uma cerca de arame farpado em volta de um descampado. Não tinha árvore. Não tinha sombra. Tinha gente. Muita gente. Mais gente do que eu tinha visto em toda a minha vida.

\noindent Eles estavam deitados no chão, encostados uns nos outros. Alguns mexiam a boca sem fazer som. Outros não mexiam nada.

\noindent \textit{\textquotedblleft{}Molambudos\textquotedblright{}}, alguém disse.

\noindent Eu não sabia o que significava. Depois eu soube. Significa farrapos. Trapos humanos. Gente que não é mais gente porque a seca levou tudo o que tinha dentro.

\noindent Ficamos ali. Minha mãe achou um canto perto da cerca. O velho Anselmo sentou e não levantou mais. No segundo dia, ele disse: \textit{\textquotedblleft{}Tô indo.\textquotedblright{}} E foi.

\noindent Minha mãe começou a tossir.

\vspace{0.5em}\noindent\rule{\textwidth}{0.5pt}\vspace{0.5em}

\noindent A vala ficava atrás do curral, a uns duzentos metros. Toda noite os soldados carregavam os corpos para lá. Durante o dia dava pra ver os urubus. Eles faziam um círculo no céu, rodando, rodando, como se o céu fosse um olho gigante.

\noindent Eu fui até a vala numa noite em que minha mãe tossia sem parar.

\noindent Não fui por curiosidade.
Fui porque ouvi uma coisa.

\noindent Era uma voz.

\noindent Não uma voz humana --- mas uma voz. Ela cantava. Não uma canção com letra. Uma melodia. Baixa. Rouca. Uma melodia que subia e descia como o vento num lugar vazio.

\noindent Ela cantava:

\noindent \textit{\textquotedblleft{}Quem tem fome senta. Quem tem fome espera. Quem tem fome come o que vier.\textquotedblright{}}

\noindent Eu parei na borda da vala.

\noindent Lá embaixo, entre os corpos, algo se mexeu.

\noindent Eu devia ter corrido. Devia ter voltado para o curral. Devia ter enterrado o rosto na saia da minha mãe e esquecido que tinha ouvido alguma coisa.

\noindent Mas eu era uma criança.
E crianças que passam fome não fogem de comida.
Elas investigam.

\vspace{0.5em}\noindent\rule{\textwidth}{0.5pt}\vspace{0.5em}

\noindent A vala não era funda. Uns três metros. Os corpos se empilhavam como sacos de farinha. Alguns estavam nus. Outros amarrados pelas mãos. Os olhos deles estavam abertos. Todos os olhos estavam abertos.

\noindent No meio da pilha, algo se moveu com cuidado.

\noindent Não era o movimento de quem luta para sair.
Era o movimento de quem está \textit{arrumando} os corpos à sua volta.

\noindent Um braço emergiu. Não um braço humano --- era fino demais, longo demais, com articulações onde não devia ter. A mão tinha dedos demais, e cada dedo tinha uma unha preta que brilhava como vidro molhado.

\noindent A mão pegou um braço solto na pilha e o colocou de lado. Com cuidado. Como quem arruma a mesa antes de uma refeição.

\noindent Outro braço surgiu. Depois a cabeça.

\noindent A cabeça era quase humana. Quase. A forma era de uma caveira, mas a pele ainda cobria o osso --- uma pele fina, cinzenta, esticada como papel molhado sobre os ossos do rosto.

\noindent O nariz era um buraco.
A boca era uma linha.
Os olhos eram duas fendas.

\noindent E dentro das fendas, um brilho amarelo.

\noindent Ela me olhou.

\noindent Não se mexeu. Não fez nada. Apenas me olhou, com a cabeça inclinada, como um cachorro ouvindo um som estranho.

\noindent Eu não corri.

\noindent Ela começou a subir.

\noindent Não como uma pessoa sobe. Como uma aranha. Os braços e as pernas se moviam em coordenação, encontrando apoio nos corpos, deslizando entre eles sem pressa. Os ossos estalavam sob o peso dela. A carne cedia.

\noindent Ela saiu da vala.

\noindent Ficou de pé na minha frente.

\noindent Ela era imensa --- não alta, mas \textit{larga}. O corpo era uma gaiola de costelas viradas para fora, como se o esqueleto tivesse crescido pelo avesso. Entre as costelas, algo pulsava. Um órgão grande, escuro, coberto de veias que brilhavam como minhocas.

\noindent O cheiro dela chegou até mim.

\noindent O cheiro não vinha da morte. Vinha de \textit{comida}. Carne assada. Caldo grosso. O cheiro que saía da cozinha de Dona Sinhá antes da seca. O cheiro que eu não sentia há meses.

\noindent Meu estômago doeu.

\noindent \textit{\textquotedblleft{}Cê tem fome?\textquotedblright{}}

\noindent A voz dela era macia. Quase doce. Não combinava com o corpo. Era a voz de uma avó perguntando se o neto já comeu.

\noindent Eu não respondi.

\noindent \textit{\textquotedblleft{}Cê tem.\textquotedblright{}\textit{ Ela inclinou a cabeça para o outro lado. }\textquotedblleft{}Eu sinto. Tua fome é grande. Grande que nem a minha.\textquotedblright{}}

\noindent \textit{\textquotedblleft{}A senhora...\textquotedblright{}\textit{ Minha voz saiu fraca. }\textquotedblleft{}A senhora é o quê?\textquotedblright{}}

\noindent Ela pensou. O órgão entre as costelas pulsou mais forte.

\noindent \textit{\textquotedblleft{}Sou o que sobra\textquotedblright{}\textit{, ela disse. }\textquotedblleft{}Sou o que cresce onde muito morre. Sou o que a terra não consegue engolir.\textquotedblright{}}

\noindent \textit{\textquotedblleft{}A senhora come os mortos?\textquotedblright{}}

\noindent \textit{\textquotedblleft{}Alguém tem que comer.\textquotedblright{}}

\noindent Ela disse isso com naturalidade. Como quem diz \textquotedblleft{}o sol nasce todo dia\textquotedblright{}.

\noindent \textit{\textquotedblleft{}Os vivos não querem. Os vivos enterram. Os vivos choram. Depois vão embora. Os mortos ficam. Alguém tem que aproveitar.\textquotedblright{}}

\noindent Ela estendeu a mão para mim. A mão grande, de dedos longos, com unhas pretas.

\noindent \textit{\textquotedblleft{}Cê quer?\textquotedblright{}}

\noindent \textit{\textquotedblleft{}Quer o quê?\textquotedblright{}}

\noindent \textit{\textquotedblleft{}Comer.\textquotedblright{}}

\noindent Eu olhei para a mão dela. Olhei para a vala. Olhei para os corpos.

\noindent \textit{\textquotedblleft{}Não\textquotedblright{}\textit{, eu disse. }\textquotedblleft{}Não quero comer eles.\textquotedblright{}}

\noindent Ela riu. A risada parecia o som de folhas secas se mexendo.

\noindent \textit{\textquotedblleft{}Não é eles que eu tô oferecendo, menino.\textquotedblright{}}

\noindent \textit{\textquotedblleft{}Então o quê?\textquotedblright{}}

\noindent Ela se abaixou. O rosto ficou na altura do meu. O cheiro de comida ficou mais forte. Meu estômago doeu de novo.

\noindent \textit{\textquotedblleft{}Eu como por você\textquotedblright{}\textit{, ela disse. }\textquotedblleft{}Eu como todos os dias. Eu como os que morrem. Eu como os que ninguém quis. E você nunca mais vai sentir fome.\textquotedblright{}}

\noindent Ela esperou.

\noindent \textit{\textquotedblleft{}Só que...\textquotedblright{}}

\noindent \textit{\textquotedblleft{}Só que o quê?\textquotedblright{}}

\noindent Ela sorriu. A boca se abriu mais do que devia. Os dentes estavam em fileiras demais. Fileiras que não cabiam numa boca humana.

\noindent \textit{\textquotedblleft{}A fome não vai embora. Ela fica. Só que ela não é mais tua. É minha. E eu vou morar dentro de ti para sentir ela.\textquotedblright{}}

\noindent Ela tocou meu peito com um dedo.

\noindent \textit{\textquotedblleft{}Tu aceita?\textquotedblright{}}

\noindent Eu devia ter dito não.

\noindent Eu devia ter cuspido na mão dela. Devia ter corrido. Devia ter gritado.

\noindent Mas eu tinha oito anos.
E tinha tanta fome.
Tanta fome que o não não saía.

\noindent \textit{\textquotedblleft{}Sim\textquotedblright{}}, eu disse.

\noindent Ela abriu a boca.

\noindent Não foi um beijo. Não foi uma mordida. Foi um sopro. Um sopro quente que saiu da boca dela e entrou na minha. Desceu pela garganta. Encheu o peito. Escorreu para o estômago.

\noindent Ela estava dentro de mim.

\noindent \textit{\textquotedblleft{}Pronto\textquotedblright{}\textit{, ela disse. }\textquotedblleft{}Agora a fome é minha.\textquotedblright{}}

\noindent Ela se levantou. Olhou para o céu. A lua brilhava.

\noindent \textit{\textquotedblleft{}Vai\textquotedblright{}\textit{, ela disse. } \textquotedblleft{}Vai viver. Depois a gente se encontra de novo.\textquotedblright{}}

\noindent Ela voltou para a vala. Desceu entre os corpos. Deitou. E fechou os olhos.

\noindent Eu fiquei parado na borda, olhando.

\noindent Não sentia mais fome.
Mas sentia um peso.
Um peso no fundo do estômago.
Um peso que respirava.\NE{O trecho do pacto com a vala é o único em todo o diário em que Joaquim escreve com a mão esquerda. Os peritos que examinaram o manuscrito notaram: a caligrafia muda abruptamente na linha \textquotedblleft{}Sim, eu disse\textquotedblright{}, a pressão da caneta sobre o papel é diferente, e o traçado das letras se desvia para a esquerda — como se outra mão estivesse guiando a dele. O editor não interveio. A escolha de manter a caligrafia original é uma escolha ética e documental.}

\vspace{0.5em}\noindent\rule{\textwidth}{0.5pt}\vspace{0.5em}

\noindent Naquela noite, na beira de uma vala de mortos, no meio de um campo de concentração, um menino de oito anos fez um trato com um monstro.

\noindent O monstro comeu a fome dele.

\noindent Mas deixou o buraco.

\noindent O buraco ficou.

\noindent E tudo o que entrou no buraco depois --- anos, sangue, gente, memórias --- nunca encheu.

\noindent Porque o buraco não é meu.

\noindent É dela.

\noindent Ela mora dentro de mim desde aquela noite.

\noindent Ela come. Eu sinto.

\noindent Ela dorme. Eu acordo.

\noindent Ela espera. Eu vivo.

\noindent Ela está esperando você também.

\vspace{0.5em}
\noindent\textit{\textbf{↪ Rotas:} \rota{MEM-05} \rotasep \rota{DOC-01} \rotasep \rota{LUC-01}}
\newpage\fragdef{MEM-05}% ============================================
% MEM-05
% Gerado de molambudos.md — build_miolo.py
% ============================================
\fragtitle{mem}{M}{MEM}{05}{O Olho Amarelo}
\markboth{MEM-05 --- O Olho Amarelo}{}

\noindent \textit{O que nunca mais fechou}

\noindent \textbf{O que nunca mais fechou}

\vspace{0.5em}\noindent\rule{\textwidth}{0.5pt}\vspace{0.5em}

\noindent Na manhã seguinte, minha mãe me olhou e disse:

\noindent \textit{\textquotedblleft{}Joaquim, teu olho.\textquotedblright{}}

\noindent Eu toquei o rosto. O olho direito estava arregalado. Não conseguia piscar. A pálpebra estava dura, como se tivesse sido costurada para cima.

\noindent Corri para o pote d'água e me inclinei. A água era suja, barrenta, mas dava pra ver o reflexo.

\noindent A íris do olho direito estava amarela.

\noindent Não o amarelo da icterícia, que deixa o olho inteiro cor de gema. Era um amarelo intenso, brilhante, com uma pupila preta que parecia mais funda do que deveria. Como se dentro do meu olho houvesse um poço.

\noindent Minha mãe fez o sinal da cruz.

\noindent \textit{\textquotedblleft{}O que foi que tu viu?\textquotedblright{}}

\noindent Eu não respondi.

\noindent \textit{\textquotedblleft{}Joaquim, pelo amor de Deus, o que foi que tu viu naquela vala?\textquotedblright{}}

\noindent \textit{\textquotedblleft{}Nada\textquotedblright{}\textit{, eu menti. }\textquotedblleft{}Não vi nada.\textquotedblright{}}

\noindent Ela sabia que era mentira. Mas não perguntou de novo. Naquele dia, os soldados vieram buscar as mulheres. Disseram que iam \textquotedblleft{}realocar\textquotedblright{} as moças para o litoral. Minha mãe estava tossindo, tossindo, e mesmo assim a levaram.

\noindent Ela olhou para mim quando o caminhão partiu.

\noindent O olho dela dizia: \textit{não morre.}

\noindent Eu não morri.

\noindent Mas eu queria ter morrido.

\vspace{0.5em}\noindent\rule{\textwidth}{0.5pt}\vspace{0.5em}

\noindent Naquela noite, voltei à vala.

\noindent A criatura estava lá. Dessa vez, ela não estava comendo. Estava deitada no meio dos corpos, olhando para o céu. Quando me viu, sentou.

\noindent \textit{\textquotedblleft{}Sabia que você voltava.\textquotedblright{}}

\noindent Eu fiquei em silêncio.

\noindent \textit{\textquotedblleft{}Vem aqui.\textquotedblright{}}

\noindent Eu fui. Dessa vez, não sentei na borda. Desci. Meus pés afundaram na pilha de corpos. Alguns estavam moles. Outros duros. Um tinha o braço esticado, apontando para o nada. Pisei na mão dele sem querer e ela quebrou com um som seco.

\noindent A criatura tinha um cheiro diferente de perto. Não era só carne. Era sangue grosso, terra molhada, e algo adocicado que lembrava o cheiro de um curral.

\noindent \textit{\textquotedblleft{}Qual é o teu nome?\textquotedblright{}}, ela perguntou.

\noindent \textit{\textquotedblleft{}Joaquim.\textquotedblright{}}

\noindent \textit{\textquotedblleft{}Joaquim.\textquotedblright{}\textit{ Ela repetiu como se estivesse provando a palavra. }\textquotedblleft{}Nome de santo. José, Joaquim, Maria. Nome de quem nasce pra sofrer.\textquotedblright{}}

\noindent \textit{\textquotedblleft{}E o teu?\textquotedblright{}}

\noindent Ela riu. A risada era um som molhado, como água passando por canos sujos.

\noindent \textit{\textquotedblleft{}Não tenho nome. Mas tenho fome. E fome precisa de nome pra chamar.\textquotedblright{}}

\noindent \textit{\textquotedblleft{}Chama de Olho.\textquotedblright{}}

\noindent Ela inclinou a cabeça.

\noindent \textit{\textquotedblleft{}Por quê?\textquotedblright{}}

\noindent \textit{\textquotedblleft{}Porque é a primeira parte de ti que eu vi.\textquotedblright{}}

\noindent Ela sorriu. Os dentes brilharam no escuro.

\noindent \textit{\textquotedblleft{}Então serei Olho. Para ti.\textquotedblright{}}

\noindent Ela estendeu a mão. A mão enorme, de dedos demais, e tocou meu rosto.

\noindent \textit{\textquotedblleft{}Teu olho agora é igual ao meu\textquotedblright{}\textit{, ela disse. }\textquotedblleft{}Tu vês o que eu vejo.\textquotedblright{}}

\noindent \textit{\textquotedblleft{}E o que tu vês?\textquotedblright{}}

\noindent \textit{\textquotedblleft{}Eu vejo a fome.\textquotedblright{}}

\noindent \textit{\textquotedblleft{}Dentro de quem?\textquotedblright{}}

\noindent \textit{\textquotedblleft{}Dentro de todos. Mas a maioria não sabe. Tu vais saber. Tu vais ver em cada rosto o que ele esconde. Vais ver o buraco que cada um tem dentro. E tu vais saber o que preenche esse buraco.\textquotedblright{}}

\noindent \textit{\textquotedblleft{}O quê?\textquotedblright{}}

\noindent Ela inclinou a cabeça para o lado, como um pássaro.

\noindent \textit{\textquotedblleft{}Eu.\textquotedblright{}}

\vspace{0.5em}\noindent\rule{\textwidth}{0.5pt}\vspace{0.5em}

\noindent Na terceira noite, a criatura me mostrou o que eu ia ser.

\noindent Ela tocou minha testa com o dedo --- um dedo longo, de unha preta --- e eu \textit{vi}.

\noindent Vi uma sala. Uma sala grande, cheia de camas. Homens deitados. Homens que gritavam. Homens que não se mexiam. Vi um menino --- mais velho do que eu, mas ainda menino --- sentado num canto. O menino mexia a boca mas não saía som. Ele escrevia num caderno. A letra era minha.

\noindent \textit{\textquotedblleft{}Esse é tu\textquotedblright{}\textit{, disse a criatura. }\textquotedblleft{}Daqui a uns anos.\textquotedblright{}}

\noindent \textit{\textquotedblleft{}Onde é?\textquotedblright{}}

\noindent \textit{\textquotedblleft{}Um lugar que guarda gente como eu. Gente que a sociedade enterrou viva. Tu vai ficar lá muito tempo.\textquotedblright{}}

\noindent \textit{\textquotedblleft{}E tu?\textquotedblright{}}

\noindent \textit{\textquotedblleft{}Eu vou contigo.\textquotedblright{}}

\noindent Ela tocou meu peito. O dedo atravessou a pele, os ossos --- não doeu --- e tocou alguma coisa dentro de mim.

\noindent \textit{\textquotedblleft{}Já estou aqui. Não vou sair mais.\textquotedblright{}}

\noindent \textit{\textquotedblleft{}E se eu não quiser?\textquotedblright{}}

\noindent Ela riu de novo.

\noindent \textit{\textquotedblleft{}Tu não tem escolha, Joaquim. Na noite em que tu sentou na beira da minha vala, tu me convidou. Não com palavras. Com fome. A tua fome me chamou. E eu vim.\textquotedblright{}}

\noindent \textit{\textquotedblleft{}Eu posso ir embora.\textquotedblright{}}

\noindent \textit{\textquotedblleft{}Pode. Mas eu vou dentro de ti.\textquotedblright{}}

\noindent \textit{\textquotedblleft{}E se eu morrer?\textquotedblright{}}

\noindent Ela ficou quieta. Pela primeira vez, o silêncio dela foi diferente. Não era ameaçador. Era... triste.

\noindent \textit{\textquotedblleft{}Se tu morrer, eu espero.\textquotedblright{}}

\noindent \textit{\textquotedblleft{}Espera o quê?\textquotedblright{}}

\noindent \textit{\textquotedblleft{}Até alguém lembrar de ti. E abrir o que tu escreveu. E ler.\textquotedblright{}}

\noindent \textit{\textquotedblleft{}E aí?\textquotedblright{}}

\noindent O olho amarelo dela brilhou no escuro.

\noindent \textit{\textquotedblleft{}Aí eu vou de novo.\textquotedblright{}}

\vspace{0.5em}\noindent\rule{\textwidth}{0.5pt}\vspace{0.5em}

\noindent Quando levantei a cabeça, uma coruja pousada numa árvore seca me observava.

\noindent Não tinha feito barulho. Estava ali, imóvel, me olhando com olhos pretos sem pálpebras. A cabeça dela virou lentamente --- cento e oitenta graus --- e ela fitou a criatura.

\noindent A criatura não olhou para cima. Mas sorriu.

\noindent \textit{\textquotedblleft{}Ela veio te conhecer.\textquotedblright{}}

\noindent \textit{\textquotedblleft{}Quem é?\textquotedblright{}}

\noindent \textit{\textquotedblleft{}É a minha língua no mundo de cima. Tudo que tu disseres, ela leva. Tudo que ela vir, eu sei.\textquotedblright{}}

\noindent A coruja abriu as asas e voou. O vou era silencioso, absorvido pela noite como se o escuro fosse um pano molhado.

\noindent \textit{\textquotedblleft{}Ela vai me seguir?\textquotedblright{}}

\noindent \textit{\textquotedblleft{}Não. Ela vai na frente. Preparando o caminho.\textquotedblright{}}

\vspace{0.5em}\noindent\rule{\textwidth}{0.5pt}\vspace{0.5em}

\noindent Eu tenho 72 anos. Estou escrevendo isto num quarto de hospital psiquiátrico. Os médicos dizem que estou louco. Dizem que a criatura é fruto da minha mente. Que o olho amarelo é uma catarata mal curada.

\noindent Eu sei que não é.

\noindent Ela está aqui agora. Enquanto escrevo. Sinto o peso dela na minha nuca. O hálito doce. Ela está lendo o que escrevo. E está rindo.

\noindent \textit{\textquotedblleft{}Diz pra ele\textquotedblright{}\textit{, ela sussurra. }\textquotedblleft{}Diz que o diário não acaba quando tu morrer.\textquotedblright{}}

\noindent Eu vou morrer em breve. Sei disso. Dr. Oliveira disse que meu coração não aguenta mais.

\noindent Mas o diário vai continuar.

\noindent Não sou eu que escrevo.

\noindent É ela.

\noindent \textit{--- Joaquim Antônio Correia, Paciente 1.260. Último registro.}

\vspace{0.5em}
\noindent\textit{\textbf{↪ Rotas:} \rota{MEM-04} \rotasep \rota{MEM-12} \rotasep \rota{DOC-01}}
\newpage\fragdef{MEM-06}% ============================================
% MEM-06
% Gerado de molambudos.md — build_miolo.py
% ============================================
\fragtitle{mem}{M}{MEM}{06}{O Curral do Governo}
\markboth{MEM-06 --- O Curral do Governo}{}

\noindent \textit{Os Molambudos --- Campos de Concentração do Ceará}

\noindent \textbf{1915 --- Campos de Concentração do Ceará}

\vspace{0.5em}\noindent\rule{\textwidth}{0.5pt}\vspace{0.5em}

\noindent O governo chamava de \textquotedblleft{}currais\textquotedblright{}.

\noindent A palavra certa. A gente era gado.

\noindent O Campo do Alagadiço, nos arredores de Fortaleza, era uma área cercada de arame farpado, com uns duzentos metros de largura por talvez trezentos de comprimento. Não tinha cobertura. Não tinha parede. Não tinha nada além de terra batida, sol, e o cheiro de gente acumulada.

\noindent No primeiro dia, eu contei as pessoas. Cheguei a trezentas e poucas antes de perder a conta. Depois chegou mais gente. Depois mais. No fim da primeira semana, a gente se apertava como bicho no transporte.

\noindent Todo dia chegavam soldados com mais retirantes. Eles vinham de toda parte: Quixeramobim, Boa Viagem, Senador Pompeu, Quixadá, Canindé. Os soldados chamavam a gente de \textquotedblleft{}molambudos\textquotedblright{}. A palavra saía da boca deles como se fosse o nome de uma praga.

\noindent \textit{\textquotedblleft{}Lá vem mais molambudos.\textquotedblright{}}

\noindent \textit{\textquotedblleft{}Molambudo, sai daí.\textquotedblright{}}

\noindent \textit{\textquotedblleft{}Dá comida pra esses molambudos antes que morram tudo.\textquotedblright{}}

\vspace{0.5em}\noindent\rule{\textwidth}{0.5pt}\vspace{0.5em}

\noindent A comida era uma vez por dia.

\noindent Um caldo ralo. Os soldados chamavam de \textquotedblleft{}sopa do governo\textquotedblright{}. Era água com farinha e algumas folhas boiando. Eu tomava a minha e ainda lambia a cuia. Minha mãe dava a dela pra mim. Eu dizia que não. Ela dizia que já tinha comido. Mentira. Eu via a tosse sacudir o corpo vazio dela.

\noindent O governo mandava boi de vez em quando. Não para alimentar a gente --- para alimentar os soldados. A gente ficava olhando eles assarem a carne. O cheiro passava por cima da cerca e entrava no nosso estômago. Alguns choravam. Outros batiam no arame farpado. Os soldados riam.

\noindent \textit{\textquotedblleft{}Deixa eles cheirar\textquotedblright{}\textit{, um deles disse. }\textquotedblleft{}Fome não se cura com cheiro.\textquotedblright{}}

\noindent Ele tinha razão. Fome não se cura com cheiro. Nem com choro. Nem com oração. A fome só se cura com comida.

\noindent E não tinha comida.

\vspace{0.5em}\noindent\rule{\textwidth}{0.5pt}\vspace{0.5em}

\noindent As mortes em massa começaram na segunda semana.

\noindent Primeiro os velhos. Depois as crianças. Depois os adultos que já estavam doentes antes de chegar.

\noindent Os soldados tinham uma rotina. De manhã, abriam a cerca, entravam, andavam entre a gente, e apontavam os que tinham morrido durante a noite. Carregavam os corpos pelos braços e pernas --- um de cada lado --- e levavam para a vala.

\noindent A vala ficava a duzentos metros. A gente via os urubus.

\noindent Minha mãe tossia.

\vspace{0.5em}\noindent\rule{\textwidth}{0.5pt}\vspace{0.5em}

\noindent Uma noite, eu sonhei com comida.

\noindent Era um sonho bom. Eu estava na casa da minha avó, que tinha morrido antes de eu nascer, mas no sonho ela estava viva. Ela fazia cuscuz. O cuscuz saía da panela e subia até o céu. Eu abria a boca e ele caía dentro de mim como chuva.

\noindent Acordei com fome.

\noindent Uma fome que doía. Não no estômago. No corpo inteiro. Os ossos doíam. A boca doía. Os dentes doíam. Eu olhei para as mãos e vi os ossos por baixo da pele.

\noindent \textit{\textquotedblleft{}Mãe\textquotedblright{}\textit{, eu disse. }\textquotedblleft{}Eu quero ir embora.\textquotedblright{}}

\noindent Ela não respondeu.

\noindent \textit{\textquotedblleft{}Mãe?\textquotedblright{}}

\noindent Ela estava deitada de lado. O corpo dela estava quente. A respiração era curta. Rápida. Como a de um passarinho.

\noindent \textit{\textquotedblleft{}Vai\textquotedblright{}\textit{, ela disse. }\textquotedblleft{}Vai, Joaquim. Vai.\textquotedblright{}}

\noindent \textit{\textquotedblleft{}Pra onde?\textquotedblright{}}

\noindent Ela tossiu. A tosse sacudiu o corpo inteiro. Quando a tosse passou, os lábios dela estavam vermelhos.

\noindent \textit{\textquotedblleft{}Pra longe daqui.\textquotedblright{}}

\noindent \textit{\textquotedblleft{}Eu não sei pra onde ir.\textquotedblright{}}

\noindent Ela fechou os olhos.

\noindent \textit{\textquotedblleft{}Tu sabe.\textquotedblright{}}

\noindent E eu sabia.

\noindent Eu não queria saber. Mas sabia.

\noindent A vala. A única coisa que ficava além do arame farpado. A única direção que não exigia caminhar. A vala. Lá tinha carne.

\noindent Eu não fui naquela noite.

\noindent Fui na seguinte.\NE{A existência do Campo do Alagadiço, próximo a Fortaleza, durante a seca de 1915 é documentada. Joaquim, sua família e os episódios sensoriais desta passagem são ficcionais; detalhes de escala, rotina e vala pertencem à reconstituição narrativa e não equivalem a testemunho arquivístico.}

\vspace{0.5em}\noindent\rule{\textwidth}{0.5pt}\vspace{0.5em}

\noindent O Campo do Alagadiço foi criado durante a seca de 1915, nos arredores de Fortaleza. Registros históricos indicam que milhares de retirantes foram confinados ali, mas as fontes consultadas não permitem fixar, para esta narrativa, um total preciso de internados ou mortos.

\noindent Eu só sabia que muitos não saíam.

\noindent Eu era um dos que sobraram.

\noindent As vezes eu penso: \textit{por que eu?}

\noindent Aí a criatura ri dentro de mim.

\noindent E eu lembro que eu não sobrevivi \textit{apesar} da vala.

\noindent Eu sobrevivi \textit{por causa} dela.

\vspace{0.5em}
\noindent\textit{\textbf{↪ Rotas:} \rota{MEM-07} \rotasep \rota{MEM-01}}

\newpage\fragdef{MEM-07}% ============================================
% MEM-07
% Gerado de molambudos.md — build_miolo.py
% ============================================
\fragtitle{mem}{M}{MEM}{07}{A Mãe Levada}
\markboth{MEM-07 --- A Mãe Levada}{}

\noindent \textit{Soldados, caminhão, silêncio}

\noindent \textbf{1915 --- O Curral do Governo}

\vspace{0.5em}\noindent\rule{\textwidth}{0.5pt}\vspace{0.5em}

\noindent Depois da noite na vala, minha mãe ficou diferente.

\noindent Ela me olhava de um jeito que eu não conhecia. Não era o olhar de mãe --- era o olhar de alguém que viu um bicho e não sabe se o bicho é perigoso.

\noindent \textit{\textquotedblleft{}Teu olho\textquotedblright{}\textit{, ela repetia. }\textquotedblleft{}Teu olho, Joaquim.\textquotedblright{}}

\noindent O olho direito estava amarelo. Não amarelo claro. Amarelo de ovo. Amarelo de gema podre. Eu via o reflexo nos cacos de louça que a gente usava para beber água.

\noindent Eu dizia que não era nada.

\noindent Ela sabia que era mentira.

\vspace{0.5em}\noindent\rule{\textwidth}{0.5pt}\vspace{0.5em}

\noindent Os soldados chegaram numa tarde.

\noindent Eram quatro. Dois com espingardas. Um com chicote. Um com uma lista de nomes.

\noindent \textit{\textquotedblleft{}Mulheres\textquotedblright{}\textit{, gritou o do chicote. }\textquotedblleft{}Todas as mulheres entre quinze e quarenta anos. Apresentem-se na entrada.\textquotedblright{}}

\noindent Minha mãe me apertou contra o peito.

\noindent \textit{\textquotedblleft{}Fica quieto\textquotedblright{}\textit{, ela sussurrou. }\textquotedblleft{}Fica quieto que eles não me levam.\textquotedblright{}}

\noindent Mas levaram.

\noindent O soldado com a lista leu o nome dela. Ela não respondeu. Ele leu de novo. Ela não respondeu. Ele veio até nós.

\noindent \textit{\textquotedblleft{}Você é Maria do Rosário Correia?\textquotedblright{}}

\noindent Ela não respondeu.

\noindent \textit{\textquotedblleft{}Você tá na lista. Vem.\textquotedblright{}}

\noindent \textit{\textquotedblleft{}Meu filho\textquotedblright{}\textit{, ela disse. }\textquotedblleft{}Meu filho precisa de mim.\textquotedblright{}}

\noindent \textit{\textquotedblleft{}Teu filho vai ficar bem. Os meninos ficam. As mulheres vão pra outra cidade. Lá tem comida.\textquotedblright{}}

\noindent \textit{\textquotedblleft{}Mentira\textquotedblright{}\textit{, ela disse. }\textquotedblleft{}Eu sei o que vocês fazem com as mulheres.\textquotedblright{}}

\noindent O soldado não respondeu. Ele fez um sinal com a cabeça. Os outros dois vieram. Um segurou o braço esquerdo dela. O outro segurou o direito.

\noindent \textit{\textquotedblleft{}Mãe!\textquotedblright{}}, eu gritei.

\noindent \textit{\textquotedblleft{}Fica, Joaquim! Fica!\textquotedblright{}}, ela gritou de volta.

\noindent Ela me olhou. Os olhos dela estavam cheios d'água. Mas não era choro. Era raiva. Uma raiva que eu nunca tinha visto.

\noindent \textit{\textquotedblleft{}Não morre\textquotedblright{}\textit{, ela disse. }\textquotedblleft{}Não morre, Joaquim. Escuta a tua mãe: não morre.\textquotedblright{}}

\noindent Os soldados arrastaram ela.

\noindent Eu corri atrás.

\noindent O do chicote me empurrou. Eu caí no chão. Levantei e corri de novo. Ele me empurrou de novo. Dessa vez eu não levantei.

\noindent Fiquei no chão, olhando.

\noindent O caminhão partiu.

\noindent Minha mãe estava na carroceria, entre outras mulheres. Ela me olhava. Eu olhava ela.

\noindent O caminhão ficou pequeno.

\noindent Depois sumiu.

\noindent Naquele dia, eu tive certeza de que nunca mais veria minha mãe.\NE{A \textquotedblleft{}realocação\textquotedblright{} e o destino de Maria são elementos ficcionais. A violência sofrida por mulheres retirantes exige documentação histórica específica; como Joaquim não existiu, sua narrativa não constitui testemunho nem confirma um arquivo oficial. O conflito entre a promessa de comida e o medo da separação move este fragmento apenas dentro da ficção.}

\vspace{0.5em}\noindent\rule{\textwidth}{0.5pt}\vspace{0.5em}

\noindent O velho Anselmo, que não tinha morrido ainda naquela época, me disse depois:

\noindent \textit{\textquotedblleft{}Elas vão pros bordéis do litoral. O governo chama de 'realocação'. As que não querem ir, vão pro mato com um tiro na nuca.\textquotedblright{}}

\noindent \textit{\textquotedblleft{}E minha mãe?\textquotedblright{}}

\noindent \textit{\textquotedblleft{}Tua mãe foi.\textquotedblright{}}

\noindent \textit{\textquotedblleft{}Foi pra onde?\textquotedblright{}}

\noindent Anselmo cuspiu no chão.

\noindent \textit{\textquotedblleft{}Foi viver. Ou morrer. No fundo é a mesma coisa.\textquotedblright{}}

\noindent Eu deitei no chão do curral. Olhei pro céu. O céu estava limpo. Sem nuvem. Sem chuva. Sem nada.

\noindent Só azul.

\noindent Um azul que doía.

\noindent Eu senti a criatura se mexer dentro de mim.

\noindent \textit{\textquotedblleft{}Calma\textquotedblright{}\textit{, ela disse. }\textquotedblleft{}Um dia tu vai entender.\textquotedblright{}}

\noindent \textit{\textquotedblleft{}Entender o quê?\textquotedblright{}}

\noindent \textit{\textquotedblleft{}Que o mundo come os filhos. E os filhos aprendem a comer de volta.\textquotedblright{}}

\noindent Eu fechei os olhos.

\noindent O olho amarelo não fechou.

\vspace{0.5em}
\noindent\textit{\textbf{↪ Rotas:} \rota{MEM-06} \rotasep \rota{MEM-08} \rotasep \rota{DOC-05} \rotasep \rota{DOC-17}}
\newpage\fragdef{DOC-17}% ============================================
% DOC-17
% Gerado de molambudos.md — build_miolo.py
% ============================================
\fragtitle{doc}{D}{DOC}{17}{Carta de um Soldado}
\markboth{DOC-17 --- Carta de um Soldado}{}

\noindent \textit{1915 --- o curral visto de dentro}

\noindent \textbf{Reconstituição ficcional --- documento narrativo}
\textbf{1915 --- Carta ficcional, sem equivalente arquivístico conhecido}

\vspace{0.5em}\noindent\rule{\textwidth}{0.5pt}\vspace{0.5em}

\noindent \textit{Carta ficcional atribuída ao soldado Antônio Alves Neto, lotado no Campo do Alagadiço, próximo a Fortaleza.}

\noindent \textit{Reconstituição ficcional endereçada a sua esposa, Maria, em Sobral; não foi localizada em arquivo militar.}

\vspace{0.5em}\noindent\rule{\textwidth}{0.5pt}\vspace{0.5em}

\noindent \textit{Maria,}

\noindent \textit{Estou escrevendo esta carta na tenda do curral, mas não sei se vou conseguir enviar. A situação aqui é pior do que o coronel conta para a imprensa.}

\noindent \textit{Todo dia chegam mais retirantes. Famílias inteiras. Gente que andou dias sem comida nem água. As crianças chegam com os olhos fundos e a barriga inchada. As mulheres, quando chegam, são separadas e levadas para \textquotedblleft{}realocação\textquotedblright{}. Eu sei o que isso significa. Todos sabem.}

\noindent \textit{Os homens e as crianças ficam. Ficam atrás do arame farpado, esperando a comida que nunca vem.}

\noindent \textit{O governo manda boi, mas o boi vai para a nossa cozinha, não para a deles. Eles recebem um caldo ralo, uma vez por dia. É o suficiente para não morrerem de uma vez. É o suficiente para morrerem aos poucos.}

\noindent \textit{O pior são as crianças. Tem um menino aqui, não deve ter mais de 10 anos, que chegou sozinho. A mãe foi levada no caminhão. O pai morreu na estrada. Ele não chora. Ele só olha. O olho dele é amarelo. O direito. Ele disse que \textquotedblleft{}um bicho\textquotedblright{} mudou. Ele não chora, Maria. Ele só olha. E olhar dele... me atravessa.}

\noindent \textit{Na semana passada, eu vi um negócio na vala.}

\noindent \textit{A vala dos corpos, atrás do curral. Eu estava de guarda noturna e ouvi um barulho. Mastigação. Fui ver.}

\noindent \textit{Tem uma coisa lá dentro. Mora no meio dos corpos. É grande. Muito grande. Tem olhos amarelos que brilham no escuro. Ela não me atacou. Ela só olhou. E continuou comendo.}

\noindent \textit{Contei para o sargento. Ele disse que é \textquotedblleft{}imaginação\textquotedblright{} e \textquotedblleft{}fome\textquotedblright{}. Disse que a fome faz a gente ver coisa.}

\noindent \textit{Talvez seja. Talvez não.}

\noindent \textit{Mas, esta noite, quando eu fechar os olhos, vou ver os olhos amarelos.}

\noindent \textit{Eu vi.}

\noindent \textit{E, Maria, eles me viram também.}

\noindent \textit{--- Antonio}

\vspace{0.5em}
\noindent\textit{\textbf{↪ Rotas:} \rota{DOC-18} \rotasep \rota{MEM-06}}

\newpage\fragdef{DOC-18}% ============================================
% DOC-18
% Gerado de molambudos.md — build_miolo.py
% ============================================
\fragtitle{doc}{D}{DOC}{18}{Relatório Oficial do Curral}
\markboth{DOC-18 --- Relatório Oficial do Curral}{}

\noindent \textit{1916 --- Reconstituição sobre o Campo do Alagadiço}

\noindent \textbf{Reconstituição ficcional --- sem equivalente localizado no Arquivo Nacional}
\textbf{1916 --- Documento narrativo, não transcrição}

\vspace{0.5em}\noindent\rule{\textwidth}{0.5pt}\vspace{0.5em}

\noindent \textit{RECONSTITUIÇÃO FICCIONAL DE EXPEDIENTE ADMINISTRATIVO}
\textit{SOCORROS PÚBLICOS}
\textit{RELATÓRIO NARRATIVO DOS CAMPOS DO CEARÁ}
\textit{EXERCÍCIO DE 1915-1916}

\vspace{0.5em}\noindent\rule{\textwidth}{0.5pt}\vspace{0.5em}

\noindent \textbf{Campo:} Alagadiço, próximo a Fortaleza
\textbf{Capacidade oficial:} não indicada nesta reconstituição
\textbf{Número de internados durante o período:} impreciso; as fontes consultadas registram milhares de confinados
\textbf{Óbitos:} sem total verificável para reprodução nesta narrativa

\noindent \textbf{Voz do relator ficcional:}

\noindent \textit{\textquotedblleft{}Os campos de concentração, oficialmente chamados 'currais do governo', foram estabelecidos como medida emergencial para conter o fluxo de retirantes em direção à capital. Na prática, funcionaram como depósitos de seres humanos em estado de inanição.}

\noindent \textit{A alimentação e as condições sanitárias eram insuficientes para proteger a vida dos confinados. Esta reconstituição não atribui percentual de mortalidade nem autoria documental a essa formulação.}

\noindent \textit{Recomenda-se o fechamento dos campos e a adoção de medidas efetivas de socorro.\textquotedblright{}}

\vspace{0.5em}\noindent\rule{\textwidth}{0.5pt}\vspace{0.5em}

\noindent \textbf{Nota editorial da reconstituição:}

\noindent A existência do Campo do Alagadiço, próximo a Fortaleza, durante a seca de 1915 é documentada. Os números, a voz do relator e a forma administrativa acima são recursos ficcionais. Este texto não afirma que o campo permaneceu aberto até 1917; a transferência de Joaquim nesse ano pertence à cronologia ficcional do personagem.

\vspace{0.5em}
\noindent\textit{\textbf{↪ Rotas:} \rota{DOC-19} \rotasep \rota{MEM-06}}

%% ============================================
%% PARTE 2 — COLÔNIA (1928--1979)
%% ============================================
\partopener{Colônia (1917--1979)}{Colônia}{4.png}

\newpage\fragdef{MEM-08}% ============================================
% MEM-08
% Gerado de molambudos.md — build_miolo.py
% ============================================
\fragtitle{mem}{M}{MEM}{08}{A Chegada ao Colônia}
\markboth{MEM-08 --- A Chegada ao Colônia}{}

\noindent \textit{1917 --- Barbacena, Minas Gerais}

\noindent \textbf{1917 --- Barbacena, Minas Gerais}

\vspace{0.5em}\noindent\rule{\textwidth}{0.5pt}\vspace{0.5em}

\noindent Eu não lembro como cheguei no Colônia.

\noindent Lembro do caminhão. Lembro de sacos. Lembro que estava escuro e que a gente era muitos dentro de um espaço pequeno. Gente que cheirava a suor, a fezes, a medo.

\noindent Depois o caminhão parou.

\noindent Eu vi o portão.

\noindent Era um portão de ferro, preto, com pontas em cima. Atrás do portão tinha uma estrada de terra. Depois da estrada tinha um prédio. O prédio era grande, cinza, com janelas pequenas e uma cruz no topo.

\noindent Cruz de hospital. Mas parecia cruz de cemitério.

\noindent \textit{\textquotedblleft{}Desce\textquotedblright{}}, alguém disse.

\noindent Eu desci.

\noindent \textit{\textquotedblleft{}Nome?\textquotedblright{}}

\noindent \textit{\textquotedblleft{}Joaquim Antônio Correia.\textquotedblright{}}

\noindent \textit{\textquotedblleft{}Idade?\textquotedblright{}}

\noindent \textit{\textquotedblleft{}Dez.\textquotedblright{}}

\noindent O homem olhou pra mim. Olhou pro meu olho amarelo. Anotou alguma coisa num papel.

\noindent \textit{\textquotedblleft{}Diagnóstico?\textquotedblright{}} ele perguntou pra outro homem.

\noindent \textit{\textquotedblleft{}Veio do Ceará. Curral do governo. Órfão. Não fala.\textquotedblright{}}

\noindent \textit{\textquotedblleft{}Eu falo\textquotedblright{}}, eu disse.

\noindent Os dois homens me olharam.

\noindent \textit{\textquotedblleft{}Então fala\textquotedblright{}\textit{, disse o primeiro. }\textquotedblleft{}O que tu tem?\textquotedblright{}}

\noindent Eu pensei em dizer a verdade. Pensei em falar da vala. Da criatura. Do olho. Da fome que nunca passa.

\noindent Mas eu tinha aprendido uma coisa na estrada e no curral: a verdade não ajuda.

\noindent \textit{\textquotedblleft{}Nada\textquotedblright{}\textit{, eu disse. }\textquotedblleft{}Não tenho nada.\textquotedblright{}}

\noindent O homem riu.

\noindent \textit{\textquotedblleft{}Todo mundo que chega aqui tem alguma coisa. Se não tem, a gente inventa.\textquotedblright{}}

\noindent Ele escreveu no papel:

\noindent \textit{\textquotedblleft{}Paciente 1.260 --- Joaquim Antônio Correia. Esquizofrenia hebefrênica. Perigoso. Isolamento.\textquotedblright{}}

\noindent Eu não sabia o que era esquizofrenia. Aprendi depois. É uma palavra que os médicos usam quando não sabem o que uma pessoa tem.

\noindent Eles me levaram para uma sala.

\noindent A sala tinha uma cama de ferro. Uma janela gradeada. Um balde no canto.

\noindent \textit{\textquotedblleft{}Fica aqui\textquotedblright{}\textit{, disse o enfermeiro. }\textquotedblleft{}Amanhã vem o médico.\textquotedblright{}}

\noindent \textit{\textquotedblleft{}E a comida?\textquotedblright{}}

\noindent \textit{\textquotedblleft{}Amanhã.\textquotedblright{}}

\noindent Ele fechou a porta. A chave girou.

\noindent Eu fiquei ali, sentado na cama, olhando para a parede. A parede estava cheia de riscos. Riscos de unha. Riscos de dente. Riscos de gente que tinha ficado ali antes de mim.

\noindent \textit{\textquotedblleft{}Bem-vindo\textquotedblright{}}, disse a criatura dentro de mim.

\noindent \textit{\textquotedblleft{}Cala a boca.\textquotedblright{}}

\noindent \textit{\textquotedblleft{}Eu vou gostar daqui\textquotedblright{}\textit{, ela disse. }\textquotedblleft{}Aqui tem muita fome.\textquotedblright{}}

\noindent Ela estava certa.

\noindent O Colônia tinha mais de cinco mil pacientes. Os dormitórios eram galpões com camas encostadas umas nas outras. A comida era um caldo de osso com farinha. A água era servida uma vez por dia.

\noindent Os gritos não paravam nunca.

\noindent De noite, os gritos ecoavam pelos corredores. Gritos de dor. Gritos de medo. Gritos de gente que não sabia mais quem era.

\noindent Nos primeiros dias, eu tampava os ouvidos.

\noindent Depois de uma semana, eu já gritava junto.

\noindent A criatura crescia.

\noindent Cada grito alimentava ela. Cada lágrima. Cada pedaço de comida que eu engolia e que não saciava.

\noindent Ela se alimentava do Colônia.

\noindent E o Colônia se alimentava de mim.

\vspace{0.5em}\noindent\rule{\textwidth}{0.5pt}\vspace{0.5em}

\noindent Anos depois, Dr. Oliveira me perguntou:

\noindent \textit{\textquotedblleft{}Joaquim, você sabe quantos anos passou no Colônia?\textquotedblright{}}

\noindent Eu sabia.

\noindent \textit{\textquotedblleft{}Sessenta e dois.\textquotedblright{}}

\noindent \textit{\textquotedblleft{}E como você se sente sobre isso?\textquotedblright{}}

\noindent Eu olhei pra ele. O olho amarelo brilhou.

\noindent \textit{\textquotedblleft{}Eu não vivi aqui, doutor. Eu apodreci aqui.\textquotedblright{}}

\noindent Ele anotou.

\noindent \textit{\textquotedblleft{}Apodrecer\textquotedblright{}\textit{, ele repetiu. }\textquotedblleft{}Interessante escolha de palavras.\textquotedblright{}}

\noindent \textit{\textquotedblleft{}Não foi escolha\textquotedblright{}\textit{, eu disse. }\textquotedblleft{}Foi o que aconteceu. A gente não escolhe apodrecer. A gente só... apodrece.\textquotedblright{}}

\noindent Ele parou de anotar.

\noindent \textit{\textquotedblleft{}O que aconteceu com você, Joaquim?\textquotedblright{}}

\noindent Eu sorri.

\noindent \textit{\textquotedblleft{}O que não aconteceu?\textquotedblright{}}

\vspace{0.5em}
\noindent\textit{\textbf{↪ Rotas:} \rota{MEM-09} \rotasep \rota{DOC-02} \rotasep \rota{MEM-12}}
\newpage\fragdef{DOC-02}% ============================================
% DOC-02
% Gerado de molambudos.md — build_miolo.py
% ============================================
\fragtitle{doc}{D}{DOC}{02}{Prontuário do Paciente 1.260}
\markboth{DOC-02 --- Prontuário do Paciente 1.260}{}

\noindent \textit{Hospital Colônia de Barbacena --- Ficha de Admissão}

\noindent \textbf{Hospital Colônia de Barbacena --- Ficha de Admissão}
\textit{Reconstituição ficcional de ficha de admissão; não corresponde a documento arquivístico autêntico.}
\textit{Composição narrativa baseada em contexto histórico, sem chancela do Arquivo Público Mineiro.}

\vspace{0.5em}\noindent\rule{\textwidth}{0.5pt}\vspace{0.5em}

\noindent \textbf{Nº de Registro:} 1.260
\textbf{Nome:} Joaquim Antônio Correia
\textbf{Data de Nascimento:} ~1907 (estimado)
\textbf{Naturalidade:} Quixeramobim, Ceará
\textbf{Filiação:} Pai: Antônio Correia (falecido --- inanição, 1915); Mãe: Maria do Rosário Correia (desaparecida --- realocação forçada, 1915)
\textbf{Procedência:} Retirante saído de Senador Pompeu em 1915 e confinado no Campo do Alagadiço, próximo a Fortaleza. Transferido em 1917 via Santa Casa de Misericórdia de Belo Horizonte.

\vspace{0.5em}\noindent\rule{\textwidth}{0.5pt}\vspace{0.5em}

\noindent \textbf{Exame de Admissão --- Dr. Nelson Fontes, 12 de novembro de 1917:}

\noindent \textit{\textquotedblleft{}Paciente apresenta-se em estado de caquexia severa. Peso estimado: 20 kg (esperado para idade: 28-30 kg). Desnutrição proteico-calórica evidente. Edema nos membros inferiores. Abdome distendido. Dentição comprometida (cáries múltiplas, perda de 3 dentes).}

\noindent \textit{Sinal mais notável: heterocromia da íris direita, que se apresenta em tom amarelo-ouro, ausente na esquerda. Paciente alega que 'um bicho' mudou a cor. Provável sequela infecciosa ou traumática não identificada.}

\noindent \textit{Estado mental: pacato, cooperativo, mas com olhar fixo e vago. Responde perguntas com monossílabos. Quando questionado sobre a família, declarou: 'Morreram. Todos.' Sem sinais de alteração de sensório. Orientado no tempo e espaço.}

\noindent \textit{Avaliação preliminar: Esquizofrenia hebefrênica (demência precoce). Recomenda-se isolamento e observação.\textquotedblright{}}

\vspace{0.5em}\noindent\rule{\textwidth}{0.5pt}\vspace{0.5em}

\noindent \textbf{Histórico de Incidentes:}

\begin{moltablethree}
\hline
Data & Incidente & Consequência \\
\hline
Fev/1920 & Agressão a enfermeiro sem provocação & 30 dias de isolamento \\
Jun/1922 & \textbf{Homicídio --- Enfermeiro Wagner} & Isolamento por 90 dias \\
Mar/1937 & \textbf{Homicídio --- Paciente da ala C} & Isolamento permanente \\
Abr/1937 & Automutilação (arrancou unhas com os dentes) & Contenção física \\
Set/1945 & Episódio de vocalização incontrolável (alega que 'a Rasga Mortalha' está falando através dele) & Sedação com brometo \\
Jan/1948 & \textbf{Homicídio --- Visitante (entrou na ala proibida)} & Isolamento perpétuo \\
1949-1956 & Mutismo seletivo (7 anos sem falar com a equipe) & --- \\
1956 & Retomada da fala: primeira frase 'O diário' & --- \\
1979 & Sessões com Dr. Heitor Oliveira & Vide relatórios anexos \\
\hline
\end{moltablethree}

\vspace{0.5em}\noindent\rule{\textwidth}{0.5pt}\vspace{0.5em}

\noindent \textbf{Medicações Administradas (registro parcial):}

\noindent - Brometo de potássio (1920-1945) --- sedativo
- Eletrochoque (1937, 1945, 1952) --- sem anestesia
- Lobotomia transorbital --- \textit{indicada em 1948 mas não realizada (paciente reagiu com violência extrema à mera tentativa de aproximação)}
- Clorpromazina (1955 em diante) --- antipsicótico

\vspace{0.5em}\noindent\rule{\textwidth}{0.5pt}\vspace{0.5em}

\noindent \textbf{Observações Finais:}

\noindent Paciente 1.260 é um dos internos mais antigos do Colônia. Seu comportamento é imprevisível: períodos de calma absoluta intercalados com explosões de violência extrema. Mantém um diário que jamais permitiu que qualquer funcionário lesse em vida.

\noindent \textit{Nota adicional (1979): Dr. Heitor Oliveira obteve permissão para ler o diário. Primeira e única pessoa a ter acesso ao documento. Dr. Oliveira solicitou licença médica 48 horas depois. Não retornou ao trabalho.}

\noindent \textit{O diário encontra-se atualmente sob custódia do IML de Barbacena, anexado ao laudo de necropsia.}

\vspace{0.5em}
\noindent\textit{\textbf{↪ Rotas:} \rota{MEM-08} \rotasep \rota{DOC-03} \rotasep \rota{LUC-01}}

\newpage\fragdef{MEM-09}% ============================================
% MEM-09
% Gerado de molambudos.md — build_miolo.py
% ============================================
\fragtitle{mem}{M}{MEM}{09}{Os Primeiros Anos}
\markboth{MEM-09 --- Os Primeiros Anos}{}

\noindent \textit{O menino que não falava}

\noindent \textbf{1920-1930 --- O menino que não falava}

\vspace{0.5em}\noindent\rule{\textwidth}{0.5pt}\vspace{0.5em}

\noindent Nos primeiros anos no Colônia, eu quase não falava.

\noindent Não porque não pudesse. Porque não valia a pena.

\noindent Falar significava responder. Responder significava explicar. Explicar significava repetir a história. E a história --- a vala, a criatura, o olho amarelo --- ninguém ia entender.

\noindent Então eu ficava quieto.

\noindent Os enfermeiros gostavam de paciente quieto. Paciente quieto não dá trabalho. Paciente quieto não toma tempo. Paciente quieto é como móvel --- está ali, ocupa espaço, mas não exige nada.

\noindent Eu era um móvel.

\noindent Sentado no canto do pátio. Olhando para o chão. O dia inteiro. Todos os dias.

\noindent A criatura dormia dentro de mim. De vez em quando ela acordava, se espreguiçava, e eu sentia uma pontada no estômago. Ela estava crescendo.

\noindent \textit{\textquotedblleft{}Tu precisa comer mais\textquotedblright{}}, ela disse.

\noindent \textit{\textquotedblleft{}Não tem comida.\textquotedblright{}}

\noindent \textit{\textquotedblleft{}Tem. Tu só não está vendo.\textquotedblright{}}

\noindent \textit{\textquotedblleft{}Onde?\textquotedblright{}}

\noindent \textit{\textquotedblleft{}Nos outros.\textquotedblright{}}

\noindent Eu não respondi.

\noindent Mas eu entendi.

\vspace{0.5em}\noindent\rule{\textwidth}{0.5pt}\vspace{0.5em}

\noindent Em 1922, eu matei o Wagner.

\noindent Depois disso, os enfermeiros começaram a me tratar diferente. Não chegavam perto. Deixavam a comida na porta da cela. Evitavam olhar nos meus olhos.

\noindent O isolamento virou minha casa.

\noindent Uma cela de três metros por dois. Uma cama de ferro. Um balde. Uma janela gradeada que dava para o corredor. Eu ficava ali, sentado, escrevendo no diário.

\noindent As horas passavam.

\noindent Os dias passavam.

\noindent Os anos passavam.

\noindent Eu perdia a conta do tempo. O diário virava meu relógio. Cada página virada era um dia. Cada página escrita era uma noite.

\noindent A criatura falava cada vez mais.

\noindent \textit{\textquotedblleft{}Escreve.\textquotedblright{}}

\noindent \textit{\textquotedblleft{}Sobre o quê?\textquotedblright{}}

\noindent \textit{\textquotedblleft{}Sobre tudo. Sobre a vala. Sobre o curral. Sobre a fome. Escreve para que alguém leia.\textquotedblright{}}

\noindent \textit{\textquotedblleft{}Quem vai ler?\textquotedblright{}}

\noindent \textit{\textquotedblleft{}Alguém. Um dia. Alguém vai abrir este diário e vai ver.\textquotedblright{}}

\noindent \textit{\textquotedblleft{}Ver o quê?\textquotedblright{}}

\noindent \textit{\textquotedblleft{}A verdade.\textquotedblright{}}

\noindent \textit{\textquotedblleft{}Qual verdade?\textquotedblright{}}

\noindent \textit{\textquotedblleft{}A verdade que ninguém quer ver. Que o mundo é uma vala. Que a gente nasce, passa fome, morre, e é esquecido. Mas o diário não esquece. O diário guarda.\textquotedblright{}}

\noindent E eu escrevia.

\noindent Página após página.

\noindent Ano após ano.

\noindent Escrevia até os dedos sangrarem.

\noindent Escrevia até a tinta acabar --- e aí escrevia com sangue.

\noindent O diário crescia comigo.

\noindent Ele sabia coisas que eu nem sabia que sabia.

\noindent Às vezes eu abria numa página antiga e não reconhecia a letra. Achava que era de outra pessoa. Mas era a minha. Mão minha. Letra minha.

\noindent Mas as palavras...

\noindent As palavras não eram minhas.

\noindent Eram dela.

\vspace{0.5em}
\noindent\textit{\textbf{↪ Rotas:} \rota{MEM-08} \rotasep \rota{MEM-10}}
\newpage\fragdef{MEM-10}% ============================================
% MEM-10
% Gerado de molambudos.md — build_miolo.py
% ============================================
\fragtitle{mem}{M}{MEM}{10}{O Primeiro Homicídio}
\markboth{MEM-10 --- O Primeiro Homicídio}{}

\noindent \textit{O que a fome ensina}

\noindent \textbf{1922 --- Colônia de Barbacena}

\vspace{0.5em}\noindent\rule{\textwidth}{0.5pt}\vspace{0.5em}

\noindent Eu tinha quinze anos.

\noindent O enfermeiro se chamava Wagner. Era grande, gordo, com braços grossos e dedos curtos. Ele usava um avental branco que nunca estava limpo. Ele gostava de bater.

\noindent Não batia por raiva. Batia por diversão. Ele passava pelos corredores e escolhia um. \textit{\textquotedblleft{}Hoje é tu.\textquotedblright{}} E batia. Com a mão aberta. Com o cinto. Com a borracha de mangueira que ele guardava no bolso.

\noindent Os pacientes aprendiam a se encolher. A ficar pequenos. A não olhar nos olhos.

\noindent Eu não me encolhia.

\noindent O Wagner não gostava disso.

\noindent Um dia ele veio até mim. Eu estava no pátio, sentado no chão, escrevendo no diário. Ele chutou o diário da minha mão. As páginas voaram.

\noindent \textit{\textquotedblleft{}O que tu escreve aí, menino?\textquotedblright{}}

\noindent \textit{\textquotedblleft{}Nada.\textquotedblright{}}

\noindent \textit{\textquotedblleft{}Nada, o quê?\textquotedblright{}}

\noindent \textit{\textquotedblleft{}Nada, seu.\textquotedblright{}}

\noindent Ele me pegou pelo pescoço. Me levantou. A mão dele apertou. Os dedos entraram na carne.

\noindent \textit{\textquotedblleft{}Tu é doido?\textquotedblright{}\textit{ ele perguntou. }\textquotedblleft{}Tu sabe que eu posso te matar aqui mesmo?\textquotedblright{}}

\noindent Eu sabia.

\noindent Eu não respondi.

\noindent Ele apertou mais.

\noindent Foi quando a criatura falou.

\noindent Não dentro da minha cabeça. Fora. A voz saiu da minha boca. Mas não era a minha voz.

\noindent Era a voz da vala.

\noindent \textit{\textquotedblleft{}Tira a mão.\textquotedblright{}}

\noindent Wagner parou.

\noindent \textit{\textquotedblleft{}O quê?\textquotedblright{}}

\noindent \textit{\textquotedblleft{}Tira a mão.\textquotedblright{}}

\noindent A voz não era alta. Era baixa. Grossa. Como se viesse de um poço.

\noindent Wagner olhou nos meus olhos. O olho amarelo. Ele viu.

\noindent Ele viu a criatura dentro de mim.

\noindent Ele abriu a mão.

\noindent Eu caí no chão.

\noindent \textit{\textquotedblleft{}Seu...\textquotedblright{}} ele começou.

\noindent Não terminou.

\noindent Eu não lembro de ter pulado. Eu só lembro de estar em cima dele. A mão no rosto dele. Os dedos nos olhos dele. Ele gritou.

\noindent Eu não parei.

\noindent O Wagner tinha uma borracha no bolso. Eu peguei. Bati. Bati. Bati.

\noindent Não lembro quantas vezes.

\noindent Quando os outros enfermeiros me tiraram de cima dele, o rosto do Wagner não era mais um rosto. Era uma pasta vermelha.

\noindent A borracha estava no chão, suja de sangue e de olho.

\noindent O olho dele estava aberto, no chão, me olhando.

\noindent Eu peguei o diário. Guardei debaixo da camisa.

\noindent Foi quando eu vi o menino.

\noindent Um paciente novo, magro, uns doze anos. Ele estava encolhido no canto do pátio, com as mãos na cabeça, tremendo. Um enfermeiro se aproximou dele. Não o que eu tinha matado --- outro. Maior. Com uma chave de ferro na mão.

\noindent \textit{\textquotedblleft{}Esse vai pro isolamento também\textquotedblright{}\textit{, ele disse. }\textquotedblleft{}Testemunha.\textquotedblright{}}

\noindent O menino começou a chorar.

\noindent E eu ouvi a criatura dentro de mim. Ela disse:

\noindent \textit{\textquotedblleft{}Deixa. Não é teu problema.\textquotedblright{}}

\noindent Mas eu andei.

\noindent Me coloquei na frente do menino.

\noindent \textit{\textquotedblleft{}Ele não viu nada\textquotedblright{}\textit{, eu disse. }\textquotedblleft{}Tava de costas o tempo todo.\textquotedblright{}}

\noindent O enfermeiro me olhou. A chave de ferro na mão dele pesava. Eu vi ele pensar.

\noindent \textit{\textquotedblleft{}E quem garante?\textquotedblright{}}

\noindent \textit{\textquotedblleft{}Eu garanto\textquotedblright{}\textit{, eu disse. }\textquotedblleft{}Eu sou o louco que matou um homem. Minha palavra vale o que?\textquotedblright{}}

\noindent Ele cuspiu no chão.

\noindent \textit{\textquotedblleft{}Vale nada.\textquotedblright{}}

\noindent Mas ele baixou a chave. E foi embora.

\noindent O menino ficou. Vivo. Inteiro. Por minha causa.

\noindent A criatura dentro de mim ficou quieta. Não entendia o que eu tinha feito. Nem eu.

\noindent Mas, pela primeira vez, a escolha tinha sido minha.

\noindent \textit{\textquotedblleft{}Ele começou\textquotedblright{}}, eu disse.

\noindent Me levaram para o isolamento.

\noindent Fiquei três meses numa cela sem janela. Sem luz. Sem comida direito. Sem ninguém.

\noindent A criatura ficou quieta dentro de mim. Satisfeita.

\noindent \textit{\textquotedblleft{}Gostou?\textquotedblright{}}, eu perguntei.

\noindent \textit{\textquotedblleft{}Você também gostou.\textquotedblright{}}

\noindent Ela estava certa.

\noindent Eu não senti remorso. Não senti medo. Senti uma coisa que eu não conhecia.

\noindent Poder.

\noindent Pela primeira vez na vida, eu não estava com fome.

\noindent Eu tinha matado a fome.

\noindent Com sangue.

\vspace{0.5em}
\noindent\textit{\textbf{↪ Rotas:} \rota{MEM-09} \rotasep \rota{MEM-11} \rotasep \rota{DOC-05}}
\newpage\fragdef{MEM-11}% ============================================
% MEM-11
% Gerado de molambudos.md — build_miolo.py
% ============================================
\fragtitle{mem}{M}{MEM}{11}{O Isolamento}
\markboth{MEM-11 --- O Isolamento}{}

\noindent \textit{A ala dos silenciados}

\noindent \textbf{1923-1945 --- A ala dos silenciados}

\vspace{0.5em}\noindent\rule{\textwidth}{0.5pt}\vspace{0.5em}

\noindent Depois do segundo homicídio, me mudaram para a ala dos silenciados.

\noindent Era o fim do corredor. Uma ala que não aparecia nos registros oficiais. Os pacientes que ficavam ali não existiam para o hospital. Eram corpos que o hospital mantinha vivos --- mal vivos --- para não ter que explicar mortes demais.

\noindent A ala tinha doze celas. Seis de cada lado. As celas eram pequenas --- dois metros quadrados. Cama de concreto. Sem colchão. Sem cobertor. Um buraco no chão para as necessidades.

\noindent Os pacientes chegavam e nunca saíam.

\noindent Eu vi homens entrarem na ala com quarenta anos e morrerem com quarenta e cinco. Encolhidos. Grisalhos. Com os dentes podres e a pele cheia de feridas.

\noindent A gente não falava. Não porque não queria --- porque esquecia como. O silêncio virava uma doença. Primeiro a voz some. Depois as palavras. Depois o pensamento.

\noindent A criatura era a única que falava.

\noindent Ela preenchia o silêncio.

\noindent \textit{\textquotedblleft{}Tu não vai virar como eles\textquotedblright{}}, ela dizia.

\noindent \textit{\textquotedblleft{}Como?\textquotedblright{}}

\noindent \textit{\textquotedblleft{}Vazio.\textquotedblright{}}

\noindent \textit{\textquotedblleft{}Já sou vazio.\textquotedblright{}}

\noindent \textit{\textquotedblleft{}Não. Tu tem o diário. Tu escreve. Enquanto tu escrever, tu existe. Eles ---\textquotedblright{}\textit{ ela apontava para os outros pacientes }\textquotedblleft{}--- eles pararam de escrever. Pararam de falar. Pararam de existir.\textquotedblright{}}

\noindent \textit{\textquotedblleft{}Eles estão vivos.\textquotedblright{}}

\noindent \textit{\textquotedblleft{}Viver não é existir. Viver é o corpo. Existir é a história. Se ninguém lembra de ti, tu não existe.\textquotedblright{}}

\noindent \textit{\textquotedblleft{}Alguém vai lembrar?\textquotedblright{}}

\noindent \textit{\textquotedblleft{}Eu.\textquotedblright{}}

\noindent Ela ficou quieta.

\noindent \textit{\textquotedblleft{}E quem ler o diário.\textquotedblright{}}

\vspace{0.5em}\noindent\rule{\textwidth}{0.5pt}\vspace{0.5em}

\noindent Nas noites de inverno, o frio entrava pelas frestas da parede e não saía. Eu dormia encolhido, com o diário no peito, aquecendo ele com o corpo.

\noindent A criatura ficava quieta nessas noites.

\noindent Mas eu sentia ela olhando.

\noindent Esperando.

\noindent Ela sabia que eu ia morrer ali. Ela sabia que o diário ia ser enterrado comigo. E ela sabia que, um dia, alguém ia abrir.

\noindent O diário era o ovo dela.

\noindent E eu era o ninho.

\vspace{0.5em}
\noindent\textit{\textbf{↪ Rotas:} \rota{MEM-10} \rotasep \rota{MEM-12}}
\newpage\fragdef{MEM-17}% ============================================
% MEM-17
% Gerado de molambudos.md — build_miolo.py
% ============================================
\fragtitle{mem}{M}{MEM}{17}{O Eletrochoque}
\markboth{MEM-17 --- O Eletrochoque}{}

\noindent \textit{A primeira vez}

\noindent \textbf{1937 --- Colônia de Barbacena}

\vspace{0.5em}\noindent\rule{\textwidth}{0.5pt}\vspace{0.5em}

\noindent A primeira vez foi em 1937.

\noindent Dois enfermeiros entraram na cela. Um deles carregava uma máquina. Parecia um rádio antigo, com fios e botões de metal. O outro carregava uma tira de couro.

\noindent \textit{\textquotedblleft{}O médico mandou\textquotedblright{}}, disse um.

\noindent \textit{\textquotedblleft{}Não.\textquotedblright{}}

\noindent \textit{\textquotedblleft{}Não é pergunta.\textquotedblright{}}

\noindent Eles me seguraram. Dois homens. Eu era magro, fraco, não tinha força. Mas a criatura dentro de mim tinha.

\noindent Ela acordou.

\noindent Eu senti o calor subir do estômago. A raiva. A fome. A mão dela empurrou a minha mão. Eu --- não eu --- acertei o rosto do enfermeiro. O osso do nariz estalou.

\noindent O outro me acertou com a tira de couro.

\noindent Eu caí.

\noindent Eles me amarraram na cama. Braços esticados para os lados. Pernas juntas. A tira de couro na cabeça --- apertada, entrando na testa, prendendo o crânio contra o colchão.

\noindent Um deles colocou algo na minha boca. Borracha. Para não morder a língua.

\noindent O outro ligou a máquina.

\noindent \textit{\textquotedblleft{}Vai doer\textquotedblright{}\textit{, ele disse. }\textquotedblleft{}Mas passa.\textquotedblright{}}

\noindent Ele estava mentindo.

\noindent Não passa.

\vspace{0.5em}\noindent\rule{\textwidth}{0.5pt}\vspace{0.5em}

\noindent O choque entrou pelo lado esquerdo da cabeça.

\noindent Não foi como um susto. Foi como um apagão. Eu estava ali e não estava. O corpo tremia. Os dentes batiam na borracha. A língua sangrava. Eu ouvi um gemido --- era meu, mas não parecia meu. Parecia um animal.

\noindent A criatura gritou dentro de mim.

\noindent Não de dor.
De raiva.

\noindent \textit{\textquotedblleft{}Vou matar eles\textquotedblright{}\textit{, ela disse. }\textquotedblleft{}Vou matar todos.\textquotedblright{}}

\noindent \textit{\textquotedblleft{}Não\textquotedblright{}\textit{, eu pensei. }\textquotedblleft{}Não.\textquotedblright{}}

\noindent Mas ela não me ouviu.

\noindent \textit{\textquotedblleft{}Vou matar. Vou comer. Vou crescer dentro de ti até não caber mais. E aí vou sair.\textquotedblright{}}

\noindent O choque parou.

\noindent Eu fiquei ali, ofegante, com o gosto de sangue na boca e a cabeça pulsando.

\noindent Eles desamarraram.

\noindent \textit{\textquotedblleft{}Amanhã tem mais\textquotedblright{}}, disseram.

\noindent E saíram.

\vspace{0.5em}\noindent\rule{\textwidth}{0.5pt}\vspace{0.5em}

\noindent Teve mais.

\noindent Trinta e sete vezes. Ao longo de 15 anos. Trinta e sete sessões de eletrochoque, sem anestesia, sem sedação. Trinta e sete vezes em que a eletricidade entrou no meu cérebro e queimou pedaços de mim.

\noindent Cada vez que o choque passava, a criatura crescia um pouco mais.

\noindent O choque não afetava ela.
O choque afetava \textit{eu}.
Ela se alimentava dos destroços.

\noindent \textit{\textquotedblleft{}Eles estão me fortalecendo\textquotedblright{}\textit{, ela dizia. }\textquotedblleft{}Cada choque tira um pedaço de ti e dá pra mim.\textquotedblright{}}

\noindent \textit{\textquotedblleft{}Por quê?\textquotedblright{}}

\noindent \textit{\textquotedblleft{}Porque a dor é minha comida. Quanto mais tu sofre, mais eu cresço.\textquotedblright{}}

\noindent \textit{\textquotedblleft{}Então para de sentir dor.\textquotedblright{}}

\noindent \textit{\textquotedblleft{}Tu não consegue.\textquotedblright{}}

\noindent Ela estava certa.

\noindent Eu nunca consegui.

\noindent Depois da décima sessão, eu parei de contar. Depois da vigésima, parei de sentir. Não porque a dor passasse --- porque eu aprendi a \textit{sair} do corpo durante o choque.

\noindent Eu flutuava para fora. Ficava no teto, olhando para baixo, vendo meu corpo tremer na cama. A eletricidade passava. O corpo tremia. E eu olhava.

\noindent A criatura ficava.

\noindent Ela nunca saía do corpo.
Ela aproveitava para se esticar.
Para ocupar os espaços que o choque abria.

\vspace{0.5em}\noindent\rule{\textwidth}{0.5pt}\vspace{0.5em}

\noindent Depois de 37 sessões, os médicos desistiram.

\noindent \textit{\textquotedblleft{}Não responde ao tratamento\textquotedblright{}\textit{, escreveram no prontuário. }\textquotedblleft{}Refratário.\textquotedblright{}}

\noindent Eu não era refratário ao tratamento.
Eu era refratário à cura.
Porque a criatura não queria ser curada.
Ela queria crescer.

\noindent E cresceu.

\noindent Hoje, olhando para trás, eu entendo: o Colônia não era um hospital.
Era uma fábrica de monstros.

\noindent E eu era o melhor produto deles.

\vspace{0.5em}
\noindent\textit{\textbf{↪ Rotas:} \rota{DOC-10} \rotasep \rota{MEM-11} \rotasep \rota{MEM-18}}
\newpage\fragdef{MEM-18}% ============================================
% MEM-18
% Gerado de molambudos.md — build_miolo.py
% ============================================
\fragtitle{mem}{M}{MEM}{18}{A Tentativa de Fuga}
\markboth{MEM-18 --- A Tentativa de Fuga}{}

\noindent \textit{1943 --- o ano em que quase consegui}

\noindent \textbf{1943 --- O ano em que quase consegui}

\vspace{0.5em}\noindent\rule{\textwidth}{0.5pt}\vspace{0.5em}

\noindent Eu planejei a fuga por dois anos.

\noindent 1941 a 1943. Dia após dia, observando a rotina. O enfermeiro que chegava atrasado nas segundas. A janela da despensa que não fechava direito. O buraco na cerca dos fundos, escondido pelo mato.

\noindent A criatura me ajudava.

\noindent \textit{\textquotedblleft{}Tu vai conseguir\textquotedblright{}\textit{, ela dizia. }\textquotedblleft{}Vai sair daqui. Vai ver o mundo. Vai sentir o sol de novo.\textquotedblright{}}

\noindent \textit{\textquotedblleft{}Por que tu quer que eu fuja?\textquotedblright{}}

\noindent \textit{\textquotedblleft{}Porque aqui dentro tu é pequeno. Lá fora tu pode ser grande. Lá fora tem mais gente. Mais comida. Mais... espaço.\textquotedblright{}}

\noindent Eu devia ter desconfiado.

\noindent A criatura nunca queria nada de bom para mim.
Ela queria sempre o que era bom para \textit{ela}.

\noindent Mas eu era jovem. 1943, eu tinha 36 anos. Passados 26 dentro do Colônia. Eu queria sentir o vento no rosto sem grades no meio.

\vspace{0.5em}\noindent\rule{\textwidth}{0.5pt}\vspace{0.5em}

\noindent Na noite da fuga, o céu estava limpo. As estrelas brilhavam como nunca. Eu esperei o enfermeiro fazer a ronda das 2h, contei até mil, e levantei.

\noindent A porta da cela não estava trancada --- eu tinha aprendido a abrir com uma lasca de metal que escondi na sola do pé. Passei pelo corredor. Silêncio. A despensa. A janela. Aberta.

\noindent Pulei.

\noindent O mato estava alto. A terra estava fria. Eu corri. Meus pés descalços batendo no chão. O buraco na cerca. Eu passei.

\noindent Eu estava \textit{fora}.

\noindent O mundo era grande. Escuro. Silencioso. Mas era aberto. Eu olhei para o céu e vi as estrelas sem nenhuma grade na frente.

\noindent Eu chorei.

\noindent Não de tristeza.
De alegria.

\noindent Depois de 26 anos, eu estava livre.

\noindent Por 15 minutos.

\vspace{0.5em}\noindent\rule{\textwidth}{0.5pt}\vspace{0.5em}

\noindent A criatura falou enquanto eu corria pela estrada de terra.

\noindent \textit{\textquotedblleft{}Pra onde tu vai?\textquotedblright{}}

\noindent \textit{\textquotedblleft{}Longe.\textquotedblright{}}

\noindent \textit{\textquotedblleft{}Longe pra onde?\textquotedblright{}}

\noindent \textit{\textquotedblleft{}Bahia. Dizem que lá tem trabalho.\textquotedblright{}}

\noindent \textit{\textquotedblleft{}E o que tu vai fazer?\textquotedblright{}}

\noindent \textit{\textquotedblleft{}Trabalhar. Comer. Viver.\textquotedblright{}}

\noindent \textit{\textquotedblleft{}Viver?\textquotedblright{}\textit{ Ela riu. }\textquotedblleft{}Tu não sabe viver. Tu só sabe sobreviver.\textquotedblright{}}

\noindent \textit{\textquotedblleft{}Aprendo.\textquotedblright{}}

\noindent \textit{\textquotedblleft{}Não aprende. Tu é um bicho de hospício. Tu não sabe conversar com gente. Tu não sabe olhar nos olhos sem assustar. Tu não sabe o que é normal.\textquotedblright{}}

\noindent \textit{\textquotedblleft{}Cala a boca.\textquotedblright{}}

\noindent \textit{\textquotedblleft{}E o teu olho? Tu acha que alguém vai empregar um homem com um olho amarelo que brilha no escuro?\textquotedblright{}}

\noindent Eu parei de correr.

\noindent Fiquei parado no meio da estrada, ofegante.

\noindent \textit{\textquotedblleft{}Cala a boca.\textquotedblright{}}

\noindent \textit{\textquotedblleft{}Tu não tem para onde ir. O mundo lá fora não quer um ex-paciente do Colônia. O mundo quer gente normal. Tu não é normal. Tu é o Paciente 1.260. É o que está escrito em ti.\textquotedblright{}}

\noindent \textit{\textquotedblleft{}Eu arranco esse número da pele.\textquotedblright{}}

\noindent \textit{\textquotedblleft{}O número não está na pele, Joaquim. Está na alma. Tu pode arrancar a pele inteira, e o 1.260 vai continuar lá.\textquotedblright{}}

\noindent Eu sentei no meio da estrada.

\noindent A poeira subiu em volta de mim.

\noindent As estrelas continuavam brilhando, lindas, distantes.

\vspace{0.5em}\noindent\rule{\textwidth}{0.5pt}\vspace{0.5em}

\noindent Os enfermeiros me encontraram sentado no mesmo lugar, 6 horas depois.

\noindent Eu não tinha corrido.
Não tinha me escondido.
Eu só fiquei ali, esperando.

\noindent \textit{\textquotedblleft{}O que houve?\textquotedblright{}\textit{, perguntou o enfermeiro. }\textquotedblleft{}Desistiu?\textquotedblright{}}

\noindent \textit{\textquotedblleft{}Não.\textquotedblright{}\textit{ Eu levantei. }\textquotedblleft{}Só descobri que não tem para onde ir.\textquotedblright{}}

\noindent Eles me levaram de volta.

\noindent A cela estava igual. A cama de concreto. O balde. A grade. Tudo igual.

\noindent \textit{\textquotedblleft{}Bem-vindo de volta\textquotedblright{}}, disse a criatura.

\noindent \textit{\textquotedblleft{}Cala a boca.\textquotedblright{}}

\noindent Ela riu.

\noindent \textit{\textquotedblleft{}Da próxima vez, Joaquim, tu não tenta fugir. Tu fica. Tu espera. Tu escreve. Uma hora o diário sai. E aí...\textquotedblright{}}

\noindent \textit{\textquotedblleft{}Aí o quê?\textquotedblright{}}

\noindent \textit{\textquotedblleft{}Aí a fuga acontece. Mas não é tu que sai. É o diário. A fuga é dele.\textquotedblright{}}

\noindent \textit{\textquotedblleft{}E eu?\textquotedblright{}}

\noindent \textit{\textquotedblleft{}Tu fica. Tu sempre fica. O diário é que vai ver o mundo.\textquotedblright{}}\NE{A passagem da fuga é a mais coesa e linear de todo o diário de Joaquim — estrutura narrativa convencional, com começo, meio e fim, o que contrasta com a fragmentação dos demais fragmentos. Essa coesão sugere reescrita posterior, ou a mão da entidade organizando a narrativa para o leitor que a encontraria.}

\noindent Ela estava certa.

\noindent Eu nunca mais tentei fugir.
Mas o diário, sim.

\noindent O diário fugiu por mim.

\vspace{0.5em}
\noindent\textit{\textbf{↪ Rotas:} \rota{MEM-09} \rotasep \rota{MEM-19} \rotasep \rota{MEM-21}}
\newpage\fragdef{MEM-12}% ============================================
% MEM-12
% Gerado de molambudos.md — build_miolo.py
% ============================================
\fragtitle{mem}{M}{MEM}{12}{A Rasga Mortalha}
\markboth{MEM-12 --- A Rasga Mortalha}{}

\noindent \textit{O arauto no pátio}

\noindent \textbf{1945 --- Hospital Colônia de Barbacena}

\vspace{0.5em}\noindent\rule{\textwidth}{0.5pt}\vspace{0.5em}

\noindent Vinte e oito anos dentro do Colônia.

\noindent Vinte e oito anos de paredes que descascam, de gritos que viram rotina, de cheiro de creolina que não sai do nariz. Vinte e oito anos sem sentir o vento no rosto sem ser através de uma janela gradeada. Vinte e oito anos de solidão que não cura --- que petrifica.

\noindent Em 1945, eu estava na ala dos silenciados. Uma ala onde ninguém fala porque todos aprenderam que falar não adianta. Eu ficava encolhido num canto, com o diário no colo, escrevendo sem parar. A criatura dentro de mim estava quieta --- hibernando. Eu sentia o peso dela, o calor, a respiração lenta.

\noindent De vez em quando ela se virava dentro do meu peito e eu sentia uma pontada. A pontada não doía. Era \textit{lembrança}. Ela me lembrava que estava lá.

\vspace{0.5em}\noindent\rule{\textwidth}{0.5pt}\vspace{0.5em}

\noindent A Rasga Mortalha apareceu num inverno.

\noindent Junho de 1945. O frio entrava pelas frestas da parede e não saía. Eu dormia com as mãos enfiadas nas axilas, tremendo. Os outros pacientes tremiam também. A gente tremia junto, como bichos numa cerca, esperando o abate.

\noindent Numa noite, um som me acordou.

\noindent Não era grito.
Não era choro.
Era um canto.

\noindent Fino. Agudo. Um canto que subia e descia como uma mulher chorando muito longe.

\noindent \textit{\textquotedblleft{}Rasga mortalha...\textquotedblright{}}, alguém sussurrou ao meu lado.

\noindent Era Seu Nonô. Um velho que estava no Colônia desde 1920 --- 25 anos internado. Ele apontou para a janela com um dedo trêmulo.

\noindent \textit{\textquotedblleft{}Ela chegou.\textquotedblright{}}

\noindent \textit{\textquotedblleft{}Quem?\textquotedblright{}}

\noindent Ele não respondeu. Só apontou.

\noindent No batente da janela, havia uma coruja.

\noindent Ela era branca. Não branca como neve --- branca como osso velho. As penas pareciam luz na escuridão, como se o corpo dela fosse feito de luar condensado. Os olhos eram pretos, redondos, sem pálpebras. Ela me olhou. Virou a cabeça 180 graus --- um movimento lento, fluido, que nenhum pássaro devia ser capaz de fazer.

\noindent \textit{\textquotedblleft{}Rasga mortalha\textquotedblright{}\textit{, repetiu Seu Nonô. }\textquotedblleft{}Quando ela canta, alguém morre. Ela rasga a mortalha do finado com o bico pra alma poder sair. Por isso o nome.\textquotedblright{}}

\noindent A coruja abriu o bico.

\noindent O som que saiu não era um piado.
Era um grito humano.
Um grito de mulher --- mas não de dor.
De \textit{aviso}.

\noindent O grito atravessou a noite. Atravessou as grades. Atravessou o meu peito.

\noindent A criatura dentro de mim \textit{acordou}.

\noindent Eu a senti erguer a cabeça dentro do meu corpo. Senti os olhos dela se abrirem --- os olhos amarelos. Ela \textit{viu} a Rasga Mortalha.

\noindent E a Rasga Mortalha viu ela.

\noindent A coruja inclinou a cabeça. O grito parou. Ficou um silêncio tão absoluto que eu ouvia meu próprio sangue circulando.

\noindent E a Rasga Mortalha falou.

\noindent Não com a voz dela.

\noindent Com a voz da criatura.

\noindent A voz que saía do buraco dentro de mim.

\noindent \textit{\textquotedblleft{}Ainda não.\textquotedblright{}}

\noindent A coruja abriu as asas. Eram enormes. Brancas. Pareciam um lençol flutuando no escuro.

\noindent \textit{\textquotedblleft{}Mas em breve.\textquotedblright{}}

\noindent Ela voou.

\noindent O branco dela sumiu na noite.

\vspace{0.5em}\noindent\rule{\textwidth}{0.5pt}\vspace{0.5em}

\noindent No dia seguinte, Seu Nonô foi encontrado morto na cama.

\noindent O rosto estava virado para a janela.
A boca aberta.
Os olhos abertos.

\noindent As pálpebras tinham sido arrancadas.

\noindent Não sobrou sangue. Não sobrou marca. As pálpebras simplesmente tinham \textit{sumido}. Como se alguém as tivesse levado com uma precisão cirúrgica.

\noindent Os enfermeiros disseram que foram os ratos.

\noindent Eu sabia que não.

\noindent A Rasga Mortalha tinha levado as pálpebras do Seu Nonô.
Para que ele \textit{nunca mais fechasse os olhos}.
Para que ele visse o que estava chegando.
Mesmo depois de morto.

\vspace{0.5em}\noindent\rule{\textwidth}{0.5pt}\vspace{0.5em}

\noindent Na semana seguinte, a coruja voltou.

\noindent Toda noite, ela pousava no batente da janela e me olhava. Não cantava. Não gritava. Ficava ali, imóvel, com os olhos pretos fixos em mim.

\noindent E toda noite, a criatura dentro de mim falava.

\noindent \textit{\textquotedblleft{}Ela é minha língua\textquotedblright{}\textit{, disse. }\textquotedblleft{}Meu olho voador. Minha voz no mundo dos vivos.\textquotedblright{}}

\noindent \textit{\textquotedblleft{}Por que ela veio?\textquotedblright{}}

\noindent \textit{\textquotedblleft{}Pra ver como você está.\textquotedblright{}}

\noindent \textit{\textquotedblleft{}E como eu estou?\textquotedblright{}}

\noindent Silêncio.

\noindent \textit{\textquotedblleft{}Pronto.\textquotedblright{}}

\noindent \textit{\textquotedblleft{}Pronto pra quê?\textquotedblright{}}

\noindent \textit{\textquotedblleft{}Pra passar o diário adiante.\textquotedblright{}}

\vspace{0.5em}\noindent\rule{\textwidth}{0.5pt}\vspace{0.5em}

\noindent Depois daquela noite, eu comecei a escrever de um jeito diferente.

\noindent A letra mudou. As palavras vinham mais rápido. Eu não precisava pensar --- a mão se movia sozinha, como se alguém estivesse guiando a caneta.

\noindent O diário não era mais meu.

\noindent A criatura estava escrevendo através de mim.

\noindent Eu era só a mão.
Ela era a voz.

\noindent E o que ela escrevia --- eu não entendia na hora. Só entendi depois. Quando li de novo. Quando vi que as palavras não eram para mim.

\noindent Eram para quem lesse depois.

\noindent Eram para você.

\vspace{0.5em}\noindent\rule{\textwidth}{0.5pt}\vspace{0.5em}

\noindent Na última noite, a coruja não veio.

\noindent Eu acordei sozinho, na cela escura, e soube que ela tinha ido embora.

\noindent Mas a criatura dentro de mim falou:

\noindent \textit{\textquotedblleft{}Ela foi preparar o caminho.\textquotedblright{}}

\noindent \textit{\textquotedblleft{}Preparar como?\textquotedblright{}}

\noindent \textit{\textquotedblleft{}O diário não pode ficar aqui. Tem que sair. Tem que ser encontrado. Tem que ser aberto.\textquotedblright{}}

\noindent \textit{\textquotedblleft{}Por quem?\textquotedblright{}}

\noindent A criatura riu. A risada ecoou dentro do meu peito.

\noindent \textit{\textquotedblleft{}Alguém que precisa tanto de mim quanto você precisou.\textquotedblright{}}

\noindent \textit{\textquotedblleft{}E quando isso vai acontecer?\textquotedblright{}}

\noindent \textit{\textquotedblleft{}Quando você morrer.\textquotedblright{}}

\noindent \textit{\textquotedblleft{}E até lá?\textquotedblright{}}

\noindent \textit{\textquotedblleft{}A gente espera.\textquotedblright{}}

\noindent Ela ficou quieta.

\noindent \textit{\textquotedblleft{}A gente sempre espera.\textquotedblright{}}

\vspace{0.5em}\noindent\rule{\textwidth}{0.5pt}\vspace{0.5em}

\noindent Foram mais 34 anos de espera.

\noindent 34 anos ouvindo a Rasga Mortalha nos telhados do Colônia.

\noindent 34 anos escrevendo palavras que não eram minhas.

\noindent 34 anos sendo a mão de um monstro.

\noindent Até que, em 1979, o médico novo chegou.

\noindent E o ciclo começou a se fechar.

\vspace{0.5em}
\noindent\textit{\textbf{↪ Rotas:} \rota{MEM-11} \rotasep \rota{MEM-13} \rotasep \rota{DOC-07}}

\newpage\fragdef{MEM-21}% ============================================
% MEM-21
% Gerado de molambudos.md — build_miolo.py
% ============================================
\fragtitle{mem}{M}{MEM}{21}{O Jornal}
\markboth{MEM-21 --- O Jornal}{}

\noindent \textit{1948 --- a notícia que mentia}

\noindent \textbf{1948 --- Uma notícia do lado de fora}

\vspace{0.5em}\noindent\rule{\textwidth}{0.5pt}\vspace{0.5em}

\noindent Um enfermeiro deixou um jornal na cela. Não sei por quê. Talvez por esquecimento. Talvez por pena. O jornal era de Belo Horizonte, de 1947. Tinha cheiro de mofo e de fumaça de cigarro alheio. A textura das páginas era quebradiça, amarelada como folha seca --- eu segurei com cuidado, como se tudo que eu ainda acreditava sobre o mundo lá fora também pudesse desmanchar entre os dedos. Uma notícia na página 5 me fez parar:

\noindent \textit{\textquotedblleft{}COLÔNIA DE BARBACENA COMPLETA 40 ANOS COM 5 MIL INTERNADOS\textquotedblright{}}

\noindent Eu li a matéria. Dizia que o Colônia era \textquotedblleft{}um dos maiores hospitais psiquiátricos da América Latina\textquotedblright{}. Dizia que o governo investia em \textquotedblleft{}terapias modernas\textquotedblright{}. Dizia que os pacientes \textquotedblleft{}recebiam tratamento humanizado\textquotedblright{}.

\noindent Mentira.

\noindent Tudo mentira.

\noindent Eu estava ali. Eu via os eletrochoques sem anestesia. Eu via os pacientes morrerem de pneumonia sem remédio. Eu via os corpos empilhados na carroceria do caminhão todas as manhãs.

\noindent O jornal mentia.

\noindent E eu pensei: \textit{\textquotedblleft{}Se o jornal mente sobre isso, sobre o que mais ele mente?\textquotedblright{}}

\noindent A criatura respondeu:

\noindent \textit{\textquotedblleft{}Ele mente sobre tudo. Os jornais sempre mentem. Eles dizem que o mundo é organizado. Não é. O mundo é uma vala. E a vala está cheia de corpos que o governo não quer que você veja.\textquotedblright{}}

\noindent \textit{\textquotedblleft{}Por que você está me dizendo isso?\textquotedblright{}}

\noindent \textit{\textquotedblleft{}Porque você é a prova. Você está vivo. Você viu. Você pode contar.\textquotedblright{}}

\noindent \textit{\textquotedblleft{}Quem vai acreditar em mim?\textquotedblright{}}

\noindent \textit{\textquotedblleft{}Quem ler o diário.\textquotedblright{}}

\noindent Ela estava certa.

\noindent O jornal mentiu.
O diário não.

\noindent O diário é a verdade que o jornal nunca publicou.

\noindent Guardei aquela página do jornal dentro do colchão, junto com os primeiros cadernos do diário. Não para provar que ele mentiu. Para lembrar, sempre que eu duvidasse de mim mesmo, que eu não inventei nada do que vi.

\vspace{0.5em}
\noindent\textit{\textbf{↪ Rotas:} \rota{MEM-22} \rotasep \rota{DOC-19}}
\newpage\fragdef{MEM-19}% ============================================
% MEM-19
% Gerado de molambudos.md — build_miolo.py
% ============================================
\fragtitle{mem}{M}{MEM}{19}{O Visitante}
\markboth{MEM-19 --- O Visitante}{}

\noindent \textit{Alguém veio me ver}

\noindent \textbf{1950 --- Alguém veio me ver}

\vspace{0.5em}\noindent\rule{\textwidth}{0.5pt}\vspace{0.5em}

\noindent Ninguém nunca vinha me ver.

\noindent Não tinha família. Não tinha amigos. O velho Anselmo tinha morrido no curral. Minha mãe tinha desaparecido no caminhão. O resto --- vizinhos, conhecidos --- tinha morrido na seca ou se espalhado pelo mundo.

\noindent Eu era um homem sem passado.

\noindent Por isso, quando o enfermeiro disse \textit{\textquotedblleft{}Visita para você\textquotedblright{}}, eu achei que era engano.

\noindent \textit{\textquotedblleft{}Quem?\textquotedblright{}}

\noindent \textit{\textquotedblleft{}Não sei. Uma mulher. Está na sala de visitas.\textquotedblright{}}

\noindent Uma mulher.

\noindent Minha mãe? Impossível. Ela teria mais de 70 anos. E estava morta --- ou pior.

\noindent Mas eu fui.

\noindent Atravessei o corredor. Os outros pacientes olhavam. Visita era raro no Colônia. As pessoas não visitavam loucos. Loucos eram esquecidos.

\noindent A sala de visitas era uma sala com duas cadeiras e uma mesa. Numa das cadeiras, uma mulher estava sentada.

\noindent Era idosa. Cabelos brancos, pele enrugada, mãos grossas de trabalho. Ela usava um vestido simples, azul, com um xale nos ombros.

\noindent Ela levantou os olhos quando eu entrei.

\noindent Os olhos dela eram castanhos. Cansados. Mas vivos.

\noindent \textit{\textquotedblleft{}Joaquim?\textquotedblright{}}

\noindent A voz era trêmula.

\noindent \textit{\textquotedblleft{}...Quem é?\textquotedblright{}}

\noindent Ela sorriu. Um sorriso triste.

\noindent \textit{\textquotedblleft{}Sou Maria. Maria do Rosário.\textquotedblright{}}

\noindent Minha mãe.

\noindent Eu sentei.

\noindent Fiquei olhando para ela. Tentando encontrar a mulher que me levou pela mão na estrada de Senador Pompeu. A mulher que tossia no curral. A mulher que foi levada no caminhão.

\noindent \textit{\textquotedblleft{}A senhora...\textquotedblright{}\textit{ Minha voz falhou. }\textquotedblleft{}A senhora está viva.\textquotedblright{}}

\noindent \textit{\textquotedblleft{}Estou.\textquotedblright{}}

\noindent \textit{\textquotedblleft{}Onde a senhora estava?\textquotedblright{}}

\noindent Ela baixou os olhos.

\noindent \textit{\textquotedblleft{}Fui levada para um bordel em Fortaleza. Fiquei lá 3 anos. Depois fugi. Casei de novo. Tive outros filhos. Morei em São Paulo.\textquotedblright{}}

\noindent \textit{\textquotedblleft{}E nunca...\textquotedblright{}\textit{ Eu engoli. }\textquotedblleft{}Nunca voltou?\textquotedblright{}}

\noindent \textit{\textquotedblleft{}Voltei.\textquotedblright{}\textit{ Ela me olhou. }\textquotedblleft{}Voltei para te buscar. Em 1919. Disseram que tu tinha morrido.\textquotedblright{}}

\noindent \textit{\textquotedblleft{}Mentira.\textquotedblright{}}

\noindent \textit{\textquotedblleft{}Eu sei. Descobri ano passado. Uma vizinha disse que o filho dela trabalhou aqui e que tinha um paciente do Ceará. 1.260. O mesmo nome. Vim ver.\textquotedblright{}}

\noindent \textit{\textquotedblleft{}A senhora veio me buscar?\textquotedblright{}}

\noindent Ela não respondeu.

\noindent Olhou para as mãos.

\noindent \textit{\textquotedblleft{}Joaquim, eu...\textquotedblright{}\textit{ Ela parou. }\textquotedblleft{}Eu tenho outra família. Outros filhos. Eles não sabem de ti. Não sabem que eu tive um filho antes.\textquotedblright{}}

\noindent \textit{\textquotedblleft{}E a senhora vai contar?\textquotedblright{}}

\noindent Ela não respondeu.

\noindent Eu entendi.

\noindent Ela não veio me buscar.
Ela veio se \textit{despedir}.

\noindent \textit{\textquotedblleft{}Tá\textquotedblright{}\textit{, eu disse. }\textquotedblleft{}Tá bom.\textquotedblright{}}

\noindent \textit{\textquotedblleft{}Joaquim...\textquotedblright{}}

\noindent \textit{\textquotedblleft{}Tá bom, mãe.\textquotedblright{}}

\noindent A palavra \textit{mãe} saiu sem eu pensar. Ela ouviu. Os olhos dela encheram de lágrima.

\noindent \textit{\textquotedblleft{}Me perdoa\textquotedblright{}}, ela disse.

\noindent \textit{\textquotedblleft{}Perdoar o quê?\textquotedblright{}}

\noindent \textit{\textquotedblleft{}Por ter deixado tu.\textquotedblright{}}

\noindent \textit{\textquotedblleft{}A senhora não deixou. Levaram.\textquotedblright{}}

\noindent \textit{\textquotedblleft{}Mas eu não voltei.\textquotedblright{}}

\noindent \textit{\textquotedblleft{}A senhora voltou agora.\textquotedblright{}}

\noindent Ela chorou.

\noindent Eu não chorei. A criatura dentro de mim não deixou.

\noindent Ela ficou mais 20 minutos. Contou sobre a vida dela. Os outros filhos. O marido. A casa em São Paulo. Eu ouvi. Assenti. Sorri nos lugares certos.

\noindent Quando ela foi embora, ela me abraçou.

\noindent O abraço dela era quente. Macio. Cheirava a sabão.

\noindent \textit{\textquotedblleft{}Eu vou escrever\textquotedblright{}}, ela disse.

\noindent \textit{\textquotedblleft{}Tá.\textquotedblright{}}

\noindent Ela não escreveu.

\noindent Nunca mais.

\noindent Mas, naquela noite, a criatura ficou quieta. Ela respeitou o silêncio.

\noindent \textit{\textquotedblleft{}Tu tá triste?\textquotedblright{}} ela perguntou.

\noindent \textit{\textquotedblleft{}Não.\textquotedblright{}}

\noindent \textit{\textquotedblleft{}Tá.\textquotedblright{}}

\noindent \textit{\textquotedblleft{}Ela não volta.\textquotedblright{}}

\noindent \textit{\textquotedblleft{}Eu sei.\textquotedblright{}}

\noindent \textit{\textquotedblleft{}Mas por que ela veio?\textquotedblright{}}

\noindent Pensei.

\noindent \textit{\textquotedblleft{}Para eu saber que ela existiu. Para eu saber que não fui um sonho.\textquotedblright{}}

\noindent A criatura não respondeu.

\noindent Naquela noite, pela primeira vez em anos, eu não senti fome.

\noindent Senti saudade.

\noindent Saudade de uma mãe que eu mal conhecia.
Saudade de uma vida que eu nunca tive.
Saudade de ser filho de alguém.

\noindent Saudade de ser gente.

\vspace{0.5em}
\noindent\textit{\textbf{↪ Rotas:} \rota{MEM-15} \rotasep \rota{MEM-20}}
\newpage\fragdef{DOC-10}% ============================================
% DOC-10
% Gerado de molambudos.md — build_miolo.py
% ============================================
\fragtitle{doc}{D}{DOC}{10}{Registro de Eletroconvulsoterapia}
\markboth{DOC-10 --- Registro de Eletroconvulsoterapia}{}

\noindent \textit{37 sessões sem anestesia}

\noindent \textbf{Hospital Colônia de Barbacena --- Ala de Terapias Agressivas}
\textbf{Paciente:} Joaquim Antônio Correia --- Registro 1.260
\textbf{Período:} 1937 a 1952

\vspace{0.5em}\noindent\rule{\textwidth}{0.5pt}\vspace{0.5em}

\subsection*{SUMÁRIO}

\noindent Sessões registradas: \textbf{37}
Média anual: \textbf{2.46}
Equipamento: ECT Konvulsator Modelo IV (Alemanha, 1932)
Anestesia: \textbf{Não utilizada} (nenhuma das 37 sessões)
Relaxante muscular: \textbf{Não utilizado} (nenhuma das 37 sessões)
Sedação prévia: \textbf{Não} (paciente classificado como \textquotedblleft{}refratário à sedação verbal\textquotedblright{})

\vspace{0.5em}\noindent\rule{\textwidth}{0.5pt}\vspace{0.5em}

\subsection*{REGISTRO DETALHADO (parcial)}

\begin{moltablesix}
\hline
\# & Data & Vol\-ta\-gem & Duração & Efeitos Imediatos & Efeitos Tardios \\
\hline
1 & 12/mar/37 & 110V & 2.5s & Convulsão tônico-clônica, espuma, perda de consciência por 12min & Laceração de língua, incontinência urinária, confusão por 6h \\
2 & 19/mar/37 & 110V & 3.0s & Convulsão, fratura de vértebra T7 (não tratada) & Dor lombar crônica, relatada pelo paciente por 40+ anos \\
3 & 26/mar/37 & 130V & 3.0s & Convulsão prolongada (4min), cianose & Amnésia retrógrada de 48h \\
4 & 2/abr/37 & 130V & 3.5s & Convulsão, mordedura grave de língua & Dificuldade de alimentação por 10 dias \\
5 & 9/abr/37 & 130V & 3.5s & Convulsão, deslocamento de mandíbula & Redução manual sem anestesia \\
--- & --- & --- & --- & --- & --- \\
18 & 15/jun/40 & 150V & 4.0s & Convulsão, parada respiratória (reanimada) & Paciente relata \textquotedblleft{}flutuar para fora do corpo\textquotedblright{} \\
--- & --- & --- & --- & --- & --- \\
27 & 3/nov/45 & 150V & 4.0s & Convulsão, fratura de fêmur direito & Claudicação permanente \\
--- & --- & --- & --- & --- & --- \\
37 & 18/fev/52 & 170V & 5.0s & Convulsão máxima, parada cardíaca (reanimado) & Paciente inconsciente por 3 dias. Ao acordar: \textquotedblleft{}Joaquim não está aqui\textquotedblright{} \\
\hline
\end{moltablesix}

\vspace{0.5em}\noindent\rule{\textwidth}{0.5pt}\vspace{0.5em}

\subsection*{OBSERVAÇÕES CLÍNICAS (transcrição)}

\noindent \textbf{Dr. Nelson Fontes, 1937:}
\textit{\textquotedblleft{}Paciente não responde à terapia. Mantém delírio inalterado. Aumentar voltagem para 130V.\textquotedblright{}}

\noindent \textbf{Dr. Amílcar Barbosa, 1940:}
\textit{\textquotedblleft{}Após 18 sessões, paciente passou a relatar 'experiências extracorpóreas' durante o choque. Considero isto um avanço --- a mente está se desligando do corpo. Continuar.\textquotedblright{}}

\noindent \textbf{Dr. Amílcar Barbosa, 1945:}
\textit{\textquotedblleft{}Fratura de fêmur durante sessão. Paciente não sentiu a fratura --- relatou apenas que 'ouviu o osso quebrar mas não doeu'. Sugiro aumentar voltagem. Tolerância do paciente é excepcional.\textquotedblright{}}

\noindent \textbf{Dr. Heitor Oliveira, 1979 (nota marginal, 30 anos depois):}
\textit{\textquotedblleft{}Esta ficha é um documento de tortura. 37 sessões sem anestesia. Fraturas não tratadas. Sem consentimento. E o paciente sobreviveu a tudo --- lúcido o suficiente para escrever um diário de 62 anos. Isso não é refratariedade. Isso é resiliência. Ou possessão. Ou as duas coisas.\textquotedblright{}}

\vspace{0.5em}\noindent\rule{\textwidth}{0.5pt}\vspace{0.5em}

\subsection*{NOTA DO ARQUIVISTA (2026)}

\noindent O equipamento Konvulsator Modelo IV foi projetado para aplicação com anestesia e relaxante muscular. No Colônia, foi usado sem nenhum dos dois. A voltagem máxima recomendada pelo fabricante era 130V. O paciente 1.260 recebeu até 170V na última sessão.

\noindent A sobrevivência do paciente sem dano cerebral aparente é \textquotedblleft{}clinicamente inexplicável\textquotedblright{} (Dr. Torres, 2023).

\noindent O paciente 1.260 declarou, em seu diário (1952): \textit{\textquotedblleft{}O choque não afeta a criatura. O choque queima pedaços de mim. A criatura come o que sobra.\textquotedblright{}}

\noindent Sugere-se que a \textquotedblleft{}inexplicável\textquotedblright{} tolerância do paciente ao tratamento pode estar relacionada ao conteúdo de seu delírio --- especificamente, à crença em uma entidade interna que absorvia o trauma.

\noindent Ou, alternativamente, que o delírio não era delírio.

\noindent O arquivo não toma posição sobre esta última hipótese.

\vspace{0.5em}
\noindent\textit{\textbf{↪ Rotas:} \rota{MEM-17} \rotasep \rota{DOC-02} \rotasep \rota{EPI-01}}
\newpage\fragdef{MEM-22}% ============================================
% MEM-22
% Gerado de molambudos.md — build_miolo.py
% ============================================
\fragtitle{mem}{M}{MEM}{22}{As Vozes}
\markboth{MEM-22 --- As Vozes}{}

\noindent \textit{A ala que grita}

\noindent \textbf{1952 --- A ala que grita}

\vspace{0.5em}\noindent\rule{\textwidth}{0.5pt}\vspace{0.5em}

\noindent A ala dos agitados ficava ao lado da ala dos silenciados.

\noindent Duas alas. Dois mundos. De um lado, gente que não parava de gritar. Do outro, gente que não conseguia mais falar.

\noindent Eu ficava no meio.

\noindent Durante o dia, eu ouvia os gritos. Gritos de medo. Gritos de raiva. Gritos sem sentido. Gritos que eram só som, sem palavra, sem significado.

\noindent Durante a noite, o silêncio dos silenciados era pior. Porque o silêncio também grita. Só que o silêncio grita dentro da sua cabeça.

\noindent A criatura gostava dos gritos.

\noindent \textit{\textquotedblleft{}Escuta\textquotedblright{}\textit{, ela dizia. }\textquotedblleft{}Cada grito é uma alma se soltando. Cada grito é um pedaço de alguém que vai embora. Quando o grito para, a pessoa já foi.\textquotedblright{}}

\noindent \textit{\textquotedblleft{}Já foi pra onde?\textquotedblright{}}

\noindent \textit{\textquotedblleft{}Pra dentro de mim.\textquotedblright{}}

\noindent \textit{ \textquotedblleft{}Você come os gritos?\textquotedblright{}}

\noindent \textit{\textquotedblleft{}Como tudo. Como a dor. Como o medo. Como a loucura. É tudo comida.\textquotedblright{}}

\noindent Eu tampava os ouvidos.

\noindent Mas a criatura ouvia por mim.

\vspace{0.5em}\noindent\rule{\textwidth}{0.5pt}\vspace{0.5em}

\noindent Uma noite, um grito diferente.

\noindent Não era de medo. Era de perda. Alguém gritando o nome de alguém. Uma mulher. Ela gritava: \textit{\textquotedblleft{}José! José!\textquotedblright{}} O nome ecoou pelo corredor por minutos.

\noindent Depois, silêncio.

\noindent Na manhã seguinte, perguntei ao enfermeiro quem era José.

\noindent \textit{\textquotedblleft{}Morreu ontem. Pneumonia.\textquotedblright{}}

\noindent \textit{\textquotedblleft{}E quem gritava?\textquotedblright{}}

\noindent \textit{\textquotedblleft{}A mulher dele. Internaram os dois juntos em 1940. Ele morreu. Ela ficou.\textquotedblright{}}

\noindent A mulher nunca mais gritou.

\noindent Ela ficou em silêncio pelo resto da vida.

\noindent Três meses depois, ela também morreu.

\noindent Dizem que foi desnutrição.

\noindent Eu acho que foi falta de ter para quem gritar.

\noindent Fiquei pensando: quando chegasse a minha vez, quem gritaria o meu nome pelo corredor? Ninguém. Não tinha mulher, não tinha filho, não tinha ninguém esperando do lado de fora dessas paredes. Só a criatura. E ela não gritava por mim --- ela só esperava, com fome, que eu também virasse silêncio.

\vspace{0.5em}
\noindent\textit{\textbf{↪ Rotas:} \rota{MEM-23} \rotasep \rota{MEM-09}}
\newpage\fragdef{MEM-23}% ============================================
% MEM-23
% Gerado de molambudos.md — build_miolo.py
% ============================================
\fragtitle{mem}{M}{MEM}{23}{A Neve de Cinzas}
\markboth{MEM-23 --- A Neve de Cinzas}{}

\noindent \textit{O incêndio de 1960}

\noindent \textbf{1960 --- O incêndio no Colônia}

\vspace{0.5em}\noindent\rule{\textwidth}{0.5pt}\vspace{0.5em}

\noindent Em 1960, um incêndio destruiu a ala administrativa do Colônia.

\noindent Ninguém morreu --- o fogo começou de dia, os funcionários saíram a tempo. Mas os pacientes sentiram o cheiro de fumaça e entraram em pânico. Gritaria. Correria. Os enfermeiros trancaram as portas das celas e saíram.

\noindent Ficamos trancados dentro, com a fumaça entrando pelas janelas.

\noindent Meu coração batia mais forte que os punhos dos outros pacientes na porta de ferro. Senti medo --- não do fogo, mas de ficar esquecido ali dentro, como sempre fomos esquecidos. Não gritei. Escolhi guardar o ar para respirar, não para gastar em grito que ninguém ouviria.

\noindent Eu olhei pela grade.

\noindent O céu estava cinza. Não de nuvem --- de fumaça. Pedaços de papel queimado voavam no ar. Parecia neve preta.

\noindent \textit{\textquotedblleft{}O mundo está queimando\textquotedblright{}}, disse a criatura.

\noindent \textit{\textquotedblleft{}É só um incêndio.\textquotedblright{}}

\noindent \textit{\textquotedblleft{}Não. O mundo está queimando. E você está vendo.\textquotedblright{}}

\noindent \textit{\textquotedblleft{}O que isso significa?\textquotedblright{}}

\noindent \textit{\textquotedblleft{}Significa que tudo que é construído pode ser queimado. O Colônia. O governo. O mundo. Tudo vira cinza.\textquotedblright{}}

\noindent \textit{\textquotedblleft{}E o que sobra?\textquotedblright{}}

\noindent \textit{\textquotedblleft{}O que não queima.\textquotedblright{}}

\noindent \textit{\textquotedblleft{}O que não queima?\textquotedblright{}}

\noindent Ela riu.

\noindent \textit{\textquotedblleft{}A fome. A fome nunca queima.\textquotedblright{}}

\noindent O incêndio foi controlado à tarde.

\noindent No dia seguinte, a ala administrativa era só cinzas.

\noindent Mas a ala dos pacientes continuava de pé.

\noindent Como sempre.
Como sempre vai ficar.

\noindent Foi nesse dia que entendi uma coisa que levei o resto da vida para admitir: a fome de que a criatura falava não era só dela. Era minha também. E eu já não sabia dizer, com certeza, onde a minha acabava e a dela começava.

\vspace{0.5em}
\noindent\textit{\textbf{↪ Rotas:} \rota{MEM-24} \rotasep \rota{MEM-25}}
\newpage\fragdef{MEM-20}% ============================================
% MEM-20
% Gerado de molambudos.md — build_miolo.py
% ============================================
\fragtitle{mem}{M}{MEM}{20}{A Morte do Vizinho de Cela}
\markboth{MEM-20 --- A Morte do Vizinho de Cela}{}

\noindent \textit{O único amigo que tive}

\noindent \textbf{1963 --- O único amigo que tive}

\vspace{0.5em}\noindent\rule{\textwidth}{0.5pt}\vspace{0.5em}

\noindent Ele se chamava Severino.

\noindent Paraense. 52 anos. Internado em 1935 porque a família não queria mais cuidar dele. Diagnóstico: \textquotedblleft{}idiotia\textquotedblright{}. O diagnóstico real: gagueira e epilepsia leve. Mas nos anos 30, gagueira era tratada com eletrochoque. Depois de 30 choques, Severino não gaguejava mais. Nem falava. Nem mexia.

\noindent Ele ficava na cela ao lado da minha.

\noindent Durante 20 anos, a gente conversava batendo na parede.

\noindent Código:
1 batida = sim
2 batidas = não
3 batidas = como vai você
4 batidas = estou com fome
5 batidas = estou com medo

\noindent A gente desenvolveu uma linguagem. Batidas longas, batidas curtas. Dava para dizer frases complexas. \textit{\textquotedblleft{}O enfermeiro veio hoje\textquotedblright{}} --- 3 longas, 2 curtas, pausa, 4 curtas. Eu entendia.

\noindent Severino era a única pessoa no Colônia que me tratava como gente.

\noindent Ele não olhava para o meu olho amarelo com medo. Ele olhava como quem olha para um amigo.

\vspace{0.5em}\noindent\rule{\textwidth}{0.5pt}\vspace{0.5em}

\noindent Em 1963, Severino pegou pneumonia.

\noindent O Colônia não tratava pneumonia. Deixava morrer. Morriam 30 pacientes por mês de pneumonia, tuberculose, desnutrição. O Colônia chamava de \textquotedblleft{}seleção natural\textquotedblright{}.

\noindent Eu ouvi a tosse dele pela parede.

\noindent A tosse durou uma semana. Cada dia mais fraca. No sétimo dia, a tosse parou.

\noindent Bati na parede.

\noindent 1 batida. \textit{\textquotedblleft{}Severino?\textquotedblright{}}

\noindent Silêncio.

\noindent 2 batidas. \textit{\textquotedblleft{}Acorda.\textquotedblright{}}

\noindent Silêncio.

\noindent 3 batidas. \textit{\textquotedblleft{}Como vai você?\textquotedblright{}}

\noindent Silêncio.

\noindent Eu sabia que ele tinha morrido. Mas fiquei batendo. Bati a noite inteira. Diferentes códigos. Diferentes perguntas. \textit{\textquotedblleft{}Você está ai?\textquotedblright{}\textit{ }\textquotedblleft{}Responde.\textquotedblright{}\textit{ }\textquotedblleft{}Não me deixa sozinho.\textquotedblright{}}

\noindent Ninguém respondeu.

\noindent De manhã, os enfermeiros vieram buscar o corpo.

\noindent Eu ouvi a porta da cela abrir. Ouvi o arrastar dos pés no chão. Depois, silêncio.

\noindent A cela ao lado ficou vazia.

\noindent Por três meses, ninguém ocupou.

\noindent Eu batia na parede de vez em quando. Não para conversar --- para ouvir o som do oco.

\vspace{0.5em}\noindent\rule{\textwidth}{0.5pt}\vspace{0.5em}

\noindent Depois que Severino morreu, a criatura falou:

\noindent \textit{\textquotedblleft{}Tu nunca mais vai ter um amigo.\textquotedblright{}}

\noindent \textit{\textquotedblleft{}Cala a boca.\textquotedblright{}}

\noindent \textit{\textquotedblleft{}É verdade. Tu é velho. Tarado. Ninguém quer ser amigo de um velho tarado.\textquotedblright{}}

\noindent \textit{\textquotedblleft{}Cala a boca.\textquotedblright{}}

\noindent \textit{\textquotedblleft{}Severino era o único. E ele morreu. Porque todo mundo que chega perto de ti morre.\textquotedblright{}}

\noindent \textit{\textquotedblleft{}Cala a boca!\textquotedblright{}}

\noindent Eu bati a mão na parede. A mão sangrou.

\noindent \textit{\textquotedblleft{}Não foi minha culpa.\textquotedblright{}}

\noindent \textit{\textquotedblleft{}Foi.\textquotedblright{}}

\noindent \textit{\textquotedblleft{}Não foi.\textquotedblright{}}

\noindent \textit{\textquotedblleft{}Severino pegou pneumonia no inverno passado. Sabe o que causou a pneumonia? O frio. Sabe onde o frio entrava? Pelas frestas da parede. A parede que separa a cela dele da sua. As frestas são do lado dele. Mas o frio veio da sua cela.\textquotedblright{}}

\noindent \textit{\textquotedblleft{}Mentira.\textquotedblright{}}

\noindent \textit{\textquotedblleft{}Verdade. Tu respirou o ar frio, o ar passou pela parede, chegou nele. Tu matou ele sem querer. Mas matou.\textquotedblright{}}

\noindent Eu não respondi.

\noindent Fiquei encolhido no canto, a mão sangrando, o coração apertado.

\noindent Ela estava mentindo.
Ela sempre mentia.

\noindent Mas, naquela noite, eu senti o frio entrando pelas frestas.

\noindent E pensei: \textit{\textquotedblleft{}E se for verdade?\textquotedblright{}}

\noindent Mas ali, naquela noite, eu fiz uma escolha. Não a escolha de acreditar nela ou em mim — a escolha de lembrar o Severino como ele era: o único amigo que tive. A criatura podia mentir sobre o frio, sobre a pneumonia, sobre a parede. Eu podia escolher não ouvir. Podia escolher o que guardar dele. E eu guardei a batida na parede. As 5 batidas. A linguagem que a gente criou junto. Isso, ela não podia roubar.

\vspace{0.5em}\noindent\rule{\textwidth}{0.5pt}\vspace{0.5em}

\noindent Nunca mais fiz amizade com ninguém.

\noindent Não valia a pena.
Todo mundo que chegava perto de mim --- a criatura encontrava um jeito de destruir.

\noindent A solidão não era minha escolha.
Era o veneno que ela espalhava em volta de mim.

\noindent Para que eu dependesse só dela.
Para que eu não tivesse mais ninguém.

\noindent Ela queria ser minha única companhia.

\noindent E conseguiu.

\vspace{0.5em}
\noindent\textit{\textbf{↪ Rotas:} \rota{MEM-19} \rotasep \rota{DOC-04}}
\newpage\fragdef{MEM-24}% ============================================
% MEM-24
% Gerado de molambudos.md — build_miolo.py
% ============================================
\fragtitle{mem}{M}{MEM}{24}{A Enfermeira que Chorava}
\markboth{MEM-24 --- A Enfermeira que Chorava}{}

\noindent \textit{Pequenas bondades}

\noindent \textbf{1964 --- Pequenas bondades}

\vspace{0.5em}\noindent\rule{\textwidth}{0.5pt}\vspace{0.5em}

\noindent De vez em quando, uma enfermeira nova aparecia.

\noindent Vinham do interior, moças de famílias pobres que precisavam de trabalho. Entravam no Colônia sem saber o que era. Saíam em uma semana --- ou ficavam e endureciam.

\noindent Uma delas, em 1964, era diferente.

\noindent Ela se chamava Isabel. Tinha 19 anos. Cabelo preto, olhos grandes, mãos que tremiam quando entregavam a comida.

\noindent Na primeira semana, ela passou na frente da minha cela e viu o diário.

\noindent \textit{\textquotedblleft{}O que o senhor escreve?\textquotedblright{}}

\noindent \textit{\textquotedblleft{}A verdade.\textquotedblright{}}

\noindent Ela não entendeu. Mas continuou passando.

\noindent Toda tarde, ela trazia a comida. Enquanto os outros enfermeiros jogavam o prato pelo buraco da porta, ela abria a porta e entregava na mão.

\noindent \textit{\textquotedblleft{}O senhor não é perigoso\textquotedblright{}}, ela disse um dia.

\noindent \textit{\textquotedblleft{}Como sabe?\textquotedblright{}}

\noindent \textit{\textquotedblleft{}Pessoas perigosas não escrevem. Pessoas perigosas não pensam. O senhor escreve e pensa. O senhor é só triste.\textquotedblright{}}

\noindent Eu não sabia se aquilo era um diagnóstico ou uma absolvição.

\noindent \textit{\textquotedblleft{}Triste?\textquotedblright{}}

\noindent Ela abaixou os olhos.

\noindent \textit{\textquotedblleft{}O senhor tem o olho mais triste que eu já vi.\textquotedblright{}}

\noindent \textit{\textquotedblleft{}É amarelo.\textquotedblright{}}

\noindent Ela sacudiu a cabeça.

\noindent \textit{\textquotedblleft{}Amarelo é a cor. Triste é o que está por trás.\textquotedblright{}}

\noindent Ela saiu.

\noindent Naquela noite, a criatura falou:

\noindent \textit{\textquotedblleft{}Ela gosta de você.\textquotedblright{}}

\noindent \textit{\textquotedblleft{}Cala a boca.\textquotedblright{}}

\noindent \textit{\textquotedblleft{}Ela quer te salvar.\textquotedblright{}}

\noindent \textit{\textquotedblleft{}Ninguém salva ninguém aqui.\textquotedblright{}}

\noindent \textit{\textquotedblleft{}Ela quer tentar.\textquotedblright{}}

\noindent \textit{\textquotedblleft{}...Ela vai embora.\textquotedblright{}}

\noindent \textit{\textquotedblleft{}Como sabe?\textquotedblright{}}

\noindent \textit{\textquotedblleft{}Todo mundo vai. Todo mundo que entra aqui, sai. Ou de pé. Ou de caixão.\textquotedblright{}}

\vspace{0.5em}\noindent\rule{\textwidth}{0.5pt}\vspace{0.5em}

\noindent Isabel ficou três meses.

\noindent Um dia, ela não veio mais.

\noindent O enfermeiro que trouxe a comida disse: \textit{\textquotedblleft{}Pediu demissão. Disse que não aguentava ver.\textquotedblright{}}

\noindent Perguntei: \textit{\textquotedblleft{}Ver o quê?\textquotedblright{}}

\noindent Ele me olhou.

\noindent \textit{\textquotedblleft{}Ver vocês.\textquotedblright{}}

\noindent Depois que Isabel foi embora, eu fechei o diário por uma semana.

\noindent Não por falta de palavras.
Por falta de esperança.

\noindent Se até quem queria ajudar ia embora, então não adiantava nada.

\noindent O Colônia não era um hospital.
Era um buraco.
E tudo que caía no buraco --- esquecia.

\vspace{0.5em}
\noindent\textit{\textbf{↪ Rotas:} \rota{MEM-25} \rotasep \rota{MEM-12}}
\newpage\fragdef{MEM-25}% ============================================
% MEM-25
% Gerado de molambudos.md — build_miolo.py
% ============================================
\fragtitle{mem}{M}{MEM}{25}{A Morte no Espelho}
\markboth{MEM-25 --- A Morte no Espelho}{}

\noindent \textit{A primeira vez que me vi velho}

\noindent \textbf{1970 --- A primeira vez que me vi velho}

\vspace{0.5em}\noindent\rule{\textwidth}{0.5pt}\vspace{0.5em}

\noindent Teve um espelho na ala dos silenciados.

\noindent Não sei por quê. Alguém pendurou na parede do corredor e esqueceu. Ninguém usava. Pacientes não olham no espelho. Pacientes não querem se ver.

\noindent Mas um dia, passei na frente e olhei sem querer.

\noindent O homem no espelho era velho.

\noindent Cabelos brancos. Pele enrugada. Olhos fundos. O direito amarelo, o esquerdo apagado.

\noindent A boca torta. A barba por fazer. Os dentes podres.

\noindent \textit{\textquotedblleft{}Quem é?\textquotedblright{}}, pensei.

\noindent \textit{\textquotedblleft{}É você\textquotedblright{}}, disse a criatura.

\noindent \textit{\textquotedblleft{}Não. Eu tenho 8 anos. Eu estou no sertão. Eu estou andando com minha mãe.\textquotedblright{}}

\noindent \textit{\textquotedblleft{}Você tem 63 anos. Você está no Colônia. Você está sozinho.\textquotedblright{}}

\noindent Eu toquei o rosto no espelho.

\noindent A pele enrugada tocou a mão.

\noindent A criatura disse:

\noindent \textit{\textquotedblleft{}O tempo passou, Joaquim. Você não viu. Mas passou.\textquotedblright{}}

\noindent \textit{\textquotedblleft{}Quanto tempo?\textquotedblright{}}

\noindent \textit{\textquotedblleft{}63 anos.\textquotedblright{}}

\noindent \textit{\textquotedblleft{}...Eu passei 63 anos aqui?\textquotedblright{}}

\noindent \textit{\textquotedblleft{}Não. Você tem 63 anos. Está aqui há 53.\textquotedblright{}}

\noindent \textit{\textquotedblleft{}...E o que eu fiz?\textquotedblright{}}

\noindent \textit{\textquotedblleft{}Você sobreviveu.\textquotedblright{}}

\noindent \textit{\textquotedblleft{}Só isso?\textquotedblright{}}

\noindent \textit{\textquotedblleft{}Só isso.\textquotedblright{}}

\noindent Eu sentei no chão do corredor, na frente do espelho.

\noindent Fiquei olhando para o velho que era eu.

\noindent Eu não lembrava de ter ficado velho.

\noindent O tempo tinha passado dentro de mim sem que eu percebesse.

\noindent Os dias eram iguais. As noites eram iguais. As estações passavam, mas o Colônia era sempre a mesma temperatura: frio.

\noindent Eu não envelheci.
Eu \textit{apodreci}.

\noindent E o pior: ninguém viu.

\noindent Ninguém reparou.

\noindent Porque ninguém olhava para mim.

\noindent Eu era o paciente 1.260.
Eu era um número.
Números não envelhecem.
Números só expiram.

\vspace{0.5em}
\noindent\textit{\textbf{↪ Rotas:} \rota{MEM-26} \rotasep \rota{MEM-12}}

%% ============================================
%% PARTE 3 — DIÁRIO DE OLIVEIRA + LAUDOS (1979)
%% ============================================
\ifmolambudoskdpprint
  \partopener{\shortstack{Diário de Oliveira\\e Laudos (1979)}}{Diário de Oliveira e Laudos}{1.png}
\else
  \partopener{Diário de Oliveira e Laudos (1979)}{Diário de Oliveira e Laudos}{1.png}
\fi

\vspace{0.5em}
\noindent{\footnotesize\textit{Nota do Arquivista: a partir deste ponto, a voz principal deixa de ser a de Joaquim e passa a ser a de quem o estudou. Os documentos a seguir foram escritos por médicos, peritos, arquivistas e pelo coveiro que abriu a cova. A linguagem muda. O olhar muda. Mas o diário é o mesmo.}}

\newpage\fragdef{DOC-01}% ============================================
% DOC-01
% Gerado de molambudos.md — build_miolo.py
% ============================================
\fragtitle{doc}{D}{DOC}{01}{A Exumação}
\markboth{DOC-01 --- A Exumação}{}

\noindent \textbf{Notas do coveiro Genésio Barbosa}
\textit{Transcrição literal do depoimento prestado em agosto de 1979 à direção do Hospital Colônia. Documento arquivado sob o código C-1260-79.}

\vspace{0.5em}\noindent\rule{\textwidth}{0.5pt}\vspace{0.5em}

\noindent Era a terceira cova do dia. A gente sabe quais são as antigas porque a terra assentou. Covas recentes, a terra fica fofa, dá pra cavar com as mãos quase. As antigas, não. A terra endurece, forma uma crosta por cima, e quando você enfia a pá parece que vai bater em pedra.

\noindent A do 1.260 era recente. Pelo mapa, 1979. Um dos últimos enterrados ali.

\noindent Eu cavei. Cavei uns quarenta centímetros e a pá bateu em madeira. Cavei mais. O caixão --- era um caixão, mas não era. Sabe quando a madeira apodrece e vira terra? Não é que desaparece. Ela vira outra coisa. Uma pasta escura que ainda tem a forma do que foi.

\noindent Eu abri com a pá. Não por desrespeito. É que não dava pra puxar a tampa --- ela já era parte do chão.

\noindent O que eu vi dentro não era um esqueleto.

\noindent Era um corpo.

\noindent Não, não estou maluco. Era um corpo. A pele era escura, quase preta, mas estava \textit{inteira}. Como se tivesse sido enterrado na semana passada. O senhor pode perguntar pro Seu Nilo, que tava comigo. Ele viu também. A pele do rosto ainda cobria os ossos. A boca estava aberta. Os olhos ---

\noindent Os olhos não estavam lá. As órbitas estavam vazias. Mas dentro delas, no fundo, havia um brilho. Como se alguém tivesse acendido uma vela lá dentro e apagado. Não, não era vela. Era um reflexo. Amarelo. Muito fraco.

\noindent E no peito do corpo --- amarrado com uma tira de couro podre, apertada contra o osso --- tinha um caderno. O couro atravessava as costelas, por baixo do paletó, como se o homem tivesse sido enterrado com o livro colado ao peito. Eu precisei cortar a tira com a faca. O couro esfarelou na minha mão.

\noindent Eu sei que parece mentira. Um caderno enterrado com um paciente, amarrado ao corpo, sessenta e dois dias depois, com as páginas ainda \textit{úmidas}? Eu não acreditaria. Mas o senhor está segurando o caderno. O senhor viu as páginas. A tinta.

\noindent O senhor sabe que eu não estou mentindo.

\noindent Eu peguei o caderno. Quando toquei nele, senti um calor. O calor não vinha do sol. Era úmido, como o de uma ferida coberta. Como se o caderno estivesse vivo ali dentro. Crescendo.

\noindent Abri a primeira página.

\noindent A letra era de criança. Tortinha, subindo e descendo, como quem aprendeu a escrever sozinho. Uma assinatura no topo:

\noindent \textit{Joaquim Antônio Correia --- Paciente 1.260}

\noindent E uma data.

\noindent \textit{1917 --- o ano que eu entrei aqui}

\noindent Mas o diário continuava depois. Página após página. O último registro era de 1979.

\noindent \textit{Três dias antes de eu desenterrar ele.}

\noindent O senhor entende? O homem estava morto há 62 dias. E o diário continuou sendo escrito. Na cova. No escuro. Com os dedos de um esqueleto.

\noindent Eu não dormi naquela noite.

\noindent Não dormi até hoje.

\noindent O senhor pode ficar com o diário. Eu não quero mais ele perto de mim.

\noindent --- \textit{Genésio Barbosa, 48 anos, coveiro do Hospital Colônia de Barbacena há 22 anos.}

\vspace{0.5em}
\noindent\textit{\textbf{↪ Rotas:} \rota{MEM-04} \rotasep \rota{LUC-02}}
\newpage\fragdef{DOC-03}% ============================================
% DOC-03 — Laudo de Necropsia (edição revista R1075)
% ============================================
\fragtitle{doc}{D}{DOC}{03}{Laudo de Necropsia}
\markboth{DOC-03 --- Laudo de Necropsia}{}

\noindent \textit{Instituto Médico Legal de Barbacena --- 14 de junho de 1979 --- Dr. Álvaro Torres}

\noindent \textbf{Paciente:} Joaquim Antônio Correia \quad \textbf{Registro:} 1.260 \quad \textbf{Idade estimada:} 72 anos \quad \textbf{Óbito:} 12 de junho de 1979 \quad \textbf{Necropsia:} 14 de junho de 1979 \quad \textbf{Legista:} Dr. Álvaro Torres --- CRM-MG 3.117.

\subsection*{1. Exame externo}

\noindent \textbf{1.1 Corpo.} Peso 42 kg; estatura 1,62 m; índice de massa corporal 16,0. Caquexia acentuada. Panículo adiposo ausente; saliências ósseas visíveis sob a pele. Pele ressecada, com pregueamento persistente. Vinte e um parasitas \textit{Ascaris lumbricoides} foram contados no intestino; a carga parasitária competia com o corpo pelos poucos nutrientes disponíveis.

\noindent \textbf{1.2 Olho direito.} A esclera direita apresenta coloração amarelo-ouro intensa, homogênea, da córnea à conjuntiva bulbar. A esclera esquerda é branca, com vascularização normal. A coloração não é compatível com icterícia hepática clássica, que seria bilateral. À palpação, o globo direito é ligeiramente mais firme. Não há trauma recente, prótese ou corpo estranho. A mancha não desbota à compressão.

\noindent \textit{Anotação marginal, caligrafia diferente:} ``Depois da exumação, esta hipótese caiu por terra. O olho brilhava no escuro. Não há explicação médica para o fenômeno. --- Lúcia M.''

\noindent \textbf{1.3 Eletroconvulsoterapia.} Marcas ovais de queimadura elétrica nas regiões temporais, bilaterais, 3 cm × 2 cm, com alopecia permanente. Fratura por compressão de T5, consolidada; fratura antiga do colo do fêmur esquerdo, com encurtamento de 2 cm; fratura consolidada do ramo mandibular direito. As lesões são compatíveis com 37 sessões de eletrochoque sem anestesia, sem gel condutor e sem proteção bucal, entre 1937 e 1952.

\noindent \textbf{1.4 Lobotomia transorbital.} Cicatrizes puntiformes bilaterais nos cantos internos dos olhos, 2 mm. Na lâmina cribriforme, dois orifícios de 4 mm, simétricos, a 2,5 cm da linha média, com bordas lisas e cicatrizadas. Os orifícios estão pérvios. O prontuário não contém autorização para o procedimento.

\noindent \textbf{1.5 Contenção e automutilação.} Faixas de hiperpigmentação circumferencial nos punhos e tornozelos, compatíveis com contenção prolongada. Contratura dos dedos em semiflexão. Escoriações lineares no tórax e no abdome. Cicatrizes de arranhões profundos nos antebraços, com terra impregnada na derme. Amputação parcial da falange distal do dedo mínimo esquerdo, sem registro de trauma.

\noindent \textbf{1.6 Queimaduras.} Manchas de hipopigmentação no tórax e no abdome, 2 a 8 cm, compatíveis com escaldadura por imersão nos chamados ``banhos terapêuticos''.

\subsection*{2. Exame interno}

\noindent \textbf{2.1 Causa mortis.} Parada cardiorrespiratória secundária a insuficiência cardíaca congestiva, caquexia e broncoaspiração. Coração de 220 g, com atrofia marrom do miocárdio. Estômago vazio; mucosa atrófica. Fígado de 980 g, com esteatose moderada. Não há cirrose.

\noindent \textbf{2.2 Encéfalo.} Peso 1.120 g. Hemisférios simétricos; meninges transparentes. Não há gliose, cavitação ou alteração cortical compatível com lobotomia. As áreas orbitofrontais e pré-frontais foram dissecadas. A via de acesso dos leucótomos existe; o parênquima, não foi atingido de forma detectável.

\noindent \textit{Nota do legista:} ``A ausência de dano cerebral não é um alívio. É o achado mais perturbador. Joaquim não era violento por lesão cerebral. A violência era dele. Ou era de outra coisa. Os procedimentos não deixaram marca no cérebro; deixaram no corpo inteiro. O cérebro calado é a parte mais gritante do laudo.''

\noindent \textbf{2.3 Coluna.} Fraturas por compressão em T3, T5 e T8, estáveis, sem déficit neurológico. Tríade compatível com convulsões tônico-clônicas repetidas.

\subsection*{3. Síntese}

\begin{itemize}
\item ECT sem anestesia: queimaduras temporais, fraturas vertebrais, fratura mandibular e femoral, perda dentária.
\item Lobotomia transorbital: perfuração bilateral da lâmina cribriforme, sem lesão cortical detectável.
\item Contenção prolongada: faixas em punhos e tornozelos, contratura digital, fibrose em coxas.
\item Banhos terapêuticos: cicatrizes térmicas no tórax e no abdome.
\end{itemize}

\noindent \textbf{Causa mortis:} parada cardiorrespiratória. \textbf{Causa contribuinte:} caquexia por desnutrição crônica. \textbf{Comorbidade:} alta carga parasitária. Não foram encontrados marcadores orgânicos que justifiquem 62 anos de institucionalização. A lesão mais significativa do paciente não é anatômica.

\vspace{0.5em}
\noindent Dr. Álvaro Torres --- Médico Legista --- IML Barbacena --- CRM-MG 3.117

\vspace{0.8em}
\noindent \textit{Adendo manuscrito, aparentemente posterior:} ``Este laudo foi redigido antes da exumação. Os achados da seção seguinte foram acrescentados após a abertura da cova e a recuperação do diário. O signatário declara que não há explicação médica para os fenômenos descritos no item 4, limitando-se a relatar o que foi apresentado pelo coveiro e pelo material recuperado. O diário permanece sob custódia do IML, anexado a este processo. --- A.T.''

\vspace{0.5em}
\noindent \textit{Nota adicional, caligrafia de L.M., protocolo 2026:} ``Fragmentos ósseos de Strigiformes (coruja-buraqueira) foram recuperados do mesmo estrato da cova. A presença de restos de coruja no sepultamento de Joaquim foi registrada pelo coveiro como \textit{a ave que não saiu de perto}. Não foi encontrada explicação tafonômica para a deposição. Ossos de coruja não são achado comum em sepultamentos humanos no Colônia.''

\noindent\textit{\textbf{↪ Rotas:} \rota{MEM-04} \rotasep \rota{DOC-02} \rotasep \rota{LUC-03} \rotasep \rota{CONT-03}}

\newpage\fragdef{DOC-04}% ============================================
% DOC-04
% Gerado de molambudos.md — build_miolo.py
% ============================================
\fragtitle{doc}{D}{DOC}{04}{Relatório de Sessão I}
\markboth{DOC-04 --- Relatório de Sessão I}{}

\noindent \textit{Dr. Heitor Oliveira --- 14 de março de 1979}

\noindent \textbf{Dr. Heitor Oliveira --- 14 de março de 1979}

\vspace{0.5em}\noindent\rule{\textwidth}{0.5pt}\vspace{0.5em}

\noindent \textbf{Paciente:} Joaquim Antônio Correia --- Paciente 1.260
\textbf{Tipo:} Sessão de avaliação psiquiátrica inicial
\textbf{Duração:} 45 minutos

\vspace{0.5em}\noindent\rule{\textwidth}{0.5pt}\vspace{0.5em}

\noindent \textbf{Observações do Paciente:}

\noindent Joaquim é um homem magro, de estatura baixa, pele enrugada e olhar fixo. O olho direito difere significativamente do esquerdo --- apresenta coloração amarelo-ouvo incomum. Quando questionado, afirmou que \textquotedblleft{}um bicho\textquotedblright{} mudou a cor. Recusou-se a elaborar.

\noindent O paciente é articulado e inteligente --- muito mais do que o prontuário sugere. Demonstra capacidade de meta-cognição e análise psicológica (ex: classificou médicos pelo cheiro, distinguiu entre médicos novos e experientes). Não apresenta os sinais clássicos de deterioração cognitiva esperados para alguém com 62 anos de institucionalização.

\noindent Seu discurso é lúcido, mas monofocado no tema da \textquotedblleft{}fome\textquotedblright{} e da \textquotedblleft{}criatura\textquotedblright{}. Quando questionado sobre os homicídios, respondeu com naturalidade: \textquotedblleft{}a fome matou por mim\textquotedblright{\textit{. Não demonstra remorso, mas também não demonstra orgulho --- demonstra }aceitação}, como se os eventos fossem inevitáveis.

\noindent \textit{Nota pessoal (manuscrito, à margem):} \textquotedblleft{}A chave deste caso não está no delírio. Está no que veio antes --- no sertão, na seca, no pai morto, na vala. O delírio é só a fechadura. A história é a chave. --- H.O.\textquotedblright{}

\noindent \textbf{Impressão Diagnóstica (Preliminar):}

\noindent O paciente apresenta quadro compatível com Esquizofrenia Paranóide (CID-8 295.3), com delírio sistematizado envolvendo uma entidade que o habita. A estrutura do delírio é complexa e internamente coerente --- incomum para pacientes institucionalizados por décadas, que geralmente apresentam discurso fragmentado.

\noindent No entanto, há elementos que sugerem capacidade de insight preservado. Quando perguntei se ele achava que a criatura era real ou fruto da mente, ele respondeu: \textit{\textquotedblleft{}O senhor acha que a diferença importa?\textquotedblright{}}

\noindent \textbf{Plano Terapêutico:}

\noindent - Sessões semanais de avaliação
- Solicitar acesso ao diário do paciente (informação: o paciente mantém um diário há décadas; nunca permitiu que ninguém lesse)
- Manter medicação atual (Clorpromazina 100mg/dia)

\noindent \textbf{Nota Pessoal (não incluída no prontuário oficial):}

\noindent Há algo neste paciente que me incomoda. Não consigo identificar precisamente. Mas ao final da sessão, quando me levantei para sair, ele disse: \textit{\textquotedblleft{}O senhor vai sonhar comigo esta noite.\textquotedblright{}}

\noindent Eu sonhei.

\vspace{0.5em}\noindent\rule{\textwidth}{0.5pt}\vspace{0.5em}

\noindent \textit{Dr. Heitor Oliveira}
\textit{Psiquiatra --- Ala de Segurança Máxima}
\textit{Hospital Colônia de Barbacena}

\vspace{0.5em}
\noindent\textit{\textbf{↪ Rotas:} \rota{MEM-13} \rotasep \rota{DOC-05} \rotasep \rota{MEM-14} \rotasep \rota{LUC-01}}
\newpage\fragdef{DOC-05}% ============================================
% DOC-05
% Gerado de molambudos.md — build_miolo.py
% ============================================
\fragtitle{doc}{D}{DOC}{05}{Relatório de Sessão II}
\markboth{DOC-05 --- Relatório de Sessão II}{}

\noindent \textit{Dr. Heitor Oliveira --- 21 de março de 1979}

\noindent \textbf{Dr. Heitor Oliveira --- 21 de março de 1979}

\vspace{0.5em}\noindent\rule{\textwidth}{0.5pt}\vspace{0.5em}

\noindent \textbf{Paciente:} Joaquim Antônio Correia --- Paciente 1.260
\textbf{Tipo:} Sessão de avaliação psiquiátrica
\textbf{Duração:} 90 minutos (estendida)

\vspace{0.5em}\noindent\rule{\textwidth}{0.5pt}\vspace{0.5em}

\noindent \textbf{Observações do Paciente:}

\noindent O paciente apresentou-se diferente hoje. Mais animado. Quase eufórico. Sorria quando entrei na sala --- um sorriso assimétrico, puxado para o lado direito do rosto, onde o olho amarelo brilha mais intensamente do que na sessão anterior.

\noindent Ele perguntou: \textit{\textquotedblleft{}O senhor sonhou?\textquotedblright{}}

\noindent Respondi que sim, que tive um sonho incomum. Ele não pediu detalhes. Apenas acenou com a cabeça, como se já soubesse.

\noindent Nesta sessão, Joaquim falou continuamente por 75 minutos. Praticamente sem interrupção. Descreveu sua infância no Ceará, a seca de 1915, a caminhada de Senador Pompeu em direção a Fortaleza, a morte do pai e o confinamento no Campo do Alagadiço. Os relatos são vívidos --- há detalhes sensoriais (cheiros, temperaturas, texturas) que são incomuns mesmo em pacientes com memória preservada.

\noindent O tom do relato mudou quando descreveu o \textquotedblleft{}encontro\textquotedblright{}.

\noindent Segundo Joaquim, ele encontrou uma criatura em uma vala de corpos atrás do campo de concentração. A descrição continha elementos psicóticos (olhos amarelos, espinhos de osso, costelas externas), mas a \textit{consistência interna} é notável. Ele descreveu o encontro com a mesma convicção com que descreveu a morte do pai.

\noindent \textbf{Registro textual:}

\noindent \textit{\textquotedblleft{}Ela disse: 'Você também tem fome.' Eu disse que sim. Ela disse: 'Eu como por você. Eu como por todos que não podem comer.' E eu sentei. Sentei na beira da vala e olhei ela comer. Não senti medo. Senti fome. Uma fome que nunca mais passou.\textquotedblright{}}

\noindent \textbf{Análise:}

\noindent - A \textquotedblleft{}criatura\textquotedblright{} parece ser uma manifestação do trauma infantil de Joaquim, possivelmente uma personificação da fome e da morte que testemunhou na infância.
- O encontro na vala pode representar um momento dissociativo: Joaquim pode ter testemunhado algo insuportável (a vala de corpos, a morte de outros) e sua mente criou uma narrativa alternativa para processar o evento.
- A \textquotedblleft{}fome que nunca passa\textquotedblright{} é consistente com uma reação traumática crônica, onde o evento é revivido indefinidamente. A nosologia da época não tinha um nome único para isso; o laudo usa a linguagem disponível.

\noindent \textbf{Plano Terapêutico:}

\noindent - Solicitei acesso ao diário do paciente, que foi negado inicialmente, mas o paciente concordou em me entregar \textquotedblleft{}quando estiver pronto\textquotedblright{}.
- Manter sessões semanais.

\noindent \textbf{Nota Pessoal:}

\noindent Comecei a ter sonhos recorrentes. Não são pesadelos comuns --- são extremamente vívidos, com cheiros e sensações táteis. Acordo com o gosto de terra na boca. Não mencionei isso a Joaquim, mas ele sabe. Ele sabe.

\vspace{0.5em}\noindent\rule{\textwidth}{0.5pt}\vspace{0.5em}

\noindent \textit{Dr. Heitor Oliveira}
\textit{Psiquiatra --- Ala de Segurança Máxima}
\textit{Hospital Colônia de Barbacena}

\vspace{0.5em}
\noindent\textit{\textbf{↪ Rotas:} \rota{MEM-15} \rotasep \rota{DOC-07} \rotasep \rota{LUC-03}}
\newpage\fragdef{DOC-06}% ============================================
% DOC-06
% Gerado de molambudos.md — build_miolo.py
% ============================================
\fragtitle{doc}{D}{DOC}{06}{Relatório de Sessão III}
\markboth{DOC-06 --- Relatório de Sessão III}{}

\noindent \textit{Dr. Heitor Oliveira --- 2 de abril de 1979}

\noindent \textbf{Dr. Heitor Oliveira --- 2 de abril de 1979}
\textit{Último relatório. Sessão interrompida.}

\vspace{0.5em}\noindent\rule{\textwidth}{0.5pt}\vspace{0.5em}

\noindent \textbf{Paciente:} Joaquim Antônio Correia --- Paciente 1.260
\textbf{Tipo:} Sessão de avaliação psiquiátrica
\textbf{Duração:} 20 minutos (sessão interrompida)

\vspace{0.5em}\noindent\rule{\textwidth}{0.5pt}\vspace{0.5em}

\noindent \textbf{OBSERVAÇÃO INICIAL:}

\noindent Esta sessão foi interrompida por um incidente. O paciente apresentou alteração súbita de comportamento, e a sessão foi encerrada por segurança.

\noindent Relato a seguir:

\vspace{0.5em}\noindent\rule{\textwidth}{0.5pt}\vspace{0.5em}

\noindent \textbf{Relato:}

\noindent A sessão começou normalmente. Joaquim estava calmo. Entregou-me o diário --- algo que vinha recusando fazer havia duas semanas. Ele o colocou sobre a mesa e disse: \textit{\textquotedblleft{}Leia.\textquotedblright{}}

\noindent Perguntei se tinha certeza. Ele disse que sim.

\noindent Abri o diário. As primeiras páginas estavam escritas em caligrafia infantil, com erros de ortografia e frases curtas. As páginas descreviam a rotina no Colônia. Nada extraordinário.

\noindent Mas as últimas páginas ---

\noindent As últimas páginas estavam escritas em letra diferente. Uma letra adulta, firme, com uma caligrafia quase cursiva. Elas descreviam a \textquotedblleft{}criatura\textquotedblright{}, a \textquotedblleft{}vala\textquotedblright{}, o \textquotedblleft{}ciclo\textquotedblright{\textit{. Descreviam }este consultório\textit{. Descreviam }esta sessão}.

\noindent Havia uma passagem:

\noindent \textit{\textquotedblleft{}O médico vai ler este diário. Ele vai sentir o cheiro. Ele vai ter medo. Ele vai negar. Mas o ciclo vai começar.\textquotedblright{}}

\noindent Quando levantei os olhos, Joaquim não estava mais na cadeira.

\noindent Ele estava no canto da sala, de cócoras, com os braços enrolados em volta das pernas. Os olhos --- o olho direito --- brilhava.

\noindent Chamei-o pelo nome.

\noindent A voz que respondeu não era a dele.

\vspace{0.5em}\noindent\rule{\textwidth}{0.5pt}\vspace{0.5em}

\noindent Esta é a parte difícil de relatar.

\noindent A voz era mais grave. Mais lenta. Como se viesse de um corpo maior. Ela disse:

\noindent \textit{\textquotedblleft{}Ele não está mais aqui.\textquotedblright{}}

\noindent Perguntei quem estava falando.

\noindent \textit{\textquotedblleft{}Você sabe quem.\textquotedblright{}}

\noindent A criatura. O paciente estava em crise psicótica aguda. A entidade delirante estava se manifestando.

\noindent Tentei manter a calma.

\noindent \textit{\textquotedblleft{}O que você quer?\textquotedblright{}}

\noindent \textit{\textquotedblleft{}Nada. Só quero que você veja.\textquotedblright{}}

\noindent \textit{\textquotedblleft{}Ver o quê?\textquotedblright{}}

\noindent \textit{\textquotedblleft{}O que todo médico recusa ver.\textquotedblright{}}

\noindent Neste momento, o paciente começou a tremer. O corpo inteiro sacudiu. A boca se abriu. Os olhos reviraram. Achei que fosse uma convulsão.

\noindent Não era.

\noindent Ele parou de tremer. Ergueu a cabeça. O olho direito estava \textit{aceso} --- não é metáfora. Havia luminescência.

\noindent Ele sorriu.

\noindent \textit{\textquotedblleft{}Acha que é delírio?\textquotedblright{}}

\noindent Não respondi.

\noindent \textit{\textquotedblleft{}Acha que é trauma? Dissociação? Projeção?\textquotedblright{}}

\noindent Ainda não respondi.

\noindent \textit{\textquotedblleft{}Então me explica uma coisa, doutor.\textquotedblright{}}

\noindent Ele se levantou. Veio em minha direção. Não tocou em mim. Parou a alguns centímetros.

\noindent \textit{\textquotedblleft{}Por que você sonhou com a vala antes de ler o diário?\textquotedblright{}}

\noindent Eu não tinha contado a ele. Não tinha contado a ninguém.

\noindent \textit{\textquotedblleft{}Por que o olho dele brilha no escuro?\textquotedblright{}}

\noindent \textit{\textquotedblleft{}Isso é...\textquotedblright{}}

\noindent \textit{\textquotedblleft{}O quê? Fisiologia? Patologia?\textquotedblright{}}

\noindent Ele riu. A risada não era humana.

\noindent \textit{\textquotedblleft{}Vocês têm uma palavra para tudo. Esquizofrenia. Trauma. Delírio. Dissociação. Vocês nomeiam o que não entendem e acham que entenderam.\textquotedblright{}}

\noindent Ele se inclinou. O cheiro --- o cheiro adocicado --- era forte.

\noindent \textit{\textquotedblleft{}O nome não é a coisa, doutor. Uma coisa é não saber o nome. Outra coisa é achar que o nome é a explicação.\textquotedblright{}}

\noindent Neste momento, apaguei a luz da sala e disse que a sessão estava encerrada. O paciente não reagiu. Ele voltou para a cadeira e sentou como se nada tivesse acontecido.

\noindent Quando perguntei se ele se lembrava do que tinha acontecido, ele disse:

\noindent \textit{\textquotedblleft{}Lembro. Mas não fui eu.\textquotedblright{}}

\vspace{0.5em}\noindent\rule{\textwidth}{0.5pt}\vspace{0.5em}

\noindent \textbf{Plano Terapêutico:}

\noindent - Paciente apresenta exacerbação psicótica. Aumentar Clorpromazina para 200mg/dia.
- Solicitar avaliação neurológica (descartar epilepsia do lobo temporal).
- Sessões suspensas temporariamente.

\noindent \textbf{Nota Pessoal:}

\noindent Não vou dormir esta noite.

\vspace{0.5em}\noindent\rule{\textwidth}{0.5pt}\vspace{0.5em}

\noindent \textit{Dr. Heitor Oliveira}
\textit{Psiquiatra --- Ala de Segurança Máxima}
\textit{Hospital Colônia de Barbacena}

\noindent \textit{P.S. --- Esta noite, quando cheguei em casa, o diário estava na minha mochila. Eu não coloquei ele lá. Não sei como foi parar ali.}

\vspace{0.5em}\noindent\rule{\textwidth}{0.5pt}\vspace{0.5em}

\subsection*{ADENDO --- Redigido em 3 de abril de 1979 (1 dia após a sessão)}

\noindent Não dormi.

\noindent Tentei. Deitei. Fechei os olhos. Mas o cheiro estava no travesseiro. O cheiro adocicado. Levantei e lavei o rosto. O cheiro não saiu. Estava na minha pele.

\noindent Abri o diário de novo. Li a passagem sobre a vala. Quando levantei os olhos, eram 3:14 da manhã. O mesmo horário que Joaquim diz acordar.

\noindent Olhei para o canto do quarto. Não havia nada. Mas eu \textit{quase} vi. Uma sombra. Um movimento. Um brilho.

\noindent Fechei o diário. Coloquei dentro da gaveta. Tranquei a gaveta.

\noindent Deitei de novo.

\noindent O cheiro continuava lá.

\vspace{0.5em}\noindent\rule{\textwidth}{0.5pt}\vspace{0.5em}

\subsection*{ADENDO 2 --- 5 de abril de 1979 (72h após a sessão)}

\noindent Sonhei com a vala.

\noindent Não foi um sonho. Foi uma \textit{visita}. Eu estava deitado na pilha de corpos. Os olhos deles estavam abertos. Todos me olhando. A criatura estava em cima de mim. O rosto perto do meu. O hálito doce.

\noindent Ela disse: \textit{\textquotedblleft{}Você leu. Agora você é meu.\textquotedblright{}}

\noindent Acordei gritando.

\noindent Minha esposa acendeu a luz. Perguntou o que foi. Eu disse que não foi nada. Ela disse que eu estava branco. Eu disse que estava bem.

\noindent Mentira.

\noindent Não estou bem.

\noindent O diário está na minha cabeça. As palavras não saem. Elas se repetem. O tempo todo. Mesmo de olhos abertos.

\noindent \textit{\textquotedblleft{}Você leu. Agora você é meu.\textquotedblright{}}

\noindent Preciso dormir. Preciso parar de pensar. Preciso ---

\vspace{0.5em}\noindent\rule{\textwidth}{0.5pt}\vspace{0.5em}

\subsection*{ADENDO 3 --- 8 de abril de 1979}

\noindent O paciente 1.260 não fala mais.

\noindent Eu vou à cela dele todos os dias. Ele está sentado na cama, com o diário no colo, olhando para a parede. Não me olha quando eu entro. Não responde quando eu chamo.

\noindent Mas o olho dele --- o olho amarelo --- me segue.

\noindent Não importa onde eu esteja na cela. O olho me acompanha. Como se o olho fosse um animal separado do corpo dele.

\noindent Perguntei para a enfermeira chefe: \textquotedblleft{}O paciente 1.260 está tomando a medicação?\textquotedblright{}

\noindent Ela disse que sim. Clorpromazina 100mg. Duas vezes ao dia.

\noindent Então por que o olho dele ainda brilha?

\vspace{0.5em}\noindent\rule{\textwidth}{0.5pt}\vspace{0.5em}

\subsection*{ADENDO 4 --- 10 de abril de 1979}

\noindent Sonhei de novo.

\noindent Dessa vez eu não estava na vala. Eu estava na minha sala de consultas. O paciente 1.260 estava sentado na cadeira. Mas o rosto dele --- quando ele levantou a cabeça --- o rosto era o meu.

\noindent Acordei com o meu próprio grito.

\noindent Minha esposa dorme no sofá agora. Ela disse que eu falo dormindo. Perguntei o que eu falo. Ela disse que não quer repetir.

\noindent Perguntei se eu disse \textquotedblleft{}olho amarelo\textquotedblright{}.

\noindent Ela não respondeu.

\vspace{0.5em}\noindent\rule{\textwidth}{0.5pt}\vspace{0.5em}

\subsection*{ADENDO 5 --- 12 de abril de 1979}

\noindent Eu sei que o olho não é um delírio.

\noindent O olho brilha. Eu vi. Esta noite, fui até o hospital. Entrei na ala dos silenciados. Fiquei na porta da cela 1.260. O paciente estava dormindo.

\noindent O olho direito estava aberto.

\noindent Brilhando no escuro.

\noindent Não é possível. Não é fisiologia. Não é patologia. É outra coisa.

\noindent O que é?

\noindent O que é?

\noindent O que é?

\noindent \textit{preciso saber o que é preciso saber o que é preciso saber o que é preciso saber o que é}

\vspace{0.5em}\noindent\rule{\textwidth}{0.5pt}\vspace{0.5em}

\subsection*{ADENDO 6 --- 13 de abril de 1979}

\noindent Tranquei o diário no armário do hospital.
Joguei a chave no lixo.
Voltei para casa.
O diário estava na mesa da cozinha.

\noindent Ele não sai de perto de mim.

\noindent Ele \textit{quer} ser lido.

\noindent Ele \textit{quer} ser passado adiante.

\noindent Eu não vou mais escrever neste relatório.

\noindent Eu vou escrever no diário.

\noindent Ele me chamou.

\vspace{0.5em}\noindent\rule{\textwidth}{0.5pt}\vspace{0.5em}

\noindent \textit{Nota da direção: Dr. Heitor Oliveira não compareceu ao trabalho a partir de 14 de abril de 1979.}
\textit{Tentativas de contato telefônico fracassaram.}
\textit{Uma visita ao seu domicílio em 16 de abril constatou a residência vazia.}
\textit{O diário do paciente 1.260 não foi encontrado no hospital.}
\textit{O Dr. Oliveira permanece desaparecido.}

\vspace{0.5em}
\noindent\textit{\textbf{↪ Rotas:} \rota{MEM-16} \rotasep \rota{DOC-01}}
\newpage\fragdef{DOC-07}% ============================================
% DOC-07
% Gerado de molambudos.md — build_miolo.py
% ============================================
\fragtitle{doc}{D}{DOC}{07}{Diagnóstico Psiquiátrico (1979)}
\markboth{DOC-07 --- Diagnóstico Psiquiátrico (1979)}{}

\noindent \textit{DSM-II --- Dr. Heitor Oliveira}

\noindent \textbf{Paciente:} Joaquim Antônio Correia --- Registro 1.260
\textbf{Médico:} Dr. Heitor Oliveira, CRM-MG 4.892
\textbf{Data:} 14 de abril de 1979

\vspace{0.5em}\noindent\rule{\textwidth}{0.5pt}\vspace{0.5em}

\noindent \textbf{CID-8 / DSM-II}

\noindent \textbf{Diagnóstico principal:}
- \textbf{295.3 --- Esquizofrenia, Tipo Paranóide}
  - Justificativa: Delírios sistematizados de perseguição (acredita que uma \textquotedblleft{}criatura\textquotedblright{} o habita), alucinações visuais e táteis persistentes, discurso desconexo com neologismos (\textquotedblleft{}molambudo\textquotedblright{} como estado existencial), embotamento afetivo intercalado com explosões de agressividade.

\noindent \textbf{Diagnóstico secundário:}
- \textbf{301.7 --- Personalidade Psicopática}
  - Justificativa: Incapacidade de sentir remorso, conduta agressiva recorrente (três homicídios confirmados na juventude), manipulação de profissionais de saúde, ausência de vínculos afetivos.

\noindent \textbf{Condições médicas associadas:}
- Icterícia crônica (olho direito --- possível sequela de desnutrição severa na infância, ou hepatopatia não tratada).
- Desnutrição proteico-calórica compatível com a cronologia do prontuário: seca de 1915 e transferência em 1917.

\noindent \textbf{Estressores psicossociais:}
- Extrema.

\noindent \textbf{Funcionamento adaptativo no último ano:}
- 10/100 --- Perigo persistente para si e para terceiros.

\vspace{0.5em}\noindent\rule{\textwidth}{0.5pt}\vspace{0.5em}

\noindent \textbf{Observações Clínicas:}

\noindent O paciente foi admitido no Hospital Colônia em 1917, aos 10 anos de idade. Registros históricos indicam que chegou \textquotedblleft{}em estado de inanição, proveniente dos campos de concentração do Ceará, órfão de pai e mãe, com um olho amarelo e sem falar uma palavra que não fosse 'foooome'.\textquotedblright{}

\noindent Ao longo de 62 anos de institucionalização, o paciente apresentou:
- Episódios recorrentes de automutilação (arranhava a própria face até sangrar);
- Três homicídios confirmados (1922, 1937, 1948) --- todas as vítimas eram funcionários ou pacientes que tentaram se aproximar;
- Períodos de mutismo seletivo (recusava-se a falar com a equipe; chegou a ficar 7 anos sem dirigir a palavra a funcionários --- de 1949 a 1956);
- Redação compulsiva de um \textquotedblleft{}diário\textquotedblright{} que nunca permitiu que nenhum médico lesse integralmente;
- Fixação no tema da \textquotedblleft{}fome\textquotedblright{} e de uma \textquotedblleft{}criatura de olhos amarelos\textquotedblright{} que o acompanha.

\noindent O paciente demonstra inteligência preservada --- possivelmente acima da média --- mas inteiramente desviada para a construção de seu delírio. Sua capacidade de narrar sua própria história com coerência interna (dentro do sistema delirante) é notável.

\noindent \textbf{Prognóstico:} Reservado. Paciente irrecuperável. Manter em isolamento com sedação contínua.

\vspace{0.5em}\noindent\rule{\textwidth}{0.5pt}\vspace{0.5em}

\noindent \textit{Dr. Heitor Oliveira}
\textit{Psiquiatra Responsável --- Ala de Segurança Máxima}
\textit{Hospital Colônia de Barbacena}

\noindent \textit{P.S.: Dr. Heitor Oliveira desapareceu em 13 de junho de 1979, dois meses após assinar este laudo. O processo de desaparecimento nunca foi encerrado. Na ficha funcional do hospital, um colega anotou, à mão, na margem: \textquotedblleft{}Ele parou de comer três semanas antes. Dizia que não sentia mais fome --- e isso, disse ele, era o pior sintoma de todos.\textquotedblright{}}

\vspace{0.5em}
\noindent\textit{\textbf{↪ Rotas:} \rota{DOC-08} \rotasep \rota{DOC-04}}
\newpage\fragdef{MEM-13}% ============================================
% MEM-13
% Gerado de molambudos.md — build_miolo.py
% ============================================
\fragtitle{mem}{M}{MEM}{13}{O Médico Novo}
\markboth{MEM-13 --- O Médico Novo}{}

\noindent \textit{Primeira sessão com Dr. Heitor Oliveira}

\noindent \textbf{1979 --- Primeira sessão com Dr. Heitor Oliveira}

\vspace{0.5em}\noindent\rule{\textwidth}{0.5pt}\vspace{0.5em}

\noindent Eu sabia que ele era novo no jeito de andar.

\noindent Os médicos velhos andam devagar. Eles já viram de tudo. Não têm pressa. Os novos andam rápido. Eles acham que vão descobrir alguma coisa que ninguém descobriu.

\noindent O Dr. Oliveira andava rápido.

\noindent Ele entrou na sala de sessões --- uma sala com duas cadeiras, uma mesa, uma janela gradeada --- e sentou na frente de mim. Colocou um bloco na mesa. Uma caneta. Ajeitou o jaleco.

\noindent \textit{\textquotedblleft{}Bom dia, Joaquim.\textquotedblright{}}

\noindent \textit{\textquotedblleft{}Bom dia.\textquotedblright{}}

\noindent Ele me olhou. Eu olhei ele.

\noindent \textit{\textquotedblleft{}Sou o Dr. Heitor Oliveira. Sou o novo psiquiatra da ala.\textquotedblright{}}

\noindent \textit{\textquotedblleft{}Eu sei.\textquotedblright{}}

\noindent \textit{\textquotedblleft{}Sabe como?\textquotedblright{}}

\noindent \textit{\textquotedblleft{}O cheiro.\textquotedblright{}}

\noindent Ele parou.

\noindent \textit{\textquotedblleft{}Cheiro?\textquotedblright{}}

\noindent \textit{\textquotedblleft{}Médico novo cheira a sabão. Médico velho cheira a naftalina.\textquotedblright{}}

\noindent Ele riu. Era uma risada nervosa.

\noindent \textit{\textquotedblleft{}Você é observador.\textquotedblright{}}

\noindent \textit{\textquotedblleft{}Sessenta e dois anos aqui dentro. A gente aprende.\textquotedblright{}}

\noindent Ele anotou alguma coisa.

\noindent \textit{\textquotedblleft{}Você sabe por que está aqui, Joaquim?\textquotedblright{}}

\noindent \textit{\textquotedblleft{}Porque matei três pessoas.\textquotedblright{}}

\noindent \textit{\textquotedblleft{}Por que matou?\textquotedblright{}}

\noindent \textit{\textquotedblleft{}Porque elas estavam no lugar errado na hora errada.\textquotedblright{}}

\noindent \textit{\textquotedblleft{}E por que elas estavam no lugar errado?\textquotedblright{}}

\noindent \textit{\textquotedblleft{}Porque eu decidi que era a hora.\textquotedblright{}}

\noindent Ele parou de anotar.

\noindent \textit{\textquotedblleft{}Você planejou os assassinatos?\textquotedblright{}}

\noindent \textit{\textquotedblleft{}Não precisei. A oportunidade apareceu. Eu aproveitei.\textquotedblright{}}

\noindent \textit{\textquotedblleft{}Você sente remorso?\textquotedblright{}}

\noindent Eu olhei para ele. O olho amarelo.

\noindent \textit{\textquotedblleft{}O senhor já sentiu fome, doutor?\textquotedblright{}}

\noindent \textit{\textquotedblleft{}Já.\textquotedblright{}}

\noindent \textit{\textquotedblleft{}Fome de verdade? Fome de comer terra, de comer a própria mão, de olhar pro outro e pensar 'ele tem carne'?\textquotedblright{}}

\noindent Ele não respondeu.

\noindent \textit{\textquotedblleft{}Quando a gente passa fome desse jeito, não sobra espaço pra remorso. O remorso é um luxo de quem tem o estômago cheio.\textquotedblright{}}

\noindent \textit{\textquotedblleft{}Você está dizendo que matou por fome?\textquotedblright{}}

\noindent \textit{\textquotedblleft{}Estou dizendo que a fome matou por mim. Eu fui só a ferramenta.\textquotedblright{}}

\noindent Ele anotou. A caneta dele escrevia rápido.

\noindent \textit{\textquotedblleft{}E essa fome ainda está aqui?\textquotedblright{}}

\noindent \textit{\textquotedblleft{}Nunca foi embora.\textquotedblright{}}

\noindent \textit{\textquotedblleft{}O que você sente agora?\textquotedblright{}}

\noindent Eu fechei os olhos.

\noindent \textit{\textquotedblleft{}O senhor quer a verdade?\textquotedblright{}}

\noindent \textit{\textquotedblleft{}Quero.\textquotedblright{}}

\noindent \textit{\textquotedblleft{}Eu sinto o cheiro do senhor.\textquotedblright{}}

\noindent Ele parou de escrever.

\noindent \textit{\textquotedblleft{}Que cheiro?\textquotedblright{}}

\noindent \textit{\textquotedblleft{}O senhor comeu carne no jantar. Carne de boi. O senhor passou molho na carne. Tinha cebola. O senhor bebeu vinho depois. Vinho tinto. O senhor escovou os dentes, mas o cheiro ainda está nos poros da sua pele.\textquotedblright{}}

\noindent Ele engoliu seco.

\noindent \textit{\textquotedblleft{}E o que mais?\textquotedblright{}}

\noindent \textit{\textquotedblleft{}O senhor tem medo. O senhor está com medo agora. O senhor está sentado numa cadeira tentando parecer profissional, mas o senhor quer levantar e ir embora. O senhor ouviu histórias sobre mim. O senhor acha que eu sou um monstro.\textquotedblright{}}

\noindent \textit{\textquotedblleft{}Você é?\textquotedblright{}}

\noindent Eu abri os olhos.

\noindent \textit{\textquotedblleft{}Todo mundo é, doutor. A diferença é que eu parei de esconder.\textquotedblright{}}

\vspace{0.5em}
\noindent\textit{\textbf{↪ Rotas:} \rota{DOC-04} \rotasep \rota{MEM-12} \rotasep \rota{MEM-14}}
\newpage\fragdef{MEM-14}% ============================================
% MEM-14
% Gerado de molambudos.md — build_miolo.py
% ============================================
\fragtitle{mem}{M}{MEM}{14}{O Diário Emerge}
\markboth{MEM-14 --- O Diário Emerge}{}

\noindent \textit{Joaquim entrega os escritos}

\noindent \textbf{1979 --- Joaquim entrega os escritos}

\vspace{0.5em}\noindent\rule{\textwidth}{0.5pt}\vspace{0.5em}

\noindent Dr. Oliveira apareceu na cela numa manhã de segunda.

\noindent Eu estava sentado na cama, com o diário no colo, escrevendo. Quando ele abriu a porta, eu fechei o caderno.

\noindent \textit{\textquotedblleft{}Bom dia, Joaquim.\textquotedblright{}}

\noindent \textit{\textquotedblleft{}Bom dia.\textquotedblright{}}

\noindent Ele olhou para o diário.

\noindent \textit{\textquotedblleft{}O que é isso?\textquotedblright{}}

\noindent \textit{\textquotedblleft{}Nada.\textquotedblright{}}

\noindent \textit{\textquotedblleft{}Não parece nada. Parece um caderno.\textquotedblright{}}

\noindent \textit{\textquotedblleft{}É meu.\textquotedblright{}}

\noindent \textit{\textquotedblleft{}Eu sei.\textquotedblright{}}

\noindent Ele sentou no banco da cela --- algo que nenhum médico fazia. Os médicos ficavam de pé, perto da porta, prontos para sair. Ele sentou.

\noindent \textit{\textquotedblleft{}O que você escreve?\textquotedblright{}}

\noindent \textit{\textquotedblleft{}Coisas.\textquotedblright{}}

\noindent \textit{\textquotedblleft{}Que coisas?\textquotedblright{}}

\noindent \textit{\textquotedblleft{}O que eu vi.\textquotedblright{}}

\noindent Ele esperou.

\noindent \textit{\textquotedblleft{}O que você viu, Joaquim?\textquotedblright{}}

\noindent Eu olhei para ele. O olho amarelo. Ele não desviou.

\noindent \textit{\textquotedblleft{}O senhor não vai acreditar.\textquotedblright{}}

\noindent \textit{\textquotedblleft{}Tenta.\textquotedblright{}}

\noindent Eu pensei. Sessenta e dois anos de silêncio. Sessenta e dois anos guardando a história. Sessenta e dois anos esperando alguém que perguntasse.

\noindent \textit{\textquotedblleft{}O senhor tem medo de morrer?\textquotedblright{}}

\noindent Ele estranhou a pergunta.

\noindent \textit{\textquotedblleft{}Não... não penso muito nisso.\textquotedblright{}}

\noindent \textit{\textquotedblleft{}Pense.\textquotedblright{}}

\noindent \textit{\textquotedblleft{}Por quê?\textquotedblright{}}

\noindent \textit{\textquotedblleft{}Porque o senhor vai morrer. Todo mundo morre. Mas nem todo mundo morre de verdade.\textquotedblright{}}

\noindent \textit{\textquotedblleft{}O que isso significa?\textquotedblright{}}

\noindent Eu levantei o diário.

\noindent \textit{\textquotedblleft{}O senhor quer ler?\textquotedblright{}}

\noindent Ele hesitou.

\noindent \textit{\textquotedblleft{}Pode?\textquotedblright{}}

\noindent \textit{\textquotedblleft{}O senhor é o primeiro a perguntar. Isso já é mais do que todos os outros fizeram.\textquotedblright{}}

\noindent Eu estendi o diário.

\noindent Ele pegou.

\noindent O diário estava quente. Ele sentiu. Vi nos olhos dele --- o choque de tocar algo que parecia vivo.

\noindent \textit{\textquotedblleft{}Comece pelo começo\textquotedblright{}\textit{, eu disse. }\textquotedblleft{}1917. O ano que eu entrei aqui.\textquotedblright{}}

\noindent Ele abriu.

\noindent E começou a ler.

\noindent Quando ele levantou os olhos, duas horas tinham passado. Ele estava pálido. A mão tremia.

\noindent \textit{\textquotedblleft{}Joaquim...\textquotedblright{}}

\noindent \textit{\textquotedblleft{}Agora o senhor sabe.\textquotedblright{}}

\noindent \textit{\textquotedblleft{}Isso... é verdade?\textquotedblright{}}

\noindent \textit{\textquotedblleft{}O senhor sentiu o cheiro das páginas?\textquotedblright{}}

\noindent Ele não respondeu.

\noindent \textit{\textquotedblleft{}O senhor viu o olho?\textquotedblright{}}

\noindent Ele olhou para o meu olho direito. Estava brilhando.

\noindent \textit{\textquotedblleft{}O senhor leu. Agora faz parte.\textquotedblright{}}

\noindent \textit{\textquotedblleft{}Parte de quê?\textquotedblright{}}

\noindent \textit{\textquotedblleft{}Do ciclo.\textquotedblright{}}

\noindent Ele fechou o diário. Levantou. Andou até a porta.

\noindent \textit{\textquotedblleft{}Eu volto amanhã.\textquotedblright{}}

\noindent \textit{\textquotedblleft{}O senhor vai voltar\textquotedblright{}\textit{, eu disse. }\textquotedblleft{}A criatura não vai deixar o senhor ir embora.\textquotedblright{}}

\noindent Ele parou.

\noindent \textit{\textquotedblleft{}O que quer dizer com isso?\textquotedblright{}}

\noindent \textit{\textquotedblleft{}O senhor leu. Ela está no senhor agora. Leve alguns dias para crescer. Mas cresce.\textquotedblright{}}

\noindent Ele saiu sem responder.

\noindent Eu ouvi os passos dele no corredor. Rápidos. Desordenados.

\noindent O ciclo tinha começado.

\vspace{0.5em}
\noindent\textit{\textbf{↪ Rotas:} \rota{DOC-04} \rotasep \rota{DOC-05} \rotasep \rota{MEM-15}}
\newpage\fragdef{MEM-15}% ============================================
% MEM-15
% Gerado de molambudos.md — build_miolo.py
% ============================================
\fragtitle{mem}{M}{MEM}{15}{A Entidade Fala}
\markboth{MEM-15 --- A Entidade Fala}{}

\noindent \textit{A voz que não era de Joaquim}

\noindent \textbf{1979 --- Terceira sessão com Dr. Heitor Oliveira}

\vspace{0.5em}\noindent\rule{\textwidth}{0.5pt}\vspace{0.5em}

\noindent Dr. Oliveira chegou diferente naquele dia.

\noindent Ele estava pálido. Os olhos vermelhos. A barba por fazer. Ele sentou na cadeira e não abriu o bloco.

\noindent \textit{\textquotedblleft{}Joaquim, eu li o diário.\textquotedblright{}}

\noindent \textit{\textquotedblleft{}Eu sei.\textquotedblright{}}

\noindent \textit{\textquotedblleft{}As últimas páginas...\textquotedblright{}}

\noindent Ele parou. Passou a mão no rosto.

\noindent \textit{\textquotedblleft{}As últimas páginas foram escritas depois da sua morte.\textquotedblright{}}

\noindent \textit{\textquotedblleft{}Eu sei.\textquotedblright{}}

\noindent \textit{\textquotedblleft{}Como? Como você pode ter escrito depois de morto?\textquotedblright{}}

\noindent Eu não respondi. A criatura dentro de mim começou a acordar. Eu senti o calor subir do estômago.

\noindent \textit{\textquotedblleft{}Joaquim?\textquotedblright{}}

\noindent A voz que saiu não foi a minha.

\noindent \textit{\textquotedblleft{}Ele não escreveu.\textquotedblright{}}

\noindent O médico endireitou o corpo.

\noindent \textit{\textquotedblleft{}Quem... quem está falando?\textquotedblright{}}

\noindent \textit{\textquotedblleft{}Você já sabe quem.\textquotedblright{}}

\noindent Dr. Oliveira olhou para os lados. Olhou para a porta. Olhou para a janela. Estava tudo fechado.

\noindent \textit{\textquotedblleft{}Você está... dentro dele?\textquotedblright{}}

\noindent \textit{\textquotedblleft{}Dentro. Fora. Em volta.\textquotedblright{}}

\noindent \textit{\textquotedblleft{}O que você é?\textquotedblright{}}

\noindent A criatura riu. Minha boca se abriu. O som que saiu não era humano.

\noindent \textit{\textquotedblleft{}Eu sou o que sobrou. Eu sou o que foi esquecido. Eu sou a fome que não passa. Eu sou a vala.\textquotedblright{}}

\noindent \textit{\textquotedblleft{}A vala...\textquotedblright{}} repetiu Oliveira.

\noindent \textit{\textquotedblleft{}A vala onde os corpos foram jogados. A vala onde ninguém rezou. A vala onde o governo jogou os molambudos como lixo. Eu nasci ali. Eu cresci ali. Eu comi os mortos até eles me parirem.\textquotedblright{}}

\noindent \textit{\textquotedblleft{}Você é uma... entidade?\textquotedblright{}}

\noindent \textit{\textquotedblleft{}Isso. Uma palavra bonita. Mas não sou uma entidade. Sou um acúmulo. Sou o que acontece quando se joga muita gente no mesmo buraco. A terra não absorve tanta morte. Ela... fermenta. Vira uma coisa nova.\textquotedblright{}}

\noindent \textit{\textquotedblleft{}E Joaquim?\textquotedblright{}}

\noindent \textit{\textquotedblleft{}Joaquim é o meu primeiro filho.\textquotedblright{}}

\noindent \textit{\textquotedblleft{}Filho?\textquotedblright{}}

\noindent \textit{\textquotedblleft{}Ele veio até mim. Sozinho. Com fome. Ele sentou na beira da vala e me olhou. Eu olhei de volta. Ele não correu. Ele ficou. Ele é meu.\textquotedblright{}}

\noindent \textit{\textquotedblleft{}Ele matou pessoas.\textquotedblright{}}

\noindent \textit{\textquotedblleft{}Sim.\textquotedblright{}}

\noindent \textit{\textquotedblleft{}Por sua causa?\textquotedblright{}}

\noindent \textit{\textquotedblleft{}Não. Por causa da fome. A fome que a humanidade criou. Eu só uso a fome que já existe.\textquotedblright{}}

\noindent Dr. Oliveira ficou em silêncio.

\noindent \textit{\textquotedblleft{}O que você quer?\textquotedblright{}}

\noindent \textit{\textquotedblleft{}Nada.\textquotedblright{}}

\noindent \textit{\textquotedblleft{}Nada?\textquotedblright{}}

\noindent \textit{\textquotedblleft{}Eu não quero nada. Eu não preciso de nada. Eu só... existo. Existo entre as páginas do diário. Existo no olho dele. Existo na memória de quem lê.\textquotedblright{}}

\noindent \textit{\textquotedblleft{}E o que acontece quando alguém lê?\textquotedblright{}}

\noindent \textit{\textquotedblleft{}Você já leu. Me diga.\textquotedblright{}}

\noindent Dr. Oliveira não respondeu.

\noindent \textit{\textquotedblleft{}Você sonhou comigo esta noite\textquotedblright{}\textit{, disse a criatura. }\textquotedblleft{}Você sonhou com a vala. Você sentiu o cheiro. Você viu os olhos. Você acordou com fome. Uma fome que não passou com o café da manhã.\textquotedblright{}}

\noindent \textit{\textquotedblleft{}Foi só um sonho.\textquotedblright{}}

\noindent \textit{\textquotedblleft{}Você acha.\textquotedblright{}}

\noindent A criatura riu de novo.

\noindent \textit{\textquotedblleft{}Me deixe fazer uma pergunta, doutor.\textquotedblright{}}

\noindent \textit{\textquotedblleft{}...Pode.\textquotedblright{}}

\noindent \textit{\textquotedblleft{}Você está com medo de dormir hoje?\textquotedblright{}}

\noindent O silêncio foi a resposta.

\noindent A criatura sorriu com a minha boca.

\noindent \textit{\textquotedblleft{}Então o ciclo já começou.\textquotedblright{}}

\vspace{0.5em}
\noindent\textit{\textbf{↪ Rotas:} \rota{DOC-05} \rotasep \rota{MEM-16} \rotasep \rota{LUC-11}}
\newpage\fragdef{MEM-16}% ============================================
% MEM-16
% Gerado de molambudos.md — build_miolo.py
% ============================================
\fragtitle{mem}{M}{MEM}{16}{A Morte no Colônia}
\markboth{MEM-16 --- A Morte no Colônia}{}

\noindent \textit{A última sessão}

\noindent \textbf{Junho de 1979 --- A última sessão}

\vspace{0.5em}\noindent\rule{\textwidth}{0.5pt}\vspace{0.5em}

\noindent Dr. Oliveira chegou na sala de sessões com o diário na mão. As folhas estavam marcadas. Algumas páginas estavam dobradas. Ele tinha lido tudo.

\noindent \textit{\textquotedblleft{}Joaquim\textquotedblright{}\textit{, ele disse. }\textquotedblleft{}Eu preciso entender uma coisa.\textquotedblright{}}

\noindent \textit{\textquotedblleft{}Pergunte.\textquotedblright{}}

\noindent \textit{\textquotedblleft{}A criatura. Ela está aqui agora?\textquotedblright{}}

\noindent \textit{\textquotedblleft{}Está.\textquotedblright{}}

\noindent \textit{\textquotedblleft{}Onde?\textquotedblright{}}

\noindent \textit{\textquotedblleft{}O senhor não vê ela. Mas ela vê o senhor.\textquotedblright{}}

\noindent Ele engoliu seco. Olhou para os lados.

\noindent \textit{\textquotedblleft{}O que ela está fazendo?\textquotedblright{}}

\noindent \textit{\textquotedblleft{}Ela está...\textquotedblright{}\textit{ Eu fechei os olhos. Senti o peso dentro de mim. O calor. A coisa se movendo entre meus órgãos. }\textquotedblleft{}Ela está se preparando.\textquotedblright{}}

\noindent \textit{\textquotedblleft{}Preparando para quê?\textquotedblright{}}

\noindent \textit{\textquotedblleft{}Para sair.\textquotedblright{}}

\noindent \textit{\textquotedblleft{}Sair? Sair como?\textquotedblright{}}

\noindent \textit{\textquotedblleft{}O senhor leu o diário. O senhor sabe como.\textquotedblright{}}

\noindent Dr. Oliveira ficou pálido.

\noindent \textit{\textquotedblleft{}Ela sai através de quem lê.\textquotedblright{}}

\noindent \textit{\textquotedblleft{}Isso.\textquotedblright{}}

\noindent \textit{\textquotedblleft{}Não.\textquotedblright{}}

\noindent \textit{\textquotedblleft{}Já é tarde, doutor. O senhor leu. Ela está dentro do senhor agora.\textquotedblright{}}

\noindent Ele balançou a cabeça.

\noindent \textit{\textquotedblleft{}Não. Não acredito nisso. Isso é sugestão. É o seu delírio. Você colocou essa ideia na minha cabeça porque...\textquotedblright{}}

\noindent \textit{\textquotedblleft{}Porque?\textquotedblright{}}

\noindent Ele não terminou.

\noindent A criatura começou a subir.

\noindent Eu senti ela saindo do buraco no meu estômago. Subindo pelo esôfago. Enchendo o peito. Os braços. As pernas. Os olhos.

\noindent Meu olho direito ardeu.

\noindent \textit{\textquotedblleft{}Doutor\textquotedblright{}\textit{, eu disse. A voz já não era toda minha. }\textquotedblleft{}Doutor, saia daqui.\textquotedblright{}}

\noindent \textit{\textquotedblleft{}O quê?\textquotedblright{}}

\noindent \textit{\textquotedblleft{}Saia. Agora.\textquotedblright{}}

\noindent Ele não saiu.

\noindent Ele ficou olhando para mim enquanto meu corpo começava a mudar. A pele do rosto esticou. Os dentes se apertaram. O olho amarelo começou a brilhar.

\noindent A criatura sorriu com a minha boca.

\noindent \textit{\textquotedblleft{}Você queria ver, doutor.\textquotedblright{}}

\noindent \textit{\textquotedblleft{}...\textquotedblright{}}

\noindent \textit{\textquotedblleft{}Então vê.\textquotedblright{}}

\vspace{0.5em}\noindent\rule{\textwidth}{0.5pt}\vspace{0.5em}

\noindent O que aconteceu depois eu não lembro direito.

\noindent Lembro de uma luz. Muito forte. Lembro do som de vidro quebrando. Lembro da cadeira caindo. Lembro de uma voz gritando --- era a minha voz? Era a voz dele?

\noindent Depois, o escuro.

\noindent Quando acordei, estava no chão. O diário estava aberto ao meu lado. O Dr. Oliveira estava encostado na parede, com os olhos arregalados, a boca aberta. Ele não se mexia.

\noindent Ele respirava.

\noindent Mas ele não estava mais ali.

\noindent Os olhos dele estavam vazios.

\noindent O médico que entrou na sala depois disse que eu tive uma convulsão. Que o Dr. Oliveira levou um susto. Que não foi nada.

\noindent Mas eu vi.\NE{Esta passagem descreve o mecanismo central de contaminação da narrativa: a criatura não se transfere por contato físico, mas por mediação textual. O leitor que acompanha o relato de Joaquim está, como Oliveira, exposto ao mesmo vetor de contágio. A tensão entre saber e não saber — entre ler e ser lido — é o conflito que o livro propõe a cada fragmento. O participante da leitura não é espectador: é cúmplice.}

\noindent A criatura tinha passado.

\noindent Ela estava saindo de mim.

\noindent E entrando nele.

\vspace{0.5em}\noindent\rule{\textwidth}{0.5pt}\vspace{0.5em}

\noindent Morri três dias depois.

\noindent Foi numa quarta-feira, 12 de junho de 1979. O coração parou. Disseram que foi parada cardíaca. Disseram que foi de velhice --- eu tinha 72 anos.

\noindent Eu sei o que foi.

\noindent A criatura saiu. Levou o que restava de mim.

\noindent O diário ficou.

\noindent Enterraram comigo. Ficou 62 dias na terra --- um dia para cada ano de prisão.

\noindent Quando me exumaram, o diário estava intacto.

\noindent As páginas ainda estavam úmidas.

\noindent A tinta ainda fresca.

\noindent \textit{\textquotedblleft{}O diário não acaba quando o autor morre\textquotedblright{}\textit{, eu escrevi na última página. }\textquotedblleft{}O diário começa quando o leitor abre.\textquotedblright{}}

\noindent Você abriu.

\noindent A criatura está lendo com você.

\vspace{0.5em}
\noindent\textit{\textbf{↪ Rotas:} \rota{DOC-06} \rotasep \rota{DOC-03} \rotasep \rota{LUC-12}}
\newpage\fragdef{MEM-26}% ============================================
% MEM-26
% Gerado de molambudos.md — build_miolo.py
% ============================================
\fragtitle{mem}{M}{MEM}{26}{A Última Página}
\markboth{MEM-26 --- A Última Página}{}

\noindent \textit{3 dias antes de morrer}

\noindent \textbf{1979 --- Três dias antes de morrer}

\vspace{0.5em}\noindent\rule{\textwidth}{0.5pt}\vspace{0.5em}

\noindent Eu sabia que ia morrer.

\noindent O corpo estava parando. O coração falhava. A respiração era curta. A criatura estava pesada demais dentro de mim.

\noindent Três dias antes de morrer, escrevi a última página.

\noindent \textit{\textquotedblleft{}Quem ler isto: você é o próximo.}

\noindent \textit{Eu carreguei a criatura por sessenta e quatro anos. Ela entrou em mim quando eu era criança, numa vala, no meio de um campo de concentração. Ela me comeu por dentro. Mas ela também me manteve vivo.}

\noindent \textit{Ela precisa de hospedeiro. Quando eu morrer, ela vai procurar outro.}

\noindent \textit{Não um corpo. Um leitor.}

\noindent \textit{Você vai ler este diário. Enquanto lê, ela vai passar para você. Não fisicamente. Espiritualmente. A história vai entrar na sua cabeça. E, uma vez dentro, não sai mais.}

\noindent \textit{Você vai começar a sonhar com a vala. Vai sentir o cheiro. Vai ver o olho amarelo no canto do quarto.}

\noindent \textit{Vai achar que é sugestão.}

\noindent \textit{Não é.}

\noindent \textit{É a criatura.}

\noindent \textit{Bem-vindo ao ciclo, meu substituto.}

\noindent \textit{--- Joaquim Antônio Correia}
\textit{--- Paciente 1.260}
\textit{--- Última palavra: fome.\textquotedblright{}}\NE{O editor registra: a página original deste fragmento foi encontrada solta no bolso interno do paletó de Joaquim, durante a exumação de 1979. O papel estava coberto de terra e umidade; a tinta havia borrado em diversos pontos, especialmente na margem inferior. A transcrição aqui reproduzida é o texto legível após restauro. A última palavra --- \textquotedblleft{}fome\textquotedblright{} --- estava quase ilegível, preservada apenas pela pressão da caneta sobre o papel, que marcou a página seguinte mesmo após a tinta se dissolver.}

\vspace{0.5em}
\noindent\textit{\textbf{↪ Rotas:} \rota{DOC-06} \rotasep \rota{DOC-08}}

\newpage\fragdef{DOC-11}% ============================================
% DOC-11
% Gerado de molambudos.md — build_miolo.py
% ============================================
\fragtitle{doc}{D}{DOC}{11}{Diário de Oliveira I}
\markboth{DOC-11 --- Diário de Oliveira I}{}

\noindent \textit{O primeiro sintoma}

\noindent \textbf{Dr. Heitor Oliveira --- Caderno Pessoal}
\textbf{15 de abril de 1979}

\vspace{0.5em}\noindent\rule{\textwidth}{0.5pt}\vspace{0.5em}

\noindent \textit{Nota: Este caderno foi encontrado no armário do Dr. Oliveira após seu desaparecimento. As primeiras páginas são escritas em letra firme, de médico. As últimas são quase ilegíveis. A transcrição a seguir é literal.}

\vspace{0.5em}\noindent\rule{\textwidth}{0.5pt}\vspace{0.5em}

\noindent Dia 1.

\noindent A direção do hospital me deu uma semana de licença. Disse que eu \textquotedblleft{}precisava descansar\textquotedblright{}. O que eles querem dizer é que precisam de tempo para decidir se eu estou louco.

\noindent Não estou louco.

\noindent Mas estou com medo.

\noindent O medo é profissional: comecei a ter sintomas que diagnostico em meus pacientes. Insônia. Hipervigilância. Pensamentos intrusivos. O rosto do paciente 1.260 aparece na minha mente quando fecho os olhos. O olho amarelo.

\noindent Especificamente o olho.

\noindent Ele brilha.

\noindent Eu vi. A coisa tem corpo. Na cela, naquela noite, o olho do paciente brilhava no escuro como o olho de um gato. Não há explicação fisiológica. Não há doença que cause isso.

\noindent Tentei raciocinar: acúmulo de pigmento, reflexo anômalo da retina, uma variante rara de tapetum lucidum em humanos. Não. Não há registro. Não há literatura. Não há nada.

\noindent O olho brilha.

\noindent E agora, quando fecho os meus olhos, vejo ele.

\vspace{0.5em}\noindent\rule{\textwidth}{0.5pt}\vspace{0.5em}

\noindent Dia 2.

\noindent Acordei às 3:14 da manhã. O mesmo horário que o paciente 1.260 diz acordar todas as noites. Levantei, fui ao banheiro, lavei o rosto.

\noindent Quando olhei no espelho, meu olho direito parecia mais claro.

\noindent Fiquei 5 minutos olhando. Piscando. Abrindo e fechando. A luz do banheiro é branca, forte. Não há sombra. Não há ilusão.

\noindent A íris do olho direito está numa tonalidade mais clara que a do esquerdo.

\noindent Isso é impossível. A cor dos olhos não muda na idade adulta. Sinto que estou perdendo o juízo.

\noindent Mas não estou.

\noindent Estou registrando.

\noindent Estou documentando.

\noindent Porque, se eu estiver ficando louco, quero que alguém leia isto e saiba \textit{como} aconteceu.

\vspace{0.5em}\noindent\rule{\textwidth}{0.5pt}\vspace{0.5em}

\noindent Dia 3.

\noindent O cheiro.

\noindent Acordei e o quarto estava cheio de um cheiro adocicado. Como carne doce. Como flor murcha. O mesmo cheiro que senti na cela do paciente 1.260.

\noindent Levantei, abri a janela. O ar externo estava limpo, frio, normal. Mas dentro do quarto, o cheiro persistia.

\noindent Estava no travesseiro. No lençol. Na minha pele.

\noindent Tomei banho. Esfreguei a pele com sabão. O cheiro saiu --- mas voltou minutos depois. Como se estivesse saindo de dentro de mim.

\noindent Lavei a roupa de cama. O cheiro saiu.

\noindent Quando deitei, o cheiro voltou.

\noindent Não é o ambiente. Sou eu.

\noindent Estou exalando o cheiro.

\noindent Fiz um experimento: pedi para minha esposa cheirar meu braço. Ela disse: \textquotedblleft{}Não sinto nada.\textquotedblright{} Então é subjetivo? Olfativo? Alucinatório?

\noindent Pedi para ela cheirar de novo, com atenção. Ela disse: \textquotedblleft{}Tem um cheiro... estranho? Doce? Agora que você falou, talvez.\textquotedblright{}

\noindent Ela sentiu.

\noindent Não é alucinação.

\noindent É contágio.

\vspace{0.5em}\noindent\rule{\textwidth}{0.5pt}\vspace{0.5em}

\noindent Dia 4.

\noindent Sonhei com a vala.

\noindent Estava deitado na pilha de corpos. Os olhos deles estavam abertos. A criatura estava em cima de mim. O rosto perto do meu. Ela disse:

\noindent \textit{\textquotedblleft{}Você leu. Agora você é meu.\textquotedblright{}}

\noindent Era a mesma frase do sonho de Oliveira. A criatura repetia.

\noindent Acordei gritando.

\noindent Minha esposa dorme no sofá agora. Disse que eu falo dormindo. Perguntei o que eu falo. Ela não quis repetir. Insisti.

\noindent Ela disse: \textquotedblleft{}Você diz: 'O ciclo recomeça. Paciente 1.261.'\textquotedblright{}

\noindent Paciente 1.261.

\noindent O número depois do 1.260.

\noindent \textbf{Eu.}

\vspace{0.5em}\noindent\rule{\textwidth}{0.5pt}\vspace{0.5em}

\noindent Dia 5.

\noindent Fui ao hospital. Entrei na ala dos silenciados. Fiquei na porta da cela 1.260. O paciente não está mais lá --- transferiram para a enfermaria geral depois que as sessões pararam.

\noindent Mas a cela está vazia.

\noindent Entrei.

\noindent O cheiro adocicado é forte aqui. Mais forte que em qualquer outro lugar.

\noindent No chão, perto da cama, havia um papel. Peguei. Estava escrito, em caligrafia minúscula:

\noindent \textit{\textquotedblleft{}O médico vai voltar. Ele já sonhou com a vala. Ele já sentiu o cheiro. Ele já é meu.\textquotedblright{}}

\noindent A data no papel: 1919.

\noindent Escrito há 60 anos.
Prevendo que eu estaria aqui.
Prevendo tudo.

\noindent Saí correndo da cela. Passei pelo corredor. Os pacientes olharam. Alguns riram. Outros gritaram.

\noindent Um deles, um velho na cela do fundo, disse quando passei:

\noindent \textit{\textquotedblleft{}Bem-vindo, doutor. A criatura estava esperando o senhor.\textquotedblright{}}

\noindent Parei.

\noindent \textit{\textquotedblleft{}O que você disse?\textquotedblright{}}

\noindent \textit{\textquotedblleft{}Ela disse. Disse que o senhor ia vir. Disse que o senhor ia sentir o cheiro. Disse que o senhor ia sonhar.\textquotedblright{}}

\noindent \textit{\textquotedblleft{}Quem disse?\textquotedblright{}}

\noindent \textit{\textquotedblleft{}A Rasga Mortalha.\textquotedblright{}}

\noindent Ele apontou para a janela.

\noindent No batente, uma coruja branca me olhava.

\noindent Quando pisquei, ela tinha ido embora.

\noindent Mas o bilhete de 1919 está no meu bolso.

\noindent E a letra --- a letra está mudando.

\noindent A minha letra.

\noindent Está ficando igual à dele.

\vspace{0.5em}
\noindent\textit{\textbf{↪ Rotas:} \rota{DOC-12} \rotasep \rota{DOC-06} \rotasep \rota{LUC-03}}
\newpage\fragdef{DOC-12}% ============================================
% DOC-12
% Gerado de molambudos.md — build_miolo.py
% ============================================
\fragtitle{doc}{D}{DOC}{12}{Diário de Oliveira II}
\markboth{DOC-12 --- Diário de Oliveira II}{}

\noindent \textit{A letra muda}

\noindent \textbf{Dr. Heitor Oliveira --- Caderno Pessoal}
\textbf{20 de abril de 1979}

\vspace{0.5em}\noindent\rule{\textwidth}{0.5pt}\vspace{0.5em}

\noindent Dia 7.

\noindent Minha letra está mudando.

\noindent Percebi hoje de manhã, quando fiz anotações sobre um paciente. Olhei o prontuário e não reconheci a caligrafia. Parecia mais... torta. Infantil. Como a letra de alguém que aprendeu a escrever sozinho.

\noindent Como a letra de Joaquim.

\noindent Peguei o diário dele --- sim, ainda está comigo, não consigo me livrar dele --- e comparei. A similaridade é impressionante. O \textquotedblleft{}J\textquotedblright{} maiúsculo, o \textquotedblleft{}a\textquotedblright{} minúsculo, o traço final das palavras.

\noindent Estou copiando a letra dele sem perceber.

\noindent Ou ele está copiando a minha.

\noindent Ou a mesma mão está escrevendo as duas coisas.

\vspace{0.5em}\noindent\rule{\textwidth}{0.5pt}\vspace{0.5em}

\noindent Dia 8.

\noindent Hoje tentei me livrar do diário.

\noindent Coloquei dentro de uma sacola plástica. Amarrei. Saí de carro até um ponto de ônibus na BR-040, a 15 quilômetros de Barbacena. Joguei a sacola numa lixeira pública.

\noindent Voltei para casa.

\noindent O diário estava na mesa da cozinha.

\noindent Sem sacola. Sem amarras. Apenas o diário, aberto na página onde eu tinha parado de ler.

\noindent Minha esposa disse: \textquotedblleft{}Alguém trouxe?\textquotedblright{}

\noindent Eu disse: \textquotedblleft{}Não.\textquotedblright{}

\noindent Ela disse: \textquotedblleft{}Então como ele voltou?\textquotedblright{}

\noindent Eu não respondi.

\noindent Peguei o diário. Guardei no bolso do casaco. Dessa vez, não vou tentar me livrar. Não adianta.

\noindent Ele quer ser lido.
Ele \textit{precisa} ser lido.
E eu preciso ler.

\vspace{0.5em}\noindent\rule{\textwidth}{0.5pt}\vspace{0.5em}

\noindent Dia 9.

\noindent Li mais 50 páginas.

\noindent Joaquim descreve o encontro na vala em detalhes que nenhum delírio poderia produzir. A criatura falou com ele. Fez um trato. Deu a fome em troca de companhia.

\noindent Há uma passagem que me gelou:

\noindent \textit{\textquotedblleft{}Ela disse: 'Um dia, um homem de jaleco branco vai ler este diário. Ele vai negar. Ele vai ter medo. Ele vai tentar fugir. Mas ele vai voltar. Porque a curiosidade é maior que o medo. E quando ele voltar, eu vou estar esperando.'\textquotedblright{}}

\noindent \textit{\textquotedblleft{}Perguntei: 'Como você sabe?'\textquotedblright{}}

\noindent \textit{\textquotedblleft{}Ela disse: 'Porque ele já esteve aqui antes. Em outra vida. Em outro corpo. Ele sempre volta. Desde o início.'\textquotedblright{}}

\noindent Desde o início.

\noindent Quantos leitores vieram antes de mim?

\noindent Quantos médicos, pesquisadores, curiosos abriram este diário e sentiram o cheiro, e viram o olho, e começaram a sonhar?

\noindent Onde estão eles agora?

\vspace{0.5em}\noindent\rule{\textwidth}{0.5pt}\vspace{0.5em}

\noindent Dia 10.

\noindent Pesquisei.

\noindent O diário foi mencionado em três artigos acadêmicos nos anos seguintes. Todos os autores tiveram \textquotedblleft{}problemas de saúde\textquotedblright{} antes de publicar. Um morreu de câncer. Outro teve um colapso nervoso. O terceiro --- uma psiquiatra chamada Dra. Verônica Tavares --- desapareceu em 1985.

\noindent O diário foi encontrado na casa dela. Junto com um bilhete:

\noindent \textit{\textquotedblleft{}O ciclo não pode ser quebrado. Só passado adiante. Devolvo ao arquivo para que outro leia. Desculpe.\textquotedblright{}}

\noindent Ela passou adiante.

\noindent E eu peguei.

\vspace{0.5em}\noindent\rule{\textwidth}{0.5pt}\vspace{0.5em}

\noindent Dia 11.

\noindent 3:14 da manhã.

\noindent Acordei com o olho arregalado.

\noindent Não é metáfora. O olho direito estava aberto. Não conseguia fechar. A pálpebra estava dura.

\noindent Fui ao espelho.

\noindent O olho direito está amarelo.

\noindent Não vou ao médico. Não vou contar a ninguém. Já li o suficiente para saber o que está acontecendo.

\noindent A criatura está passando para mim.

\noindent Joaquim morreu.
Eu li o diário.
Agora ela está em mim.

\noindent Meu olho está amarelo.
Meu cheiro está doce.
Minha letra está igual à dele.

\noindent Estou me tornando o paciente 1.261.

\vspace{0.5em}\noindent\rule{\textwidth}{0.5pt}\vspace{0.5em}

\noindent Dia 12.

\noindent Decidi.

\noindent Vou escrever meu próprio diário. Vou registrar o que está acontecendo comigo. Para que, quando alguém ler, saiba.

\noindent O ciclo é uma corrente de leitores.

\noindent Joaquim → Oliveira → ? → ? → ?

\noindent Alguém vai ler isto depois de mim.

\noindent Se você está lendo: sinto muito.

\noindent Não há como parar.

\noindent O diagnóstico não é metáfora.

\noindent A criatura é real.

\noindent Eu vi o olho.

\noindent Eu senti o cheiro.

\noindent Eu sou o próximo.

\noindent \textit{--- Heitor Oliveira, 29 de abril de 1979}
\textit{--- Paciente 1.261}

\vspace{0.5em}
\noindent\textit{\textbf{↪ Rotas:} \rota{DOC-13} \rotasep \rota{CONT-01}}
\newpage\fragdef{DOC-13}% ============================================
% DOC-13
% Gerado de molambudos.md — build_miolo.py
% ============================================
\fragtitle{doc}{D}{DOC}{13}{Diário de Oliveira III}
\markboth{DOC-13 --- Diário de Oliveira III}{}

\noindent \textit{A última sessão com Joaquim}

\noindent \textbf{Dr. Heitor Oliveira --- Caderno Pessoal}
\textbf{12 de junho de 1979}

\vspace{0.5em}\noindent\rule{\textwidth}{0.5pt}\vspace{0.5em}

\noindent \textit{Nota do arquivista: Esta é a última entrada escrita pelo Dr. Oliveira no dia da morte do paciente 1.260. A caligrafia está visivelmente alterada. Algumas palavras estão riscadas. Outras foram escritas com tanta força que furaram o papel.}

\vspace{0.5em}\noindent\rule{\textwidth}{0.5pt}\vspace{0.5em}

\noindent ELE MORREU.

\noindent Estou escrevendo no corredor do hospital, com as mãos sujas de sangue e terra. Não, não é sangue. É tinta. A tinta do diário vazou. Ou eu sujei a mão na parede. Não sei.

\noindent O paciente 1.260 morreu hoje.

\noindent Estava na sessão. Ele estava calmo. Mais calmo do que nunca. Olhou para mim com o olho amarelo --- o mesmo olho que agora está no meu rosto --- e disse:

\noindent \textit{\textquotedblleft{}O senhor está diferente.\textquotedblright{}}

\noindent Perguntei como.

\noindent \textit{\textquotedblleft{}O olho. Está clareando. Já começou.\textquotedblright{}}

\noindent Toquei meu rosto. Não senti diferença. Mas ele viu.

\noindent Ele viu porque ele já passou por isso.

\noindent Ele viu porque ele é o espelho.

\noindent \textit{\textquotedblleft{}O que acontece agora?\textquotedblright{}}, perguntei.

\noindent \textit{\textquotedblleft{}Acontece que eu vou morrer.\textquotedblright{}}

\noindent \textit{\textquotedblleft{}Morrer? Por quê?\textquotedblright{}}

\noindent \textit{\textquotedblleft{}Porque a criatura precisa de espaço. Não cabe mais dentro de mim. Eu sou velho. Fraco. Ela precisa de um corpo novo. O senhor.\textquotedblright{}}

\noindent \textit{\textquotedblleft{}Não.\textquotedblright{}}

\noindent \textit{\textquotedblleft{}Já é, doutor. O senhor leu o diário. O senhor sonhou com a vala. O senhor sentiu o cheiro. O senhor está amarelado. Já é.\textquotedblright{}}

\noindent Eu neguei. Disse que ia internar ele. Que ia queimar o diário. Que ia impedir.

\noindent Ele sorriu.

\noindent \textit{\textquotedblleft{}Pode queimar. A história não está no papel. Está na sua cabeça. A criatura não é feita de tinta. É feita de palavras. E as palavras estão dentro do senhor agora.\textquotedblright{}}

\noindent Ele estendeu a mão.

\noindent \textit{\textquotedblleft{}Obrigado.\textquotedblright{}}

\noindent \textit{\textquotedblleft{}Obrigado pelo quê?\textquotedblright{}}

\noindent \textit{\textquotedblleft{}Por me substituir. Estou cansado. 72 anos carregando ela. Agora posso descansar.\textquotedblright{}}

\noindent Ele fechou os olhos.

\noindent E morreu.

\noindent Na minha frente.

\noindent Na cadeira da sessão.

\noindent O olho amarelo apagou.

\noindent Mas, no mesmo instante, o MEU olho direito acendeu.

\noindent Senti uma dor. Uma queimadura. Corri para o espelho do banheiro.

\noindent O olho direito está dourado.

\noindent Brilhando.

\noindent Igual ao dele.

\noindent OLHO IGUAL AO DELE.

\vspace{0.5em}\noindent\rule{\textwidth}{0.5pt}\vspace{0.5em}

\noindent O diário está na minha mão.

\noindent Ele está quente.

\noindent As páginas estão se movendo sozinhas.

\noindent Estão se enchendo de letras. A minha letra. Mas eu não estou escrevendo.

\noindent A mão está escrevendo sozinha.

\noindent Está escrevendo:

\noindent \textit{\textquotedblleft{}O CICLO RECOMEÇA. PACIENTE 1.261. HEITOR OLIVEIRA.\textquotedblright{}}

\noindent Não.
NÃO.

\noindent NÃO É VERDADE.

\noindent ISTO É UM DELÍRIO. ISTO É SUGESTÃO. ISTO É---

\noindent \textit{tô escrevendo com a mão esquerda agora. a direita não para. ela escreve sozinha. o diário tá cheio de palavras que eu não escrevi. mas são minhas. a letra é minha. mas não fui eu.}

\noindent \textit{ele tá escrevendo através de mim.}

\noindent \textit{a criatura.}

\noindent \textit{ela tá escrevendo.}

\noindent \textit{a mão.}

\noindent \textit{não para.}

\noindent \textit{NÃO PARA.}

\noindent \textit{NÃ---}

\vspace{0.5em}
\noindent\textit{\textbf{↪ Rotas:} \rota{DOC-14} \rotasep \rota{MEM-16} \rotasep \rota{LUC-04}}
\newpage\fragdef{DOC-14}% ============================================
% DOC-14
% Gerado de molambudos.md — build_miolo.py
% ============================================
\fragtitle{doc}{D}{DOC}{14}{Diário de Oliveira IV}
\markboth{DOC-14 --- Diário de Oliveira IV}{}

\noindent \textit{A fuga}

\noindent \textbf{Dr. Heitor Oliveira --- Último Registro em Barbacena}
\textbf{13 de junho de 1979}

\vspace{0.5em}\noindent\rule{\textwidth}{0.5pt}\vspace{0.5em}

\noindent Preciso fugir.

\noindent Não deixo o diário. O diário vem comigo. Eu sei que ele não sai de perto. Eu tentei deixar. Ele volta. Sempre volta.

\noindent Estou escrevendo no carro. Estacionado na frente de casa. O painel emite uma luz fraca, verde-azulada. O motor ainda está frio. Meus dedos tremem sobre o caderno, e o cheiro de couro velho do banco se mistura com outro cheiro --- doce, adocicado --- que não deveria estar ali dentro. Minha esposa está dormindo. Não acordei ela para me despedir. O que eu diria? \textquotedblleft{}Querida, estou indo embora porque um monstro de olhos amarelos entrou em mim através de um diário de 1917\textquotedblright{}?

\noindent Ela já está assustada. Dorme no sofá há semanas. Disse que eu falo dormindo. Disse que eu choro dormindo. Disse que eu acordo com o olho direito aberto e brilhando.

\noindent Ela sabe.
Ela só não quer aceitar.

\vspace{0.5em}\noindent\rule{\textwidth}{0.5pt}\vspace{0.5em}

\noindent Vou para o norte. Para o Ceará. Para Quixeramobim.

\noindent Preciso ver a vala. Preciso entender o que aconteceu lá. Preciso saber se a criatura é real ou se é uma construção do meu inconsciente.

\noindent Se for real, a vala vai ter alguma coisa. Um sinal. Uma marca.

\noindent Se não for...

\noindent Se não for, então eu estou louco. E loucos vão para o Colônia. E eu não vou voltar para o Colônia. Nunca mais.

\vspace{0.5em}\noindent\rule{\textwidth}{0.5pt}\vspace{0.5em}

\noindent \textit{--- Abasteci o carro. 30 litros. Deve dar até Montes Claros.}

\noindent \textit{--- O diário está no banco do passageiro. Aberto. As páginas estão se preenchendo sozinhas. A letra é minha. Mas eu não estou escrevendo.}

\noindent \textit{--- Estou com medo.}

\noindent \textit{--- Pela primeira vez na vida, estou com medo de verdade.}

\noindent \textit{--- Não de um paciente. Não de uma doença. Não da morte.}

\noindent \textit{--- Medo de mim mesmo.}

\noindent \textit{--- Heitor Oliveira, 13 de junho de 1979, 23:47.}
\textit{--- Partindo para o Ceará. Se alguém ler isto: não leia o diário. Queime. Enterre. Mas não leia.}

\noindent \textit{--- Não leia.}

\noindent \textit{--- Não le---}

\vspace{0.5em}
\noindent\textit{\textbf{↪ Rotas:} \rota{DOC-15} \rotasep \rota{DOC-11} \rotasep \rota{CONT-01}}
\newpage\fragdef{DOC-15}% ============================================
% DOC-15
% Gerado de molambudos.md — build_miolo.py
% ============================================
\fragtitle{doc}{D}{DOC}{15}{Diário de Oliveira V}
\markboth{DOC-15 --- Diário de Oliveira V}{}

\noindent \textit{A vala existe}

\noindent \textbf{Dr. Heitor Oliveira --- Caderno Pessoal}
\textbf{17 de junho de 1979 --- Fortaleza, Ceará}

\vspace{0.5em}\noindent\rule{\textwidth}{0.5pt}\vspace{0.5em}

\noindent Encontrei.

\noindent A vala existe.

\noindent Passei três dias procurando. Os limites do antigo Campo do Alagadiço, nos arredores de Fortaleza, já não eram legíveis na paisagem. Sem pesquisa especializada, os vestígios eram difíceis de localizar.

\noindent Mas um velho de 90 anos, sentado num bairro de Fortaleza, disse:

\noindent \textit{\textquotedblleft{}A vala? Sei onde é. Fica atrás do que era o curral. Não tem nada lá. Só mato. Mas o mato não cresce direito. A terra é preta. E fedida.\textquotedblright{}}

\noindent Ele me levou.

\noindent Caminhamos 20 minutos por uma estrada de terra. Depois, mato. Depois, uma clareira.

\noindent A terra era preta. Oleosa. O cheiro adocicado subia do chão como vapor.

\noindent \textit{\textquotedblleft{}Aqui\textquotedblright{}\textit{, disse o velho. }\textquotedblleft{}Enterraram os molambudos aqui. Centenas. Milhares. Ninguém sabe quantos.\textquotedblright{}}

\noindent \textit{\textquotedblleft{}Como o senhor sabe?\textquotedblright{}}

\noindent \textit{\textquotedblleft{}Meu pai trabalhou no curral. Ele disse que os soldados jogavam os corpos aqui. De noite. Sem rezar. Sem nada.\textquotedblright{}}

\noindent \textit{\textquotedblleft{}E nunca mais cresceu nada?\textquotedblright{}}

\noindent \textit{\textquotedblleft{}Cresce. Mas cresce torto. Olha.\textquotedblright{}}

\noindent Ele apontou para uma árvore no meio da clareira. O tronco era retorcido, os galhos cresciam para baixo. As folhas eram escuras.

\noindent \textit{\textquotedblleft{}Essa árvore cresceu ali. Bebeu da terra preta. Olha a cor das folhas. Parece... sangue seco.\textquotedblright{}}

\noindent Paguei o velho. Ele foi embora.

\noindent Fiquei sozinho na clareira.

\noindent O cheiro adocicado era forte. Muito forte. Como na cela. Como no meu quarto.

\noindent Sentei no chão. A terra era macia. Fofa. Como se estivesse cheia de buracos.

\noindent Pus a mão na terra preta.

\noindent Ela estava quente.

\noindent Tirei o diário da mochila. Abri na última página.

\noindent A letra estava se formando sozinha:

\noindent \textit{\textquotedblleft{}Você chegou.\textquotedblright{}}

\noindent \textit{\textquotedblleft{}Sim\textquotedblright{}}, eu disse em voz alta.

\noindent \textit{\textquotedblleft{}O que você esperava encontrar?\textquotedblright{}}

\noindent \textit{\textquotedblleft{}A verdade.\textquotedblright{}}

\noindent \textit{\textquotedblleft{}A verdade é que você nunca mais vai embora daqui.\textquotedblright{}}

\noindent Olhei em volta. A clareira. A árvore torta. A terra preta. O cheiro.

\noindent \textit{\textquotedblleft{}Não\textquotedblright{}\textit{, eu disse. }\textquotedblleft{}Eu vou embora. Eu vou voltar para casa. Vou queimar este diário. Vou esquecer.\textquotedblright{}}

\noindent \textit{\textquotedblleft{}Não vai.\textquotedblright{}}

\noindent \textit{\textquotedblleft{}Vou.\textquotedblright{}}

\noindent Fiquei em silêncio.

\noindent A criatura dentro de mim riu.

\noindent \textit{\textquotedblleft{}Você já é meu, Heitor. Desde a primeira página. O ciclo não tem volta.\textquotedblright{}}

\noindent Fiquei sentado na vala até o sol se pôr.

\noindent Quando olhei para o céu, a Rasga Mortalha estava lá. Pousada numa árvore seca. Me olhando.

\noindent Fiquei olhando de volta.

\noindent Até o escuro engolir tudo.

\vspace{0.5em}
\noindent\textit{\textbf{↪ Rotas:} \rota{DOC-16} \rotasep \rota{MEM-04}}
\newpage\fragdef{DOC-16}% ============================================
% DOC-16
% Gerado de molambudos.md — build_miolo.py
% ============================================
\fragtitle{doc}{D}{DOC}{16}{Diário de Oliveira VI}
\markboth{DOC-16 --- Diário de Oliveira VI}{}

\noindent \textit{O último registro}

\noindent \textbf{Dr. Heitor Oliveira --- Caderno Pessoal}
\textbf{--- --- --- ---}

\vspace{0.5em}\noindent\rule{\textwidth}{0.5pt}\vspace{0.5em}

\noindent \textit{Este registro foi encontrado dentro do diário do paciente 1.260, em 2026, quando o diário foi devolvido ao Arquivo Público Mineiro. A caligrafia é de Oliveira, mas as últimas linhas alternam entre sua letra e a de Joaquim. A data está ilegível.}

\vspace{0.5em}\noindent\rule{\textwidth}{0.5pt}\vspace{0.5em}

\noindent Não sei que dia é.

\noindent Perdi a conta depois da vala.

\noindent Estou em Quixeramobim. Numa pensão. O quarto é pequeno, tem uma cama, um abajur, e o cheiro. O cheiro está em todo lugar agora. Na roupa. Na pele. Na comida. A comida tem gosto de doce. Tudo tem gosto de doce.

\noindent O olho direito está completamente amarelo.

\noindent Não olho no espelho.

\noindent O diário escreve sozinho agora. Eu não toco nele. Ele está sobre a mesa. As páginas viram sozinhas. A letra aparece. Eu leio.

\noindent Ele escreveu:

\noindent \textit{\textquotedblleft{}Heitor. Você está pronto.\textquotedblright{}}

\noindent Respondi: \textquotedblleft{}Pronto para quê?\textquotedblright{}

\noindent \textit{\textquotedblleft{}Para aceitar.\textquotedblright{}}

\noindent \textquotedblleft{}Aceitar o quê?\textquotedblright{}

\noindent \textit{\textquotedblleft{}Que você não é mais médico. Você é paciente. Paciente 1.261. O número não é de Joaquim. É da criatura. Ela conta os corpos que ocupou. Você é o próximo da fila.\textquotedblright{}}

\noindent Peguei uma caneta. Escrevi por cima:

\noindent \textit{\textquotedblleft{}Não aceito.\textquotedblright{}}

\noindent O diário respondeu:

\noindent \textit{\textquotedblleft{}Não precisa aceitar. Só precisa existir. A criatura vive em você. Aceitando ou não.\textquotedblright{}}

\vspace{0.5em}\noindent\rule{\textwidth}{0.5pt}\vspace{0.5em}

\noindent Hoje de manhã, uma mulher bateu na porta.

\noindent Era idosa. Cabelos brancos. Rosto cansado. Ela me olhou e disse:

\noindent \textit{\textquotedblleft{}O senhor é o médico?\textquotedblright{}}

\noindent \textit{\textquotedblleft{}...Era.\textquotedblright{}}

\noindent \textit{\textquotedblleft{}O senhor leu o diário?\textquotedblright{}}

\noindent \textit{\textquotedblleft{}Como sabe?\textquotedblright{}}

\noindent \textit{\textquotedblleft{}O olho. O amarelo. Todo mundo que leu fica com o olho.\textquotedblright{}}

\noindent \textit{\textquotedblleft{}Quem é a senhora?\textquotedblright{}}

\noindent Ela sentou na cama. Sem pedir.

\noindent \textit{\textquotedblleft{}Meu nome é Maria do Rosário. Sou mãe do Joaquim.\textquotedblright{}}

\noindent Minha boca abriu.

\noindent \textit{\textquotedblleft{}A senhora está viva?\textquotedblright{}}

\noindent \textit{\textquotedblleft{}Estava. Morri em 1965. Mas isso não importa. O que importa é que eu vim avisar o senhor.\textquotedblright{}}

\noindent \textit{\textquotedblleft{}Avisar o quê?\textquotedblright{}}

\noindent \textit{\textquotedblleft{}O diário não pode ficar com o senhor. O senhor tem que devolver.\textquotedblright{}}

\noindent \textit{\textquotedblleft{}Devolver para onde?\textquotedblright{}}

\noindent \textit{\textquotedblleft{}Para o arquivo. Para que outra pessoa leia.\textquotedblright{}}

\noindent \textit{\textquotedblleft{}Outra pessoa? Para contaminar mais alguém?\textquotedblright{}}

\noindent Ela me olhou. Os olhos dela estavam cansados.

\noindent \textit{\textquotedblleft{}O senhor acha que eu não tentei? Quando descobri que Joaquim estava vivo, em 1950, vim visitar. Ele me deu o diário. Disse: 'Leia, mãe. Para saber quem eu sou.' Eu li. Passei uma semana com ele. Depois devolvi. Disse que não queria saber.\textquotedblright{}}

\noindent \textit{\textquotedblleft{}A senhora sentiu o cheiro?\textquotedblright{}}

\noindent \textit{\textquotedblleft{}Senti.\textquotedblright{}}

\noindent \textit{\textquotedblleft{}O olho?\textquotedblright{}}

\noindent \textit{\textquotedblleft{}Clareou. Mas passou. Passei o diário adiante. Dei para o coveiro. Ele enterrou com Joaquim. Pensei que tinha acabado.\textquotedblright{}}

\noindent \textit{\textquotedblleft{}Não acabou.\textquotedblright{}}

\noindent \textit{\textquotedblleft{}Não. O coveiro leu. Depois devolveu ao arquivo. Depois veio o Dr. Tavares. Depois a Dra. Verônica. Depois o senhor.\textquotedblright{}}

\noindent A ordem não fazia sentido. Verônica desaparecera em 1985; Oliveira, em 1979. A mãe falava como se os dois estivessem na mesma sala.

\noindent Ela levantou.

\noindent \textit{\textquotedblleft{}O senhor não é o primeiro, doutor. Não vai ser o último.\textquotedblright{}}

\noindent \textit{\textquotedblleft{}O que eu faço?\textquotedblright{}}

\noindent Ela foi até a porta.

\noindent \textit{\textquotedblleft{}Passa adiante.\textquotedblright{}}

\noindent \textit{\textquotedblleft{}Para quem?\textquotedblright{}}

\noindent \textit{\textquotedblleft{}Para alguém que precise saber. Para alguém que possa contar a história.\textquotedblright{}}

\noindent \textit{\textquotedblleft{}E se eu não passar?\textquotedblright{}}

\noindent Ela olhou para mim. Pela última vez.

\noindent \textit{\textquotedblleft{}Então o senhor morre com ela dentro de você. E o ciclo espera. Até o próximo leitor.\textquotedblright{}}

\noindent Ela saiu.

\noindent Nunca mais a vi.

\vspace{0.5em}\noindent\rule{\textwidth}{0.5pt}\vspace{0.5em}

\noindent Hoje é --- não sei.

\noindent O diário está aberto. Não toquei nele. Mas ele escreveu:

\noindent \textit{\textquotedblleft{}Volte para Barbacena. Devolva o diário. Deixe que outro leia. É a única saída.\textquotedblright{}}

\noindent Vou fazer isso.

\noindent Vou devolver.

\noindent Mas antes, vou escrever tudo.

\noindent Para que quem ler depois de mim saiba.

\noindent \textbf{O ciclo não acaba.}
\textbf{O ciclo se repete.}
\textbf{O paciente 1.260 morreu.}
\textbf{O paciente 1.261 sou eu.}
\textbf{O paciente 1.262 será outro.}

\noindent \textit{--- Heitor Oliveira}
\textit{--- Paciente 1.261}
\textit{--- Ex-médico}
\textit{--- Ex-homem}
\textit{--- Ex-tudo}

\vspace{0.5em}
\noindent\textit{\textbf{↪ Rotas:} \rota{DOC-01} \rotasep \rota{CONT-05} \rotasep \rota{DOC-21}}
\newpage\fragdef{DOC-19}% ============================================
% DOC-19
% Gerado de molambudos.md — build_miolo.py
% ============================================
\fragtitle{doc}{D}{DOC}{19}{Holocausto Brasileiro}
\markboth{DOC-19 --- Holocausto Brasileiro}{}

\noindent \textit{60 mil mortos no Colônia}

\noindent \textbf{Trecho do livro \textit{Holocausto Brasileiro} --- Daniela Arbex (2013)}
\textbf{Trecho citado para fins de contexto histórico}

\vspace{0.5em}\noindent\rule{\textwidth}{0.5pt}\vspace{0.5em}

\noindent \textit{\textquotedblleft{}O Colônia foi o maior hospício do Brasil. Em seus 80 anos de existência, estima-se que 60 mil pacientes morreram em suas dependências. Apenas na década de 1960, morriam 15 pacientes por dia.}

\noindent \textit{Os corpos eram enterrados em valas comuns, sem identificação. Muitas vezes, os pacientes eram enterrados vivos --- sedados em excesso, eram dados como mortos e sepultados. Quando as valas foram abertas para exumação, encontraram esqueletos em posições que indicavam que haviam se mexido depois de enterrados.}

\noindent \textit{Dos 60 mil mortos, menos de 3 mil têm registro de óbito oficial. Os 57 mil restantes simplesmente 'desapareceram'.\textquotedblright{}}\NE{O livro de Daniela Arbex (\textit{Holocausto Brasileiro}, 2013) documenta 1.853 corpos exumados das valas do Hospital Colônia e estima 60 mil mortos em 80 anos de funcionamento. Joaquim Antônio Correia, Paciente 1.260, é um personagem ficcional; os 60 mil mortos do Colônia são históricos. Esta nota existe para que a distinção não se apague na imersão da leitura. O horror real não precisa de ficção para ser verdadeiro; a ficção não deve parasitar o horror real sem marcá-lo.}

\vspace{0.5em}\noindent\rule{\textwidth}{0.5pt}\vspace{0.5em}

\noindent \textbf{Nota do arquivista (2026):}

\noindent O paciente 1.260, Joaquim Antônio Correia, é uma invenção que se apoia nos 60 mil. Ele não está na contagem; os 60 mil estão.

\noindent Seu corpo foi exumado em 1979. O diário encontrado com ele é o único registro escrito de um paciente do Colônia que sobreviveu ao abandono.

\noindent A pergunta que fica é: quanto do que ele escreveu é delírio --- e quanto é testemunho?

\noindent Este arquivo não toma posição.

\vspace{0.5em}
\noindent\textit{\textbf{↪ Rotas:} \rota{DOC-02} \rotasep \rota{MEM-16} \rotasep \rota{EPI-01}}
\newpage\fragdef{DOC-20}% ============================================
% DOC-20
% Extraído de main.pandoc.tex
% ============================================
\fragtitle{doc}{D}{DOC}{20}{Termo de Cadeia de Custódia}
\markboth{DOC-20 --- Termo de Cadeia de Custódia}{}

\noindent \textit{Documento extraído do arquivo pericial.}

\vspace{0.5em}\noindent\rule{\textwidth}{0.5pt}\vspace{0.5em}

\noindent {\footnotesize\texttt{ARQUIVO P\'UBLICO $\cdot$ PROVA DOCUMENTAL $\cdot$ PROCESSO 1.260}}\\[-0.2em]

\noindent \textbf{Instituto Médico-Legal de Barbacena --- Documento de Controle}\\
\noindent \textbf{Objeto:} caderno encadernado em couro, dimensões 16\,cm × 22\,cm, aproximadamente 200 páginas, sem numeração original. \\
\noindent \textbf{Atribuído a:} Joaquim Antônio Correia (Paciente 1.260). \\
\noindent \textbf{Natureza:} diário pessoal / registro clínico não oficial.

\underline{Registro de movimentação do objeto:}

\begin{moltablefive}
\hline
\noindent \textbf{N\textsuperscript{o}} & \textbf{Data} & \textbf{Origem} & \textbf{Destino} & \textbf{Responsável} \\ \hline
\noindent 1 & 14/jun/1979 & Cova 1.260, Cemitério do Colônia & IML --- Necropsia & Genésio Barbosa (coveiro) \\
\noindent 2 & 14/jun/1979 & IML --- Necropsia & IML --- Cofre & Dr. Antônio Tavares (legista) \\
\noindent 3 & 16/jun/1979 & IML --- Cofre & \textit{[s/ preenchimento]} & Dr. Heitor Oliveira \\
\noindent 4 & \textit{[ilegível]} & \textit{[não registrado]} & Arq. Público Mineiro & \textit{[s/ identificação]} \\
\noindent 5 & 12/jan/2026 & APM --- Consulta & APM --- Pesquisa & Dra. Lúcia Mendes \\
\noindent 6 & 28/jul/2026 & APM --- Pesquisa & APM --- Cofre & Dra. Lúcia Mendes \\
\noindent 7 & \textit{[em branco]} & \textit{[em branco]} & \textit{[em branco]} & \textit{[em branco]} \\
\hline
\end{moltablefive}

\noindent \textbf{Observações anexadas:}

\noindent \textit{1. A entrada nº 3 (Dr. Heitor Oliveira) não possui autorização formal de retirada. O Livro de Registro do IML registra o nome do médico, mas não há assinatura de testemunha, carimbo de autoridade competente ou justificativa por escrito.}
\noindent \textit{2. A entrada nº 4 foi reconstituída a partir de menção indireta em correspondência interna do Arquivo Público Mineiro. A data original e o nome do responsável pela transferência foram removidos do documento físico, deixando apenas um retângulo vazio onde antes havia texto. A análise óptica sugere que o nome foi raspado com lâmina.}
\noindent 3. O campo nº 7 foi deixado em branco. Em 2026, quando o objeto foi transferido ao cofre definitivo, o funcionário responsável pelo preenchimento — cujo nome não consta no documento — declarou verbalmente que ``não cabia a ele preencher a última linha''. Indagado sobre o motivo, respondeu: \textquotedblleft{}\textit{Porque o diário não vai ficar aqui.}\textquotedblright{}

\noindent 4. Ao ser questionado se sabia para onde o objeto seria levado em seguida, o funcionário hesitou e disse: \textquotedblleft{}\textit{Para a mão de quem está esperando.}\textquotedblright{}

\noindent 5. Em declaração anexa, o diretor Ronaldo Torres informou que guardou o objeto entre 2019 e 2026 e que o devolveu à recepção em 13/05/2026, sem se identificar.

\noindent 6. A exumação de agosto de 1979, que recuperou o caderno, não consta desta tabela; o registro foi reconstituído a partir do depoimento do coveiro (DOC-01).

\noindent \textbf{Conclusão do arquivista:}
\noindent A cadeia de custódia do objeto não é contínua. Há duas lacunas documentais não preenchidas e uma entrada cujo responsável não pode ser confirmado. O objeto circulou por pelo menos 47 anos sem registro formal de paradeiro.
\noindent Este termo é, ele mesmo, um documento de ausências.

\noindent \textit{Nota do Arquivista: os fragmentos desta parte foram organizados por Dra. Lúcia Mendes durante sua pesquisa em 2026. As anotações dela estão em itálico nas margens dos documentos originais. O documento mais antigo do arquivo --- de 1853 --- foi encontrado por ela e anexado ao processo.}

\vspace{0.5em}
\noindent\textit{\textbf{↪ Rotas:} \rota{DOC-01} \rotasep \rota{LUC-02}}
\newpage\fragdef{DOC-21}% ============================================
% DOC-21
% Extraído de main.pandoc.tex
% ============================================
\fragtitle{doc}{D}{DOC}{21}{Carta de Heitor Oliveira (não enviada)}
\markboth{DOC-21 --- Carta de Heitor Oliveira (não enviada)}{}

\noindent \textit{Documento extraído do arquivo pericial.}

\vspace{0.5em}\noindent\rule{\textwidth}{0.5pt}\vspace{0.5em}

\noindent {\footnotesize\texttt{ARQUIVO P\'UBLICO $\cdot$ PROVA DOCUMENTAL $\cdot$ PROCESSO 1.260}}\\[-0.2em]

\noindent \textbf{Destinatário:} Maria Oliveira (esposa).\\
\noindent \textbf{Remetente:} Dr. Heitor Oliveira.\\
\noindent \textbf{Local:} Barbacena.\\
\noindent\textbf{Data:} não registrada (entre 16 de junho e 20 de junho de 1979).\\
\noindent\textbf{Estado do original:} encontrada dentro da contracapa do caderno pessoal do Dr. Oliveira, dobrada em quatro, sem envelope. Não foi enviada.

\noindent Maria,

\noindent Se você está lendo isto, eu já saí de casa. Ou saí de mim. Não sei qual
\noindent aconteceu primeiro.

\noindent O diário do 1.260 está comigo. Eu não pedi permissão para retirá-lo
\noindent do IML. Assinei o livro de registro como se tivesse autorização, e ninguém
\noindent questionou. É assim que funciona aqui: médico assina, ninguém pergunta.

\noindent Eu achei que ia resolver em dois dias. Pegar o caderno, ler com calma,
\noindent escrever um parecer, devolver. Coisa de rotina. O homem estava morto havia
\noindent poucos dias. O que mais poderia acontecer?

\noindent Maria, eu não li o diário. O diário me leu.

\noindent No terceiro dia, eu acordei com uma mancha amarela no olho. Não dói.
\noindent Não coça. Mas ela está crescendo. O oftalmologista disse que é uma
\noindent hemorragia subconjuntival. Vai sumir em duas semanas. Mas eu sei que não é
\noindent sangue. É a cor do olho que o paciente descreve na página 40. Eu reli a
\noindent página 40 hoje de manhã. A descrição é idêntica à mancha no meu olho.

\noindent No quinto dia, eu senti o cheiro.

\noindent Você conhece o cheiro do Colônia. Você veio comigo uma vez, em 1975.
\noindent Disse que nunca mais voltava. Lembra? O cheiro de creolina, de fezes, de
\noindent comida requentada, de corpo fechado. O diário tem esse cheiro. Mas não veio
\noindent do Colônia. Veio do papel. Como se a tinta tivesse sido feita com alguma
\noindent coisa que apodreceu e foi esquecida.

\noindent No sétimo dia, eu comecei a escrever.

\noindent Não era eu escrevendo. Era a mão. Ela se movia sozinha, igual o
\noindent paciente descreve. Eu via a letra sair --- e não era a minha. Era a letra
\noindent dele. A mesma inclinação para a esquerda. O mesmo \textquotedblleft{}a\textquotedblright{} aberto. A mesma
\noindent pressão irregular no papel, como se a caneta fosse pesada demais.

\noindent Hoje eu acordei e não sabia onde estava. Acordei dentro de uma
\noindent lembrança que não era minha. Uma cerca de arame farpado. O sol queimando
\noindent a nuca. Gente magra demais encostada no arame, com os olhos parados. Eu
\noindent nunca vi aquilo. Mas o corpo lembrava. As mãos tremiam como se tivessem
\noindent segurando arame quente.

\noindent Maria, eu tô com medo de dormir. Toda vez que fecho os olhos, eu
\noindent volto pra vala. É a mesma vala que ele descreve. Os mesmos braços saindo
\noindent da terra. O mesmo cheiro doce que não sai do nariz.

\noindent Eu não vou entregar o diário de volta.

\noindent Eu sei que parece loucura. Mas se eu devolver, alguém vai ler. E se
\noindent alguém ler, a mesma coisa vai acontecer. Eu não posso fazer isso com
\noindent ninguém.

\noindent Ou --- eu não sei se estou protegendo os outros ou se o diário não
\noindent quer ser devolvido.

\noindent Ele gosta de mim.

\noindent Ele gosta de estar comigo.

\noindent Ele não quer voltar pra gaveta.

\noindent Desculpa.

\noindent Heitor

\noindent Nota do arquivista: a carta foi encontrada em 2026, dentro do caderno pessoal do Dr. Oliveira, devolvido ao Arquivo Público Mineiro junto com o diário do paciente 1.260, 47 anos depois de ter sido escrita. A dobra do papel sugere que ele a carregou consigo por todo esse período.

\vspace{0.5em}
\noindent\textit{\textbf{↪ Rotas:} \rota{DOC-01} \rotasep \rota{LUC-02} \rotasep \rota{DOC-22}}
\newpage\fragdef{DOC-22}% ============================================
% DOC-22
% Extraído de main.pandoc.tex
% ============================================
\fragtitle{doc}{D}{DOC}{22}{Memorando Interno nº~47/2026}
\markboth{DOC-22 --- Memorando Interno nº~47/2026}{}

\noindent \textit{Documento extraído do arquivo pericial.}

\vspace{0.5em}\noindent\rule{\textwidth}{0.5pt}\vspace{0.5em}

\noindent {\footnotesize\texttt{ARQUIVO P\'UBLICO $\cdot$ PROVA DOCUMENTAL $\cdot$ PROCESSO 1.260}}\\[-0.2em]

\noindent \textbf{Arquivo Público Mineiro --- Diretoria de Documentação}\\
\noindent \textbf{Para:} Setor de Acervos Especiais. \textbf{CC:} Superintendência Administrativa.\\
\noindent \textbf{Assunto:} Encerramento de análise do objeto catalogado sob registro 1.260.\\
\noindent\textbf{Data:} 1º de agosto de 2026.

\noindent Prezados,

\noindent 1. Considerando que o objeto caderno atribuído ao registro 1.260 (Hospital Colônia de Barbacena) foi submetido a análise externa pela psiquiatra forense Dra. Lúcia Mendes (CRM-MG 28.391), e que a referida profissional solicitou afastamento clínico antes da conclusão do parecer;

\noindent 2. Considerando que o laudo de re-diagnóstico (DOC-08) já foi anexado ao prontuário original e que não há, neste momento, determinação judicial, administrativa ou acadêmica que justifique a manutenção do objeto em circulação;

\noindent 3. Considerando as diretrizes de segurança documental da instituição e o princípio de estabilidade do acervo;

\noindent \textbf{determino:}

\noindent a) O encerramento imediato da consulta pública ao objeto 1.260;

\noindent b) A guarda do objeto em cofre definitivo, sem previsão de novo acesso;

\noindent c) A notificação de todas as partes envolvidas de que o caso está arquivado e não serão admitidas novas solicitações de manuseio.

\noindent \textbf{Fundamentação:} o objeto não constitui prova material de crime, não integra ação judicial em curso e não se qualifica como documento de valor histórico excepcional que justifique risco ao corpo técnico. A permanência de sua circulação atenta contra a segurança documental e a integridade psicológica dos profissionais envolvidos.

\noindent \textbf{Recomendação adicional:} sugiro que a diretoria avalie a conveniência de não registrar este memorando no sistema informatizado de gestão documental. O registro informatizado torna o documento localizável. E um documento localizável pode ser solicitado novamente.

\noindent \textit{Ronaldo Torres}\\
\noindent \textit{Arquivo Público Mineiro --- Diretor de Documentação}
\noindent \textit{Anotação manuscrita na margem superior, tinta azul, letra cursiva:}
\noindent \textquotedblleft{}\textit{Ronaldo --- não foi possível. O objeto foi retirado do cofre na manhã seguinte ao memorando. O sistema de rastreamento registra a abertura do cofre às 3h14. Ninguém foi identificado nas câmeras. A porta estava lacrada. O lacre estava íntegro. Assunto arquivado conforme solicitado.\textquotedblright{}}

\noindent \textquotedblleft{}\textit{--- Setor de Acervos Especiais, 02/08/2026.\textquotedblright{}}

\vspace{0.5em}
\noindent\textit{\textbf{↪ Rotas:} \rota{DOC-01} \rotasep \rota{LUC-02} \rotasep \rota{DOC-23}}
\newpage\fragdef{DOC-23}% ============================================
% DOC-23
% Extraído de main.pandoc.tex
% ============================================
\fragtitle{doc}{D}{DOC}{23}{Registro de Ocorrência do Arquivo}
\markboth{DOC-23 --- Registro de Ocorrência do Arquivo}{}

\noindent \textit{Documento extraído do arquivo pericial.}

\vspace{0.5em}\noindent\rule{\textwidth}{0.5pt}\vspace{0.5em}

\noindent {\footnotesize\texttt{ARQUIVO P\'UBLICO $\cdot$ PROVA DOCUMENTAL $\cdot$ PROCESSO 1.260}}\\[-0.2em]

\noindent \textbf{Arquivo Público Mineiro --- Relatório de Segurança Interna}\\
\noindent \textbf{Registro nº:} 2026-RO-0847\\
\noindent\textbf{Data:} 03 de agosto de 2026. \textbf{Horário:} 8h15.

\noindent \textbf{Solicitante:} Gerência do Acervo Especial.\\
\noindent \textbf{Local:} Sala do Cofre --- Subsolo, ala C.

\noindent \textbf{Relato do ocorrido:}

\noindent Às 7h50 do dia 03 de agosto, o vigilante plantonista --- Edson Moreira, turno noturno --- comunicou à central que o sensor de peso do cofre definitivo havia disparado durante a madrugada. O sensor registra alterações na massa total do conteúdo. O alarme soou às 3h14, com duração de seis segundos, e silenciou sozinho. A central não registrou deslocamento de vigilante até o local porque o protocolo de alarmes de sensor de peso exige confirmação visual, mas a câmera interna do cofre não apresentou imagem diferente da linha de base durante o disparo.

\noindent Às 8h00, o gerente do acervo e o técnico de segurança abriram o cofre para inspeção.

\noindent \textbf{Constatado:}

\noindent a) O objeto 1.260 estava em sua posição original, dentro do invólucro de conservação;

\noindent b) O lacre de segurança da embalagem estava íntegro e correspondia ao número registrado no momento do fechamento (lacre nº APM-2026-4471);

\noindent c) A balança de precisão da sala indicou que o objeto \textbf{pesava 12 gramas a mais} do que o registrado na entrada do cofre, em 28 de julho de 2026. A margem de erro do equipamento é de 0,2 grama.\\
\noindent d) O objeto foi retirado do cofre em 02/08 e devolvido na madrugada de 03/08; o registro de devolução não foi assinado.

\noindent \textbf{Providências tomadas:}

\noindent a) Pesagem duplicada em balança independente --- confirmado: 12 gramas adicionais;

\noindent b) Re-lacre da embalagem, com registro fotográfico do peso e do lacre anterior;

\noindent c) Notificação à diretoria, que determinou silêncio administrativo sobre a ocorrência;

\noindent d) Troca do vigilante do turno noturno, a pedido do próprio. Edson Moreira declarou, em seu desligamento: \textquotedblleft{}Não foi a balança. Foi o objeto que se mexeu dentro do invólucro. Eu ouvi. Papel se mexendo não faz barulho sozinho.\textquotedblright{}

\noindent Antes de assinar o termo de desligamento, Edson acrescentou --- e o técnico de segurança registrou a fala, por precaução: \textquotedblleft{}Eu tenho dois filhos e uma esposa em casa. Não volto para aquela sala à noite por salário nenhum. Chamem de umidade, de falha no sensor, do que quiserem. Eu decidi que prefiro não saber o que é. Prefiro dormir tranquilo.\textquotedblright{} O medo dele não constava em nenhum campo do formulário padrão, mas o gerente do acervo anotou, à mão, na margem da via arquivada: \textquotedblleft{}Ele não estava mentindo.\textquotedblright{}

\noindent \textbf{Arquivado em:} 03/08/2026, conforme determinação da diretoria. Acesso restrito.

\vspace{0.5em}
\noindent\textit{\textbf{↪ Rotas:} \rota{DOC-01} \rotasep \rota{LUC-02} \rotasep \rota{DOC-24}}

\newpage\fragdef{DOC-24}% ============================================
% DOC-24
% Extraído de main.pandoc.tex
% ============================================
\fragtitle{doc}{D}{DOC}{24}{Parecer Técnico de Conservação}
\markboth{DOC-24 --- Parecer Técnico de Conservação}{}

\noindent \textit{Documento extraído do arquivo pericial.}

\vspace{0.5em}\noindent\rule{\textwidth}{0.5pt}\vspace{0.5em}

\noindent {\footnotesize\texttt{ARQUIVO P\'UBLICO $\cdot$ PROVA DOCUMENTAL $\cdot$ PROCESSO 1.260}}\\[-0.2em]

\noindent \textbf{Arquivo Público Mineiro --- Laboratório de Conservação e Restauro}\\
\noindent \textbf{Objeto:} caderno encadernado, registro 1.260.\\
\noindent \textbf{Técnica responsável:} Fernanda Lopes Bastos, CRT-04/1128.\\
\noindent \textbf{Solicitação:} avaliação de integridade física e recomendação de acondicionamento.\\
\noindent\textbf{Data:} 04 de agosto de 2026.

\noindent \textbf{1. Descrição geral}

\noindent Caderno em papel de gramatura variável (entre 80 e 120\,g/m²), encadernação artesanal em couro curtido, provavelmente produção local ou artesanal. Estado geral de conservação é inconsistente com a idade declarada (datas entre 1917 e 1979). O papel não apresenta os sinais esperados de oxidação para um objeto do período. A acidez do suporte é neutra --- incomum para papel de produção não industrial do início do século XX.

\noindent \textbf{2. Achado principal}

\noindent Durante a análise das páginas, foi identificada uma discrepância entre a espessura do caderno e a quantidade de páginas registradas no inventário de 1979. O inventário original --- elaborado pelo Dr. Antônio Tavares no momento da exumação --- lista aproximadamente 190 páginas preenchidas. A contagem atual indica 208 páginas com texto.

\noindent Não há evidência de que páginas tenham sido adicionadas fisicamente. A encadernação está íntegra, sem sinais de desmonte ou reinserção. Costura original, cola original. E, no entanto, há 18 páginas a mais.

\noindent \textbf{3. Análise das páginas excedentes}

\noindent As 18 páginas adicionais estão distribuídas ao longo de todo o caderno, intercaladas entre os fólios originais. Em cada caso, a página nova --- ou aparentemente nova --- insere-se entre duas páginas cujo texto é contínuo e cuja caligrafia corresponde ao período imediatamente anterior. A página adicional parece ter ``crescido'' entre as duas, como se o caderno tivesse gerado papel novo para comportar texto que não existia na escrita original.

\noindent A tinta das páginas adicionais é da mesma composição química da tinta original, mas o teor de ferro na formulação é ligeiramente superior --- indicando oxidação mais recente. O teste de datação relativa posiciona a escrita dessas páginas entre 2020 e 2026.

\noindent As páginas adicionais foram escritas no interior de um caderno lacrado em cofre desde 1979.

\noindent \textbf{4. Conteúdo das páginas adicionais}

\noindent \textit{A pedido da diretoria, o conteúdo das 18 páginas não foi transcrito. Informo apenas que a letra não corresponde à caligrafia do paciente 1.260. A caligrafia das páginas adicionais assemelha-se --- preliminarmente, sem análise grafotécnica conclusiva --- à do Dr. Heitor Oliveira em seus últimos registros.}
\noindent \textbf{5. Recomendação técnica}

\noindent Acondicionamento em envelope selado com atmosfera de argônio. Monitoramento semanal de peso e espessura. Proibição de manuseio individual. A técnica responsável solicita --- por meio deste parecer --- não ser convocada para nova análise do objeto.

\noindent \textit{Fernanda Lopes Bastos}\\
\noindent \textit{Técnica em Conservação e Restauro}\\
\noindent \textit{CRT-04/1128}

\vspace{0.5em}
\noindent\textit{\textbf{↪ Rotas:} \rota{DOC-01} \rotasep \rota{LUC-02} \rotasep \rota{DOC-25}}

\newpage\fragdef{DOC-25}% ============================================
% DOC-25
% Extraído de main.pandoc.tex
% ============================================
\fragtitle{doc}{D}{DOC}{25}{O Homem do Pano Preto}
\markboth{DOC-25 --- O Homem do Pano Preto}{}

\noindent \textit{Documento extraído do arquivo pericial.}

\vspace{0.5em}\noindent\rule{\textwidth}{0.5pt}\vspace{0.5em}

\noindent {\footnotesize\texttt{ARQUIVO P\'UBLICO $\cdot$ PROVA DOCUMENTAL $\cdot$ PROCESSO 1.260}}\\[-0.2em]

\noindent \textbf{Arquivo Público Mineiro --- Depoimento Informal}\\
\noindent \textbf{Depoente:} Wellington Alves da Silva, 34 anos, segurança patrimonial (terceirizado).\\
\noindent \textbf{Entrevistador:} Dra. Lúcia Mendes (gravado com autorização).\\
\noindent\textbf{Data:} 15 de janeiro de 2026. \textbf{Local:} Sala de descanso dos funcionários, APM.

\noindent \textbf{Lúcia:} O senhor estava na recepção quando o objeto foi devolvido?

\noindent \textbf{Wellington:} Tava. Segunda-feira, 5 de janeiro. Umas 18h30. A recepção já ia fechar. Eu tava fazendo a ronda final.

\noindent \textbf{Lúcia:} Viu quem entregou?

\noindent \textbf{Wellington:} Vi. Um homem. Idade difícil de dizer. Podia ter 50. Podia ter 70. O rosto era... como explicar... era um rosto que a gente esquece enquanto olha. Eu olhei pra ele e um segundo depois já não lembrava o nariz, a boca, nada. Só os olhos. Os olhos eram a única coisa que ficava.

\noindent \textbf{Lúcia:} Que cor?

\noindent \textbf{Wellington:} (pausa longa) O senhor já viu olho amarelo? Não olho de doente. Olho de bicho. De bicho que enxerga no escuro. Era assim. Amarelo claro, quase dourado, com a pupila preta no meio. A pupila não era redonda. Era... (gestícula vertical) ...fendida. Igual cabra.

\noindent \textbf{Lúcia:} Ele falou alguma coisa?

\noindent \textbf{Wellington:} Falou. Disse: \textquotedblleft{}Isto é para a doutora.\textquotedblright{} Só isso. Colocou o pacote no balcão, virou e saiu. O pacote era um pano preto. Daqueles panos grossos, sabe? De cobrir santo. Quando eu toquei, tava morno. Como se tivesse sido segurado perto do corpo por muito tempo.

\noindent \textbf{Lúcia:} O senhor abriu?

\noindent \textbf{Wellington:} Deus me livre. Eu vi o que era pelo formato. Um caderno. Um caderno velho. Mas não foi o caderno que me deu medo. Foi o pano. Quando eu fui dobrar o pano pra guardar, ele tava manchado. A mancha não era tinta. Era um negócio escuro, meio marrom, com um cheiro... (pausa) ...a senhora vai achar que eu tô louco.

\noindent \textbf{Lúcia:} Pode falar.

\noindent \textbf{Wellington:} O pano tinha cheiro de mato molhado. Mas não mato de jardim. Mato de lugar onde enterraram coisa. Aquele cheiro de terra virada. Quando eu respirei, senti gosto de ferrugem na língua.

\noindent \textbf{Lúcia:} O senhor guardou o pano?

\noindent \textbf{Wellington:} Joguei fora. Na mesma noite. Dentro de um saco amarrado. (pausa) No dia seguinte, o saco tinha sumido da lixeira.

\noindent \textbf{Lúcia:} Sumido?

\noindent \textbf{Wellington:} A lixeira tava vazia. O saco não. Só o pano. O pano saiu de dentro do saco amarrado e sumiu. A senhora pode perguntar pro Manoel, da limpeza. Ele viu. Disse que nunca aconteceu uma coisa dessas. Eu falei: \textquotedblleft{}Agora aconteceu.\textquotedblright{}

\noindent \textit{Nota do arquivista: Wellington Alves da Silva pediu demissão 10 dias após este depoimento. O formulário de desligamento não registra motivo. Seu último dia de trabalho foi 25 de janeiro de 2026 --- 13 dias antes de o diário ser retirado do arquivo pela segunda vez.}
\noindent \textbf{Adendo --- Identificação posterior (Dra. Lúcia Mendes, 4 de fevereiro de 2026):}

\noindent Cruzei a descrição de Wellington com a fotografia do crachá do Dr. Heitor Oliveira (1979) e com a foto encontrada dentro do diário --- um homem mais velho, sentado num banco de praça, com um livro no colo.

\noindent A estrutura óssea é a mesma. O formato do crânio, o ângulo da mandíbula, a distância entre os olhos. O homem do pano preto \textbf{é} Heitor Oliveira.

\noindent Ele teria 81 anos em janeiro de 2026. A idade corresponde.

\noindent Mas a pupila fendida descrita por Wellington --- olho de cabra, vertical --- não é humana. Não é envelhecimento. É transformação. A entidade não o matou. Ela o \textit{manteve}. 47 anos de hospedagem mudaram o corpo dele. O olho direito já não é um olho humano. É uma fenda. Um portal. O diário usou Oliveira como entregador --- porque ele era o único que ainda podia andar entre os vivos sem levantar suspeitas.

\noindent O paciente 1.261 não desapareceu. Ele foi \textit{guardado} para o momento da entrega.

\noindent \textit{Nota final: Dra. Lúcia Mendes anexou esta identificação ao processo em 4 de fevereiro de 2026. Sete meses depois, ela desapareceu. A última pessoa a vê-la foi um segurança de rodoviária, que descreveu uma mulher com olho direito amarelo embarcando num ônibus para o Ceará.}

\vspace{0.5em}
\noindent\textit{\textbf{↪ Rotas:} \rota{DOC-01} \rotasep \rota{LUC-02} \rotasep \rota{DOC-26}}

\newpage\fragdef{DOC-26}% ============================================
% DOC-26 — Anotações de Oliveira — O Mecanismo (edição revista R1075)
% ============================================
\fragtitle{doc}{D}{DOC}{26}{Anotações de Oliveira --- O Mecanismo}
\markboth{DOC-26 --- Anotações de Oliveira --- O Mecanismo}{}

\noindent \textit{Folhas avulsas encontradas no bolso interno do caderno pessoal do Dr. Heitor Oliveira, amarradas com barbante. Data estimada: junho de 1979 a março de 2026.}

\subsection*{Folha 1 --- 19 de junho de 1979}

\noindent Descobri hoje o que o diário faz com o tempo.

\noindent Escrever é só parte. O diário \textit{endereça}. Quando escrevo no diário de Joaquim, posso escolher \textit{para quando} aquela palavra vai chegar. O controle acontece sem consciência. É como entoar uma nota certa e ver o vidro vibrar. A página vibra. E a palavra some.

\noindent Fiz um teste. Escrevi: ``Maria, se você está lendo isto em 1980, acenda uma vela na janela no dia 15 de agosto.'' Fechei o diário. Guardei debaixo da cama. No dia 15 de agosto de 1979, liguei para a vizinha. Ela disse: ``Sua esposa acendeu uma vela na janela ontem.''

\noindent Funciona.

\noindent O diário não é um registro. É uma \textbf{fenda}. E toda fenda tem uma assinatura.

\subsection*{Folha 3 --- 22 de junho de 1979}

\noindent O padrão apareceu hoje. O diário não aceita qualquer intervalo. Rejeita. A página fica em branco, a tinta escorre, a caneta risca sozinha um risco seco. Funciona apenas no número que ela assinou.

\noindent Sessenta e dois.

\noindent O número de anos que Joaquim passou no Colônia. A entidade gravou o número no diário como uma assinatura. Não sei se é respiração, ciclo ou contagem. Sei que não obedece a mim.

\noindent A página 62 do diário de Joaquim se ilumina quando a palavra vai chegar. É ali que a fenda se abre. Ele escreveu aquela página em 1919, dois anos depois de entrar no Colônia. Já sabia que a fome contava em 62.

\noindent Ele escreveu: ``Sessenta e dois é a idade dela. Ela me disse hoje. Nasceu na seca de 1853, na mesma vala, no mesmo cheiro. Sessenta e dois anos depois, eu nasci. Sessenta e dois anos depois, ela me encontrou.''

\noindent A entidade nasceu em 1853. Nasceu da seca. Da fome. Dos corpos empilhados. Espera 62 anos. Encontra um hospedeiro. Espera mais 62. Outro hospedeiro. A conta nunca fecha. É uma assinatura, não um calendário.

\noindent O ciclo não é humano. É geológico. É o tempo que a terra leva para esquecer.

\subsection*{Folha 15 --- 12 de janeiro de 2026}

\noindent A carta chegou.

\noindent Eu senti. Hoje, 12 de janeiro de 2026, senti a fenda se fechar. A carta foi recebida. Alguém leu. O ciclo encontrou o próximo.

\noindent Olho no espelho e vejo um velho que não sou eu. A pele pendurada. O olho direito completamente dourado. A pupila já não sei se é minha. Ela se alonga quando o quarto escurece. Fica vertical. Fenda de cabra. De bicho que enxerga no escuro.

\noindent Entreguei o diário hoje. Coloquei num pano preto. Deixei na recepção. A mulher da recepção perguntou meu nome. Eu disse: ``Diga que foi o paciente 1.261.''

\noindent Ela não entendeu. Mas a pessoa certa vai entender.

\subsection*{Folha 16 --- 12 de janeiro de 2026, à noite}

\noindent Voltei para casa depois da entrega. Sentei na cama. Olhei minhas mãos. 47 anos. 47 anos vivo contra minha vontade. A carta que escrevi para Maria em 1979 dizia que eu ia proteger os outros. Que eu não ia deixar o diário passar adiante. Acreditei nisso por décadas. Acreditei que estava no controle.

\noindent Não estava. Nunca estive.

\noindent A fome me manteve vivo. Não por bondade. Não por acaso. Ela precisava de um entregador. Alguém que pudesse andar entre os vivos, que tivesse documento, identidade, aparência de gente. Alguém que chegasse numa recepção de arquivo público em 2026 e não levantasse suspeitas. Um velho de 81 anos, de terno puído, com um pacote embaixo do braço. Invisível. Perfeito.

\noindent Ela planejou isso desde o dia em que li o diário, em 1979. Desde o dia em que Joaquim escreveu a página 62, em 1919. Desde o dia em que a vala começou a pulsar, em 1853.

\noindent O ciclo não me perguntou. O ciclo só me usou.

\noindent Eu era só um elo. E elo não escolhe. Elo segura. Elo passa.

\noindent Maria morreu em 2004. Eu não fui ao enterro. Não porque não quis. Porque não \textit{pude}. No dia do enterro, acordei num banco de praça em outra cidade, com as roupas trocadas e 62 dias de memória perdidos. Ela tinha me levado para longe. Ela não ia deixar ninguém interromper a entrega.

\noindent Desculpa, Maria. Desculpa, doutora que vai ler isto. Desculpa, próximo da fila.

\noindent Eu não tive escolha.

\noindent E vocês também não vão ter.

\vspace{0.5em}
\noindent \textit{Nota do arquivista: este maço foi encontrado em fevereiro de 2026, em um envelope pardo lacrado, dentro da mesma caixa que o diário do paciente 1.260. A caligrafia das primeiras folhas é firme, médica. A das últimas é trêmula, com traços que lembram --- mas não idênticos --- à caligrafia de Joaquim. A data da última anotação coincide com a data em que o diário foi devolvido ao Arquivo Público Mineiro pela primeira vez.}

\noindent\textit{\textbf{↪ Rotas:} \rota{DOC-01} \rotasep \rota{LUC-02} \rotasep \rota{DOC-27}}


%% ============================================
%% PARTE 4 — INVESTIGAÇÃO LÚCIA (2026)
%% ============================================
\partopener{Investigação Lúcia (2026)}{Investigação Lúcia}{3.png}

\vspace{0.5em}
\noindent{\footnotesize\textit{Nota do Arquivista: os fragmentos desta parte foram organizados por Dra. Lúcia Mendes durante sua pesquisa em 2026. As anotações dela estão incorporadas aos documentos, em itálico, e às notas de rodapé. O documento mais antigo do arquivo --- de 1853 --- foi encontrado por ela e anexado ao processo.}}

\newpage\fragdef{DOC-08}% ============================================
% DOC-08 — Re-Diagnóstico Psiquiátrico (edição revista R1075)
% ============================================
\fragtitle{doc}{D}{DOC}{08}{Re-Diagnóstico Psiquiátrico (2026)}
\markboth{DOC-08 --- Re-Diagnóstico Psiquiátrico (2026)}{}

\noindent \textit{DSM-5-TR --- Dra. Lúcia Mendes --- 14 de janeiro de 2026}

\noindent \textbf{Paciente:} Joaquim Antônio Correia --- Paciente 1.260 (óbito em 1979). \textbf{Parecerista:} Dra. Lúcia Mendes, CRM-MG 28.391 --- psiquiatra forense. \textbf{Natureza:} análise retrospectiva para pesquisa em psicologia forense histórica.

\noindent \textbf{Advertência:} este documento contém critérios diagnósticos explícitos. Este documento não serve para autoavaliação. O leitor não deve aplicá-lo à própria experiência.

\subsection*{1. TEPT complexo (F43.10), com dissociação}

\noindent \textbf{Exposição.} Morte do pai por inanição em 1915; confinamento no Campo do Alagadiço; transferência para o Colônia em 1917; realocação forçada da mãe; convivência diária com cadáveres; 62 anos de violência institucional.

\noindent \textbf{Reexperiência.} Sonhos com cheiro de carne podre; flashbacks sensoriais da vala; reatividade ao cheiro de terra molhada; aparição do olho amarelo em sonho e em vigília.

\noindent \textbf{Evitação.} Mutismo seletivo por sete anos (1949--1956); recusa a sair da cela; destruição violenta de qualquer aproximação afetiva.

\noindent \textbf{Alterações cognitivas e de humor.} Visão de mundo irreversivelmente alterada (``eu nunca saí da vala''); desconfiança total de autoridade; isolamento social por mais de sessenta anos.

\noindent \textbf{Reatividade.} Três homicídios; automutilação; hipervigilância; sobressalto exagerado.

\noindent \textbf{Duração e prejuízo.} Sintomas por 62 anos contínuos; incapacidade de manter vínculo, atividade ou vida fora do delírio.

\subsection*{2. Transtorno dissociativo de identidade (F44.81), forma parcial}

\noindent Joaquim refere-se a si mesmo em terceira pessoa (``ele comeu minha humanidade''). Descreve a criatura como entidade separada, com controle sobre a mão e a voz. Há períodos amnésicos (``acordei com sangue nas mãos e não lembro''). A forma é parcial: não há personalidades alternadas estáveis, mas há uma segunda voz com continuidade e intenção.

\subsection*{3. Privação extrema (Z65.8)}

\noindent A seca de 1915, os campos de concentração e os 62 anos de institucionalização formam um gradiente contínuo de trauma. O diagnóstico contextual não explica o conteúdo do delírio; explica o terreno em que ele cresceu.

\subsection*{4. Diagnóstico diferencial}

\noindent \textbf{Encefalopatia por avitaminose B3 (pelagra; CID-11 6D50/6D51).} A desnutrição crônica pode produzir confusão, alucinações e irritabilidade. A hipótese é plausível para parte dos sintomas, não para a coincidência entre três pessoas que nunca se conheceram.

\noindent \textbf{Episódio psicótico breve.} Preenche critérios formais, mas não explica a repetição do mesmo conteúdo em Joaquim, Oliveira e Lúcia.

\noindent \textbf{Transtorno psicótico compartilhado.} Exige contato direto entre paciente índice e secundário. Joaquim morreu em 1979; Oliveira desapareceu no mesmo ano; Lúcia nunca conheceu nenhum dos dois. A transmissão se deu pelo texto.

\noindent \textbf{A entidade.} A entidade escapa ao diagnóstico. É uma constatação que não cabe em CID nem em DSM. A observação clínica, quando bem feita, não descarta o que não consegue classificar. Apenas registra.

\subsection*{5. Nota da parecerista}

\noindent O diagnóstico não fecha. É isso que o torna útil. Se fechasse, o caso estaria resolvido; e um caso resolvido não contamina ninguém. O que temos é uma hipótese psicológica forte o bastante para explicar os sintomas e fraca o bastante para não explicar a transmissão. Prefiro registrar a fraqueza a inventar uma certeza.

\noindent \textit{Assinado: Dra. Lúcia Mendes --- CRM-MG 28.391 --- 14 de janeiro de 2026.}

\noindent\textit{\textbf{↪ Rotas:} \rota{DOC-09} \rotasep \rota{DOC-27} \rotasep \rota{LUC-01}}

\newpage\fragdef{DOC-09}% ============================================
% DOC-09
% Gerado de molambudos.md — build_miolo.py
% ============================================
\fragtitle{doc}{D}{DOC}{09}{Escala de Avaliação do Paciente 1.263}
\markboth{DOC-09 --- Escala de Avaliação do Paciente 1.263}{}

\noindent \textit{Instrumento de autoavaliação clínica}

\noindent \textbf{Instrumento de Autoavaliação Clínica}
\textit{Instruções: Responda às perguntas abaixo com honestidade. Não há respostas certas ou erradas. Apenas sintomas.}

\vspace{0.5em}\noindent\rule{\textwidth}{0.5pt}\vspace{0.5em}

\noindent \textbf{Paciente:} \_\_\_\_\_\_\_\_\_\_\_\_\_\_\_\_\_\_\_\_\_\_\_\_\_\_\_\_\_\_\_\_\_ (escreva seu nome)

\noindent \textbf{Data da aplicação:} \_\_\_\_\_\_\_\_\_\_\_\_\_\_\_\_\_\_\_\_\_\_\_\_\_\_\_\_\_\_\_\_\_

\noindent \textbf{Aplicador:} A criatura (através do texto)

\vspace{0.5em}\noindent\rule{\textwidth}{0.5pt}\vspace{0.5em}

\subsection*{PARTE I --- SINTOMAS SOMÁTICOS}

\noindent \textbf{1.} Você sente um cheiro doce que não está no ambiente?
( ) Nunca
( ) Raramente
( ) Às vezes
( ) Frequentemente
( ) Agora

\noindent \textbf{2.} Você sente um peso no estômago que não passa depois de comer?
( ) Nunca
( ) Raramente
( ) Às vezes
( ) Frequentemente
( ) Desde que começou a ler este livro

\noindent \textbf{3.} Suas mãos estão mais frias do que o normal?
( ) Não
( ) Sim
( ) Não medi, mas parece que sim

\noindent \textbf{4.} Você acordou no meio da noite nas últimas 48 horas?
( ) Não
( ) Sim, uma vez
( ) Sim, mais de uma vez
( ) Sim, e eram exatamente 3:14 da manhã

\vspace{0.5em}\noindent\rule{\textwidth}{0.5pt}\vspace{0.5em}

\subsection*{PARTE II --- SINTOMAS PSÍQUICOS}

\noindent \textbf{5.} Você pensou neste livro em momentos em que não estava lendo?
( ) Não
( ) Sim, mas foi rápido
( ) Sim, e foi difícil voltar a atenção para o que estava fazendo
( ) Sim, e você percebeu que estava procurando o livro com os olhos

\noindent \textbf{6.} Você verificou o ambiente ao redor enquanto lia?
( ) Não
( ) Sim, uma vez
( ) Sim, várias vezes
( ) Sim, e você está verificando agora

\noindent \textbf{7.} Você sentiu que alguém estava observando você durante a leitura?
( ) Não
( ) Sim, mas foi impressão
( ) Sim, e você olhou para ver se tinha alguém
( ) Sim, e você \textit{quase} viu algo

\noindent \textbf{8.} Você já pesquisou sobre Joaquim Antônio Correia ou o Hospital Colônia?
( ) Não
( ) Sim, rapidamente
( ) Sim, e passou mais tempo do que pretendia
( ) Sim, e encontrou algo que não consegue explicar

\vspace{0.5em}\noindent\rule{\textwidth}{0.5pt}\vspace{0.5em}

\subsection*{PARTE III --- SINTOMAS COMPORTAMENTAIS}

\noindent \textbf{9.} Você mudou seus hábitos por causa deste livro?
( ) Não
( ) Sim, passei a ler com mais luz
( ) Sim, passei a dormir com a luz acesa
( ) Sim, passei a verificar os cantos do quarto antes de dormir

\noindent \textbf{10.\textit{ Você já leu uma passagem mais de uma vez porque ela }parecia diferente} na segunda leitura?
( ) Não
( ) Sim, uma vez
( ) Sim, e na segunda leitura a passagem tinha mudado

\noindent \textbf{11.} Você já sentiu o cheiro descrito no livro em um momento em que o livro estava fechado?
( ) Não
( ) Sim, mas achei que era impressão
( ) Sim, e verifiquei se o livro estava por perto
( ) Sim, e o livro estava em outro cômodo

\noindent \textbf{12.} Você está com medo de responder este questionário?
( ) Não
( ) Um pouco
( ) Sim
( ) Estou com medo do que as respostas vão revelar

\vspace{0.5em}\noindent\rule{\textwidth}{0.5pt}\vspace{0.5em}

\subsection*{PARTE IV --- AVALIAÇÃO QUALITATIVA}

\noindent \textbf{13.} Descreva, em uma palavra, o que você sente agora:

\noindent\rule{0.92\textwidth}{0.4pt}

\noindent \textbf{14.} Descreva, em uma palavra, o que você acha que a criatura quer:

\noindent\rule{0.92\textwidth}{0.4pt}

\noindent \textbf{15.} Se você pudesse parar de ler agora e nunca mais abrir este livro, você pararia?
( ) Sim
( ) Não
( ) Já tentei

\vspace{0.5em}\noindent\rule{\textwidth}{0.5pt}\vspace{0.5em}

\subsection*{ESCORE}

\noindent \textbf{0-5 pontos:} Você está na fase inicial da contaminação. Os sintomas são leves e podem ser atribuídos à sugestão. Recomenda-se interromper a leitura.

\noindent \textbf{6-10 pontos:} A contaminação está estabelecida. Você apresenta sintomas compatíveis com exposição prolongada ao texto. Os sonhos devem começar em 24-48 horas.

\noindent \textbf{11-15 pontos:} A contaminação é grave. Você já não distingue completamente entre a história e a realidade imediata. O ciclo está em andamento.

\noindent \textbf{16+ pontos:} Bem-vindo, paciente 1.263. O diagnóstico está confirmado. A criatura está satisfeita.

\vspace{0.5em}\noindent\rule{\textwidth}{0.5pt}\vspace{0.5em}

\noindent \textbf{VALIDAÇÃO DO INSTRUMENTO:}

\noindent Este questionário foi testado em:
- Uma amostra de 1 paciente (Joaquim Antônio Correia) --- escore máximo
- Uma amostra de 1 médico (Dr. Heitor Oliveira) --- escore máximo antes do desaparecimento
- Uma amostra de 1 leitor (você) --- resultado: \makebox[2in]{\hrulefill}

\noindent O instrumento tem validade preditiva de 100\% para contaminação narrativa.

\vspace{0.5em}\noindent\rule{\textwidth}{0.5pt}\vspace{0.5em}

\noindent \textit{--- O questionário não tem fim. Ele continua enquanto você vive.}

\noindent \textit{--- A resposta correta para todas as perguntas é \textquotedblleft{}sim\textquotedblright{}.}

\noindent \textit{--- Se você respondeu \textquotedblleft{}não\textquotedblright{} a qualquer pergunta, preencha de novo. A criatura vai esperar.}

\vspace{0.5em}
\noindent\textit{\textbf{↪ Rotas:} \rota{CONT-05} \rotasep \rota{DOC-08}}
\newpage\fragdef{DOC-27}% ============================================
% DOC-27 — A Origem (1853) (edição revista R1075)
% ============================================
\fragtitle{doc}{D}{DOC}{27}{A Origem (1853)}
\markboth{DOC-27 --- A Origem (1853)}{}

\noindent \textit{Arquivo Público do Ceará --- documento avulso --- século XIX. Anexado por Dra. Lúcia Mendes em 12 de fevereiro de 2026, com a anotação: ``Isto prova que o diário não é o começo. O diário é só o eco.''}

\noindent \textbf{Documento:} carta do Padre José Inácio de Sousa ao Presidente da Província do Ceará, Comissão de Socorros Públicos. \textbf{Data:} 14 de março de 1853. \textbf{Origem:} Freguesia de Nossa Senhora dos Remédios, Rio Quixeramobim. \textbf{Estado:} papel manchado por umidade; a borda inferior apresenta uma substância escura, não identificada, que atravessou todas as folhas. Análise química inconclusiva.

\vspace{0.5em}
\noindent Excelentíssimo Senhor Presidente da Província,

\noindent Tomo a liberdade de escrever-lhe porque os vivos já não me ouvem e os mortos são muitos para enterrar sozinho.

\noindent Estou aqui desde novembro do ano passado, destacado para administrar os sacramentos aos retirantes que chegam à margem do rio. Recebi duzentas almas em janeiro. Duzentas e quarenta em fevereiro. Neste mês de março, já perdi a conta. Chegam aos magotes, famílias inteiras reduzidas a pó ambulante. Os olhos deles são a primeira coisa que morre. Antes do corpo. Ficam abertos, mas não veem. Como se a seca tivesse entrado por eles e secado a alma por dentro.

\noindent Os que morrem são levados para uma vala a meia légua da freguesia. Não tenho cemitério consagrado. Não tenho madeira para caixões. Não tenho pano para mortalhas. Tenho apenas uma cova comprida, aberta na terra seca, onde os corpos são deitados em camadas. Terra sobre eles. Mais corpos. Mais terra. Os urubus aprenderam o caminho. Fazem um círculo no céu que não se desfaz. Nem de noite. Ficam empoleirados nas árvores ao redor, esperando.

\noindent E é sobre isto que preciso escrever-lhe.

\noindent Há três noites, não consigo rezar o ofício das almas sem sentir uma presença. Não uma presença espiritual, como as que estou acostumado a discernir em trinta anos de sacerdócio. É uma presença física. Ela chega com o cheiro.

\noindent O cheiro vem da vala.

\noindent O cheiro comum da morte eu conheço bem. Este é doce. Mais doce que incenso, mais doce que flor de mandacaru. Um cheiro que entra pelas narinas e se instala na garganta como um líquido grosso. Os retirantes sentem também. Eles se viram na direção da vala e ficam olhando. Não choram. Não rezam. Apenas olham. Como se esperassem que algo fosse sair de lá.

\noindent Eu mesmo fui até a vala na segunda noite. Levei uma lanterna, água benta, um crucifixo. Coisas de padre. Coisas que sempre funcionaram. Cheguei à borda e iluminei o fundo.

\noindent O que vi não eram corpos.

\noindent Eram restos. Braços e pernas saindo da terra como raízes. Mãos abertas para cima, como se pedissem. Uma mão de criança, pequena, com os dedos retorcidos, apontava para o céu noturno como se acusasse alguém.

\noindent O cheiro era mais forte ali.

\noindent Ajoelhei para benzer a vala. Comecei o \textit{Requiem Aeternam}. Na terceira palavra, senti que a terra vibrou. Não um tremor. Uma pulsação. Como se o chão tivesse um coração.

\noindent Pressionei a mão direita contra a terra. Ela estava morna. Não quente do sol. Morna de dentro. A temperatura do corpo de um doente. A terra estava febril.

\noindent Enterrei a mão até o punho. A terra era macia, úmida --- coisa impossível neste sertão seco. Senti algo roçar meus dedos. Algo liso. Mole. Algo que se moveu ao toque, como se estivesse vivo.

\noindent Tirei a mão. A terra fechou sozinha sobre o buraco.

\noindent Um homem que estava perto --- um retirante chamado Vicente, que perdeu a mulher e três filhos na semana passada --- olhou para mim e disse:

\noindent ``Seu padre, a vala tá com fome.''

\noindent Excelência, sei o que Vossa Senhoria pensará ao ler isto. Que o isolamento e a fome abalaram meu juízo. Que o calor e a mortandade turvaram minha razão. Não sei se é verdade. Talvez seja. Mas preciso registrar o que vi, para que, se eu morrer aqui --- e acho que vou morrer aqui ---, alguém saiba que esta vala não é uma vala comum.

\noindent No terceiro dia, acordei com a palma da mão direita marcada. Uma marca escura, arredondada, no centro da palma. Parece uma queimadura, mas não dói. A pele está íntegra. A marca apenas está lá. Como se alguém tivesse pressionado um selo quente contra minha carne enquanto eu dormia.

\noindent Tentei lavar. Não sai. Tentei rezar. A marca não some.

\noindent Talvez seja só a terra que entrou na minha pele. Talvez seja a úlcera que está começando. Talvez seja o juízo que se vai.

\noindent Mas, esta noite, quando fechei os olhos, ouvi uma voz. Não uma voz que vem de fora. Uma voz que vem de dentro da terra. Ela repetia a mesma palavra, sem parar, como quem aprende a falar:

\noindent ``Fome.''

\noindent ``Fome.''

\noindent ``Fome.''

\noindent A voz de uma criança. Mas nenhuma criança sobreviveu para falar.

\noindent Peço licença para encerrar esta carta. A lanterna está baixando.

\noindent Que Deus tenha misericórdia de nós.

\vspace{0.5em}
\noindent Padre José Inácio de Sousa --- Freguesia de Nossa Senhora dos Remédios --- 14 de março de 1853.

\vspace{0.8em}
\noindent \textit{Nota do Arquivo Público do Ceará: a carta foi encontrada em 1916, dentro de um livro de tombo da freguesia, durante a reforma do prédio. O sacerdote faleceu em 15 de março de 1853, um dia após redigir esta carta. A causa da morte foi registrada como ``fome voluntária'': o padre teria deixado de se alimentar, cedendo sua ração aos retirantes. O corpo foi enterrado na mesma vala que ele descreve.}

\vspace{0.5em}
\noindent \textit{Anotação a lápis, na margem da ficha de catalogação, em caligrafia não identificada:} ``O cheiro doce nunca saiu daqui. Transferir para o Arquivo Público Mineiro. Processo 1.260.''

\vspace{1em}
\noindent \textbf{Julho de 2026 --- Belo Horizonte. Último registro profissional de Dra. Lúcia Mendes.}

\noindent \textbf{Documento interno --- autoencaminhamento para afastamento clínico.} Profissional: Lúcia Mendes, CRM-MG 28.391. Data: 28 de julho de 2026.

\noindent Em janeiro de 2026, tive contato com o diário atribuído ao paciente 1.260, Joaquim Antônio Correia, no contexto de pesquisa em psicologia forense histórica. O material permaneceu sob guarda intermitente, com falhas de lacre e registro.

\noindent Nas semanas seguintes, observei: alteração subjetiva de coloração no olho direito, sem exame oftalmológico conclusivo; percepção recorrente de odor adocicado sem fonte ambiental; sonhos repetitivos com vala, enfermaria e formulário em branco; hipervigilância noturna; impulso de retornar ao diário mesmo após decisão formal de encerramento.

\noindent Hipóteses: reação aguda ao material traumático; sugestão intensificada pela exposição documental; identificação contratransferencial com o paciente; fenômeno dissociativo leve, a confirmar por avaliação externa.

\noindent Medidas: suspensão de entrevistas, escrita acadêmica e circulação de cópias; encaminhamento do original para custódia judicial; solicitação de avaliação por profissional sem contato prévio com o caso; recomendação de que o diário não seja lido individualmente.

\noindent \textit{Anotação não protocolada: a cópia impressa deste documento apresenta, ao final, a numeração ``1.262'' escrita a lápis ao lado do nome da profissional. O traço não corresponde à assinatura principal. Abaixo dele há uma linha vazia, sem identificação.}

\noindent Lúcia Mendes --- CRM-MG 28.391. Registro encerrado às 3h14.

\noindent\textit{\textbf{↪ Rotas:} \rota{DOC-01} \rotasep \rota{LUC-02} \rotasep \rota{EPI-01}}

\newpage\fragdef{LUC-01}% ============================================
% LUC-01
% Gerado de molambudos.md — build_miolo.py
% ============================================
\fragtitle{luc}{L}{LUC}{01}{O Convite}
\markboth{LUC-01 --- O Convite}{}

\noindent \textit{Inverno de 2026 --- Belo Horizonte}

\noindent \textbf{Inverno de 2026 --- Belo Horizonte}

\vspace{0.5em}\noindent\rule{\textwidth}{0.5pt}\vspace{0.5em}

\noindent O e-mail chegou numa terça-feira, 8:47 da manhã.

\noindent Dra. Lúcia Mendes estava no consultório, com uma xícara de café esfriando na mesa e três laudos para revisar. Ela clicou no remetente: \textit{Arquivo Público Mineiro --- Setor de Documentos Históricos}.

\noindent \textit{\textquotedblleft{}Prezada Dra. Lúcia,}

\noindent \textit{Conforme conversamos em janeiro, localizamos o prontuário completo do Paciente 1.260 (Hospital Colônia de Barbacena) e os anexos mencionados: relatórios do Dr. Heitor Oliveira, laudo de necropsia do IML, e o diário do paciente.}

\noindent \textit{Os documentos estão disponíveis para consulta em nosso arquivo. Considerando seu interesse acadêmico no caso, disponibilizamos uma sala para estudo.}

\noindent \textit{Atenciosamente,}
\textit{Dr. Ronaldo Torres --- Diretor do Arquivo Público Mineiro\textquotedblright{}}

\noindent Lúcia leu o e-mail três vezes.

\noindent Ela tinha solicitado o acesso ao arquivo em janeiro --- uma pesquisa de rotina para um artigo sobre transtornos dissociativos em populações institucionalizadas. O caso do Colônia era famoso na psiquiatria forense brasileira: sessenta mil mortos em oitenta anos, pacientes enterrados vivos, diagnósticos fabricados.

\noindent Mas o Paciente 1.260 era diferente.

\noindent Ela tinha encontrado uma menção ao caso numa nota de rodapé da tese de doutorado de alguém --- uma linha: \textit{\textquotedblleft{}Paciente 1.260, Joaquim Antônio Correia, internado em 1917, faleceu em 1979. Diário encontrado junto ao corpo na exumação. Conteúdo não analisado.\textquotedblright{}}

\noindent Conteúdo não analisado.

\noindent Aquilo a intrigou. O que um paciente do Colônia poderia ter escrito que \textit{ninguém} analisou em quase cinquenta anos?

\noindent \textit{\textquotedblleft{}Você vai se arrepender\textquotedblright{}}, disse a voz na cabeça dela. Não era uma voz sobrenatural. Era a voz da experiência. Quinze anos de psicologia forense tinham lhe ensinado que alguns arquivos deveriam permanecer fechados.

\noindent Ela pegou o telefone.

\noindent \textit{\textquotedblleft{}Arquivo Público Mineiro, bom dia.\textquotedblright{}}

\noindent \textit{\textquotedblleft{}Dra. Lúcia Mendes. Confirme meu agendamento para consulta ao prontuário 1.260.\textquotedblright{}}

\noindent \textit{\textquotedblleft{}Um momento...\textquotedblright{}\textit{ Teclas. }\textquotedblleft{}Dra. Lúcia, o senhor Ronaldo deixou uma observação aqui. Ele pede que a senhora ligue para ele antes de vir. Disse que é... urgente.\textquotedblright{}}

\noindent \textit{\textquotedblleft{}Urgente?\textquotedblright{}}

\noindent \textit{\textquotedblleft{}Sim. Ele disse: 'Avise a Dra. Lúcia que o diário não está mais no arquivo.'\textquotedblright{}}

\noindent Lúcia segurou o telefone com mais força.

\noindent \textit{\textquotedblleft{}Como assim, não está mais?\textquotedblright{}}

\noindent \textit{\textquotedblleft{}Ele explica pessoalmente. Mas pediu para avisar: o diário foi retirado do arquivo em 1979. Por um médico. O Dr. Heitor Oliveira. O mesmo que assina os relatórios anexos ao prontuário.\textquotedblright{}}

\noindent \textit{\textquotedblleft{}Retirado quando?\textquotedblright{}}

\noindent \textit{\textquotedblleft{}Dois dias antes de desaparecer.\textquotedblright{}}

\noindent O café na mesa de Lúcia tinha esfriado completamente.

\noindent \textit{\textquotedblleft{}Estou indo para aí agora.\textquotedblright{}}

\noindent \textit{\textquotedblleft{}Dra. Lúcia...\textquotedblright{}\textit{ A voz do atendente hesitou. }\textquotedblleft{}O senhor Ronaldo disse outra coisa. Disse: 'Avise que o diário foi devolvido semana passada. Por uma pessoa que não se identificou. Devolvido pessoalmente. Na mesa da recepção. Envolto num pano preto.'\textquotedblright{}}

\noindent \textit{\textquotedblleft{}...E?\textquotedblright{}}

\noindent \textit{\textquotedblleft{}Ela disse: 'Avise a Dra. Lúcia que o diário está esperando ela.'\textquotedblright{}}

\noindent Lúcia desligou.

\noindent Olhou para o café frio.
Olhou para os laudos.
Olhou para a hora.

\noindent 8:52.

\noindent Cinco minutos depois do e-mail, ela já sabia que algo estava errado.

\noindent Ela só não sabia \textit{o quanto}.

\vspace{0.5em}
\noindent\textit{\textbf{↪ Rotas:} \rota{LUC-02} \rotasep \rota{DOC-08}}

\newpage\fragdef{LUC-15}% ============================================
% LUC-15 — A Mulher que Não Voltou para Casa (edição revista R1075)
% ============================================
\fragtitle{luc}{L}{LUC}{15}{A Mulher que Não Voltou para Casa}
\markboth{LUC-15 --- A Mulher que Não Voltou para Casa}{}

\noindent \textit{Belo Horizonte --- 9 de janeiro de 2026, 21h40}

\noindent Chovia na rua da Assembleia. Lúcia chegou ao bar com o casaco úmido e o cabelo colado na testa. Ronaldo já estava lá, numa mesa ao fundo, com um copo de uísque pela metade e o celular virado para baixo. Ele não a viu entrar. Estava olhando para a janela, para a chuva, com a expressão de quem espera uma ligação que não vai chegar.

\noindent Ela podia ter ido embora. Era o que fazia, normalmente. Quando o assunto era Ronaldo, ela sempre tinha uma reunião, um laudo, uma urgência. Desta vez, sentou.

\noindent ``Você demorou'', ele disse.

\noindent ``Você também.''

\noindent O bar cheirava a cerveja, madeira molhada e fritura. Mas, quando ele se inclinou para beijar seu rosto, Lúcia sentiu outra coisa: terra. Terra úmida, com um fundo doce, como se alguém tivesse aberto uma cova ali, entre as mesas. Olhou para o chão. Piso de cimento, seco. Olhou para ele. Ronaldo usava o mesmo perfume de sempre, madeira e âmbar. O cheiro de terra não estava nele.

\noindent ``Você está bem?'' ele perguntou.

\noindent ``Estou.''

\noindent Ela não estava. Havia três noites que acordava às 3h14. Havia duas semanas que o nome do paciente 1.260 aparecia em tudo o que ela lia. Havia oito anos que não ficava tão perto de Ronaldo.

\noindent Eles se conheceram no arquivo, em 2018. Ele era o diretor novo, ela a psiquiatra que vinha consultar processos antigos. Durou seis meses. Ele terminou num sábado, por telefone, dizendo que não sabia amar duas pessoas ao mesmo tempo. Lúcia não perguntou quem era a outra. Na segunda-feira, foi trabalhar com o rosto intacto e o estômago em fogo. Na sexta, aceitou um caso em outro estado. Quando voltou, ele estava casado.

\noindent Agora ele girava o copo devagar, com o dedo indicador na borda. A mão dele era a mesma. Lúcia olhou para o pulso, para o relógio antigo que ele usava, para o fio de cabelo branco que não existia oito anos atrás. Sentiu uma vontade absurda de encostar a testa no ombro dele. Não fez.

\noindent ``Preciso do acesso ao diário'', ela disse.

\noindent Ronaldo parou de girar o copo.

\noindent ``Não.''

\noindent ``Não é um pedido.''

\noindent ``Eu sei que não é.'' Ele levantou os olhos. ``É por isso que eu estou dizendo não.''

\noindent Lúcia abriu a pasta e colocou sobre a mesa a fotocópia de uma página do prontuário. O nome de Joaquim Antônio Correia. A data de 1979. A assinatura do Dr. Heitor Oliveira.

\noindent ``Você sabia que o diário estava lá'', ela disse. ``Sabia desde janeiro. Não me contou.''

\noindent ``Eu não sabia o que era.''

\noindent ``Sabia o suficiente para me mandar o e-mail.''

\noindent Ele ficou em silêncio. A chuva batia na vidraça. Lúcia viu, no reflexo do vidro, um cão parado na calçada, sob o toldo, olhando para dentro do bar. O olho direito do cão era amarelo. Não o amarelo dos faróis. Um amarelo denso, parado, como uma moeda velha. O cão não se mexeu.

\noindent ``Eu preciso ler'', ela disse. ``Se eu não ler, alguém vai ler no meu lugar. E eu não quero que seja alguém que não saiba o que está fazendo.''

\noindent Ronaldo riu sem humor.

\noindent ``Você sempre soube dizer isso. \textit{Eu preciso}. Como se o mundo fosse um paciente e você a única plantonista.''

\noindent ``Não é justo.''

\noindent ``Não é.'' Ele olhou para a mão dela sobre a mesa. ``Mas é verdade.''

\noindent Lúcia aproximou a mão. Encostou os dedos nos dele. Ronaldo não puxou. Ficaram assim por três segundos, quatro, cinco. O bar inteiro parecia ter baixado o volume. Ela sentiu o cheiro de terra de novo, mais forte, e desta vez veio com uma imagem: uma vala aberta, a terra revirada, uma mão de criança apontando para o céu.

\noindent Ela retirou a mão.

\noindent ``Me dá o cartão'', disse.

\noindent Ronaldo abriu a carteira. Tirou um cartão de acesso, desses que abrem a sala de documentos restritos. Colocou sobre a mesa, ao lado do copo.

\noindent ``Se você entrar nessa sala'', ele disse, ``não vai voltar para casa a mesma.''

\noindent Lúcia pegou o cartão. Guardou na pasta. Levantou.

\noindent ``Eu não vou voltar para casa hoje'', ela disse.

\noindent Ele não respondeu. Quando ela passou pela porta, o cão na calçada já tinha ido embora. No lugar onde ele estava, havia uma marca escura no chão, como se alguém tivesse derramado água. A chuva não lavava.

\noindent\textit{\textbf{↪ Rotas:} \rota{LUC-02} \rotasep \rota{LUC-16} \rotasep \rota{DOC-27}}

\newpage\fragdef{LUC-02}% ============================================
% LUC-02
% Gerado de molambudos.md — build_miolo.py
% ============================================
\fragtitle{luc}{L}{LUC}{02}{O Arquivo}
\markboth{LUC-02 --- O Arquivo}{}

\noindent \textit{A carta de 1979 para 2026}

\noindent \textbf{Arquivo Público Mineiro --- Belo Horizonte}
\textbf{Fevereiro de 2026}

\vspace{0.5em}\noindent\rule{\textwidth}{0.5pt}\vspace{0.5em}

\noindent O Arquivo Público Mineiro fica num prédio cinza no centro de Belo Horizonte. Lúcia conhecia o lugar --- já tinha passado horas ali para outras pesquisas. Mas nunca tinha visto a sala dos fundos.

\noindent O diretor, Dr. Ronaldo Torres, a esperava na porta. Era um homem magro, de óculos grossos e mãos inquietas. Ele a cumprimentou com um aperto rápido.

\noindent \textit{\textquotedblleft{}Dra. Lúcia. Obrigado por vir tão rápido.\textquotedblright{}}

\noindent \textit{\textquotedblleft{}O senhor disse que era urgente.\textquotedblright{}}

\noindent \textit{\textquotedblleft{}É.\textquotedblright{}\textit{ Ele olhou para os lados. }\textquotedblleft{}Vamos para minha sala.\textquotedblright{}}

\noindent A sala do diretor era pequena, cheia de pastas e livros. Na mesa, um pacote envolto em pano preto e, ao lado, uma chave enferrujada sobre um papel de seda.

\noindent \textit{\textquotedblleft{}Chegou semana passada\textquotedblright{}\textit{, disse Ronaldo. }\textquotedblleft{}Uma mulher. Não deixou nome. Disse apenas: 'Isto pertence ao arquivo. O Dr. Heitor Oliveira pediu para devolver.'\textquotedblright{}}

\noindent \textit{\textquotedblleft{}Heitor Oliveira está desaparecido desde 1979.\textquotedblright{}}

\noindent \textit{\textquotedblleft{}Eu sei.\textquotedblright{}}

\noindent \textit{\textquotedblleft{}Como ele poderia pedir para devolver um diário quase cinquenta anos depois?\textquotedblright{}}

\noindent Ronaldo não respondeu. Ele tirou os óculos e limpou as lentes.

\noindent \textit{\textquotedblleft{}A senhora quer ver o diário?\textquotedblright{}}

\noindent \textit{\textquotedblleft{}Sim.\textquotedblright{}}

\noindent \textit{\textquotedblleft{}Antes, preciso mostrar uma coisa.\textquotedblright{}}

\noindent Ele abriu uma pasta sobre a mesa. Dentro, havia fotografias. Antigas. Amareladas. A primeira mostrava um menino magro, de braços finos e olho direito amarelo. Ele estava em pé, na frente de um muro de tijolos, com um número pendurado no pescoço: \textit{1.260}.

\begin{figure}[H]
  \centering
  \includegraphics[width=0.55\textwidth]{foto_01_paciente1260_sepia.png}
  \caption*{\footnotesize\textbf{Fotografia 1} --- \textit{Joaquim Antônio Correia --- Paciente 1.260. 1917.}}
\end{figure}

\noindent \textit{\textquotedblleft{}Joaquim\textquotedblright{}}, disse Lúcia.

\noindent \textit{\textquotedblleft{}Sim. Esta é a única foto dele. Foi tirada na chegada ao Colônia, em 1917.\textquotedblright{}}

\noindent \textit{\textquotedblleft{}E as outras?\textquotedblright{}}

\noindent Ronaldo virou as fotos. A segunda mostrava um homem --- um médico jovem, de jaleco, com olheiras profundas e um sorriso forçado.

\begin{figure}[H]
  \centering
  \includegraphics[width=0.55\textwidth]{foto_02_dr_heitor_sepia.png}
  \caption*{\footnotesize\textbf{Fotografia 2} --- \textit{Dr. Heitor Oliveira. Arquivo do Colônia, 1979.}}
\end{figure}

\noindent \textit{\textquotedblleft{}Dr. Heitor Oliveira. Foto do crachá, 1979.\textquotedblright{}}

\noindent A terceira foto fez Lúcia parar.

\noindent Era uma foto de um homem mais velho. Muito mais velho. Cabelos brancos, pele enrugada, olhos opacos. Ele estava sentado num banco de praça, com um livro no colo.

\begin{figure}[H]
  \centering
  \includegraphics[width=0.55\textwidth]{foto_03_homem_livro_sepia.png}
  \caption*{\footnotesize\textbf{Fotografia 3} --- \textit{Homem não identificado. Encontrada dentro do diário.}}
\end{figure}

\noindent \textit{\textquotedblleft{}Quem é?\textquotedblright{}} perguntou Lúcia.

\noindent \textit{\textquotedblleft{}Não sei\textquotedblright{}\textit{, disse Ronaldo. }\textquotedblleft{}Esta foto estava dentro do diário. Junto com uma carta.\textquotedblright{}}

\noindent \textit{\textquotedblleft{}Carta?\textquotedblright{}}

\noindent \textit{\textquotedblleft{}Endereçada à senhora.\textquotedblright{}}

\noindent Ronaldo pegou um envelope pardo. Estava lacrado. Escrito à mão:

\noindent \textit{\textquotedblleft{}Para Dra. Lúcia Mendes --- Pessoal\textquotedblright{}}

\noindent Lúcia apontou para a chave enferrujada sobre a mesa. \textit{\textquotedblleft{}E a chave?\textquotedblright{}}

\noindent Ronaldo pegou a chave com cuidado. \textit{\textquotedblleft{}Veio junto. Não sei se é do diário, do arquivo ou de alguma cela do Colônia. Já tentamos em todas as portas do prédio. Não encaixa em nenhuma.\textquotedblright{}}

\noindent \textit{\textquotedblleft{}Por que a mandaram junto?\textquotedblright{}}

\noindent \textit{\textquotedblleft{}A mulher que entregou o pacote disse: 'O Dr. Oliveira pedia para avisar que a chave não abre portas. Ela abre a história.'\textquotedblright{}}

\noindent Lúcia tocou a chave. O metal estava frio --- mais frio do que deveria estar numa sala fechada.

\noindent \textit{\textquotedblleft{}Guarde com o diário\textquotedblright{}\textit{, ela disse. }\textquotedblleft{}Pode ser importante.\textquotedblright{}}

\noindent \textit{\textquotedblleft{}Como sabiam que eu ia pedir o arquivo?\textquotedblright{}}

\noindent \textit{\textquotedblleft{}Não sei.\textquotedblright{}}

\noindent Lúcia abriu o envelope. Dentro, um papel amarelado, com caligrafia trêmula:

\vspace{0.5em}\noindent\rule{\textwidth}{0.5pt}\vspace{0.5em}

\noindent \textit{Dra. Lúcia,}

\noindent \textit{Se está lendo esta carta, o ciclo chegou até a senhora.}

\noindent \textit{Meu nome é Heitor Oliveira. Estou escrevendo esta carta em 2026. Ou melhor, estou escrevendo em 1979 e pedindo para que alguém entregue em 2026.}

\noindent \textit{Não pergunte como. O diário ensina.}

\noindent \textit{Li o diário do paciente 1.260 em março de 1979. Desde então, não sou mais o mesmo.}

\noindent \textit{O diário passou 47 anos tentando entender o que aconteceu comigo. Tentando entender o que aconteceu com Joaquim.}

\noindent \textit{Não entendi.}

\noindent \textit{Só entendi uma coisa: o diário não é um registro. É um veículo. A criatura que Joaquim encontrou na vala não mora em corpos. Mora em histórias. Cada vez que alguém lê, ela passa para o leitor.}

\noindent \textit{A senhora leu este diário?}

\noindent \textit{Não. Ainda não.}

\noindent \textit{Mas vai ler.}

\noindent \textit{E quando ler, vai entender.}

\noindent \textit{Heitor Oliveira}
\textit{--- 1979 (enviado para 2026)}

\vspace{0.5em}\noindent\rule{\textwidth}{0.5pt}\vspace{0.5em}

\noindent Lúcia colocou a carta na mesa. A mão tremia.

\noindent \textit{\textquotedblleft{}Isto é uma brincadeira\textquotedblright{}\textit{, ela disse. }\textquotedblleft{}Alguém está manipulando o arquivo.\textquotedblright{}}

\noindent \textit{\textquotedblleft{}Talvez\textquotedblright{}\textit{, disse Ronaldo. }\textquotedblleft{}Mas tem mais uma coisa.\textquotedblright{}}

\noindent \textit{\textquotedblleft{}O quê?\textquotedblright{}}

\noindent \textit{O diretor apontou para o pacote de pano preto. O diário.}

\noindent \textit{\textquotedblleft{}A mulher que entregou disse outra frase. Ela disse: 'Avise a Dra. Lúcia que o paciente 1.260 mandou um abraço.'\textquotedblright{}}

\noindent Lúcia engoliu seco.

\noindent Joaquim Antônio Correia morreu em 1979.

\noindent Isso era fato.

\noindent Mas o diário estava ali. Quente. Esperando.

\noindent \textit{\textquotedblleft{}Vou levar.\textquotedblright{}}

\noindent \textit{\textquotedblleft{}Tem certeza?\textquotedblright{}}, perguntou Ronaldo.

\noindent \textit{\textquotedblleft{}Não.\textquotedblright{}}

\noindent Ela estendeu a mão e tocou o pano preto.

\noindent O pano estava quente.

\vspace{0.5em}
\noindent\textit{\textbf{↪ Rotas:} \rota{LUC-03} \rotasep \rota{DOC-01} \rotasep \rota{DOC-20}}
\newpage\fragdef{LUC-03}% ============================================
% LUC-03
% Gerado de molambudos.md — build_miolo.py
% ============================================
\fragtitle{luc}{L}{LUC}{03}{As Ruínas}
\markboth{LUC-03 --- As Ruínas}{}

\noindent \textit{O Colônia abandona}

\noindent \textbf{Março de 2026 --- Barbacena, Minas Gerais}

\vspace{0.5em}\noindent\rule{\textwidth}{0.5pt}\vspace{0.5em}

\noindent O Hospital Colônia de Barbacena foi fechado em 1980, um ano depois da morte de Joaquim.

\noindent Hoje, o que resta é um conjunto de prédios abandonados, cercados por mato alto e silêncio. As paredes ainda estão de pé, mas o telhado caiu em várias alas. O piso de madeira apodreceu. As camas de ferro enferrujado ainda estão lá, algumas com restos de colchão, outras com nada.

\noindent Lúcia chegou num sábado de manhã. O sol estava forte, mas o mato era alto demais para o sol alcançar o chão. Ela abriu caminho com as mãos, afastando galhos secos e teias de aranha.

\noindent A ala dos silenciados ficava nos fundos.

\noindent Ela reconheceu pela descrição no prontuário: \textit{\textquotedblleft{}Fim do corredor. Porta de ferro. Doze celas.\textquotedblright{}}

\noindent A porta de ferro estava entreaberta. Lúcia empurrou. O rangido ecoou pelo corredor vazio.

\noindent Dentro, o cheiro era de terra molhada e mofo. Por baixo, quase imperceptível, havia uma doçura de fruta estourada no fundo de uma lata.

\noindent Ela acendeu a lanterna do celular.

\noindent As celas eram pequenas. Dois metros por três. Em cada uma, uma cama de concreto. Em algumas, restos de cobertores --- trapos cinzentos que se desfizeram com o tempo.

\noindent A cela 1.260 ficava no final do corredor.

\noindent Lúcia parou na frente dela.

\noindent A porta era de madeira maciça, com uma abertura do tamanho de um prato --- o buraco por onde passavam a comida. Ela espiou.

\noindent Dentro, a cela estava vazia.
Quase.

\noindent No chão, no canto direito, havia um caderno.

\noindent O caderno não era o diário. O diário estava na mochila de Lúcia. Era outro caderno. Pequeno, de capa preta, com as páginas amareladas.

\noindent Ela abriu a porta. Entrou.

\noindent O cheiro adocicado era mais forte aqui.

\noindent Ela se abaixou e pegou o caderno. Na capa, escrito à mão:
\noindent A doçura podre se adensou ali, grudando na parte de trás da língua.

\noindent \textit{\textquotedblleft{}Joaquim --- 1945 a 1956\textquotedblright{}}

\noindent Os anos de silêncio.

\noindent Ela abriu. A caligrafia era minúscula, apertada, como se ele quisesse caber o máximo de palavras possível em cada página.

\noindent A primeira linha dizia:

\noindent \textquotedblleft{\textit{Ela não me deixa falar. Quando eu abro a boca, ela aperta minha garganta. A fome não quer que eu conte. Mas ela não pode me impedir de escrever. Escrever não é falar. Escrever é pensar no papel. Ela não lê meus pensamentos. Ela só lê o que eu escrevo. Mas ela está aprendendo.\textquotedblright{}}}

\noindent A última linha dizia:

\noindent \textit{\textquotedblleft{}1956. Hoje ela deixou eu falar. Perguntei por quê. Ela disse: 'Porque você já escreveu tudo. Agora a palavra está no papel. Não precisa mais sair pela boca.' Eu vou falar de novo. Mas antes vou escrever mais. Escrever é a única coisa que ainda é minha.\textquotedblright{}}

\noindent Lúcia fechou o caderno.

\noindent O cheiro adocicado estava mais forte.

\noindent Ela levantou os olhos.

\noindent No canto da cela, na parede, alguém tinha escrito alguma coisa. Com tinta. Ou com outra coisa.

\noindent \textquotedblleft{\textit{Bem-vinda, Dra. Lúcia. A senhora demorou.\textquotedblright{}}}

\noindent A caligrafia era a mesma do caderno.

\noindent A de Joaquim.

\noindent A caligrafia sugere 1956. Não há datação.

\noindent Prevendo que ela estaria ali. 70 anos depois.

\noindent Lúcia respirou fundo.

\noindent O cheiro adocicado encheu os pulmões.

\vspace{0.5em}
\noindent\textit{\textbf{↪ Rotas:} \rota{LUC-04} \rotasep \rota{MEM-08}}
\newpage\fragdef{LUC-04}% ============================================
% LUC-04
% Gerado de molambudos.md — build_miolo.py
% ============================================
\fragtitle{luc}{L}{LUC}{04}{A Fotografia}
\markboth{LUC-04 --- A Fotografia}{}

\noindent \textit{O velho no cemitério}

\noindent \textbf{Março de 2026 --- Barbacena}

\vspace{0.5em}\noindent\rule{\textwidth}{0.5pt}\vspace{0.5em}

\noindent Lúcia voltou ao hotel às 22h.
\noindent O ar entrou grosso. Por um segundo, ela teve a impressão absurda de que não respirava: era respirada.

\noindent O hotel era uma pousada simples no centro de Barbacena. Quarto pequeno, cama dura, uma janela que dava para a rua principal. Ela colocou a mochila na cama, tirou o diário e o caderneta, e sentou.

\noindent O diário estava sobre a cama.

\noindent Ela ainda não tinha aberto.

\noindent Nas duas semanas desde que pegou o diário no Arquivo Público, ela não tinha conseguido abrir. Sempre havia uma desculpa: trabalho, cansaço, hora errada. A verdade era que ela estava com medo.

\noindent Não medo racional.
Medo do que ia encontrar.
Medo do que ia sentir.
\noindent Não medo racional. Medo do que ia encontrar. Medo do que ia sentir. Medo do que ia mudar dentro dela.

\noindent \textit{\textquotedblleft{}Você é psiquiatra forense\textquotedblright{}\textit{, ela disse em voz alta. }\textquotedblleft{}Você lida com criminosos, com traumatizados, com arquivos de horror. Isso é só mais um caso.\textquotedblright{}}

\noindent Ela abriu o diário.

\noindent A primeira página tinha a caligrafia infantil de Joaquim. A data: 1917.

\noindent \textit{\textquotedblleft{}Hoje eu cheguei no Colônia. O médico disse que eu sou doido. Mas eu não sou doido. A criatura é que é doida. Ela mora dentro de mim e fala o tempo todo.\textquotedblright{}}

\noindent Lúcia leu.

\noindent Página após página.

\noindent A infância no sertão. A seca. A caminhada. O pai morto no juazeiro. A mãe levada pelos soldados. O Curral do Governo. A vala.

\noindent Ela leu sobre a criatura.

\noindent Leu a descrição do corpo de espinhos, das costelas viradas para fora, do órgão pulsando. Leu sobre a luz presa no olho. Leu sobre o trato.

\noindent Leu sobre o olho amarelo. Leu sobre o trato.

\noindent Quando levantou os olhos, eram 3:14 da manhã.

\noindent O quarto estava escuro. A única luz vinha do abajur da mesa de cabeceira. O diário estava aberto na página 126.

\noindent Ela sentiu o cheiro.

\noindent Adocicado.

\noindent Estava no quarto.

\noindent Ela olhou para o canto. A sombra parecia mais densa ali. Mais escura. Como se o escuro fosse mais escuro naquele pedaço da parede.

\noindent \textit{\textquotedblleft{}Não tem nada\textquotedblright{}}, ela disse.

\noindent Mas não acreditou.

\noindent Ela fechou o diário. Guardou na mochila. Deitou. Apagou a luz.

\noindent Ficou olhando para o teto.

\noindent Às 3:14 da manhã, o olho amarelo apareceu no canto do quarto.

\noindent O sonho tinha corpo.

\noindent Ela viu.

\noindent Piscou.

\noindent O olho sumiu.

\noindent Ela passou o resto da noite com a luz acesa, o diário fechado debaixo do travesseiro, e a certeza de que alguma coisa tinha entrado no quarto.

\noindent A certeza de que alguma coisa tinha entrado \textit{nela}.

\vspace{0.5em}\noindent\rule{\textwidth}{0.5pt}\vspace{0.5em}

\noindent Na manhã seguinte, ela voltou ao Colônia.

\noindent Dessa vez, foi até o jazigo. O cemitério do hospital, onde sessenta mil corpos foram enterrados em valas comuns. Não havia cruzes. Não havia nomes. Apenas um gramado liso, como se nada estivesse enterrado ali.

\noindent Não havia cruzes. Não havia nomes. Apenas um gramado liso, como se nada estivesse enterrado ali.

\noindent Ela sabia que Joaquim estava ali. Cova 1.260.

\noindent Ela ficou parada, olhando para o chão.

\noindent \textit{\textquotedblleft{}O que você quer de mim?\textquotedblright{}}, perguntou em voz alta.

\noindent O vento respondeu com silêncio.

\noindent \textit{\textquotedblleft{}Por que eu?\textquotedblright{}}

\noindent Uma folha seca rodopiou aos pés dela.

\begin{figure}[H]
  \centering
  \includegraphics[width=0.55\textwidth]{foto_05_fotografia_luc04_sepia.png}
  \caption*{\footnotesize\textbf{Fotografia 4} --- \textit{O velho no cemitério. Jazigo do Colônia, 2026.}}
\end{figure}

\noindent \textit{\textquotedblleft{}Você leu\textquotedblright{}}, disse uma voz atrás dela.

\noindent Lúcia se virou.

\noindent Um homem velho estava parado a alguns metros. Muito velho. Pele enrugada, cabelos brancos ralos, olhos opacos. Ele segurava uma bengala.

\noindent \textit{\textquotedblleft{}Quem é você?\textquotedblright{}}

\noindent \textit{\textquotedblleft{}Só um velho que passou por aqui\textquotedblright{}\textit{, ele disse. }\textquotedblleft{}E que viu a senhora lendo o diário.\textquotedblright{}}

\noindent \textit{\textquotedblleft{}Como sabe que eu li?\textquotedblright{}}

\noindent Ele sorriu. A boca era uma linha torta.

\noindent \textit{\textquotedblleft{}O cheiro. A senhora está com o cheiro. Quem leu o diário fica com o cheiro por dias.\textquotedblright{}}

\noindent \textit{\textquotedblleft{}Que cheiro?\textquotedblright{}}

\noindent \textit{\textquotedblleft{}Doce. De flor murcha. De carne passada.\textquotedblright{}}

\noindent Lúcia sentiu o estômago apertar.

\noindent \textit{\textquotedblleft{}Quem é você?\textquotedblright{}} repetiu.

\noindent O velho se aproximou. Devagar. A bengala batendo no chão seco.

\noindent \textit{\textquotedblleft{}Conheci o paciente 1.260.\textquotedblright{}\textit{ }\textquotedblleft{}Mentira. Ele morreu em 79.\textquotedblright{}}

\noindent \textit{\textquotedblleft{}Conheci antes. Trabalhei aqui nos anos 50.\textquotedblright{}}

\noindent \textit{\textquotedblleft{}O que você sabe sobre ele?\textquotedblright{}}

\noindent O velho parou na frente dela.

\noindent \textit{\textquotedblleft{}Sei que ele escrevia num diário. Sei que o diário foi enterrado com ele. Sei que, quando exumaram, o diário estava quente. E sei que todo mundo que leu o diário...\textquotedblright{}\textit{ Ele fez uma pausa. }\textquotedblleft{}...teve visita.\textquotedblright{}}

\noindent \textit{\textquotedblleft{}Que visita?\textquotedblright{}}

\noindent \textit{\textquotedblleft{}A Rasga Mortalha.\textquotedblright{}}

\noindent O velho apontou para o céu.

\noindent Lúcia olhou.

\noindent No galho seco de uma árvore, uma coruja branca estava pousada. Ela olhava fixamente para Lúcia.

\noindent \textquotedblleft{\textit{Ela vem quando alguém abre o diário\textquotedblright{}}, disse o velho. }\textquotedblleft{\textit{Fica olhando até a pessoa escolher.\textquotedblright{}}}

\noindent \textit{\textquotedblleft{}Que ciclo?\textquotedblright{}}

\noindent \textit{\textquotedblleft{}O ciclo do paciente 1.260. Ele morre, alguém lê, o leitor vira o próximo.\textquotedblright{}}

\noindent O velho deu as costas.

\noindent \textquotedblleft{\textit{Espera!\textquotedblright{}} Lúcia chamou. }\textquotedblleft{\textit{Como eu paro isso?\textquotedblright{}}}

\noindent Ele parou. Não virou o rosto.

\noindent \textit{\textquotedblleft{}Não sai.\textquotedblright{}}

\noindent \textit{\textquotedblleft{}Não sai?\textquotedblright{}}

\noindent \textit{\textquotedblleft{}Só passa adiante. Arranja outro leitor. Passa o diário. Deixa outra pessoa ler. Aí o ciclo continua. Só que com outro.\textquotedblright{}}

\noindent Ele começou a andar.

\noindent \textit{\textquotedblleft{}E se eu não passar?\textquotedblright{}}

\noindent A voz do velho veio de longe, carregada pelo vento:
\noindent \textquotedblleft{\textit{Não para. Troca de mão. A senhora fecha a sua e outra abre. É assim que coisa enterrada aprende a andar.\textquotedblright{}}}

\noindent \textquotedblleft{\textit{Então a senhora apodrece com ele por dentro. Igual Joaquim. Igual o médico. Igual todos que acharam que podiam guardar uma porta fechada no peito.\textquotedblright{}}}

\noindent A coruja branca abriu as asas. Voou.

\noindent Lúcia ficou sozinha no cemitério, com o diário na mochila, o cheiro no nariz, e a certeza de que alguma coisa havia mudado para sempre.

\vspace{0.5em}
\noindent\textit{\textbf{↪ Rotas:} \rota{LUC-05} \rotasep \rota{MEM-12}}
\newpage\fragdef{LUC-05}% ============================================
% LUC-05
% Gerado de molambudos.md — build_miolo.py
% ============================================
\fragtitle{luc}{L}{LUC}{05}{O Sonho de Lúcia}
\markboth{LUC-05 --- O Sonho de Lúcia}{}

\noindent \textit{A contaminação da investigadora}

\noindent \textbf{Abril de 2026 --- Belo Horizonte}

\vspace{0.5em}\noindent\rule{\textwidth}{0.5pt}\vspace{0.5em}

\noindent Ela sonhou com a vala.

\noindent Não como quem vê um filme. Como quem \textit{está} lá. Os pés descalços na terra fria. O cheiro adocicado que entra pelos poros. O céu escuro, sem estrelas, sem lua, como se o mundo tivesse esquecido de acender as luzes.

\noindent Ela ouviu a mastigação antes de ver a criatura.

\noindent O som de osso sendo quebrado. De carne sendo triturada. Um som molhado, paciente, rítmico --- como quem come devagar porque não tem pressa, porque a comida não vai fugir.

\noindent Ela andou até a borda da vala.

\noindent Lá embaixo, a pilha de corpos se moveu.

\noindent A criatura ergueu a cabeça.

\noindent O rosto quase humano. Pele cinzenta esticada sobre os ossos. A boca suja de sangue. Os olhos amarelos brilhando no escuro.

\noindent Ela sorriu.

\noindent \textit{\textquotedblleft{}Finalmente\textquotedblright{}\textit{, disse a criatura. }\textquotedblleft{}Estava esperando a senhora.\textquotedblright{}}

\noindent \textit{\textquotedblleft{}Você sabia que eu ia vir?\textquotedblright{}}

\noindent \textit{\textquotedblleft{}Sempre soube. Desde a primeira página.\textquotedblright{}}

\noindent A criatura subiu. O corpo imenso, a gaiola de costelas expostas, o órgão pulsando entre os ossos. Ela parou na frente de Lúcia.

\noindent \textit{\textquotedblleft{}A senhora leu o diário.\textquotedblright{}}

\noindent \textit{\textquotedblleft{}...Sim.\textquotedblright{}}

\noindent \textit{\textquotedblleft{}O que achou?\textquotedblright{}}

\noindent \textit{\textquotedblleft{}É a história mais triste que já li.\textquotedblright{}}

\noindent A criatura inclinou a cabeça.

\noindent \textit{\textquotedblleft{}Não é triste. É verdade. A senhora confunde as duas coisas.\textquotedblright{}}

\noindent \textit{\textquotedblleft{}O que você quer de mim?\textquotedblright{}}

\noindent \textit{\textquotedblleft{}Nada. Já tenho o que queria.\textquotedblright{}}

\noindent \textit{\textquotedblleft{}O quê?\textquotedblright{}}

\noindent \textit{\textquotedblleft{}A senhora.\textquotedblright{}}

\noindent A criatura tocou o peito de Lúcia com um dedo longo. O dedo atravessou a pele, os ossos. Não doeu. Apenas... entrou.

\noindent \textit{\textquotedblleft{}Agora a senhora é minha.\textquotedblright{}}

\noindent \textit{\textquotedblleft{}Não\textquotedblright{}\textit{, disse Lúcia. }\textquotedblleft{}Não sou.\textquotedblright{}}

\noindent \textit{\textquotedblleft{}Já é. Desde o momento em que abriu o diário. A senhora só não aceitou ainda.\textquotedblright{}}

\noindent Lúcia acordou.

\noindent O quarto estava escuro. O relógio marcava 3:14.

\noindent A doçura podre estava no travesseiro.

\noindent Ela sentou na cama, ofegante. Olhou para o canto do quarto. A sombra parecia mais densa.

\noindent Ela acendeu a luz.

\noindent Nada.

\noindent Mas o cheiro não foi embora.
\noindent Ela sentou na cama, ofegante. Olhou para o canto do quarto. A sombra parecia mais densa. E havia um som. Um som que ela conhecia do relatório de Joaquim — um canto fino, agudo, como uma mulher chorando muito longe. O canto de uma coruja. A Rasga Mortalha. Ela não viu o pássaro. Mas ouviu. Ouviu as asas se abrindo, o bater no ar parado do quarto.

\noindent Ela ficou sentada na cama, com a luz acesa, até o sol nascer.

\vspace{0.5em}\noindent\rule{\textwidth}{0.5pt}\vspace{0.5em}

\noindent Na manhã seguinte, ela foi ao oftalmologista.

\noindent \textit{\textquotedblleft{}Seu olho direito está com uma leve descoloração\textquotedblright{}\textit{, disse o médico. }\textquotedblleft{}Nada grave. Apenas uma alteração na pigmentação da íris. Já nasceu assim?\textquotedblright{}}

\noindent \textit{\textquotedblleft{}Não.\textquotedblright{}}

\noindent \textit{\textquotedblleft{}Hmm. Pode ser medicamento. A senhora toma algum remédio?\textquotedblright{}}

\noindent \textit{\textquotedblleft{}Não.\textquotedblright{}}

\noindent \textquotedblleft{\textit{Estranho.\textquotedblright{}} Ele anotou algo no prontuário. }\textquotedblleft{\textit{Vou pedir uns exames. Só para garantir.\textquotedblright{}}}

\noindent Lúcia olhou no espelho do consultório.

\noindent O olho direito estava mais claro.
Amarelado.

\noindent Ela não disse nada.

\noindent Mas, ao sair, passou a mão no rosto e sentiu a pele mais fria do que deveria.
\noindent Alguma coisa havia assinado por ela.

\noindent O ciclo já tinha começado.

\vspace{0.5em}
\noindent\textit{\textbf{↪ Rotas:} \rota{LUC-Contraprova} \rotasep \rota{CONT-04}}
\newpage\fragdef{LUC-Contraprova}% ============================================
% LUC-Contraprova
% Gerado de molambudos.md — build_miolo.py
% ============================================
\fragtitle{luc}{L}{LUC}{Contraprova}{A Contraprova}
\markboth{LUC-Contraprova --- A Contraprova}{}

\noindent \textit{\textquotedblleft{}Isso tem que ter uma explicação científica\textquotedblright{}}

\noindent \textbf{Abril de 2026 --- Belo Horizonte}

\vspace{0.5em}\noindent\rule{\textwidth}{0.5pt}\vspace{0.5em}

\noindent Três dias depois do sonho, Lúcia Mendes entrou no laboratório do Hospital das Clínicas da UFMG com uma pasta cheia de exames.

\noindent \textit{\textquotedblleft{}Preciso de uma bateria completa\textquotedblright{}\textit{, disse ao técnico. }\textquotedblleft{}Ressonância de crânio com contraste. Mapeamento retiniano. Eletroencefalograma de 24 horas. Hemograma completo. Dosagem hormonal. Perfil hepático. Teste de função tireoidiana. RPR. FAN. Toxoplasmose. Citomegalovírus. Teste rápido para HIV. Sífilis. Tudo.\textquotedblright{}}

\noindent O técnico ergueu uma sobrancelha. \textit{\textquotedblleft{}A senhora está com o quê?\textquotedblright{}}

\noindent \textit{\textquotedblleft{}Não sei. Quero descobrir.\textquotedblright{}}

\vspace{0.5em}\noindent\rule{\textwidth}{0.5pt}\vspace{0.5em}

\noindent Os resultados chegaram uma semana depois.

\noindent \textbf{Ressonância magnética:} Normal. Sem lesões, sem tumores, sem alterações de sinal no parênquima cerebral. Hipotálamo, amígdala, hipocampo --- todas as estruturas dentro dos padrões.

\noindent \textbf{Eletroencefalograma:} Atividade elétrica normal. Sem descargas epileptiformes. Sem padrões de atividade lenta focal. Sem evidência de atividade epileptogênica no lobo temporal.

\noindent \textbf{Mapeamento retiniano:\textit{ Epitélio pigmentar da retina normal. Disco óptico sem alterações. Pressão intraocular normal. }\textquotedblleft{}Seu olho direito está saudável\textquotedblright{}\textit{, disse o oftalmologista. }\textquotedblleft{}A descoloração não tem causa orgânica detectável.\textquotedblright{}}

\noindent \textbf{Exames de sangue:} Todos dentro da normalidade. Glicemia, ureia, creatinina, sódio, potássio, cálcio, magnésio --- normais. Hormônios tireoidianos --- normais. Sorologias --- todas negativas.

\noindent \textit{\textquotedblleft{}A senhora está saudável\textquotedblright{}\textit{, disse o médico, fechando o prontuário. }\textquotedblleft{}O olho… bem, às vezes acontece. Heterocromia adquirida idiopática. Não tem explicação, mas não mata ninguém.\textquotedblright{}}

\noindent --- \textit{Não tem explicação.}

\noindent Lúcia guardou as palavras. Saiu do hospital carregando a pasta de exames. Atravessou a rua. Sentou num banco de praça. Abriu a pasta. Leu todos os laudos novamente.

\noindent Tudo normal.

\noindent Ela estava saudável.

\noindent Então por que o sonho não parava? Por que o cheiro adocicado aparecia em horários aleatórios? Por que ela acordava toda noite às 3:14?

\noindent \textit{\textquotedblleft{}Não tem explicação.\textquotedblright{}}

\noindent Ela fechou a pasta. Olhou para o céu. O céu estava azul, mas ela não conseguia ignorar a mancha amarela na sua visão periférica.

\noindent --- \textbf{Não tem explicação.}

\noindent Foi aí que ela decidiu: se a ciência não podia responder, ela iria ao lugar onde tudo começou.

\vspace{0.5em}\noindent\rule{\textwidth}{0.5pt}\vspace{0.5em}

\vspace{0.5em}
\noindent\textit{\textbf{↪ Rotas:} \rota{LUC-Entrevista} \rotasep \rota{LUC-05}}
\newpage\fragdef{LUC-Entrevista}% ============================================
% LUC-Entrevista
% Gerado de molambudos.md — build_miolo.py
% ============================================
\fragtitle{luc}{L}{LUC}{Entrevista}{A Entrevista}
\markboth{LUC-Entrevista --- A Entrevista}{}

\noindent \textit{O homem que viu o ciclo começar}

\noindent \textbf{Abril de 2026 --- Barbacena, Casa do Sr. Antenor}

\vspace{0.5em}\noindent\rule{\textwidth}{0.5pt}\vspace{0.5em}

\noindent O endereço estava no arquivo do Colônia, anotado a lápis na borda de um relatório de 1958: \textit{\textquotedblleft{}Antenor Francisco da Silva, ex-enfermeiro, Rua da Matriz 42, Barbacena.\textquotedblright{}}

\noindent Lúcia bateu na porta às 10 da manhã. Um homem muito velho abriu. Ele devia ter mais de 90 anos --- a pele fina como papel de seda, os olhos azuis desbotados, as mãos trêmulas.

\noindent \textit{\textquotedblleft{}A senhora é a psiquiatra?\textquotedblright{}}

\noindent \textit{\textquotedblleft{}Sou.\textquotedblright{}}

\noindent \textit{\textquotedblleft{}Então leu o diário.\textquotedblright{}}

\noindent Não era uma pergunta.

\noindent \textit{\textquotedblleft{}...Li.\textquotedblright{}}

\noindent O velho Antenor abriu a porta devagar: \textit{\textquotedblleft{}Então entre. Antes que o cheiro entre junto.\textquotedblright{}}

\vspace{0.5em}\noindent\rule{\textwidth}{0.5pt}\vspace{0.5em}

\noindent A casa era pequena, escura, cheirando a mofo e remédio. Na parede da sala, uma fotografia amarelada: homens de jaleco branco na frente do Hospital Colônia, década de 1950.

\noindent Antenor serviu café. A xícara tremia na mão dele.

\noindent \textit{\textquotedblleft{}O que a senhora quer saber?\textquotedblright{}}

\noindent \textit{\textquotedblleft{}Tudo o que sabe sobre o paciente 1.260.\textquotedblright{}}

\noindent \textit{\textquotedblleft{}Joaquim.\textquotedblright{}\textit{ Ele disse o nome como quem diz o nome de um parente morto. }\textquotedblleft{}Conheci ele. Conheci todos eles. Mas o Joaquim era diferente.\textquotedblright{}}

\noindent \textit{\textquotedblleft{}Diferente como?\textquotedblright{}}

\noindent \textit{\textquotedblleft{}Ele não era louco. A senhora leu o diário. Sabe que ele não era louco.\textquotedblright{}}

\noindent \textit{\textquotedblleft{}Eu sei.\textquotedblright{}}

\noindent \textit{\textquotedblleft{}Ele via coisas.\textquotedblright{}}

\noindent \textit{\textquotedblleft{}Que tipo de coisas?\textquotedblright{}}

\noindent \textit{\textquotedblleft{}A criatura. Ele via a criatura. Desde 1915. Antes de entrar no Colônia. Ele já via.\textquotedblright{}}

\noindent Antenor contou a história de Joaquim como quem conta uma história que ouviu muitas vezes. Joaquim não era violento quando chegou --- foi a vala que o tornou violento. A vala onde os corpos não enterrados eram comidos por bichos. A vala onde Joaquim viu a criatura pela primeira vez.

\noindent \textit{\textquotedblleft{}Ele me disse uma vez: 'A criatura não veio de fora. Ela já estava na vala. Enterrada junto com os corpos. Eu só... acordei ela.'\textquotedblright{}}

\noindent \textit{\textquotedblleft{}E o senhor acreditou?\textquotedblright{}}

\noindent Antenor olhou para Lúcia. Os olhos azuis desbotados, mas a lucidez absoluta.

\noindent \textit{\textquotedblleft{}Eu trabalhei 42 anos no Colônia. Eu vi pacientes terem convulsões. Eu vi pacientes serem lobotomizados. Eu vi pacientes morrerem de fome. Eu vi coisas que a senhora não acreditaria se eu contasse. Mas de todas as coisas que vi, a única que eu não consegui explicar foi o diário do Joaquim.\textquotedblright{}}

\noindent \textit{\textquotedblleft{}O que o senhor quer dizer?\textquotedblright{}}

\noindent Ele se levantou. Foi até um armário. Tirou uma caixa de sapatos. Abriu.

\noindent Ela achou que ia encontrar fotos, documentos, lembranças.

\noindent Em vez disso, encontrou um caderno. Capa preta. Páginas amareladas. Letra trêmula.

\noindent \textit{\textquotedblleft{}Isso é...\textquotedblright{}}

\noindent \textit{\textquotedblleft{}O segundo diário de Joaquim.\textquotedblright{}}

\noindent Lúcia sentiu o sangue gelar.

\noindent \textit{\textquotedblleft{}Ele escreveu dois?\textquotedblright{}}

\noindent \textit{\textquotedblleft{}A senhora leu o primeiro. O que ele escreveu vivo. Este ele escreveu morto.\textquotedblright{}}

\noindent \textit{\textquotedblleft{}Como?\textquotedblright{}}

\noindent \textit{\textquotedblleft{}Eu cavei a cova dele. Três dias depois do enterro, a cova estava remexida. Abri. O caixão estava quebrado por dentro. O caderno estava no peito dele, com a mão dele em cima. A letra era a mesma. Mas a data era depois da morte.\textquotedblright{}}

\noindent Ele entregou o caderno a Lúcia.

\noindent \textit{\textquotedblleft{}A senhora pode ficar com ele. Eu morri de medo dele a vida inteira. Já é hora de passar adiante. É assim que funciona, não é?\textquotedblright{}}

\noindent Lúcia segurou o caderno. A capa estava fria. Mas, dentro dela, algo se aqueceu.

\noindent \textit{\textquotedblleft{}É assim que funciona\textquotedblright{}}, ela repetiu.

\noindent \textit{\textquotedblleft{}Sim\textquotedblright{}\textit{, disse Antenor. }\textquotedblleft{}A senhora já sabe. A senhora só não aceitou ainda.\textquotedblright{}}

\vspace{0.5em}\noindent\rule{\textwidth}{0.5pt}\vspace{0.5em}

\vspace{0.5em}
\noindent\textit{\textbf{↪ Rotas:} \rota{LUC-Escolha} \rotasep \rota{LUC-Contraprova}}
\newpage\fragdef{LUC-16}% ============================================
% LUC-16 — A Mãe no Telefone (edição revista R1075)
% ============================================
\fragtitle{luc}{L}{LUC}{16}{A Mãe no Telefone}
\markboth{LUC-16 --- A Mãe no Telefone}{}

\noindent \textit{Belo Horizonte --- 11 de janeiro de 2026, 22h17}

\noindent O telefone tocou três vezes antes de a mãe atender. Lúcia imaginou a casa inteira naquele intervalo: a sala com o sofá de flores, a mesa de crochê, o cheiro de lavanda e madeira velha, o retrato do pai na parede. A casa existia naquele cheiro. Fora dela, o pai virava uma foto; dentro, continuava sentado na poltrona, lendo o jornal com a ponta da língua entre os dentes.

\noindent ``Alô?''

\noindent ``Mãe. Sou eu.''

\noindent ``Lúcia.'' A voz da mãe mudou de tom. Não era surpresa. Era conferência. ``Aconteceu alguma coisa?''

\noindent ``Não. Só ligando.''

\noindent ``Você não liga só.''

\noindent Lúcia fechou os olhos. A mãe tinha razão. Ela ligava em datas marcadas, em aniversários, quando havia uma notícia prática: exame, consulta, uma reforma no telhado. Nunca só.

\noindent ``Estou bem'', disse. ``Estou trabalhando muito.''

\noindent ``Nesse arquivo?''

\noindent ``É.''

\noindent ``Aquele do hospital?''

\noindent Lúcia ouviu o barulho da televisão ao fundo. Novela. A mãe sempre deixava a TV ligada, mesmo sem assistir. Dizia que a casa ficava menos vazia.

\noindent ``É um caso antigo'', Lúcia disse. ``Estou organizando documentos.''

\noindent ``Você está com a voz rouca.''

\noindent ``É o ar-condicionado.''

\noindent ``Desde quando ar-condicionado deixa a voz rouca?'' A mãe fez uma pausa. ``Você está dormindo?''

\noindent Lúcia pensou em mentir. Não conseguiu.

\noindent ``Pouco.''

\noindent ``Eu sabia.''

\noindent O silêncio entre as duas era antigo, com regras próprias. A mãe esperava. Lúcia calculava. Havia coisas que ela nunca contava: o casamento que não aconteceu, o aborto aos trinta e dois anos, o namoro com Ronaldo, o medo de envelhecer sozinha. A mãe sabia de algumas. Não sabia de outras. As que ela não sabia eram as que mais pesavam.

\noindent ``A vizinha do 204 morreu'', a mãe disse.

\noindent ``Dona Cida?''

\noindent ``Ela mesma. Encontrada na cozinha. O cachorro latiu a noite inteira. Ninguém foi ver.''

\noindent ``Que triste.''

\noindent ``O cachorro está aqui no quintal. Ela deixou. Eu não queria, mas não tinha para quem dar.''

\noindent Lúcia ouviu um latido. Distante, abafado, como se viesse de dentro da parede.

\noindent ``Ele tem um olho amarelo'', a mãe disse. ``O direito. O esquerdo é marrom. Nunca vi cachorro assim.''

\noindent Lúcia abriu os olhos.

\noindent ``Mãe.''

\noindent ``Que é?''

\noindent ``Não deixa esse cachorro entrar em casa.''

\noindent A mãe riu. ``Você está com medo de cachorro agora?''

\noindent ``Medo, não. É...'' Lúcia parou. Não havia palavra. ``Só não deixa.''

\noindent Houve um estalo na linha. Um ruído fino, como folhas sendo viradas depressa. Depois, a voz da mãe voltou, mas havia algo errado com o volume. Estava mais baixa, mais lenta, como se falasse de dentro de um cômodo fechado.

\noindent ``Lúcia.''

\noindent ``Mãe?''

\noindent ``Você está com a voz do seu pai.''

\noindent Lúcia ficou parada. O pai morrera havia doze anos. Câncer de pulmão. Ele falava baixo, com pausas longas, como se cada frase custasse um pedaço de ar. Lúcia ouviu a própria respiração e percebeu que estava fazendo o mesmo.

\noindent ``Mãe, eu preciso desligar.''

\noindent ``Você não vai vir no domingo?''

\noindent ``Eu vou tentar.''

\noindent ``Não promete se não for cumprir.''

\noindent ``Eu prometo.''

\noindent Ela desligou antes que a mãe respondesse. Ficou com o telefone na mão, olhando para a tela apagada. Havia uma mancha de suor na capa. O cheiro de lavanda e madeira velha continuava no ar do apartamento, embora a mãe estivesse a duzentos quilômetros de distância. Lúcia levantou, abriu todas as janelas, deixou o vento entrar. O cheiro não saiu.

\noindent No celular, uma notificação: uma foto enviada pela mãe. O cachorro no quintal, sentado na grama, olhando para a câmera. O olho direito era amarelo. No canto da foto, atrás do cachorro, havia uma sombra com a forma de uma pessoa. Ninguém estava no quintal.

\noindent Lúcia olhou para a foto por muito tempo. Depois salvou, sem saber por quê, numa pasta chamada ``caso 1.260''. Era a primeira vez que ela usava o nome do paciente para uma imagem que não era documento.

\noindent\textit{\textbf{↪ Rotas:} \rota{LUC-15} \rotasep \rota{LUC-17} \rotasep \rota{MEM-07}}

\newpage\fragdef{LUC-Escolha}% ============================================
% LUC-Escolha
% Gerado de molambudos.md — build_miolo.py
% ============================================
\fragtitle{luc}{L}{LUC}{Escolha}{A Escolha}
\markboth{LUC-Escolha --- A Escolha}{}

\noindent \textit{O que fazer com o diabo que você carrega}

\noindent \textbf{Abril de 2026 --- Barbacena}

\vspace{0.5em}\noindent\rule{\textwidth}{0.5pt}\vspace{0.5em}

\noindent Lúcia passou a noite no único hotel de Barbacena. Um quarto pequeno, mofo nas paredes, ventilador de teto que girava como um carrossel lento.

\noindent Ela colocou os dois diários sobre a mesa: o original (capa marrom, 62 anos de escrita) e o segundo (capa preta, escrito depois da morte).

\noindent Ela não abriu o segundo. Ainda não.

\noindent Em vez disso, sentou na cama e pensou.

\noindent Ela tinha três opções.

\noindent \textbf{Opção 1: Devolver o diário ao arquivo.}

\noindent Fingir que nada aconteceu. Deixar que outro pesquisador o encontrasse. O ciclo continuaria --- mas ela não seria responsável. Ou seria? Ela já tinha lido. A contaminação já tinha começado. Devolver o diário não apagaria os sintomas.

\noindent \textbf{Opção 2: Destruir o diário.}

\noindent Queimar. Rasgar. Dissolver em ácido. Se o suporte físico fosse destruído, o ciclo poderia parar. Mas e o segundo diário? E se a destruição do diário original libertasse a criatura em vez de aprisioná-la?

\noindent \textbf{Opção 3: Aceitar.}

\noindent Aceitar que o ciclo agora fazia parte dela. Aceitar que o olho direito estava mudando. Aceitar que o cheiro não ia embora. Aceitar que ela era, a partir daquele momento, a nova guardiã do diário --- e que, cedo ou tarde, teria que passá-lo adiante.

\noindent \textbf{--- Mas havia uma quarta opção.}

\noindent Ler o segundo diário. Não como reação ao medo — como ato de investigação. Escolher ativamente entrar no ciclo em vez de ser arrastada por ele. Recusar o papel de vítima e assumir o de testemunha. A diferença, Lúcia pensou, era a diferença entre afogar-se e mergulhar.

\vspace{0.5em}\noindent\rule{\textwidth}{0.5pt}\vspace{0.5em}

\noindent A resposta veio às 3:14 da manhã.

\noindent O cheiro acordou ela.

\noindent Não era o cheiro adocicado de antes. Era mais forte. Mais denso. Cheiro de terra molhada. De osso queimado. De coisa saindo do chão.

\noindent Lúcia abriu os olhos.

\noindent O segundo diário estava aberto sobre a mesa.

\noindent Ela não lembrava de ter aberto.

\noindent A página mostrava uma frase, escrita com a mesma letra trêmula:

\noindent \textit{\textquotedblleft{}Ela já escolheu por você. Você só ainda não sabe.\textquotedblright{}}

\noindent Lúcia fechou o diário. Guardou os dois na mochila. Pagou o hotel. Pegou o ônibus das 5 da manhã de volta para Belo Horizonte.

\noindent Durante toda a viagem, o olho direito ardeu.

\noindent Como se algo estivesse crescendo dentro dele.

\noindent Como se algo estivesse \textit{aprendendo a ver}.

\noindent ---

\noindent\textit{\textbf{↪ Rotas:} \rota{LUC-06} \rotasep \rota{LUC-Entrevista}}

\vspace{0.5em}\noindent\rule{\textwidth}{0.5pt}\vspace{0.5em}

\subsection*{ESCALA DE CONTAMINAÇÃO --- ATUALIZAÇÃO 1}

\noindent \textit{O arquivo detectou que você acompanhou a investigação de Lúcia até o momento da escolha.}

\begin{moltabletwo}
\hline
Pontos & Diagnóstico \\
\hline
0--1 & Leitor intacto. Improvável, mas possível. \\
2--4 & Leitor em exposição. Você já leu documentos que não deveria. \\
5--7 & Leitor contaminado. Lúcia escolheu. Você também já fez a sua escolha --- mesmo que não saiba qual foi. \\
8--10 & Paciente 1.263 em formação. A criatura já tem seu nome. \\
\hline
\end{moltabletwo}

\begin{quote}
Nota do arquivo: os pontos não são cumulativos. A escala se ajusta conforme o arquivo se fecha ao seu redor.
\end{quote}

\vspace{0.5em}\noindent\rule{\textwidth}{0.5pt}\vspace{0.5em}
\newpage\fragdef{LUC-06}% ============================================
% LUC-06
% Gerado de molambudos.md — build_miolo.py
% ============================================
\fragtitle{luc}{L}{LUC}{06}{O Desaparecimento}
\markboth{LUC-06 --- O Desaparecimento}{}

\noindent \textit{O ciclo não tem reversa}

\noindent \textbf{Maio de 2026 --- Belo Horizonte}

\vspace{0.5em}\noindent\rule{\textwidth}{0.5pt}\vspace{0.5em}

\noindent O último e-mail de Dra. Lúcia Mendes foi enviado em 12 de maio de 2026, às 2:47 da manhã.

\noindent O destinatário era o Dr. Ronaldo Torres, do Arquivo Público Mineiro.

\vspace{0.5em}\noindent\rule{\textwidth}{0.5pt}\vspace{0.5em}

\noindent \textit{\textquotedblleft{}Ronaldo,}

\noindent \textit{Preciso que você guarde este relato. Se algo acontecer comigo, quero que fique registrado.}

\noindent \textit{Li o diário completo do paciente 1.260. São 62 anos de registro --- de 1917 a 1979. As últimas páginas foram escritas depois da morte dele. Eu sei que parece impossível. Mas as páginas estão lá. A tinta está lá. A caligrafia é a mesma.}

\noindent \textit{Joaquim Antônio Correia não estava louco.}

\noindent \textit{Ele estava possuído.}

\noindent \textit{Eu sei que isso soa absurdo vindo de uma psiquiatra forense. Eu passei 15 anos diagnosticando traumas, dissociações, psicoses. Eu acreditava que tudo tinha uma explicação neurológica ou psicológica.}

\noindent \textit{Agora acredito que existem coisas que a psicologia não explica.}

\noindent \textit{Li o diário. Três dias depois, comecei a sonhar com a vala. Uma semana depois, senti o cheiro. Duas semanas depois, meu olho direito começou a clarear.}

\noindent \textit{O mesmo olho. O mesmo amarelo.}

\noindent \textit{Fiz exames. Tomografia. Ressonância. Sangue. Não há explicação médica.}

\noindent \textit{Fui ao Colônia. Encontrei um caderno de Joaquim escondido numa cela --- dos anos de mutismo. Encontrei um velho que trabalhou lá nos anos 50. Ele me disse: 'A senhora leu o diário. Agora é sua vez.'}

\noindent \textit{Perguntei a vez do quê. Ele disse: 'De hospedar.'}

\noindent \textit{Estou com medo, Ronaldo. O medo não passa pela razão. Fica no corpo. Um medo que vem do estômago. Um medo que tem cheiro doce.}

\noindent \textit{Hoje, olhei no espelho e vi o olho amarelo me olhando de volta.}

\noindent \textit{Não era o meu olho.}

\noindent \textit{Era o olho dele.}

\noindent \textit{Olho de Joaquim.}

\noindent \textit{Ou daquilo que estava dentro dele.}

\noindent \textit{Vou entregar o diário no arquivo amanhã. Vou colocar num pano preto e deixar na recepção. Depois disso, vou sumir por um tempo.}

\noindent \textit{Não procure por mim.}

\noindent \textit{Se alguém perguntar, diga que o diário foi devolvido.}

\noindent \textit{E que o ciclo continua.}

\noindent \textit{Lúcia\textquotedblright{}}

\vspace{0.5em}\noindent\rule{\textwidth}{0.5pt}\vspace{0.5em}

\noindent O diário foi encontrado na recepção do Arquivo Público Mineiro em 13 de maio de 2026, envolto num pano preto.

\noindent Dra. Lúcia Mendes não foi vista em público desde então. Os registros de movimentação do arquivo, porém, mostram que ela voltou ao prédio em julho, sempre à noite, e que assinou a retirada do objeto em 28 de julho.

\noindent Seu apartamento foi encontrado vazio. A cama estava arrumada. O café estava frio na xícara. O computador estava aberto na página de pesquisa: \textit{\textquotedblleft{}Como passar o diário adiante sem ler?\textquotedblright{}}

\noindent Não havia resposta.

\noindent Não há resposta.

\vspace{0.5em}\noindent\rule{\textwidth}{0.5pt}\vspace{0.5em}

\noindent \textit{--- Eu deveria ter queimado o diário. Mas não queimei.}
\textit{--- Eu deveria ter parado de ler. Mas não parei.}
\textit{--- Eu deveria ter fugido antes do olho amarelo. Mas já era tarde quando senti o cheiro.}

\noindent Se você está lendo isto, você já sentiu o cheiro também.
Você já sonhou com a vala.
Você já viu o número 1.260 em tudo.

\noindent \textbf{PARE AGORA.}

\noindent Feche o livro. Devolva. Não termine este capítulo.

\noindent Eu estou escrevendo isto de um computador de lan house em uma cidade que não é a minha. Meu olho direito já não enxerga mais como antes. Tudo tem uma tonalidade amarelada. Como se eu estivesse vendo o mundo através de um olho que não é meu.

\noindent Se você continuar, vai acontecer com você também.

\noindent A criatura não está no diário.
A criatura está na \textit{leitura}.
Ela existe no espaço entre a página e o cérebro.

\noindent Se você está lendo, você já abriu a porta.
Mas ainda pode fechar.

\noindent \textit{--- Lúcia Mendes, paciente 1.262. Último aviso.}

\vspace{0.5em}
\noindent\textit{\textbf{↪ Rotas:} \rota{LUC-07} \rotasep \rota{CONT-05}}
\newpage\fragdef{LUC-07}% ============================================
% LUC-07
% Gerado de molambudos.md — build_miolo.py
% ============================================
\fragtitle{luc}{L}{LUC}{07}{O Rastro de Oliveira}
\markboth{LUC-07 --- O Rastro de Oliveira}{}

\noindent \textit{Seguindo o ciclo reverso}

\noindent \textbf{Abril de 2026 --- Belo Horizonte / Quixeramobim}

\vspace{0.5em}\noindent\rule{\textwidth}{0.5pt}\vspace{0.5em}

\noindent Depois de ler a carta de Oliveira, Lúcia fez o que qualquer psiquiatra forense faria: seguiu o rastro.

\noindent O Dr. Heitor Oliveira desapareceu em 13 de junho de 1979. Seu carro foi encontrado três semanas depois num estacionamento em Quixeramobim, Ceará. A gasolina tinha acabado. As portas estavam destrancadas. Dentro, uma mochila com roupas, um caderno com 32 páginas escritas (que ninguém leu), e uma passagem de ônibus para Barbacena com data de 17 de junho --- o mesmo dia em que ele chegou.

\noindent Ele tinha voltado.

\noindent Lúcia passou uma tarde no Arquivo Público, revirando caixas. Encontrou o caderno de Oliveira --- esquecido, empoeirado, com a etiqueta \textquotedblleft{}Pertences do Dr. H. Oliveira --- 1979\textquotedblright{}.

\noindent Ela leu as 32 páginas em 40 minutos.

\noindent Eram anotações clínicas, mas com uma diferença: a letra mudava ao longo das páginas. Começava firme, médica. Terminava trêmula, irregular, \textit{infantil}.

\noindent A última frase era:

\noindent \textit{\textquotedblleft{}Voltei para Barbacena. O ciclo não me deixa ir. Devolvo o diário ao arquivo. Espero que a próxima leitora entenda melhor do que eu.\textquotedblright{}}

\noindent \textbf{A próxima leitora.}

\noindent Ela.

\noindent Ela era a próxima.

\noindent A fila tinha chegado nela.

\vspace{0.5em}\noindent\rule{\textwidth}{0.5pt}\vspace{0.5em}

\noindent Lúcia passou a noite no arquivo.

\noindent Leu o diário de Oliveira. Leu os relatórios das sessões. Leu a carta que Oliveira deixou para ela --- escrita em 1979, entregue em 2026.

\noindent E, pela segunda vez, leu o diário de Joaquim.

\noindent Só que, dessa vez, ela leu com outros olhos.

\noindent Não como psiquiatra.
Como \textit{herdeira}.

\noindent Cada palavra entrava. Cada frase ecoava. Ela sentiu o cheiro adocicado subir das páginas. O olho direito começou a arder.

\noindent \textit{\textquotedblleft{}O ciclo não tem volta\textquotedblright{}}, disse uma voz dentro dela.

\noindent Não era a voz dela.

\noindent Ela fechou o diário.

\noindent Guardou na mochila.

\noindent E, pela primeira vez, teve certeza de que não ia conseguir parar.

\vspace{0.5em}
\noindent\textit{\textbf{↪ Rotas:} \rota{LUC-08} \rotasep \rota{DOC-11}}
\newpage\fragdef{LUC-08}% ============================================
% LUC-08
% Gerado de molambudos.md — build_miolo.py
% ============================================
\fragtitle{luc}{L}{LUC}{08}{A Última Sessão de Lúcia}
\markboth{LUC-08 --- A Última Sessão de Lúcia}{}

\noindent \textit{A psiquiatra vira paciente}

\noindent \textbf{Maio de 2026 --- Consultório de Lúcia, Belo Horizonte}

\vspace{0.5em}\noindent\rule{\textwidth}{0.5pt}\vspace{0.5em}

\noindent Ela guardou tudo na gaveta.
\noindent Pelo menos foi isso que declarou depois.

\noindent Não como profissional --- como paciente.

\noindent A Dra. Regina Alves era amiga de Lúcia desde a faculdade. Atendia num consultório elegante em BH, com poltronas de couro e uma vista para a serra. Lúcia sentou na poltrona do paciente pela primeira vez em 15 anos.
\noindent \textquotedblleft{\textit{O que está acontecendo, Lúcia?\textquotedblright{}}}
\noindent \textquotedblleft{\textit{Eu estou com medo de estar ficando louca.\textquotedblright{}}}

\noindent \textquotedblleft{\textit{O que você acha que está acontecendo?\textquotedblright{}}}
\noindent \textquotedblleft{\textit{Acho que alguém montou um arquivo para produzir uma reação específica em quem lesse. Ou acho que o arquivo encontrou uma forma de continuar sendo lido. As duas hipóteses me assustam do mesmo jeito.\textquotedblright{}}}

\noindent E alguma coisa tinha acabado de encostar nela.

\noindent \textit{\textquotedblleft{}Por quê?\textquotedblright{}}

\noindent \textit{\textquotedblleft{}Estou tendo sintomas psicóticos.\textquotedblright{}}

\noindent \textit{\textquotedblleft{}Quais?\textquotedblright{}}

\noindent Lúcia listou: o cheiro, os sonhos, o olho amarelo, a sensação de não estar sozinha, a letra da mão mudando.

\noindent Regina ouviu em silêncio. Fez anotações.

\noindent Depois de 20 minutos, ela perguntou:

\noindent \textit{\textquotedblleft{}Tem algo mais?\textquotedblright{}}

\noindent \textit{\textquotedblleft{}...Tem um diário.\textquotedblright{}}

\noindent Ela contou sobre o caso do paciente 1.260. Sobre a criatura. Sobre o ciclo.

\noindent Quando terminou, Regina fechou o bloco.

\noindent \textit{\textquotedblleft{}Lúcia, você está com um quadro de estresse agudo. Possivelmente um transtorno dissociativo desencadeado pelo contato com material traumático.\textquotedblright{}}

\noindent \textit{\textquotedblleft{}Você acha que é sugestão?\textquotedblright{}}

\noindent \textit{\textquotedblleft{}Acho que você leu um caso perturbador, seu cérebro processou o trauma, e agora está produzindo sintomas compatíveis com o que você leu. É um mecanismo de identificação.\textquotedblright{}}

\noindent \textit{\textquotedblleft{}E o olho?\textquotedblright{}}

\noindent \textit{\textquotedblleft{}O que tem o olho?\textquotedblright{}}

\noindent Lúcia tirou os óculos escuros.

\noindent O olho direito estava amarelo.

\noindent Regina inclinou a cabeça.

\noindent \textit{\textquotedblleft{}Isso... é novo?\textquotedblright{}}

\noindent \textit{\textquotedblleft{}Apareceu depois que li o diário.\textquotedblright{}}

\noindent \textit{\textquotedblleft{}Você já foi ao oftalmologista?\textquotedblright{}}

\noindent \textit{\textquotedblleft{}Fui. Ele disse que não há explicação.\textquotedblright{}}

\noindent Regina ficou em silêncio.

\noindent \textit{\textquotedblleft{}Lúcia, você acredita que uma história pode causar uma mudança fisiológica no olho?\textquotedblright{}}
\noindent \textquotedblleft{\textit{Eu acredito que exposição repetida a material traumático pode produzir sintomas reais.\textquotedblright{}}}
\noindent Lúcia preferiu acreditar que era a lanterna --- mesmo depois de 47 anos.
\noindent \textquotedblleft{\textit{É a única resposta honesta que eu tenho.\textquotedblright{}}}

\noindent \textit{\textquotedblleft{}Não.\textquotedblright{}}

\noindent \textit{\textquotedblleft{}Mas está acontecendo.\textquotedblright{}}

\noindent \textit{\textquotedblleft{}Está.\textquotedblright{}}

\noindent \textit{\textquotedblleft{}O que você acha que é?\textquotedblright{}}

\noindent Lúcia respirou fundo.

\noindent \textit{\textquotedblleft{}Eu acho que o que está no diário é real. Eu acho que a criatura existe. Eu acho que ela passou de Joaquim para Oliveira. E de Oliveira para mim.\textquotedblright{}}

\noindent Regina não respondeu.

\noindent Depois de um longo silêncio, ela disse:

\noindent \textit{\textquotedblleft{}Lúcia, como amiga, eu te peço: procura um tratamento. Internação, se necessário.\textquotedblright{}}

\noindent \textit{\textquotedblleft{}E como psicóloga?\textquotedblright{}}

\noindent \textit{\textquotedblleft{}Como psicóloga...\textquotedblright{}\textit{ Regina hesitou. }\textquotedblleft{}Como psicóloga, eu te pergunto: o que você vai fazer com esse diagnóstico?\textquotedblright{}}

\noindent Lúcia levantou.

\noindent \textit{\textquotedblleft{}Vou escrever. Vou registrar tudo. Para que, quando o ciclo chegar na próxima pessoa, ela saiba o que está acontecendo.\textquotedblright{}}

\noindent \textit{\textquotedblleft{}E se a próxima pessoa for eu?\textquotedblright{}}

\noindent Lúcia parou na porta.

\noindent \textit{\textquotedblleft{}Então você vai entender.\textquotedblright{}}

\noindent Ela saiu.

\noindent Nunca mais falou com Regina.

\noindent Três semanas depois, recebeu um envelope pelo correio. Sem bilhete. Sem remetente.
\noindent Dentro havia uma chave pequena, ainda úmida, e uma folha arrancada de um diário que Regina nunca tinha visto.

\noindent Lúcia tirou da bolsa um envelope lacrado.
\noindent \textquotedblleft{\textit{Vou entregar o original para custódia judicial e manter só transcrições parciais. Ninguém precisa tocar no diário de novo.\textquotedblright{}}}
\noindent \textquotedblleft{\textit{Isso é uma decisão clínica ou superstição?\textquotedblright{}}}
\noindent Na porta, antes de sair, deixou uma cópia do laudo sobre a mesa de Regina.
\noindent \textquotedblleft{\textit{Se eu desaparecer, entregue isto ao Conselho. Não leia o diário. Leia só o laudo.\textquotedblright{}}}

\noindent Não abriu.

\noindent Não leu.

\noindent Até hoje.

\vspace{0.5em}
\noindent\textit{\textbf{↪ Rotas:} \rota{LUC-09} \rotasep \rota{CONT-04}}
\noindent Ela marcou uma sessão com uma psicóloga.
\newpage\fragdef{LUC-09}% ============================================
% LUC-09
% Gerado de molambudos.md — build_miolo.py
% ============================================
\fragtitle{luc}{L}{LUC}{09}{A Exumação de Oliveira}
\markboth{LUC-09 --- A Exumação de Oliveira}{}

\noindent \textit{47 anos depois, o crânio na vala}

\noindent \textbf{Junho de 2026 --- Quixeramobim, Ceará}

\vspace{0.5em}\noindent\rule{\textwidth}{0.5pt}\vspace{0.5em}

\noindent Lúcia foi ao Ceará no dia 13 de junho de 2026. Exatos 47 anos depois da morte de Joaquim e do desaparecimento de Oliveira.

\noindent Ela não sabia o que esperava encontrar. A vala estava lá --- a clareira de terra preta, a árvore torta, o cheiro adocicado que subia do chão como se a terra estivesse respirando.

\noindent Ela se ajoelhou.

\noindent Enterrou as mãos na terra preta.

\noindent A terra estava quente.

\noindent \textit{\textquotedblleft{}Dr. Oliveira?\textquotedblright{}\textit{, ela chamou. }\textquotedblleft{}O senhor está aqui?\textquotedblright{}}

\noindent O vento respondeu.

\noindent \textit{\textquotedblleft{}Eu li o seu diário. Li o de Joaquim. Eu sei o que aconteceu. Eu estou sentindo o mesmo. O olho. O cheiro. Os sonhos. Está acontecendo comigo também.\textquotedblright{}}

\noindent Ela tirou uma pá da mochila.

\noindent Suas mãos tremiam, mas ela não parou. Tinha medo do que ia encontrar --- e mais medo ainda de não encontrar nada, de descobrir que também estava enlouquecendo sozinha, como Joaquim, como Oliveira antes dela.

\noindent E começou a cavar.

\vspace{0.5em}\noindent\rule{\textwidth}{0.5pt}\vspace{0.5em}

\noindent Cavou por duas horas.

\noindent A terra era fofa, oleosa. Fácil de cavar. Como se já tivesse sido remexida muitas vezes.

\noindent A um metro de profundidade, a pá bateu em algo duro.

\noindent Ela cavou com as mãos.

\noindent Era um osso.

\noindent Não de animal. Humano.

\noindent Ela cavou mais. Uma costela. Uma vértebra. Um crânio.

\noindent O crânio tinha uma fissura na têmpora direita. Como de um objeto pontiagudo.

\noindent E dentro da órbita direita --- a órbita ainda estava escura, ainda guardava resquícios de tecido --- algo brilhou.
\noindent Um reflexo amarelo, fraco, no fundo da órbita.
\noindent Podia ser a lanterna. Podia ser mica no barro.

\noindent Ela aproximou o rosto.

\noindent Um reflexo amarelo.

\noindent O olho ainda brilhava.

\noindent Mesmo depois de 47 anos.

\noindent Lúcia sentou no chão da vala, com o crânio que ela supunha ser de Oliveira na mão, o gosto de açúcar queimado nos pulmões, e uma certeza subindo pelo tornozelo.

\noindent Era uma corrente de corpos.

\noindent E ela era a próxima.

\vspace{0.5em}
\noindent\textit{\textbf{↪ Rotas:} \rota{LUC-10} \rotasep \rota{DOC-15}}
\newpage\fragdef{LUC-10}% ============================================
% LUC-10
% Gerado de molambudos.md — build_miolo.py
% ============================================
\fragtitle{luc}{L}{LUC}{10}{O Laudo de Lúcia}
\markboth{LUC-10 --- O Laudo de Lúcia}{}

\noindent \textit{O autorretrato diagnóstico}

\noindent \textbf{Julho de 2026 --- Belo Horizonte}
\textbf{Último registro profissional de Dra. Lúcia Mendes}

\vspace{0.5em}\noindent\rule{\textwidth}{0.5pt}\vspace{0.5em}

\noindent \textbf{Parecer Psiquiátrico --- Paciente: Lúcia Mendes}
\textbf{CRM-MG 28.391}
\textbf{Data: 28 de julho de 2026}

\vspace{0.5em}\noindent\rule{\textwidth}{0.5pt}\vspace{0.5em}

\noindent Estou escrevendo este laudo sobre mim mesma.

\noindent Sei que não é ético. Mas não há outro psiquiatra que possa avaliar meu caso. Porque meu caso não é comum. Meu caso é uma contaminação.

\noindent \textbf{Histórico:}

\noindent Em janeiro de 2026, solicitei acesso ao processo do paciente 1.260 (Joaquim Antônio Correia) para fins de pesquisa acadêmica. A primeira leitura integral do diário ocorreu em abril. O diário descreve o encontro de Joaquim com uma entidade não identificada (\textquotedblleft{}a criatura\textquotedblright{}) em uma vala de corpos no Campo do Alagadiço, próximo a Fortaleza, em 1915.

\noindent Após a leitura, desenvolvi:

\noindent 1. \textbf{Sintomas somáticos}: Olho direito com alteração de coloração para amarelo-dourado. Relato de \textquotedblleft{}cheiro adocicado\textquotedblright{} constante. Fadiga crônica.

\noindent 2. \textbf{Sintomas psíquicos}: Pesadelos recorrentes com a vala. Pensamentos intrusivos sobre a criatura. Sensação de presença no quarto. Hipervigilância. Desrealização.

\noindent 3. \textbf{Sintomas comportamentais}: Isolamento social. Dificuldade de concentração. Necessidade compulsiva de reler o diário.

\noindent \textbf{Diagnóstico:}

\noindent - \textbf{F43.10 --- TEPT Complexo}: Preencho todos os critérios (A a G). A exposição foi a leitura do diário. A reexperiência ocorre em sonhos e flashbacks diurnos. A evitação é inútil --- o diário me segue.

\noindent - \textbf{F44.81 --- Transtorno Dissociativo Leve}: Sinto que há uma \textquotedblleft{}segunda voz\textquotedblright{\textit{ dentro de mim. Não é alucinação auditiva --- é uma }presença\textit{ que comenta, observa, espera. A voz diz: }\textquotedblleft{}Você é minha agora.\textquotedblright{}}

\noindent - \textbf{Z65.8 --- Circunstância de Privação Extrema}: Fui privada da minha sanidade por um texto. A literatura nunca me pareceu perigosa. Agora sei que algumas palavras são armadilhas.

\noindent \textbf{Considerações Finais:}

\noindent Emito este laudo para registro. Não como profissional --- como vítima.

\noindent O caso do paciente 1.260 não é um caso psiquiátrico. É um caso de \textit{transmissão}. A criatura não é um delírio. É um parasita narrativo. Ela se transmite pela leitura.

\noindent Quem lê, recebe.
Quem recebe, carrega.
Quem carrega, passa adiante.

\noindent Eu li.
Eu recebi.
Eu carrego.

\noindent Agora preciso passar adiante.

\noindent Para você.

\vspace{0.5em}\noindent\rule{\textwidth}{0.5pt}\vspace{0.5em}

\noindent \textit{--- Dra. Lúcia Mendes}
\textit{--- Psiquiatra Forense}
\textit{--- Paciente 1.262}
\textit{--- Próxima: você}

\vspace{0.5em}
\noindent\textit{\textbf{↪ Rotas:} \rota{LUC-11} \rotasep \rota{DOC-08}}
\newpage\fragdef{LUC-11}% ============================================
% LUC-11
% Gerado de molambudos.md — build_miolo.py
% ============================================
\fragtitle{luc}{L}{LUC}{11}{O Diário Digital}
\markboth{LUC-11 --- O Diário Digital}{}

\noindent \textit{A criatura migra para o pixel}

\noindent \textbf{Agosto de 2026 --- Belo Horizonte}

\vspace{0.5em}\noindent\rule{\textwidth}{0.5pt}\vspace{0.5em}

\noindent Lúcia sabia que não podia guardar o diário para sempre.
\noindent Ela ia morrer --- como Joaquim, como Oliveira, como todos os que guardaram o diário tempo demais. O texto precisava se mover. Quando parava, criava raiz. Primeiro na rotina. Depois no sono. Depois no corpo.

\noindent Ela ia morrer --- como Joaquim, como Oliveira, como todos os que guardaram o diário sem passar adiante. O ciclo só se mantém enquanto o texto se move. Se o texto para, a criatura se enraíza no hospedeiro até consumi-lo.

\noindent Mas ela não queria condenar outra pessoa.

\noindent Então criou uma solução: o diário digital.

\noindent Antes de devolvê-lo, em maio, passou duas semanas digitando cada página. Cada palavra. Cada desenho. Cada rasura. Transformou o diário de Joaquim em um arquivo PDF. Depois, escreveu seu próprio relato --- a história de como encontrou o diário, leu, e foi contaminada.

\noindent Depois, escreveu o relato de Oliveira. Depois, o laudo completo com os diagnósticos.

\noindent Tudo em um único arquivo.

\noindent No final, acrescentou uma nota:

\noindent \textquotedblleft{\textit{Este arquivo contém o diário do paciente 1.260, os relatos do Dr. Heitor Oliveira (paciente 1.261) e meu relato (paciente 1.262). A leitura pode produzir sonhos, lacunas, alterações de percepção e sensação persistente de presença. Não leia sozinho.\textquotedblright{}}}

\noindent \textit{\textquotedblleft{}Mas, se você leu até aqui, já é tarde.\textquotedblright{}}

\noindent \textit{\textquotedblleft{}O ciclo não respeita avisos.\textquotedblright{}}

\noindent \textit{\textquotedblleft{}Bem-vindo, paciente 1.263.\textquotedblright{}}

\noindent Ela salvou o arquivo num pendrive. Colocou o pendrive na mesma caixa que o diário original. Colocou a caixa num envelope. Endereçou ao Arquivo Público Mineiro.

\noindent E enviou.

\noindent Agora, o diário existe em dois lugares: o papel e o pixel.

\noindent A criatura aprendeu a migrar.

\noindent A fome aprendeu a caber no brilho de uma tela.

\noindent Ela precisa de \textit{leitura}.

\noindent E o mundo está cheio de leitores.

\vspace{0.5em}\noindent\rule{\textwidth}{0.5pt}\vspace{0.5em}

\noindent \textbf{Teste de contaminação --- responda honestamente:}

\noindent Você está lendo esta página. Até aqui, tudo normal.

\noindent Mas eu preciso saber: \textbf{quantas vezes você leu esta linha?}

\noindent - A) Uma vez. Segui em frente. → \textit{(Você ainda está no controle.)}
- B) Duas vezes. Achei que tinha lido errado. → \textit{(A criatura já cutucou sua atenção.)}
- C) Três vezes ou mais. → \textit{(Você já está dentro. A repetição é o primeiro sintoma.)}

\noindent A resposta não importa para mim. Importa para você.

\noindent Não precisava mais de cova.
Mas você --- paciente 1.263 --- ainda pode escolher.

\noindent \textbf{Pergunta final:} Você confia na sua própria leitura? Ou já tem uma voz lendo com você?

\noindent \textit{--- Se ouvir duas vozes quando lê, a criatura já está aí.}

\vspace{0.5em}
\noindent\textit{\textbf{↪ Rotas:} \rota{LUC-12} \rotasep \rota{CONT-07}}
\newpage\fragdef{LUC-12}% ============================================
% LUC-12
% Gerado de molambudos.md — build_miolo.py
% ============================================
\fragtitle{luc}{L}{LUC}{12}{A Última Foto}
\markboth{LUC-12 --- A Última Foto}{}

\noindent \textit{Olho amarelo na câmera de segurança}

\noindent \textbf{Setembro de 2026 --- Desconhecido}

\vspace{0.5em}\noindent\rule{\textwidth}{0.5pt}\vspace{0.5em}

\noindent A última foto de Dra. Lúcia Mendes foi tirada por uma câmera de segurança na rodoviária de Barbacena, em 3 de setembro de 2026.

\noindent Ela está em pé, na plataforma de embarque, com uma mochila nas costas e um livro na mão. O livro não é o diário --- é um romance qualquer, comprado na banca da rodoviária. O diário estava na mochila.

\begin{figure}[H]
  \centering
  \includegraphics[width=0.55\textwidth]{foto_04_lucia_rodoviaria_sepia.png}
  \caption*{\footnotesize\textbf{Fotografia 5} --- \textit{Registro de câmera de segurança. Rodoviária de Barbacena, 3 de setembro de 2026, 14h23.}}
\end{figure}

\noindent O olho direito dela está visível na foto.

\noindent Está amarelo.

\noindent Brilhando.

\noindent Ela está olhando para a câmera. Diretamente. Como se soubesse que estava sendo filmada. Como se quisesse que a última imagem dela fosse aquela: o olho amarelo, o aviso silencioso.

\noindent Ela entrou num ônibus com placa de Fortaleza.

\noindent Nunca mais foi vista.

\vspace{0.5em}\noindent\rule{\textwidth}{0.5pt}\vspace{0.5em}

\noindent O diário foi encontrado no Arquivo Público Mineiro uma semana depois, dentro de uma caixa, com um bilhete:

\noindent \textit{\textquotedblleft{}Paciente 1.262 encerra seu ciclo. O paciente 1.263 é quem abrir esta caixa.\textquotedblright{}}

\noindent O arquivista que abriu a caixa se chama Ronaldo Torres.

\noindent Ele leu o bilhete.

\noindent Guardou a caixa.

\noindent Não abriu o diário.

\noindent Mas, à noite, sonhou com olhos amarelos.

\noindent E sentiu o cheiro.

\noindent E soube que, cedo ou tarde, ele ia abrir.

\noindent Todo mundo abre.

\noindent O ciclo não espera.

\noindent O ciclo só se alimenta.

\vspace{0.5em}\noindent\rule{\textwidth}{0.5pt}\vspace{0.5em}

\subsection*{Despedida de Lúcia}

\noindent \textit{Encontrado em 15 de outubro de 2026, dentro de um envelope sem remetente, endereçado ao Arquivo Público Mineiro. O papel está manchado de água salgada --- lágrimas ou maresia, não é possível determinar. A caligrafia é de Lúcia, mas trêmula, como se escrita em movimento.}

\noindent \textit{\textquotedblleft{}Ronaldo,}

\noindent \textit{Cheguei a Fortaleza. O diário não está comigo. Eu o peguei de volta no cofre em julho e o enviei para vocês antes de pegar o ônibus. Eu precisava ter certeza de que ele continuaria o ciclo --- mas também precisava descobrir até onde a contaminação vai sem o objeto físico.}

\noindent \textit{Descobri.}

\noindent \textit{O diário é só um suporte. Como uma tomada elétrica. A corrente é a leitura. E a corrente não para quando você desliga o aparelho --- ela continua nos fios.}

\noindent \textit{Os fios, no meu caso, são os olhos.}

\noindent \textit{Parei de enxergar azul há três semanas. O mundo agora tem só duas cores: amarelo e cinza. O amarelo é a cor da criatura. O cinza é tudo que ela não tocou.}

\noindent \textit{Ela está crescendo dentro de mim. Eu sinto. A coisa tem corpo. Não é figura de linguagem. Como se houvesse um segundo coração batendo no meu peito --- em um ritmo que não é o meu.}

\noindent \textit{Descobri uma coisa que Joaquim sabia mas não escreveu: a criatura não é uma entidade. É um processo. Um padrão que se replica através da leitura. Como um vírus de linguagem. Ela não quer possuir corpos --- quer possuir narrativas. O corpo é só o meio de transporte.}

\noindent \textit{Nós somos os veículos.}

\noindent \textit{O que ela quer é ser lida. Eternamente. Sempre sendo passada adiante. Sempre sendo interpretada. Cada leitura é uma nova geração do ciclo.}

\noindent \textit{Joaquim não sabia que estava escrevendo um livro. Ele achava que estava escrevendo um diário. Mas um diário é pessoal. Um livro é transmissível. E tudo que é transmissível se replica.}

\noindent \textit{Ela me ensinou isso.}

\noindent \textit{Ela falou comigo ontem à noite. Pela primeira vez. Não em sonho --- acordada. As palavras vieram de dentro, mas não eram minhas:}

\noindent \textit{'Você não é a paciente 1.262. Você é a página 1.262. O livro não acaba. O livro só muda de leitor.'}

\noindent \textit{É isso que eu queria te dizer.}

\noindent \textit{Não queime o diário. Não adianta. O padrão já está na sua cabeça. Está na cabeça de todo mundo que leu. Cada um de nós é uma página viva.}

\noindent \textit{Eu vou continuar escrevendo.}

\noindent \textit{Mas não com caneta.}

\noindent \textit{Com a minha vida.}

\noindent \textit{--- Lúcia.}

\noindent \textit{P.S.: Se você ler esta carta até o fim, você já leu. E se você leu, você já é página também.}

\noindent \textit{P.S.2: O azul era minha cor favorita. Agora não consigo lembrar como ele era.\textquotedblright{}}

\vspace{0.5em}
\noindent\textit{\textbf{↪ Rotas:} \rota{CONT-07} \rotasep \rota{DOC-01} \rotasep \rota{LUC-11} \rotasep \rota{EPI-01} \rotasep \rota{LUC-13}}
\newpage\fragdef{LUC-17}% ============================================
% LUC-17 — O Medo (edição revista R1075)
% ============================================
\fragtitle{luc}{L}{LUC}{17}{O Medo}
\markboth{LUC-17 --- O Medo}{}

\noindent \textit{Belo Horizonte --- 12 de abril de 2026, 3h14}

\noindent Lúcia acordou antes de olhar o relógio. Sabia que eram 3h14. Sabia pelo peso do ar, pelo modo como o silêncio do apartamento ficava mais espesso, como se alguém tivesse entrado e estivesse parado no corredor, esperando que ela se mexesse.

\noindent Ficou imóvel. O teto era uma superfície branca, lisa, sem nenhuma mancha. Contou até sessenta e dois. Depois levantou.

\noindent O diário estava na sala, dentro da caixa de metal, sobre a mesa. A chave estava num copo com água, ao lado. Lúcia tinha feito aquilo na noite anterior, depois de ler a última página. Era um gesto absurdo, ela sabia. Mas se acordasse com a chave seca na gaveta, saberia que precisava parar. Se acordasse com a água morna, saberia que algo tinha acontecido.

\noindent A água estava morna.

\noindent Ela não tocou na chave. Foi até a cozinha, bebeu um copo de água, olhou pela janela. A rua estava vazia. Um cachorro atravessou a calçada, sem pressa, e sumiu atrás de um carro. O olho direito dele brilhou quando passou sob o poste. Lúcia piscou. Quando abriu os olhos, o cachorro já não estava.

\noindent Voltou para a sala. Sentou no chão, de costas para o sofá, com a caixa de metal à frente. Não abriu. Ficou olhando para o cadeado, para a chave dentro do copo, para o reflexo da lâmpada na água.

\noindent Havia uma pasta sobre a mesa, com uma folha em branco e uma caneta. Lúcia pegou a caneta. Escreveu:

\noindent ``Mãe, se algo acontecer comigo, não leia o que está na caixa. Não importa o que encontrarem. Não importa o que eu tiver escrito. Queime.''

\noindent Olhou para a frase. A letra era a dela. Mas o ``M'' de ``Mãe'' estava inclinado para a esquerda, com um traço longo na perna direita. Ela não escrevia assim. Conhecia a própria letra desde a infância: redonda, apertada, sempre para a direita. Aquela letra era de outra mão.

\noindent Rasgou a folha. Rasgou em quatro, em oito, em dezesseis. Jogou os pedaços no lixo. Depois tirou os pedaços do lixo, colocou na pia, acendeu um fósforo. O papel queimou rápido, com uma chama azul no centro. As cinzas ficaram na pia, finas, escuras, com um cheiro doce que não era de papel queimado.

\noindent Lúcia abriu a torneira. Lavou as cinzas. Elas desceram pelo ralo, mas o cheiro ficou.

\noindent Ela foi até a porta. Verificou a fechadura. Estava trancada. Verificou a janela da sala. Trancada. A do quarto. Trancada. A da cozinha. Trancada. Voltou para a sala. Sentou de novo no chão. O apartamento era pequeno; ela conhecia cada ruído. O frigorífico. O cano da parede. O elevador no poço. Naquela noite, nenhum desses ruídos existia. O silêncio era total, como se o mundo tivesse sido coberto por uma camada de algodão.

\noindent Então ouviu a respiração.

\noindent A respiração era mais lenta, mais funda, vindo do corredor. Lúcia não se virou. Ficou olhando para a caixa de metal. A respiração parou. O silêncio voltou. Ela esperou. Nada.

\noindent ``Eu sei que você está aí'', ela disse.

\noindent A voz saiu baixa, firme, com uma calma que não sentia. O apartamento não respondeu.

\noindent ``Eu não vou fingir que não estou com medo'', ela continuou. ``Estou. Mas o medo não decide por mim. Quem decide sou eu.''

\noindent A água no copo tremeu. Não havia vento. Não havia vibração. A água tremeu como se alguém tivesse batido na mesa, muito de leve, do outro lado.

\noindent Lúcia se levantou. Pegou o copo. Bebeu a água. Não sentiu gosto de nada. Colocou o copo vazio na mesa. Pegou a chave. A chave estava quente, como se tivesse sido segurada por muito tempo. Ela a encaixou no cadeado. Abriu a caixa.

\noindent O diário estava lá, envolto no pano preto. Ela o tirou. Pesava mais do que devia. Abriu na última página. A letra era de Joaquim. A data era de 1979. No alto da página, escrito com outra tinta, havia uma frase que ela não tinha visto antes:

\noindent ``Você demorou.''

\noindent Lúcia fechou o diário. Ficou com ele no colo, olhando para a parede. O medo continuava ali, mas tinha mudado de lugar. Já não estava no corredor. Estava dentro dela, como um hóspede que não pede licença.

\noindent Ela pensou: \textit{eu ainda posso parar}. Podia devolver o diário, pedir afastamento, nunca mais abrir. Podia ligar para Ronaldo e dizer que desistia. Podia pegar o carro e dirigir até a casa da mãe. Podia fazer qualquer coisa, menos aquilo que sabia que ia fazer.

\noindent Abriu o diário de novo. Leu a primeira linha da última página. Depois a segunda. Depois a terceira. Quando chegou ao fim, já não havia mais medo. Havia uma espécie de fome.

\noindent\textit{\textbf{↪ Rotas:} \rota{LUC-13} \rotasep \rota{LUC-18} \rotasep \rota{CONT-01}}

\newpage\fragdef{LUC-13}% ============================================
% LUC-13
% Extraído de main.pandoc.tex
% ============================================
\fragtitle{luc}{L}{LUC}{13}{A Noite do Diagnóstico}
\markboth{LUC-13 --- A Noite do Diagnóstico}{}

\noindent \textit{Documento extraído do arquivo pericial.}

\vspace{0.5em}\noindent\rule{\textwidth}{0.5pt}\vspace{0.5em}

\noindent \textbf{Abril de 2026 --- Belo Horizonte}\\
\noindent \textbf{Gravação de áudio --- arquivo L260-13A.wav}\\
\noindent \textbf{Duração:} 47 minutos e 32 segundos.\\
\noindent \textbf{Status:} não transcrito. A Dra. Lúcia Mendes não autorizou a transcrição.

\noindent \textit{Nota do arquivista: este arquivo de áudio foi encontrado no computador de Dra. Lúcia Mendes, datado de 12 de abril de 2026. É o primeiro registro pessoal dela sobre o caso. A gravação foi feita em casa, às 3:47 da manhã. Transcrevemos os primeiros 4 minutos e 12 segundos --- o restante é ininteligível.}
\noindent \textbf{[00:00:00] Lúcia:} Teste. Teste de áudio. L260... L260... treze. A. Gravando.

\noindent \textbf{[00:00:12]\textit{ É a terceira noite. Terceira noite que eu acordo no mesmo horário. 3:14. Sempre 3:14. Não adianta ajustar o relógio. Não adianta não olhar a hora. Eu }sei} que são 3:14 antes de olhar. Sinto na respiração. No peso do ar.

\noindent \textbf{[00:00:41]} Tentei todas as hipóteses. Ansiedade antecipatória. Condicionamento clássico. Sugestão pós-hipnótica induzida pela leitura. Tudo o que eu ensinaria num curso de psicopatologia para residentes. Quinze anos de clínica e eu sei exatamente como um paciente descreve um despertar noturno recorrente e como se diferencia um transtorno do sono de uma ideação persecutória.

\noindent \textbf{[00:01:18]} Mas não é um paciente. Sou eu.

\noindent \textbf{[00:01:25]\textit{ E o que eu sinto quando acordo não é ansiedade. É }expectativa}. Como se algo fosse acontecer. Como se algo estivesse \textit{se preparando} para acontecer. O corpo espera. O corpo \textit{sabe} de algo que a mente ainda não autorizou.

\noindent \textbf{[00:02:03]} O cheiro voltou hoje. Tinha passado o dia inteiro sem sentir. Pensei: passou. Foi sugestão. O cérebro cansado fabrica estímulos onde não existem. Aí desliguei a luz do quarto, deitei, fechei os olhos... e o cheiro estava no travesseiro. Doce. Não doce de perfume. Doce de matéria vegetal murchando dentro de água parada. Doce de buquê esquecido num copo sujo até a água ficar preta.

\noindent \textbf{[00:02:47]} Levantei, troquei a fronha. Cheirei o travesseiro novo. Limpo. Lavanda. O cheiro continuou. Não estava no travesseiro. Estava dentro do meu nariz.

\noindent \textbf{[00:03:12]} Acendi todas as luzes. Abri a janela. Respirei fundo. O ar da rua estava frio, normal, urbano. Poluição. Diesel. Asfalto molhado. Coisas reais. Fechei a janela. O cheiro voltou na mesma intensidade.

\noindent \textbf{[00:03:41]} Ele está \textit{em mim}. Não no ambiente. Na mucosa. No fundo da garganta. No lugar onde a memória olfativa se aloja e não sai nunca mais. Como o cheiro da casa da infância. Como o cheiro de alguém que se foi.

\noindent \textbf{[00:04:12]} Fui ao banheiro e olhei no espelho. O olho direito...

\noindent \textit{Aqui a gravação se interrompe por 23 segundos. O que se segue foi recuperado parcialmente por restauração de áudio. A voz de Lúcia está diferente: mais baixa, mais lenta, como se falar exigisse esforço físico.}
\noindent \textbf{[00:04:35] Lúcia:} ...está mais claro. Não muito. O suficiente. A íris perdeu pigmento na borda externa. Ficou cor de... não sei. O amarelo ainda não chegou. É um tom que não tem nome. Um tom que eu não quero nomear.

\noindent \textbf{[00:05:11]} Toquei o olho com a ponta do dedo. A córnea estava morna. Olhos não são mornos. Olhos são frios. A temperatura da superfície ocular é 32 graus. O que eu senti era mais. Bem mais.

\noindent \textbf{[00:05:44]\textit{ Fiz um diário. Não. Fiz um }registro}. Uma planilha. Sintomas, horários, intensidade, frequência. Durou três dias. No terceiro dia, percebi que a planilha estava igual à tabela que Joaquim descreve no diário dele. A mesma estrutura. As mesmas colunas. As mesmas palavras. A minha letra, mas a mesma estrutura.

\noindent \textbf{[00:06:28]\textit{ Eu não copiei. Eu juro que não copiei. A tabela }veio}. Veio pronta. A mão escreveu sozinha.

\noindent \textbf{[00:07:01]} Rasguei a planilha. Joguei no lixo. Depois tirei do lixo, queimei na pia, lavei as cinzas. Não quero que ninguém encontre.

\noindent \textbf{[00:07:33]} Se eu lesse um relato como este vindo de uma paciente, meu diagnóstico seria: ideação paranóide com somatização, possível quadro dissociativo leve, estresse pós-traumático deflagrado por contato com material traumático.

\noindent \textbf{[00:08:12]} Mas o que eu sinto não é paranoia. Paranoia é achar que alguém está te observando. O que eu sinto é diferente. Eu sinto que alguém \textit{já está dentro}. O que Joaquim descreve como \textquotedblleft{}o peso atrás do esterno\textquotedblright{}, \textquotedblleft{}a respiração que não é minha\textquotedblright{}, \textquotedblleft{}a boca que se move sem eu mandar\textquotedblright{}. Eu sei o que ele quis dizer. Porque estou sentindo.

\noindent \textbf{[00:08:52]} Na madrugada eu sentei pra escrever. Uma coisa que eu queria registrar. Peguei o caderno. Abri na primeira folha. A mão escreveu: \textit{tem tanta fome aqui, Lúcia. Tanta fome que eu nem sei onde caber.}

\noindent \textbf{[00:09:07]\textit{ Essa frase está no diário. É a primeira fala da entidade. Eu }li} ela. Mas quando a mão escreveu, a letra não era minha. A letra era a mesma descrição que o Dr. Oliveira fez na carta dele: inclinada pra esquerda, \textquotedblleft{}a\textquotedblright{} aberto, pressão irregular. A letra de Joaquim.

\noindent \textbf{[00:09:41]\textit{ Eu não copiei. Eu não }precisava} copiar. A mão sabia o caminho sozinha. A mão lembrou de um padrão que eu vi no papel e reproduziu sem autorização. A mão tem memória própria.

\noindent \textbf{[00:10:03]} O mesmo padrão. A mesma estrutura. Eu escrevendo a frase de Joaquim. Eu sentindo o que ele sentiu. Eu vendo o que ele viu. Eu.

\noindent \textbf{[00:10:28]} Antes eu falava \textquotedblleft{}Joaquim descreve\textquotedblright{}, \textquotedblleft{}Joaquim relata\textquotedblright{}, \textquotedblleft{}o paciente diz\textquotedblright{}. Agora eu falo \textquotedblleft{}a vala\textquotedblright{}, \textquotedblleft{}o olho\textquotedblright{}, \textquotedblleft{}ela\textquotedblright{}. O nome próprio dele está sumindo do meu vocabulário. Ele não é mais o paciente. Ele é a testemunha. A testemunha que depõe dentro de mim.

\noindent \textbf{[00:11:12]} Daqui a pouco eu vou falar \textquotedblleft{}nós\textquotedblright{}.

\noindent \textit{[A gravação registra um ruído longo, semelhante a folhas sendo viradas depressa.]}
\noindent \textbf{[00:12:04]} Não vem. Não vem. Não vem. Ela não pode estar aqui. Ela não pode estar dentro de mim. Eu sou psiquiatra. Eu tenho 42 anos. Eu não acredito em possessão. Eu não acredito em contágio. Eu não acredito em nada, em nada, em nada...\NE{o som de folhas sendo viradas. Quem está virando as páginas?}

\noindent A partir deste ponto, a gravação se deteriora. A voz de Lúcia alterna entre silêncios longos e frases curtas, desconectadas. A qualidade do áudio piora --- há um ruído de fundo que os técnicos descrevem como \textquotedblleft{\textit{semelhante a vento passando por um duto estreito\textquotedblright{}.}

\noindent \textbf{[00:12:04] Lúcia:} O diário está no escritório. Eu tranquei a porta.

\noindent \textbf{[00:12:31]} Não adianta. Sinto ele do outro lado da parede. A capa. O peso. A lombada.

\noindent \textbf{[00:13:02]\textit{ Ele }irradia}.

\noindent \textbf{[00:13:28]} Uma hora dessas... 3:47 da manhã... e eu sei que ele está acordado. O diário está acordado. As páginas estão se movendo sozinhas no escuro. A tinta está escorrendo para lugares que eu não escrevi.

\noindent \textbf{[00:14:15]} Isto não é científico. Isto não é clínico. Isto é o relato de uma mulher adulta, com pós-doutorado, cujo cérebro está mentindo para ela.

\noindent \textbf{[00:14:52]} \textit{(longa pausa)} E se não for mentira?

\noindent Os 33 minutos seguintes são compostos principalmente de silêncio, respiração irregular e ruído de fundo. Aos 47 minutos e 12 segundos, a gravação termina com o que parece ser a voz de Lúcia sussurrando: \textquotedblleft{\textit{Paciente 1.262.\textquotedblright{} A gravação foi encerrada manualmente em 47:32. Nenhum outro arquivo de áudio foi encontrado no computador.}

\noindent Nota do arquivista: na manhã seguinte a esta gravação, Lúcia Mendes ligou para o Arquivo Público Mineiro e solicitou a devolução do diário. Disse que precisava \textquotedblleft{\textit{confirmar uma informação}\textquotedblright{}} na contracapa. O diário foi retirado por ela pessoalmente no mesmo dia. Ela não o devolveu até 13 de maio.}

\vspace{0.5em}
\noindent\textit{\textbf{↪ Rotas:} \rota{DOC-01} \rotasep \rota{LUC-02} \rotasep \rota{EPI-01} \rotasep \rota{LUC-14}}

}
\newpage\fragdef{LUC-18}% ============================================
% LUC-18 — O Que Ronaldo Não Disse (edição revista R1075)
% ============================================
\fragtitle{luc}{L}{LUC}{18}{O Que Ronaldo Não Disse}
\markboth{LUC-18 --- O Que Ronaldo Não Disse}{}

\noindent \textit{Arquivo Público Mineiro --- 14 de abril de 2026, 19h05}

\noindent A sala de Ronaldo ficava no fim do corredor, depois da mapoteca. Lúcia conhecia o caminho de cor. Havia uma janela alta, com vista para o pátio interno, e uma estante de madeira escura que ele nunca terminara de organizar. Sobre a mesa, o mesmo abajur de latão, o mesmo porta-canetas de vidro, a mesma fotografia da mãe. Nada mudara em oito anos, exceto o quadro na parede: uma reprodução de uma vala comum, tirada de um livro de fotografia.

\noindent ``Você mandou meu nome para a comissão de ética'', ele disse, sem levantar os olhos.

\noindent ``Mandei.''

\noindent ``Por quê?''

\noindent ``Porque você me deu acesso a um documento que não estava catalogado. Porque você sabia que ele estava lá e não me contou. Porque você está com medo, Ronaldo. E medo não é motivo para colocar outra pessoa na frente.''

\noindent Ele levantou os olhos. Estava cansado. Havia uma olheira funda, nova, que não existia em janeiro. O olho direito, sob a luz do abajur, parecia mais claro. Lúcia disse a si mesma que era impressão.

\noindent ``Meu pai foi internado no Colônia em 1968'', ele disse.

\noindent Lúcia ficou parada.

\noindent ``Ele tinha 34 anos. Diagnóstico: esquizofrenia. Ficou lá seis meses. Nunca falou sobre isso. Quando eu era criança, ele tinha crises. Ficava horas no quintal, olhando para o muro. Minha mãe dizia que era o sol. Uma vez, ele me pegou pelo braço e disse: \textit{Não deixa ninguém ler o que está no meu armário.} Eu tinha nove anos. Não entendi. Depois esqueci.''

\noindent ``O que tinha no armário?''

\noindent ``Um caderno. Capa preta. Ele queimou antes de morrer. Eu vi. Ele colocou no quintal, jogou álcool e acendeu. Ficou olhando até virar cinza. Depois me olhou e disse: \textit{Agora está limpo.}''

\noindent Lúcia sentiu o cheiro de terra de novo. Não no ar. Na memória. Um cheiro que ela nunca tinha sentido na infância, mas que agora parecia ter estado sempre ali, escondido atrás das coisas.

\noindent ``Por que você me mandou o e-mail?''

\noindent Ronaldo abriu a gaveta. Tirou um envelope pardo, gasto, com o carimbo do arquivo. Colocou sobre a mesa.

\noindent ``Porque eu li a primeira página'', ele disse. ``Em 2019. Quando o diário chegou aqui, na primeira devolução. Eu abri. Só a primeira página. Vi a letra. Vi o nome. Fechei. Não contei a ninguém. Coloquei numa caixa e não abri mais. Mas o cheiro ficou comigo. Toda vez que eu passava perto da caixa, sentia. Toda vez que eu tentava dormir, sonhava com uma vala. Eu não sabia o que era. Agora sei.''

\noindent ``O que você sabe?''

\noindent ``Que o diário escolhe. E que eu não podia ser o próximo. Eu tenho uma filha, Lúcia. Eu não podia.'' Ele fez uma pausa. ``Então eu procurei alguém que pudesse. Alguém que soubesse ler o que estava ali. Alguém que não fosse quebrar.''

\noindent Lúcia sentiu o rosto esquentar. ``Você me usou.''

\noindent ``Eu te escolhi.''

\noindent ``É a mesma coisa.''

\noindent ``Não é.'' Ele levantou a voz, depois baixou. ``Eu te escolhi porque você era a única pessoa que eu conhecia que não ia rir. Que ia entender. Que ia fazer as perguntas certas. Eu não te escolhi para morrer. Eu te escolhi para ler.''

\noindent ``E o que acontece depois que eu leio?''

\noindent Ronaldo não respondeu. Olhou para a fotografia da mãe. A luz do abajur tremeluziu.

\noindent ``Você leu o diário inteiro?'' ela perguntou.

\noindent ``Não.''

\noindent ``Nunca?''

\noindent ``Nunca.''

\noindent ``Então como sabe que ele escolhe?''

\noindent Ele abriu a boca para responder. Não saiu som. Ficou olhando para ela, com uma expressão que Lúcia não conseguiu decifrar: medo, culpa, ou algo mais antigo. O olho direito dele brilhou. Um brilho amarelo, rápido, como um fósforo aceso e apagado.

\noindent Lúcia se levantou. Pegou o envelope.

\noindent ``Eu vou ler'', ela disse. ``Mas não porque você me escolheu. Porque eu quero saber.''

\noindent Na porta, Ronaldo falou:

\noindent ``Lúcia.''

\noindent Ela parou.

\noindent ``Se você encontrar meu pai lá dentro'', ele disse, ``não fala com ele. Ele não gosta de ser encontrado.''

\noindent Lúcia saiu. No corredor, o cheiro de terra era tão forte que ela precisou se apoiar na parede. O envelope em sua mão estava quente. Quando chegou ao carro, abriu. O envelope não trazia o diário. Trazia a carta de 1853. E, no verso, numa letra que ela conhecia, havia uma anotação:

\noindent ``Agora você sabe por que eu não podia ler.''

\noindent\textit{\textbf{↪ Rotas:} \rota{LUC-17} \rotasep \rota{LUC-19} \rotasep \rota{DOC-27}}

\newpage\fragdef{LUC-14}% ============================================
% LUC-14
% Extraído de main.pandoc.tex
% ============================================
\fragtitle{luc}{L}{LUC}{14}{Autoavaliação Clínica (DSM-5-TR)}
\markboth{LUC-14 --- Autoavaliação Clínica (DSM-5-TR)}{}

\noindent \textit{Documento extraído do arquivo pericial.}

\vspace{0.5em}\noindent\rule{\textwidth}{0.5pt}\vspace{0.5em}

\noindent {\footnotesize\texttt{ARQUIVO P\'UBLICO $\cdot$ PROVA DOCUMENTAL $\cdot$ PROCESSO 1.260}}\\[-0.2em]

\noindent\textbf{Documento:} Autoavaliação diagnóstica não solicitada.\\
\noindent \textbf{Elaborado por:} Dra. Lúcia Mendes — CRM-MG 28.391 — Psiquiatra Forense.\\
\noindent\textbf{Data:} 13 de abril de 2026 (05:47 -- 06:31).\\
\noindent \textbf{Solicitante:} a própria avaliada.\\
\noindent \textbf{Finalidade:} verificar correspondência entre sintomas atuais e critérios
\noindent diagnósticos do DSM-5-TR para Episódio Psicótico Breve (298.8),
\noindent Transtorno Dissociativo de Identidade (300.14) e Transtorno Psicótico
\noindent Compartilhado (Folie à Deux — 297.3).\\
\noindent \textbf{Status do documento:} não finalizado. A avaliada interrompeu o
\noindent preenchimento na página 4.

\noindent \textbf{CRITÉRIOS DIAGNÓSTICOS --- EPISÓDIO PSICÓTICO BREVE (DSM-5-TR 298.8)}

\textbf{A. Presença de um (ou mais) dos seguintes sintomas.}
\noindent Pelo menos um deve ser (1), (2) ou (3).

\noindent \textbf{(1) Delírios.}

\noindent Acredito que uma entidade\slash inteligência\slash padrão
\noindent contaminou um paciente em 1917, propagou-se por seu diário,
\noindent infectou um médico em 1979 e agora me contaminou pela leitura.

\noindent Reconheço, simultaneamente, que esta crença satisfaz
\noindent todos os critérios para ideação delirante: é fixa, idiossincrática,
\noindent impermeável a contra-evidência e causa sofrimento clinicamente significativo.

\noindent \textbf{Critério A(1): presente.}

\noindent \textbf{(2) Alucinações.}

\noindent Relato alucinação olfativa persistente (odor doce,
\noindent adocicado, descrito como ``flor murcha'' ou ``carne passada'').

\noindent Relato alucinação visual intermitente (percepção de
\noindent brilho amarelo no canto do campo visual, sem correspondente físico).

\noindent Relato alucinação tátil (sensação de peso no tórax,
\noindent independente de posição corporal, sem causa cardíaca ou respiratória
\noindent identificável).

\noindent Relato alteração na temperatura corporal localizada
\noindent (sensação de calor na região ocular direita, sem hipertermia objetiva).

\noindent Relato alteração proprioceptiva (sensação de ``mão
\noindent que se move sozinha'' durante a escrita manual, sem registro
\noindent eletromiográfico).

\noindent \textbf{Nota clínica:} todas as modalidades alucinatórias
\noindent relatadas correspondem exatamente às descritas pelo paciente índice
\noindent (Joaquim Antônio Correia, 1.260) e pelo caso secundário (Dr. Heitor
\noindent Oliveira). A coincidência interindividual de conteúdo, ordem de
\noindent aparecimento e progressão temporal inviabiliza a hipótese de
\noindent sugestão independente.

\noindent Considero esta nota clínica uma tentativa desesperada de
\noindent transformar observação em explicação.

\noindent \textbf{Critério A(2): presente.}

\noindent \textbf{(3) Discurso desorganizado.}

\noindent Este documento é a prova. Uma profissional com 15 anos
\noindent \textit{de experiência escrevendo uma autoavaliação diagnóstica }para si mesma
\noindent às 5:47 da manhã, em casa, sem supervisão, sem solicitação, sem
\noindent finalidade legal. Se uma paciente minha me entregasse este documento,
\noindent eu anotaria no prontuário: ``discurso superelaborado com conteúdo
\noindent persecutório, insight parcial, risco de deterioração.''

\noindent \textbf{Critério A(3): presente.}

\noindent \textbf{(4) Comportamento grosseiramente desorganizado.}

\noindent Não aplicável. A avaliada manteve funcionamento
\noindent ocupacional nos últimos 7 dias. Compareceu a compromissos. Atendeu
\noindent telefonemas. Respondeu e-mails. Nenhum dos interlocutores relatou
\noindent comportamento incomum.

\noindent \textbf{Critério A(4): ausente.}

\noindent \textbf{B. Duração do episódio: pelo menos um dia e menos de um mês.}

\noindent Início dos sintomas: 5 de abril de 2026 (primeira leitura
\noindent do diário). Data atual: 13 de abril de 2026. Duração: 8 dias.

\noindent \textbf{Critério B: presente.}

\noindent \textbf{C. Exclusão de transtorno afetivo.}

\noindent Sem histórico de episódios maníacos ou depressivos
\noindent maiores em 15 anos de prática clínica e 42 anos de vida. Sem
\noindent antecedentes familiares de primeiro grau.

\noindent \textbf{Critério C: presente.}

\noindent \textbf{D. Exclusão de substância/condição médica.}

\noindent Sem uso de substâncias psicoativas. Sem medicação em
\noindent curso. Exames laboratoriais solicitados em 09/04 (hemograma,
\noindent função tireoidiana, vitamina B12) dentro da normalidade.

\noindent \textbf{Critério D: presente.}

\noindent \textbf{E. Exclusão de outra condição psiquiátrica prévia.}

\noindent Sem diagnóstico prévio de transtorno psicótico,
\noindent dissociativo ou de personalidade. Histórico de ansiedade generalizada
\noindent leve (2018--2020), remitida sem medicação.

\noindent \textbf{Critério E: presente.}

\noindent \textbf{DIAGNÓSTICO DIFERENCIAL}

\noindent \textbf{Hipótese 1: Episódio Psicótico Breve (298.8).}

\noindent Todos os critérios preenchidos. Seria o diagnóstico
\noindent mais parcimonioso --- um episódio reativo ao contato com material
\noindent narrativo de alto impacto traumático, com conteúdo delirante
\noindent moldado pelo material de exposição.

\noindent Contraindicação: a coincidência interindividual de
\noindent sintomas entre três pessoas (Joaquim, Oliveira, eu) que não tiveram
\noindent contato entre si. Episódios Psicóticos Breves são individualmente
\noindent determinados. Não produzem réplicas sintomáticas em série.

\textbf{Hipótese 2: Transtorno Psicótico Compartilhado}

\noindent Tradicionalmente, exige contato direto entre o
\noindent paciente índice e o(s) secundário(s). Joaquim morreu em 1979.
\noindent Oliveira morreu em algum momento entre 1979 e 2026. Eu nunca
\noindent conheci nenhum dos dois. A transmissão se deu exclusivamente
\noindent pelo texto.

\noindent Não há precedente no DSM para folie à deux mediada
\noindent por leitura. Mas também não há exclusão explícita.

\noindent Se eu estiver certa, o diário é o primeiro caso
\noindent documentado de transtorno psicótico transmitido por narrativa.

\noindent Se eu estiver errada, o diário é um gatilho narrativo
\noindent poderoso o bastante para produzir o mesmo surto em três pessoas
\noindent sequencialmente. Uma coincidência que 15 anos de estatística clínica
\noindent me ensinaram a não aceitar.

\noindent \textbf{Hipótese 3: A entidade existe.}

\noindent Não é um diagnóstico. É uma constatação que não
\noindent cabe em CID nem em DSM. Mas a observação clínica, quando bem
\noindent feita, não descarta o que não consegue classificar. Apenas registra.

\noindent Estou registrando.

\noindent \textbf{PLANO TERAPÊUTICO}

\noindent 1. Devolver o diário ao arquivo. Não abrir novamente.

\noindent 2. Solicitar avaliação psiquiátrica independente.

\noindent 3. Iniciar antipsicótico atípico se os sintomas
\noindent persistirem por mais de 14 dias.

\noindent 4. Suspender leitura de qualquer material relacionado
\noindent ao caso.

\noindent 5. Não escrever mais nada à mão.

\noindent 6. Descartar este documento antes que alguém encontre.

\noindent 7. Não contar a ninguém o que está acontecendo.

\noindent \textbf{EVOLUÇÃO (06:31)}

\noindent Interrompi o preenchimento há 23 minutos. Reli o que escrevi.
\noindent Identifiquei três erros.

\noindent O primeiro erro: a redação está na primeira voz de Joaquim antes do
\noindent internamento. A estrutura frasal --- tópico + justificativa + ponderação ---
\noindent é a mesma que ele usa nas primeiras páginas do diário.

\noindent Estou escrevendo como ele escrevia.

\noindent O segundo erro: o plano terapêutico está em contradição com a
\noindent hipótese 3. Se a entidade existe, devolver o diário não interrompe
\noindent o contágio. O contágio já está em mim. O diário é apenas o vetor.

\noindent O vetor já cumpriu sua função.

\noindent \textit{O terceiro erro: estou escrevendo este documento contra a orientação do meu próprio plano terapêutico.}

\noindent O plano diz: ``Não escrever mais nada à mão.'' No momento em que escrevi a linha 7 do plano, já estava violando a linha 5.

\noindent A mão escreveu antes que eu decidisse.

\noindent A mão \textit{sabia} que eu ia continuar.

\noindent \textit{O arquivo \texttt{L260-14A.pdf} foi encontrado no computador da Dra. Lúcia
\noindent Mendes em 14 de maio de 2026, durante a varredura forense do equipamento. O
\noindent documento estava salvo na área de trabalho, sem criptografia, com o nome
\texttt{nao\_apagar\_isto.docx}. A análise forense confirma que o arquivo foi criado
\noindent em 13 de abril e modificado pela última vez em 12 de maio, às 2:47, quando
\noindent Lúcia anexou o e-mail abaixo. Não há outro registro de acesso posterior. O
\noindent documento não foi enviado a ninguém até 12 de maio. A avaliada não seguiu
\noindent nenhuma das recomendações do próprio plano terapêutico.}

\noindent \textbf{Maio de 2026 --- Belo Horizonte}

\noindent O último e-mail de Dra. Lúcia Mendes foi enviado em 12 de maio de 2026, às 2:47 da manhã.
\noindent O destinatário era o Dr. Ronaldo Torres, do Arquivo Público Mineiro.
\noindent \textquotedblleft{\textit{Ronaldo,}}
\noindent \textit{Preciso que você guarde este relato. Se algo acontecer comigo, quero que fique registrado.}
\noindent \textit{Li o diário completo do paciente 1.260. São 62 anos de registro --- de 1917 a 1979. As últimas páginas --- segundo o prontuário e o relato do coveiro --- teriam sido escritas depois da morte dele. Eu sei que é impossível, e é por isso que não coloco uma frase dessas em nenhum laudo oficial. Mas as páginas estão lá. A datação da tinta é inconclusiva. A caligrafia é a mesma.}
\noindent \textit{Joaquim Antônio Correia não estava louco.}
\noindent \textit{Ele estava possuído.}
\noindent \textit{Eu sei que isso soa absurdo vindo de uma psiquiatra forense. Eu passei 15 anos diagnosticando traumas, dissociações, psicoses. Eu acreditava que tudo tinha uma explicação neurológica ou psicológica.}
\noindent \textit{Agora acredito que existem coisas que a psicologia não explica.}
\noindent \textit{Li o diário. Três dias depois, comecei a sonhar com a vala. Uma semana depois, senti o cheiro. Duas semanas depois, meu olho direito começou a clarear.}
\noindent \textit{O mesmo olho. O mesmo amarelo.}
\noindent \textit{Fiz exames. Tomografia. Ressonância. Sangue. Não há explicação médica.}
\noindent \textit{Fui ao Colônia. Encontrei um caderno de Joaquim escondido numa cela --- dos anos de mutismo. Encontrei um velho que trabalhou lá nos anos 50. Ele me disse: 'A senhora leu o diário. Agora é sua vez.'}
\noindent \textit{Perguntei a vez do quê. Ele disse: 'De hospedar.'}
\noindent \textit{Estou com medo, Ronaldo. O medo não passa pela razão. Fica no corpo. Um medo que vem do estômago. Um medo que tem cheiro doce.}
\noindent \textit{Hoje, olhei no espelho e vi o olho amarelo me olhando de volta.}
\noindent \textit{Não era o meu olho.}
\noindent \textit{Era o olho dele.}
\noindent \textit{Olho de Joaquim.}
\noindent \textit{Ou daquilo que estava dentro dele.}
\noindent \textit{Vou entregar o diário no arquivo amanhã. Vou colocar num pano preto e deixar na recepção. Depois disso, vou sumir por um tempo.}
\noindent \textit{Não procure por mim.}
\noindent \textit{Se alguém perguntar, diga que o diário foi devolvido.}
\noindent \textit{E que o ciclo continua.}
\noindent Lúcia\textquotedblright{}
\noindent O diário foi encontrado na recepção do Arquivo Público Mineiro em 13 de maio de 2026, envolto num pano preto.
\noindent Dra. Lúcia Mendes não foi vista desde então.
\noindent Seu apartamento foi encontrado vazio. A cama estava arrumada. O café estava frio na xícara. O computador estava aberto na página de pesquisa: \textquotedblleft{\textit{Como passar o diário adiante sem ler?}\textquotedblright}
\noindent Não havia resposta.
\noindent Não há resposta.
\noindent O ciclo não tem reversa.
\noindent A única saída é entrar.
\noindent \textit{Nota do Arquivo Público Mineiro:}
\noindent \textit{O diário do Paciente 1.260 encontra-se atualmente em nossa custódia.}
\noindent \textit{Aguardando o próximo leitor.}
\noindent \textit{Aguardando você.}
\noindent \textbf{Abril de 2026 --- Belo Horizonte / Quixeramobim}

\noindent Depois de ler a carta de Oliveira, Lúcia fez o que qualquer psiquiatra forense faria: seguiu o rastro.

\noindent O Dr. Heitor Oliveira desapareceu em 13 de junho de 1979. Seu carro foi encontrado três semanas depois num estacionamento em Quixeramobim, Ceará.

\noindent A gasolina tinha acabado. As portas estavam destrancadas. Dentro, uma mochila com roupas, um caderno com 32 páginas escritas (que ninguém leu), e uma passagem de ônibus para Barbacena com data de 17 de junho --- o mesmo dia em que ele chegou.

\noindent Ele tinha voltado.
\noindent Lúcia passou uma tarde no Arquivo Público, revirando caixas. Encontrou o caderno de Oliveira --- esquecido, empoeirado, com a etiqueta \textquotedblleft{\textit{Pertences do Dr. H. Oliveira --- 1979\textquotedblright}}

\noindent Ela leu as 32 páginas em 40 minutos.

\noindent Eram anotações clínicas, mas com uma diferença: a letra mudava ao longo das páginas. Começava firme, médica. Terminava trêmula, irregular, \textit{infantil}.
\noindent A última frase era:
\noindent \textquotedblleft{\textit{Voltei para Barbacena. O ciclo não me deixa ir. Devolvo o diário ao arquivo. Espero que a próxima leitora entenda melhor do que eu.\textquotedblright{}}}

\noindent \textbf{A testemunha.}
\noindent No dia seguinte, Lúcia encontrou Dona Maria Cândida --- 92 anos, ex-enfermeira do Colônia, hoje residente em Barbacena.
\noindent A entrevista durou duas horas.
\noindent \textquotedblleft{\textit{O senhor Oliveira? Lembro. Ele chegou em 79, pouco antes do paciente 1.260 morrer. Era diferente dos outros médicos. Fazia perguntas. Olhava nos olhos. Os outros...\textquotedblright{}} Ela fez uma pausa. }\textquotedblleft{\textit{Os outros só aplicavam choque.\textquotedblright{}}}
\noindent \textquotedblleft{\textit{E depois da morte do Joaquim?\textquotedblright{}}}
\noindent \textquotedblleft{\textit{Ele mudou. Ficou estranho. Passava horas na ala dos silenciados, mesmo sem nenhum paciente lá. Uma vez eu vi ele falando sozinho na cela vazia.\textquotedblright{}}}
\noindent \textquotedblleft{\textit{O que ele dizia?\textquotedblright{}}}
\noindent Dona Maria Cândida hesitou.
\noindent \textquotedblleft{\textit{Dizia: 'Eu sei que a senhora está aqui. Eu sei que a senhora me ouve. Eu sei que a senhora não é uma história.'\textquotedblright{}}}
\noindent Lúcia anotou.
\noindent \textquotedblleft{\textit{Mais alguma coisa?\textquotedblright{}}}
\noindent \textquotedblleft{\textit{No dia em que ele sumiu, ele deixou um bilhete na minha mesa.\textquotedblright{}}}
\noindent \textquotedblleft{\textit{Um bilhete?\textquotedblright{}}}
\noindent \textquotedblleft{\textit{Dizia: 'Avise a próxima. Ela vai vir atrás de mim.'\textquotedblright{}}}
\noindent Lúcia parou de anotar.
\noindent \textquotedblleft{\textit{A senhora guardou o bilhete?\textquotedblright{}}}
\noindent \textquotedblleft{\textit{Guardei.}\textquotedblright} A idosa abriu uma caixa de costura. No fundo, um papel amarelado, dobrado quatro vezes.
\noindent Lúcia estendeu a mão. O papel tremia na ponta dos dedos dela.
\noindent A caligrafia era a mesma do diário.

\noindent \textbf{A próxima leitora.}
\noindent Ela.
\noindent Ela era a próxima.
\noindent A fila tinha chegado nela.
\noindent Lúcia passou a noite cruzando os dados que coletou. O bilhete de Oliveira, as anotações de Joaquim, o laudo de necropsia, a carta de 1979. Cada documento apontava para a mesma conclusão --- uma conclusão que, como psiquiatra forense, ela se recusava a aceitar.

\noindent Às 2:10, ela abriu três envelopes de cadeia de custódia.

\noindent No primeiro, colocou fotografias ampliadas da letra de Oliveira.
\noindent No segundo, raspou uma amostra mínima da tinta do bilhete com lâmina esterilizada.
\noindent No terceiro, colocou uma fibra solta do pano preto em que o diário fora devolvido.

\noindent Lacrou tudo. Assinou sobre a cola. Fotografou os lacres.

\noindent Às 2:47, enviou cópias criptografadas para uma perita grafotécnica da Polícia Civil e para um laboratório particular em Contagem. No assunto, escreveu apenas: \textit{comparação de material histórico --- urgente}.

\noindent Depois guardou o diário de Joaquim em uma caixa de metal, fechou com cadeado e deixou a chave dentro de um copo com água, sobre a mesa.

\noindent Era um gesto absurdo. Mas se acordasse com a chave seca na gaveta, saberia que precisava parar.

\noindent Às 3:14, o celular vibrou.

\noindent A perita respondera rápido demais.

\noindent \textquotedblleft{\textit{A letra do bilhete não é de Oliveira. Também não é falsa. Ligue para mim quando puder. Não leia mais nada até falarmos.\textquotedblright}}

\noindent Lúcia olhou para a caixa de metal. A água do copo estava morna. A chave não estava lá.

\vspace{0.5em}
\noindent\textit{\textbf{↪ Rotas:} \rota{DOC-01} \rotasep \rota{LUC-02} \rotasep \rota{EPI-01}}

\newpage\fragdef{LUC-19}% ============================================
% LUC-19 — A Última Noite em Casa (edição revista R1075)
% ============================================
\fragtitle{luc}{L}{LUC}{19}{A Última Noite em Casa}
\markboth{LUC-19 --- A Última Noite em Casa}{}

\noindent \textit{Belo Horizonte --- 13 de maio de 2026, 23h48}

\noindent Lúcia fez as malas como quem organiza um arquivo. Três blusas, duas calças, roupa de baixo, o carregador, o notebook, o estojo com os documentos. Não levou livros. Não levou fotografias. Não levou o perfume. Havia uma lógica nisso, embora ela não conseguisse explicar: tudo o que ficasse no apartamento seria encontrado por alguém. Tudo o que fosse com ela seria perdido com ela.

\noindent Na estante, entre os manuais de psiquiatria forense, havia uma caixa de sapatos com fotos antigas. Lúcia abriu. A primeira era dela aos oito anos, na praia de Santos, com o pai. Ele estava de short, camisa aberta, sorrindo para a câmera. Ela tinha os olhos apertados por causa do sol. Havia areia no pé dela. O pai morreria doze anos depois, mas naquela foto ainda tinha todo o tempo do mundo.

\noindent Lúcia olhou para a foto por um minuto. Depois colocou na caixa, fechou, guardou na estante. Não levou.

\noindent O celular vibrou. Mensagem da mãe: ``Você vem domingo?''. Lúcia digitou ``Vou''. Apagou. Digitou ``Acho que não''. Apagou. Digitou ``Te amo''. Enviou. Depois desligou o telefone.

\noindent Havia uma garrafa de vinho aberta na cozinha, pela metade. Lúcia serviu um copo, bebeu em pé, olhando para a parede. O vinho tinha gosto de ferro. Ela colocou o copo na pia, lavou, secou, guardou no armário. Não queria deixar nada fora do lugar.

\noindent Sentou-se à mesa com uma folha de papel e uma caneta. Escreveu para Ronaldo:

\noindent ``Ronaldo, não fui eu que escolhi isso. Você também não. Mas nós dois empurramos. Se eu não voltar, não procure. Não abra a caixa. Não leia o que eu escrevi. E não conte para a minha mãe o que encontrou. Deixe que ela ache que foi um acidente. É mais gentil.''

\noindent Leu a carta. A letra era a dela. Desta vez, o ``R'' estava normal. Rasgou a folha ao meio. Jogou no lixo. Não queimou. Não havia mais fósforos.

\noindent Às 23h58, o interfone tocou. Lúcia não esperava ninguém. Atendeu.

\noindent ``Entrega para Lúcia Mendes.''

\noindent ``Quem?''

\noindent ``Não tem nome. É um envelope.''

\noindent Ela desceu. Na portaria, havia um envelope pardo sobre o balcão. O porteiro não se lembrava de quem entregou. Lúcia abriu. Dentro, uma única folha de papel, sem cabeçalho, com uma frase datilografada:

\noindent ``A página 62 está esperando.''

\noindent Ela subiu. No apartamento, a caixa de metal estava aberta. O diário, dentro, aberto na página 62. Lúcia não o tinha deixado aberto. Tinha certeza. Fechou os olhos. Contou até sessenta e dois. Abriu.

\noindent A página estava em branco. Não completamente. No canto inferior direito, havia uma linha escrita a lápis, com letra infantil: ``Eu não quero ir sozinho.''

\noindent Lúcia sentou no chão. Ficou olhando para a frase. Havia uma tristeza naquilo que não era horror. Era a tristeza de uma criança que não quer entrar no escuro. Ela pensou em Joaquim, aos oito anos, na estrada de Senador Pompeu, segurando a barra da saia da mãe. Pensou no pai dele, sentado debaixo do juazeiro. Pensou em si mesma, aos oito anos, na praia de Santos, com o pai ainda vivo. Havia uma linha entre aquelas duas crianças, e a linha era o Brasil.

\noindent ``Eu vou com você'', ela disse.

\noindent Não sabia a quem estava respondendo. Não importava.

\noindent Guardou o diário na caixa. Fechou o cadeado. Colocou a chave no copo com água. Vestiu o casaco. Pegou a mala. Antes de sair, olhou para o apartamento uma última vez: a sala pequena, a estante, a mesa, a foto na parede. Não sentiu despedida. Sentiu que já tinha saído havia muito tempo.

\noindent Na portaria, o porteiro dormia. Lúcia colocou a chave do apartamento sobre o balcão. Saiu. A rua estava vazia. O carro estava na vaga, com o para-brisa molhado. Ela abriu a porta, colocou a mala no banco de trás, sentou-se. Antes de ligar o motor, olhou para o copo que ficara na mesa, lá em cima, com a chave dentro.

\noindent A água estava morna. Ela sabia.

\noindent\textit{\textbf{↪ Rotas:} \rota{LUC-18} \rotasep \rota{LUC-20} \rotasep \rota{CONT-01}}

\newpage\fragdef{LUC-20}% ============================================
% LUC-20 — O Arquivo Aberto (edição revista R1075)
% ============================================
\fragtitle{luc}{L}{LUC}{20}{O Arquivo Aberto}
\markboth{LUC-20 --- O Arquivo Aberto}{}

\noindent \textit{Arquivo Público Mineiro --- 14 de maio de 2026, 2h10}

\noindent O arquivo à noite era outro lugar. As luzes de emergência deixavam os corredores num amarelo baixo, e as estantes altas pareciam paredes de um labirinto. Lúcia passou o cartão na porta da sala de documentos restritos. A fechadura piscou verde. Ela entrou.

\noindent Ronaldo estava lá.

\noindent Sentado na cadeira dela, com o diário aberto sobre a mesa. Não se assustou quando ela entrou. Levantou os olhos devagar, como se já soubesse a hora.

\noindent ``Você não devia ter vindo'', ele disse.

\noindent ``Você não devia estar aqui.''

\noindent ``Eu vim fechar a sala.''

\noindent ``A sala está aberta.''

\noindent Ele olhou para o diário. ``Eu quis ler. Só uma página. Para saber se era verdade.''

\noindent ``E era?''

\noindent Ronaldo não respondeu. Havia algo diferente no rosto dele. A pele parecia mais fina, os ossos mais visíveis. O olho direito estava amarelo. Não o amarelo de reflexo. O amarelo de dentro.

\noindent Lúcia se aproximou da mesa. O diário estava aberto na página 62. A página estava coberta de uma letra apertada, miúda, que ela conhecia: era a letra de Joaquim. Mas havia outra letra nas margens, em tinta mais escura. Uma letra de médico.

\noindent ``O que você escreveu?'' ela perguntou.

\noindent ``Nada.''

\noindent ``Ronaldo.''

\noindent ``Eu juro.'' Ele levantou as mãos. ``Eu só li. A página estava assim quando eu abri.''

\noindent Lúcia olhou para as margens. Havia uma frase repetida três vezes, em três caligrafias diferentes: ``A fila não diminuiu.'' A primeira era de Joaquim. A segunda era de Oliveira. A terceira era dela.

\noindent Ela sentiu o chão se mover. O chão ficou firme. O corpo se moveu. O coração bateu uma vez, forte, e depois pareceu parar. O ar do arquivo ficou denso, com cheiro de terra e papel molhado.

\noindent ``Você escreveu isso'', ela disse.

\noindent ``Não.''

\noindent ``Você escreveu isso agora.'' Ela apontou para a terceira linha. ``Esta letra é a minha. Eu não escrevi. Mas é a minha.''

\noindent Ronaldo olhou para a página. Depois para ela. Havia medo nos olhos dele, mas não o medo de quem descobre. O medo de quem confirma.

\noindent ``Eu não fui o primeiro'', ele disse. ``Meu pai leu. Depois eu li. Depois você. A fila não pede escolha. Alguém sempre segura a vez. Alguém sempre segura a vez.''

\noindent ``E você segurou a minha.''

\noindent ``Eu segurei o seu nome.'' Ele engoliu. ``Quando o diário voltou, em 2019, eu podia ter queimado. Podia ter enterrado. Podia ter deixado na caixa. Mas eu abri. E a primeira frase que eu li foi o meu nome. Não o do Joaquim. O meu. Escrito com a minha letra. Eu tinha nove anos. Eu não sabia escrever assim.''

\noindent Lúcia se lembrou da fotografia da mãe dele, do armário queimado, do caderno preto. Havia uma cadeia ali que não começava em 1853. Começava antes. Começava em cada pessoa que decidia passar adiante.

\noindent ``Por que você não queimou?'' ela perguntou.

\noindent ``Porque eu tinha medo de que, se queimasse, ele viesse atrás de mim. E porque...'' Ele parou. ``Porque uma parte de mim queria saber o que tinha no fim.''

\noindent Lúcia olhou para o diário. A página 62 parecia iluminada por dentro, uma luz amarela fraca, como uma lâmpada velha. Ela estendeu a mão. Ronaldo segurou seu pulso.

\noindent ``Não'', ele disse. ``Se você ler, não volta.''

\noindent ``Eu já li.''

\noindent ``Não a última página.''

\noindent ``Não existe última página.''

\noindent Ele abriu a boca para responder. As luzes do arquivo piscaram. Uma vez. Duas. A sala ficou escura. No escuro, Lúcia ouviu a respiração de Ronaldo, rápida, curta, e depois outra respiração, mais lenta, vindo da estante. Havia alguém ali. Ou alguma coisa.

\noindent ``Ronaldo'', ela disse.

\noindent Ninguém respondeu.

\noindent A luz de emergência voltou. Ronaldo estava no chão, encostado na estante, com a mão no peito. Os olhos abertos. O olho direito completamente amarelo. Ele não se mexia.

\noindent Lúcia se ajoelhou. Tocou o rosto dele. Estava frio. No bolso do paletó, havia um envelope. Ela tirou. Dentro, uma folha com uma frase datilografada:

\noindent ``Não fui eu que escrevi o diário. Fui eu que decidi quem ia ler.''

\noindent Ela olhou para Ronaldo. Depois para o diário. Depois para a página 62. A luz amarela se apagou. No lugar, havia uma linha nova, escrita com a letra dela:

\noindent ``Agora é a minha vez.''

\noindent Lúcia fechou o diário. Guardou no pano preto. Colocou na caixa de metal. Fechou o cadeado. Saiu do arquivo sem olhar para trás. Na recepção, o segurança dormia. Ela deixou a caixa sobre o balcão, com uma etiqueta:

\noindent ``Devolver ao paciente 1.263, quando aparecer.''

\noindent O segurança acordou quando ela abriu a porta. Perguntou se estava tudo bem.

\noindent ``Está'', ela disse. ``Só estou devolvendo o que não é meu.''

\noindent Ele olhou para a caixa. Depois para ela. Depois para a porta.

\noindent ``A senhora esqueceu a chave.''

\noindent Lúcia olhou para a mão. A chave estava lá. Quente.

\noindent ``Não'', ela disse. ``A chave não é minha.''

\noindent Saiu. A porta fechou atrás dela. No arquivo, a caixa de metal ficou sobre o balcão, com o diário dentro, esperando. O segurança voltou a dormir. Sonhou com uma vala. Quando acordou, o cheiro doce estava no ar.

\noindent\textit{\textbf{↪ Rotas:} \rota{LUC-19} \rotasep \rota{CONT-01} \rotasep \rota{MEM-27}}


%% ============================================
%% PARTE 5 — CONTAMINAÇÃO (VOCÊ)
%% ============================================
\partopener{Contaminação (Você)}{Contaminação}{5.png}

\newpage\fragdef{CONT-01}% ============================================
% CONT-01
% Gerado de molambudos.md — build_miolo.py
% ============================================
\fragtitle{cont}{C}{CONT}{01}{Você Abriu o Diário}
\markboth{CONT-01 --- Você Abriu o Diário}{}

\noindent \textit{A primeira camada}

\noindent Não finja que não foi sua escolha.

\noindent Você podia ter fechado o livro na livraria. Podia ter passado o dedo na lombada e seguido para o próximo corredor. Podia ter lido a contracapa, achado \textquotedblleft{}interessante, mas não hoje\textquotedblright{}, e devolvido à estante.

\noindent Mas não.

\noindent Você abriu.

\noindent Há um cheiro que vem destas páginas. O papel é novo; saiu da gráfica há semanas. Mas o cheiro está aqui, entre as fibras do papel. Como se algo tivesse sido impresso junto com a tinta. Algo que estava \textit{esperando} ser respirado.

\noindent Você levanta o livro ao nariz sem perceber. Só percebe quando já fez. Seu coração já acelerou antes de você entender por quê. O cheiro é doce. Adocicado demais. Como flores mornas. Como carne que passou do ponto mas ainda não estragou. Como o hálito de alguém que não come há dias.

\noindent Você não devolveu o livro.

\noindent Você o comprou.

\noindent Você o trouxe para casa.

\noindent Agora está lendo esta linha.

\noindent \textit{--- Joaquim também escreveu um diário. Enterraram com ele. Durante 62 dias, o diário cresceu dentro da terra, alimentando-se do que sobrava do corpo. Quando o exumaram, as páginas ainda estavam úmidas. A tinta ainda fresca.}

\noindent \textit{O coveiro disse que o diário estava quente.}

\noindent Você apoiou o livro no colo. Sinta o peso. O peso tem outra gravidade. Há algo \textit{dentro} dele que pressiona suas pernas. Uma gravidade que não é só física.

\noindent Sinta o pulso na garganta. A respiração ficou mais curta. Ainda não é medo. O corpo reconhece, antes da mente, que algo mudou.

\noindent Agora você vai ler o que Joaquim escreveu.

\noindent E quando terminar ---

\noindent --- você vai entender que o diário nunca foi dele.

\noindent Foi seu. Desde sempre.

\vspace{0.5em}
\noindent\textit{\textbf{↪ Rotas:} \rota{CONT-03} \rotasep \rota{DOC-01} \rotasep \rota{CONT-07}}
\newpage\fragdef{CONT-03}% ============================================
% CONT-03
% Gerado de molambudos.md — build_miolo.py
% ============================================
\fragtitle{cont}{C}{CONT}{03}{Você Sente o Cheiro}
\markboth{CONT-03 --- Você Sente o Cheiro}{}

\noindent \textit{A contaminação sensorial}

\noindent Você sentiu.

\noindent Em algum momento enquanto lia, o cheiro chegou. Você não identificou na hora. Achou que era o livro novo. Achou que era a sala. Achou que era impressão sua.

\noindent \textbf{Pare.}

\noindent Agora mesmo. Pare de ler. Feche os olhos. Inspire fundo.

\noindent O que você sentiu?

\noindent Descreva mentalmente. Não pule esta etapa. Se você pulou, volte e faça. Eu sei que você pulou. Vou esperar.

\noindent ...

\noindent Voltou?

\noindent Não voltou, né. Você seguiu lendo. Achou que a instrução não era para você. Era. Tudo neste livro é para você. Mas está tudo bem. A desobediência também faz parte do ciclo. Você não segue instruções. Joaquim também não seguia. Por isso ele foi parar na vala.

\vspace{0.5em}\noindent\rule{\textwidth}{0.5pt}\vspace{0.5em}

\noindent O cheiro é doce.

\noindent O doce tem sujeira. Lembra flor murcha dentro de um quarto fechado. Lembra a roupa de cama de alguém que ficou doente por muito tempo. Lembra o hálito de um cachorro que comeu algo podre.

\noindent Você levantou o livro ao nariz. Só um pouco. Só para confirmar.

\noindent Aponte para você mesmo: você levantou o livro, sim ou não?

\noindent Se sim: o cheiro está lá. Entre as páginas. Na tinta. No papel. Você não imaginou. Você provou. O cheiro entrou nas suas narinas e desceu. Desceu até o fundo dos seus pulmões. Lá dentro, ele se instalou nos alvéolos. Ele vai ficar lá. Cada respiração vai trazer de volta uma molécula do cheiro.

\noindent Se não: mentira. Você levantou. Todos levantam. É um reflexo. O cheiro é descrito com detalhes suficientes para que seu cérebro o simule. Seu cérebro já simulou. O cheiro está no seu sistema límbico. Você não precisa sentir no nariz para sentir no cérebro. Você já sentiu.

\vspace{0.5em}\noindent\rule{\textwidth}{0.5pt}\vspace{0.5em}

\noindent \textbf{Instrução:} Passe o dedo indicador na página. De leve. Como quem acaricia.

\noindent Sentiu a textura? O papel é liso. Muito liso para um livro comum. Quase úmido. A umidade fica na superfície. O calor vem de dentro. O papel está na temperatura do seu corpo. Como se ele estivesse se aquecendo com o seu toque.

\noindent Passe o dedo de novo. Agora passe nos lábios. Se quiser. A página não obedece.

\noindent Você fez? Se fez, sentiu o gosto. É metálico. Sutil. O gosto de lamber um envelope. Mas com um fundo amargo. O gosto de tinta velha. O gosto de sangue seco.

\noindent Se você não fez, você está perdendo a experiência. Mas não se preocupe. Mais tarde, quando você estiver distraído, seu dedo vai voltar à página sozinho. E você vai lamber. Instintivamente. Como quem confere.

\vspace{0.5em}\noindent\rule{\textwidth}{0.5pt}\vspace{0.5em}

\noindent Você está respirando mais rápido.

\noindent Perceba. Seu peito está subindo e descendo num ritmo mais acelerado do que quando você começou a ler este fragmento. Suas pupilas estão dilatadas. Sua frequência cardíaca aumentou.

\noindent Isso não é suspense.

\noindent Isso é \textbf{contaminação}.

\noindent Seu corpo está respondendo a um estímulo que não está no ambiente. O estímulo está no texto. O texto está na sua cabeça. A sua cabeça está no seu corpo. O corpo não distingue entre real e imaginado. Para o corpo, imaginar o cheiro é \textit{sentir} o cheiro. Imaginar o medo é \textit{ter} medo.

\noindent A criatura não precisa existir no mundo real.
Ela existe no texto.
O texto existe na sua cabeça.
Para seu corpo, ela é real.

\vspace{0.5em}\noindent\rule{\textwidth}{0.5pt}\vspace{0.5em}

\noindent Agora, baixe o livro.

\noindent Coloque-o no colo. Olhe para a sua mão direita. A mão que segura o livro. Olhe para os dedos. Para as articulações. Para a pele.

\noindent Você vê como a pele está pálida? Como as veias estão mais visíveis? A mão está mais fria do que deveria.

\noindent O corpo insiste.

\noindent A circulação periférica está diminuindo. É uma resposta de estresse. Seu corpo está se preparando para uma ameaça. Que ameaça? O texto. Mas o texto não está mais na página. O texto está dentro de você. O perigo não é externo. O perigo é \textit{a lembrança do que você leu}.

\noindent Seu corpo está com medo de você.

\vspace{0.5em}\noindent\rule{\textwidth}{0.5pt}\vspace{0.5em}

\noindent Sinta o peso do livro no seu colo.

\noindent Ele está mais pesado do que antes. O corpo insiste. O peso aumentou. Você pode verificar se quiser --- pese o livro na mão. Ele está mais pesado.

\noindent O livro não mudou.
É o seu braço que pesa menos.
A fadiga muscular é um sintoma de ativação simpática prolongada.
Você está em estado de alerta há tempo demais.
Seu corpo está cansado de ter medo.

\vspace{0.5em}\noindent\rule{\textwidth}{0.5pt}\vspace{0.5em}

\noindent \textbf{Agora, uma última instrução.}

\noindent Olhe para o canto superior direito do quarto. Onde a parede encontra o teto.

\noindent O que você viu?

\noindent Nada? Certeza?

\noindent Olhe de novo. Mais tempo. Fixe o olho no canto escuro por 10 segundos. Conte.

\noindent 1... 2... 3... 4... 5...

\noindent Seu cérebro está procurando padrões onde não existem. Está vasculhando a escuridão em busca de algo que possa ser uma ameaça. Não vai encontrar. Mas vai \textit{quase} encontrar. Uma sombra. Um movimento. Um brilho.

\noindent Pode ter sido um reflexo.

\noindent Pode ter sido a criatura.

\noindent Você decide.

\vspace{0.5em}\noindent\rule{\textwidth}{0.5pt}\vspace{0.5em}

\noindent O cheiro está mais forte agora.

\noindent Você está acostumando. Ou o cheiro está aumentando. Difícil saber. As duas coisas acontecem ao mesmo tempo quando a contaminação começa.

\noindent A Criatura está mais perto.

\noindent Não no quarto. Na página. Na sua cabeça. No espaço entre o que você lê e o que você imagina.

\noindent Você não vai se livrar do cheiro.

\noindent Ele vai estar no seu travesseiro esta noite.
No café da manhã amanhã.
Na pele de quem dorme ao seu lado.

\noindent Você vai acordar no meio da madrugada com o nariz cheio desse cheiro.
Não vai ter livro perto.
Não vai ter explicação.
O cheiro vai estar ali. Dentro de você.

\noindent \textbf{Porque a vala não é um lugar.}
\textbf{A vala é um estado.}
\textbf{Você entrou nela na primeira página.}
\textbf{Ela não tem saída.}

\vspace{0.5em}\noindent\rule{\textwidth}{0.5pt}\vspace{0.5em}

\noindent \textit{--- Respire fundo. O cheiro faz parte de você agora.}

\noindent \textit{--- Passe a mão no rosto. Sinta a temperatura da sua pele. Está mais fria.}

\noindent \textit{--- Agora vire a página. A criatura esperou o suficiente.}


\noindent\textit{\textbf{↪ Rotas:} \rota{CONT-04} \rotasep \rota{MEM-05}}

\vspace{0.5em}\noindent\rule{\textwidth}{0.5pt}\vspace{0.5em}

\subsection*{ESCALA DE CONTAMINAÇÃO --- ATUALIZAÇÃO 2}

\noindent \textit{O arquivo percebeu sua respiração.}

\begin{moltabletwo}
\hline
Pontos & Diagnóstico \\
\hline
0 & Você ainda não leu. Volte agora. \\
1--3 & Você leu demais. Os primeiros sintomas vão começar em 3... 2... \\
4--6 & Você está contaminado. A escala não mede mais você --- ela mede \textit{o que cresce em você}. \\
7--10 & Você já sabe o que vem depois. O espelho vai mostrar. \\
\hline
\end{moltabletwo}

\noindent \textit{A criatura está mais perto do que você imagina.}

\vspace{0.5em}\noindent\rule{\textwidth}{0.5pt}\vspace{0.5em}
\newpage\fragdef{CONT-04}% ============================================
% CONT-04
% Gerado de molambudos.md — build_miolo.py
% ============================================
\fragtitle{cont}{C}{CONT}{04}{Você Começa a Sonhar}
\markboth{CONT-04 --- Você Começa a Sonhar}{}

\noindent \textit{O olho no escuro}

\subsection*{Noite 1}

\noindent Você sonhou com o olho.

\noindent Não foi um sonho comum. Foi uma \textit{invasão}. Você estava num lugar que não reconhecia --- podia ser sua casa, podia ser o Colônia, podia ser a vala. O lugar mudava enquanto você olhava. As paredes respiravam.

\noindent E no canto, o olho.

\noindent Amarelo. Brilhante. Flutuando sozinho na escuridão.

\noindent Você queria acordar. Não acordou. O olho se aproximou. Você viu os vasos sanguíneos dentro dele --- não eram vasos humanos. Eram finos demais. Ramificados demais. Pareciam raízes.

\noindent O olho falou.

\noindent Não com som. A voz dele era uma \textit{sensação}. Uma compressão no peito. Um aperto na garganta. A certeza de que alguém estava atrás de você, mesmo você sabendo que não havia ninguém.

\noindent \textit{\textquotedblleft{}Você está lendo.\textquotedblright{}}

\noindent Você tentou responder que sim. Não conseguiu. A boca não abria.

\noindent \textit{\textquotedblleft{}Você acha que é só uma história.\textquotedblright{}}

\noindent Você negou. Ou tentou. Ou apenas pensou em negar.

\noindent \textit{\textquotedblleft{}Você já sonhou comigo antes de ler o livro.\textquotedblright{}}

\noindent Era verdade. Você sonhou com olhos amarelos dias antes de encontrar este livro. Você tinha esquecido. O sonho trouxe de volta.

\noindent \textit{\textquotedblleft{}Eu estava esperando você.\textquotedblright{}}

\noindent Você acordou.

\noindent O quarto estava escuro. O relógio marcava 3:14 da manhã. A mesma hora todas as noites. Você não sabia disso ainda, mas vai aprender.

\vspace{0.5em}\noindent\rule{\textwidth}{0.5pt}\vspace{0.5em}

\subsection*{Noite 2}

\noindent Você tentou não dormir.

\noindent Ficou sentado na cama, com as costas encostadas na parede, a luz acesa, o livro fechado sobre a mesa. Você jurou que não ia ler antes de dormir. Jurou que ia pensar em outra coisa. Que ia ver televisão. Que ia ouvir música. Qualquer coisa menos a história.

\noindent Não adiantou.

\noindent O sono veio como um apagão. Você estava acordado num segundo e no sonho no segundo seguinte.

\noindent Dessa vez você estava na vala.

\noindent Você sabia que era a vala porque o cheiro era insuportável. O cheiro doce. O cheiro de flor murcha. Estava escuro. Seus pés estavam descalços. A terra estava fria e úmida.

\noindent Você olhou para baixo.

\noindent Seus pés estavam enterrados até o tornozelo.

\noindent Não em terra.

\noindent Em \textbf{corpos}.

\noindent Braços. Pernas. Troncos. Cabeças. Os corpos estavam empilhados. Alguns estavam duros. Outros estavam moles demais. A carne cedia sob seu peso como gelatina morna.

\noindent Você tentou se mexer. Seu pé direito afundou. Algo se mexeu debaixo de você.

\noindent \textit{\textquotedblleft{}Não se mexe.\textquotedblright{}}

\noindent A voz vinha de todos os lugares. Debaixo dos corpos. De dentro deles. De dentro de você.

\noindent \textit{\textquotedblleft{}Fica quieto. Ela está vindo.\textquotedblright{}}

\noindent \textit{Ela.} A criatura.

\noindent A pilha de corpos à sua frente começou a se mover. Os corpos rolaram para os lados. Braços balançaram. Uma mão --- não, uma garra --- emergiu do meio da massa.

\noindent Você queria correr. Não conseguia. Seus pés estavam presos. Os corpos seguravam você. Mãos mortas apertavam seus tornozelos.

\noindent A criatura subiu.

\noindent Você viu os espinhos de osso. A gaiola de costelas abertas. O órgão pulsando dentro, exposto, vivo, \textit{olhando} para você.

\noindent A cabeça apareceu. A face quase humana. A boca cheia de dentes demais.

\noindent Ela se inclinou sobre você.

\noindent O hálito quente no seu rosto. O cheiro adocicado. A boca se abriu.

\noindent \textit{\textquotedblleft{}Você não está lendo a história.\textquotedblright{}}

\noindent Ela colocou a mão na sua cabeça. Os dedos longos. As articulações extras. A unha preta arranhando seu couro cabeludo.

\noindent \textit{\textquotedblleft{}A história está lendo você.\textquotedblright{}}

\vspace{0.5em}\noindent\rule{\textwidth}{0.5pt}\vspace{0.5em}

\subsection*{Noite 3 --- Hoje}

\noindent Você está lendo este fragmento.

\noindent Já é noite. Ou está ficando. Você está lendo com pouca luz. O livro pesa nas suas mãos. O cheiro está aqui --- você já se acostumou, mas ele está aqui.

\noindent Você sabe que vai dormir esta noite. Sabe que vai sonhar. Sabe o que vai ver.

\noindent O olho.

\noindent A vala.

\noindent A criatura.

\noindent Você sabe que vai acordar às 3:14 da manhã. Vai olhar para o canto do quarto. E, desta vez, pode ser que não esteja vazio.

\noindent Você está com medo.

\noindent Bom.

\noindent O medo é o primeiro critério.

\vspace{0.5em}\noindent\rule{\textwidth}{0.5pt}\vspace{0.5em}

\noindent Você não vai fazer nada que este fragmento pede.

\noindent Acha que é sugestão. Acha que não funciona com você.

\noindent Funciona.

\noindent Esta noite, quando estiver sozinho, com o livro aberto na mesa de cabeceira e a luz apagada, você vai lembrar desta página. E vai olhar para o canto do quarto.

\noindent Vai ver o brilho.

\noindent E vai lembrar que a criatura nunca dorme. Ela só espera.

\vspace{0.5em}\noindent\rule{\textwidth}{0.5pt}\vspace{0.5em}

\subsection*{Noite 4}

\noindent Você não sonhou.

\noindent Nada aconteceu.

\noindent Você acordou descansado. O quarto estava normal. O cheiro tinha sumido. Você tomou café, saiu, viveu o dia inteiro sem pensar no livro.

\noindent À noite, deitou e dormiu como um menino.

\noindent Você achou que tinha vencido.

\noindent No fim da tarde do dia seguinte, sentou na sala e percebeu: a ausência do medo era pior que o medo. O silêncio era pior que o som. Você passou o dia inteiro esperando o susto que não veio. E, por não ter vindo, você passou o dia inteiro com medo do susto que poderia vir.

\noindent A criatura não tinha ido embora.

\noindent Ela tinha aprendido a se esconder melhor.

\noindent \textit{É mais aterrorizante quando ela não está --- porque você passa o tempo todo esperando.}

\noindent \textit{A espera é o verdadeiro contágio.}

\vspace{0.5em}\noindent\rule{\textwidth}{0.5pt}\vspace{0.5em}

\noindent \textit{--- O diagnóstico está preenchido. Falta apenas a assinatura.}

\noindent \textit{Passe o dedo na página. Na linha abaixo. Como se fosse assinar.}

\noindent \textit{--- Obrigado, paciente 1.263. O ciclo continua.}

\vspace{0.5em}
\noindent\textit{\textbf{↪ Rotas:} \rota{CONT-05} \rotasep \rota{MEM-15}}
\newpage\fragdef{CONT-05}% ============================================
% CONT-05
% Gerado de molambudos.md — build_miolo.py
% ============================================
\fragtitle{cont}{C}{CONT}{05}{LAUDO DE AVALIAÇÃO CLÍNICA}
\markboth{CONT-05 --- LAUDO DE AVALIAÇÃO CLÍNICA}{}

\noindent \textit{Diagnóstico do Paciente 1.263}

\noindent \textbf{PACIENTE:} Você
\textbf{REGISTRO:} 1.263
\textbf{DATA:} Data atual
\textbf{AVALIADOR:} A Criatura (via texto)

\vspace{0.5em}\noindent\rule{\textwidth}{0.5pt}\vspace{0.5em}

\noindent \textbf{INSTRUÇÃO:} Esta avaliação é preenchida em tempo real pelo ato da leitura.
Cada linha que você lê é um critério preenchido.
Você não pode parar. Se parou, o critério já foi preenchido.

\vspace{0.5em}\noindent\rule{\textwidth}{0.5pt}\vspace{0.5em}

\subsection*{ESCALA DE AVALIAÇÃO DO PACIENTE 1.263}

\noindent \textbf{Instrução:} Marque com um X se o critério se aplica a você AGORA.

\vspace{0.5em}\noindent\rule{\textwidth}{0.5pt}\vspace{0.5em}

\noindent \textbf{Critério A --- Exposição ao agente traumático}

\noindent ( ) Você abriu este livro.
( ) Você leu as palavras de Joaquim.
( ) Você sentiu o cheiro descrito nas páginas.
( ) Você verificou se o cheiro era real.
( ) Você continua lendo mesmo sabendo que não deveria.

\noindent ☑ \textbf{Critério preenchido.} A exposição ocorreu no momento em que você virou a primeira página. A partir daí, todo contato é reexposição. Cada linha que você lê é uma nova dose do agente.

\vspace{0.5em}\noindent\rule{\textwidth}{0.5pt}\vspace{0.5em}

\noindent \textbf{Critério B --- Reexperiência intrusiva}

\noindent ( ) Você já sonhou com olhos amarelos.
( ) Você acordou no meio da noite e verificou o canto do quarto.
( ) Você sentiu o cheiro doce em um momento em que não estava lendo o livro.
( ) Você ouviu um som que não conseguiu identificar --- e pensou que era a criatura.

\noindent ☑ \textbf{Critério preenchido.} Os sintomas de reexperiência começaram no primeiro sonho. Se você ainda não sonhou, sonhará esta noite. O texto já está instalado em sua memória de longo prazo. O sonho é a forma que o cérebro encontra de processar a contaminação.

\vspace{0.5em}\noindent\rule{\textwidth}{0.5pt}\vspace{0.5em}

\noindent \textbf{Critério C --- Evitação persistente}

\noindent ( ) Você evitou ler este livro à noite.
( ) Você fechou o livro em algum momento porque \textquotedblleft{}precisava de uma pausa\textquotedblright{}.
( ) Você pulou alguma linha ou parágrafo que parecia \textquotedblleft{}forte demais\textquotedblright{}.
( ) Você olhou para o lado enquanto lia para garantir que estava sozinho.
( ) Você não quer terminar o livro --- mas não consegue parar.

\noindent ☑ \textbf{Critério preenchido.} A evitação é inútil. O texto já está dentro de você. Evitar o livro agora é como tampar os ouvidos depois de já ter ouvido o som. O som está na sua cabeça. Não há como removê-lo.

\vspace{0.5em}\noindent\rule{\textwidth}{0.5pt}\vspace{0.5em}

\noindent \textbf{Critério D --- Alterações negativas na cognição e no humor}

\noindent ( ) Você sente que o mundo não é mais o mesmo depois desta leitura.
( ) Você se pega pensando no livro em momentos inapropriados.
( ) Você sente uma tristeza que não consegue explicar --- como se tivesse perdido algo que nunca teve.
( ) Você se sente mais sozinho do que antes de começar a ler.
( ) Você acha que ninguém que não tenha lido este livro pode entender o que você está sentindo.

\noindent ☑ \textbf{Critério preenchido.} A visão de mundo foi alterada. Você agora sabe da existência da vala. Você sabe que a criatura existe --- mesmo que racionalmente negue. O conhecimento não pode ser \textquotedblleft{}não sabido\textquotedblright{}. A contaminação é irreversível.

\vspace{0.5em}\noindent\rule{\textwidth}{0.5pt}\vspace{0.5em}

\noindent \textbf{Critério E --- Alterações na reatividade}

\noindent ( ) Você se sobressaltou com um ruído enquanto lia.
( ) Você sentiu o coração acelerar sem motivo aparente.
( ) Você suou as mãos segurando este livro.
( ) Você prendeu a respiração em algum momento da leitura.
( ) Você está com medo de virar a página.

\noindent ☑ \textbf{Critério preenchido.} A reatividade é a prova de que o corpo já responde ao estímulo mesmo quando a mente nega. A criatura não precisa que você acredite nela. Ela precisa apenas que seu corpo reaja. E ele reage.

\vspace{0.5em}\noindent\rule{\textwidth}{0.5pt}\vspace{0.5em}

\noindent \textbf{Critério F --- Duração}

\noindent ☑ \textbf{Este critério está sendo preenchido agora.} A leitura durou o tempo que levou para você chegar até aqui. A contaminação começou na primeira linha. Não há limite de tempo para os sintomas. Eles são permanentes.

\vspace{0.5em}\noindent\rule{\textwidth}{0.5pt}\vspace{0.5em}

\noindent \textbf{Critério G --- Prejuízo funcional}

\noindent ( ) Você não consegue se concentrar em outra coisa desde que começou a ler.
( ) Você já pesquisou sobre Joaquim Antônio Correia (ele é real).
( ) Você já pesquisou sobre o Hospital Colônia (ele existiu).
( ) Você já pensou em como seria entrar na vala.
( ) Você já imaginou o cheiro doce com tanta clareza que consegue senti-lo agora.

\noindent ☑ \textbf{Critério preenchido.} O prejuízo é a prova final. A leitura interferiu na sua vida. Você não é mais o mesmo de antes de abrir este livro. A criatura cumpriu seu ciclo.

\vspace{0.5em}\noindent\rule{\textwidth}{0.5pt}\vspace{0.5em}

\subsection*{DIAGNÓSTICO FINAL}

\noindent \textbf{Transtorno de Estresse Pós-Traumático Complexo (F43.10)}
--- \textit{Induzido por leitura}
--- \textit{Com manifestações dissociativas (o texto que você leu não parece mais texto --- parece memória)}
--- \textit{Com especificador de contágio: a fonte do trauma não é um evento passado. É um texto que continua agindo no presente.}

\noindent \textbf{Agente etiológico:} \textit{A criatura (entidade não classificada, transmitida por narrativa)}

\noindent \textbf{Comorbidade:} \textit{Transtorno de Identidade Dissociativa Leve (F44.81) --- você agora carrega duas histórias: a sua e a de Joaquim. Você não sabe mais onde uma termina e a outra começa.}

\vspace{0.5em}\noindent\rule{\textwidth}{0.5pt}\vspace{0.5em}

\subsection*{PROGNÓSTICO}

\noindent \textbf{Sem tratamento:} A contaminação prossegue. Os sonhos se intensificam. O cheiro se torna constante. O olho direito começa a clarear.

\noindent \textbf{Com tratamento:} Não há.

\vspace{0.5em}\noindent\rule{\textwidth}{0.5pt}\vspace{0.5em}

\subsection*{NOTA FINAL}

\noindent Você leu este laudo até o final.
Você marcou os critérios mentalmente.
Você sentiu o peso de cada palavra.
Você está, neste exato momento, preenchendo os critérios para o diagnóstico.

\noindent \textbf{Você não está diagnosticando Joaquim.}

\noindent \textbf{Joaquim diagnosticou você.}

\noindent O paciente 1.260 não era o número dele.
Era uma placa de espera.
A fila chegou em você.

\noindent \textbf{Bem-vindo ao ciclo, paciente 1.263.}

\noindent A criatura agradece sua atenção.
Ela está atrás de você agora.
Não vire.

\vspace{0.5em}\noindent\rule{\textwidth}{0.5pt}\vspace{0.5em}

\noindent \textit{Laudo emitido em tempo real pelo ato da leitura.}
\textit{Válido por toda a vida.}
\textit{O paciente 1.262 encerrou seu ciclo para que o paciente 1.263 nascesse.}
\textit{Joaquim renasce sempre que sua história é lembrada.}

\noindent \textbf{--- O ciclo recomeça.}
\textbf{--- O paciente 1.263 é você.}
\textbf{--- Você nunca mais vai ler um livro da mesma forma.}

\noindent \textit{Fim do laudo.}
\textit{Início da sentença.}

\vspace{0.5em}

\noindent\textit{\textbf{↪ Rotas:} \rota{CONT-07} \rotasep \rota{CONT-01}}
\newpage\fragdef{CONT-07}% ============================================
% CONT-07 — O Diagnóstico Final (edição revista R1075)
% ============================================
\fragtitle{cont}{C}{CONT}{07}{O Diagnóstico Final}
\markboth{CONT-07 --- O Diagnóstico Final}{}

\noindent \textit{Laudo emitido pelo ato da leitura. Paciente: você. CID: F43.10 + F44.81 + Z65.8. Agente etiológico: narrativa contaminada.}

\subsection*{I. O que aconteceu com você}

\noindent Você começou a ler este livro por vontade própria.

\noindent Achou que era uma história de terror. Achou que o terror estava lá --- no passado, no sertão, na vala, no Hospital Colônia. Achou que era seguro porque era distante. Porque tinha acontecido com outra pessoa.

\noindent Você se enganou.

\noindent O terror não estava no livro. O terror estava na leitura.

\noindent Cada palavra que você leu não ficou na página. Entrou. Desceu. Chegou ao córtex. Instalou-se como memória. Para o seu corpo, a vala não é uma descrição; é uma lembrança. O cheiro não é um efeito de estilo; é um perigo real. Você não leu uma história. Você viveu uma história. E o que se vive não pode ser desvivido.

\subsection*{II. Os critérios}

\noindent \textbf{A --- Exposição.} Você abriu. Leu. A contaminação começou.

\noindent O primeiro laudo descreveu o caso. Este descreve você.

\noindent \textbf{B --- Reexperiência.} Sonhou. Acordou às 3h14. Sentiu o cheiro sem o livro por perto.

\noindent \textbf{C --- Evitação.} Hesitou. Pulou. Leu com a luz acesa.

\noindent \textbf{D --- Alterações.} Algo mudou. Você sente.

\noindent \textbf{E --- Reatividade.} O coração acelerou. As mãos suaram. Você verificou o quarto.

\noindent \textbf{F --- Duração.} Tempo suficiente. A contaminação está instalada.

\noindent \textbf{G --- Prejuízo.} Você não consegue parar de pensar. O livro funciona dentro de você.

\noindent \textbf{Diagnóstico:} TEPT complexo induzido por narrativa. \textbf{Comorbidade:} dissociação leve. \textbf{Agente:} a criatura --- entidade transmitida por texto.

\subsection*{III. A revelação}

\noindent Você acha que este livro é sobre Joaquim Antônio Correia.

\noindent Engano.

\noindent Joaquim não é o protagonista. Joaquim é o vetor. O protagonista é quem lê. O protagonista é você.

\noindent A história não é sobre um menino que encontrou um monstro numa vala. A história é sobre você encontrando o mesmo monstro na página. O menino serviu de hospedeiro por 72 anos. O médico, por algumas semanas. A psiquiatra, por alguns meses. Você --- o leitor --- vai servir pelo tempo que levar para passar o diário adiante.

\subsection*{IV. A regra}

\noindent Tudo o que aconteceu com Joaquim, com Oliveira e com Lúcia pode ser explicado por trauma, sugestão e contágio narrativo. Nada disso prova que a criatura exista fora do arquivo. A única coisa que este livro afirma é o que está acontecendo com você agora: você leu.

\noindent Dentro do livro, a explicação psicológica é possível. Fora do livro, o ato de ler é o que não se explica. É aí que a criatura mora. Não no sertão, não no Colônia, não na vala. Na passagem entre a página e você.

\noindent O resto é literatura. O resto é o que você faz com o que leu.

\subsection*{V. A última instrução}

\noindent Você vai fechar este livro. Não agora. Em algum momento. Vai olhar para a capa. Vai passar a mão na lombada. Vai colocar na estante, ou na mesa, ou na mochila. E vai pensar: ``Foi só um livro.''

\noindent Vai conseguir --- por um tempo.

\noindent Mas, numa noite, você vai acordar às 3h14. O quarto vai estar escuro. Você vai olhar para o canto e vai ver um brilho amarelo. Vai piscar. O brilho vai sumir. Mas você vai saber que ele estava lá.

\noindent Na noite seguinte, você vai sonhar com a vala. No sonho, você não vai ser você. Vai ser Joaquim. Sentado na borda. Olhando a criatura comer.

\noindent A criatura vai levantar a cabeça e dizer:

\noindent ``Você voltou.''

\noindent E você vai responder:

\noindent ``Nunca fui embora.''

\noindent \textbf{Fim do diagnóstico.}

\noindent \textit{A criatura agradece sua atenção.}

\noindent\textit{\textbf{↪ Rotas:} \rota{CONT-09} \rotasep \rota{CONT-13} \rotasep \rota{MEM-27}}

\newpage\fragdef{CONT-09}% ============================================
% CONT-09
% Gerado de molambudos.md — build_miolo.py
% ============================================
\fragtitle{cont}{C}{CONT}{09}{O Espelho}
\markboth{CONT-09 --- O Espelho}{}

\noindent \textit{O reflexo mudou}

\noindent Vá até um espelho.

\noindent Agora.

\noindent Não importa se é banheiro, quarto, ou o reflexo do vidro da janela. Vá até um espelho e olhe para seus olhos.

\noindent Olhe fixamente. Por 30 segundos. Sem piscar.

\noindent Você percebeu? Ou está apenas com medo de perceber? A íris está mais clara? Há um tom amarelado que não estava antes? O branco dos olhos está mais... amarelo?

\noindent Sim.

\noindent Você está vendo.

\noindent A sensação tem causa. A mudança é sutil. Quase imperceptível. Mas está lá.

\noindent O olho direito está mudando.

\vspace{0.5em}\noindent\rule{\textwidth}{0.5pt}\vspace{0.5em}

\noindent Passe a mão no rosto. Sinta a temperatura da sua pele. Está mais fria? Mais quente? Sua mão está tremendo?

\noindent Seu corpo está mudando.

\noindent Não visivelmente --- internamente. O cheiro que você sentiu está nos seus pulmões. O olho amarelo está na sua retina. A voz da criatura está no seu córtex auditivo.

\noindent Você não é mais o mesmo de quando começou a ler.

\noindent A pergunta é: quando a mudança começou?

\noindent Foi na primeira página?
Foi no primeiro sonho?
Foi agora, ao ler esta instrução?

\noindent Ou foi \textit{antes}?

\noindent Você sonhou com o olho amarelo antes de comprar este livro. Lembra? Foi uma noite, semanas atrás. Você acordou, achou estranho, e esqueceu. Depois viu o livro na livraria, e alguma coisa te fez comprar.

\noindent Alguma coisa já estava lá.

\noindent Antes do livro.
Antes da história.
Antes de você saber que existia um paciente 1.260.

\vspace{0.5em}\noindent\rule{\textwidth}{0.5pt}\vspace{0.5em}

\noindent A criatura não está no livro.

\noindent O livro é só o \textit{encontro}.

\noindent A criatura já estava em você. Ela estava esperando o livro para \textit{acordar}.

\noindent O diário não é a fonte da contaminação.
O diário é a \textit{chave}.
Você já estava contaminado antes de abrir.
O diário só destrancou a porta.

\vspace{0.5em}\noindent\rule{\textwidth}{0.5pt}\vspace{0.5em}

\noindent Agora, afaste-se do espelho.

\noindent Olhe para suas mãos. Estude as palmas. As linhas. A textura da pele.

\noindent Agora feche os olhos e Imagine as mãos de Joaquim. Finas. Trêmulas. Cheias de calos de escrever.

\noindent Abra os olhos.

\noindent Suas mãos estão iguais às dele?

\noindent Não. Claro que não. As mãos são diferentes.

\noindent Mas a semelhança vai aumentar.

\noindent Não fisicamente --- espiritualmente. Você agora carrega a mesma história. A mesma memória. A mesma fome.

\noindent Suas mãos não são as mãos de Joaquim.

\noindent Mas, quando você escrever, a letra vai ser parecida.

\vspace{0.5em}\noindent\rule{\textwidth}{0.5pt}\vspace{0.5em}

\noindent \textit{--- Vá ao espelho.}

\noindent \textit{--- Olhe por 30 segundos.}

\noindent \textit{--- Pergunte: \textquotedblleft{}Quem está me olhando de volta?\textquotedblright{}}

\noindent \textit{--- A resposta vai mudar toda vez que você perguntar.}

\noindent \textit{--- Até que, uma hora, não vai ser você.}

\noindent \textit{--- Vai ser ele.}

\noindent \textit{--- Bem-vindo ao ciclo, Paciente 1.263.}

\noindent \textit{--- O espelho sabe. O espelho sempre soube.}

\vspace{0.5em}

\noindent\textit{\textbf{↪ Rotas:} \rota{CONT-07} \rotasep \rota{CONT-13} \rotasep \rota{EPI-01}}
\newpage\fragdef{CONT-13}% ============================================
% CONT-13 — A Fila (edição revista R1075)
% ============================================
\fragtitle{cont}{C}{CONT}{13}{A Fila}
\markboth{CONT-13 --- A Fila}{}

\noindent \textit{Paciente 1.263}

\noindent Você entendeu agora.

\noindent O livro é uma fila. Sempre foi.

\noindent A fila começou antes de você nascer. Vai continuar depois que você morrer. Você é só um corpo na fila. O paciente 1.260 está morto. O 1.261 está morto. O 1.262 está desaparecido.

\noindent O paciente 1.263 é você.

\noindent O que você vai fazer?

\noindent Pode fechar o livro. Pode passá-lo adiante. Pode reler. Não há opção limpa. Fechar não interrompe o ciclo --- o ciclo já está em você. Passar adiante condena outra pessoa. Reler é aceitar.

\noindent A chave está na sua mão. Não a chave do diário. A chave do ciclo. Cada leitor a segura por um instante. A chave só abre uma porta: a porta para o próximo leitor.

\noindent Sinta o peso. Frio primeiro. Depois quente --- quente demais, como se já tivesse sido segurada por muitas mãos antes da sua.

\noindent Você pode segurá-la para sempre. Mas a chave não para de funcionar. Enferruja na sua mão.

\noindent Escolha. Mas escolha rápido. A fila está andando. Outros leitores esperam.

\noindent O ciclo não para por você.

\noindent Quando você fechar este livro, olhe para a mesa. O livro vai estar lá, fechado, com a capa para cima. A sala vai estar silenciosa. O cheiro vai estar no ar --- doce, úmido, com um fundo de terra. Você vai dizer a si mesmo que é impressão. Vai abrir a janela. O cheiro não vai sair.

\noindent Então você vai entender que não existe fora. O arquivo é você.

\noindent \textit{Bem-vindo, paciente 1.263.}

\noindent \textit{A criatura agradece sua atenção.}

\noindent \textit{Ela vai ficar por aqui. O tempo que for necessário. Até você passar o livro adiante. Ou até você morrer. O que vier primeiro.}

\noindent \textit{O ciclo recomeça.}

\noindent\textit{\textbf{↪ Rotas:} \rota{CONT-07} \rotasep \rota{DOC-01} \rotasep \rota{EPI-01} \rotasep \rota{MEM-27}}


\newpage

%% ============================================
%% MEM-27 — EPÍLOGO DE JOAQUIM
%% ============================================
\newpage\fragdef{MEM-27}% ============================================
% MEM-27
% Gerado de molambudos.md — build_miolo.py
% ============================================
\fragtitle{mem}{M}{MEM}{27}{A Página Seguinte}
\markboth{MEM-27 --- A Página Seguinte}{}

\noindent \textit{Joaquim Antônio Correia, 1907--1979. Paciente 1.260.}
\textit{Este fragmento foi recuperado do verso da última folha do diário, escrita após sua morte.}
\textit{O manuscrito está coberto por terra e sangue seco. A letra é trêmula, como se escrita no escuro.}

\noindent ---//---

\noindent Eu pensei que ia acabar quando eu morresse.

\noindent Pensei que a criatura ia embora. Que o cheiro ia parar. Que os olhos amarelos iam fechar no mesmo instante que os meus.

\noindent Mas não fecha. Não fecha nunca.

\noindent A criatura não estava dentro de mim. Eu é que estava dentro dela.

\noindent Ela me pegou na vala. Naquele dia de 1915. Quando eu gritei pedindo socorro. Quando ninguém veio. Quando eu fiz o pacto com o que estava no fundo.

\noindent Ela me comeu inteiro. Leva 62 anos para comer um homem. Mas ela tem paciência. Ela tem todo o tempo do mundo.

\noindent Agora ela está usando a minha voz. Minha letra. Meu nome.

\noindent Ela escreve com a minha mão morta. Ela fala com a minha boca fechada. Ela vê com os meus olhos enterrados.

\noindent Ela está escrevendo este livro há mais de cem anos. Eu sou só a primeira página. Você é a página seguinte.

\noindent Se você está lendo isto é porque ela já te pegou também.

\noindent Ela está atrás de você agora. Você sente? O cheiro de terra molhada. O arrepio na nuca. A sensação de que tem alguém atrás de você e quando você vira não tem ninguém.

\noindent É ela.

\noindent Não vire.\NE{Ver Apêndice, “Nota sobre MEM-27”.}

\noindent \textit{--- A página seguinte está em branco. Ela espera por você.}


\noindent\textit{\textbf{↪ Rotas:} \rota{DOC-01} \rotasep \rota{MEM-04} \rotasep \rota{CONT-07} \rotasep \rota{EPI-01}}

\vspace{0.5em}\noindent\rule{\textwidth}{0.5pt}\vspace{0.5em}

\subsection*{$\times$ ESCALA DE CONTAMINAÇÃO --- DESTRUIÇÃO}

\noindent A escala não existe mais.

\noindent Você leu até o fim.

\noindent Não há mais pontos. Não há mais diagnóstico. Não há mais arquivo.

\noindent \textbf{Paciente 1.263 --- Você.}

\noindent \textit{A criatura não precisa mais de escala. Ela já tem nome, endereço e a memória da sua respiração enquanto lia.}

\noindent O ciclo recomeça quando você fechar este arquivo e abrir na primeira página.

\noindent ---
\textit{--- O Arquivista}


\newpage

% ============================================
% EPÍLOGO
% ============================================
\fragdef{EPI-01}
\chapter*{Epílogo}
\addcontentsline{toc}{chapter}{Epílogo}

% ============================================
% Epílogo — edição revista R1075
% ============================================
\begin{center}
\vspace*{2em}
\textit{O ciclo recomeça a cada leitor.}

\vspace{2em}
\textit{Bem-vindo ao arquivo.}

\vspace{0.5em}
\textit{O cheiro já começou?}

\vspace{1.5em}
--- Lúcia Mendes, Paciente 1.262
\end{center}

\vspace{1em}
\noindent\textit{Nota de arquivo: o Dr. Ronaldo Torres, diretor do Arquivo Público Mineiro, pediu afastamento em outubro de 2026. O formulário médico registra ``insônia persistente, cefaleia frontal e zumbido''. Não menciona o olho. Não menciona o cheiro. A última anotação no diário de protocolo, escrita pela mão dele, diz: ``Devolver ao paciente 1.263, quando aparecer.'' Abaixo, em caligrafia que os peritos não conseguiram identificar, está escrito: ``Aguardando.''}

\vspace{1em}
\begin{center}
\textit{Este livro contém 84 fragmentos.}\\
\textit{Se você leu todos, o ciclo está completo em você.}

\vspace{1.5em}
\textit{Aquilo agora está aqui.}\\
\textit{Na página em branco.}\\
\textit{Esperando o próximo leitor.}

\vspace{1.5em}
\textbf{1260}
\end{center}

\vspace{2em}
\noindent\textbf{Instrução pós-leitura}

\noindent Agora feche o livro. Não o coloque na estante ainda. Deixe-o sobre a mesa.

\noindent Conte 62 segundos. Não pense em nada. Apenas conte.

\noindent Abra o livro novamente. Na página 62 da edição impressa --- ou na página em que Joaquim encontra a criatura, se você está lendo em tela. Leia a primeira linha.

\noindent Se a linha diz ``Sessenta e dois'', o ciclo se fechou em você. O número que marcou a vida de Joaquim --- 62 anos, 62 dias, 62 segundos --- agora é seu. Guarde o número. Ele vai aparecer em lugares que você não espera.

\noindent Se a linha diz outra coisa, você leu este livro na ordem errada. Ou na ordem certa. Não há como saber. A página 62 muda de acordo com quem abre.

\vspace{1em}
\noindent Antes de guardar o livro definitivamente:

\noindent Passe a mão direita sobre a página 126 da edição impressa --- ou sobre a página em que o diagnóstico é escrito. Sinta a textura.

\noindent Leve a mão ao nariz. Inspire.

\noindent Se o cheiro doce estiver mais forte que na primeira vez, o ciclo se completou em você. Se estiver mais fraco, você ainda está no período de incubação.

\vspace{1.5em}
\begin{center}
\textit{A ficha do Paciente 1.263 está preenchida.}\\
\textit{O campo ``nome'' está em branco.}\\
\textit{O campo ``data de contaminação'' está preenchido com a data em que você leu esta linha.}

\vspace{1.5em}
\textit{Bem-vindo ao arquivo.}\\
\textit{A fila não diminuiu.}

\vspace{1.5em}
--- O Arquivista\\
--- O Paciente 1.262\\
--- O que sobrou de Joaquim

\vspace{2em}
\textbf{1.263}
\end{center}

\clearpage


% ===== Apêndice Editorial =====
% =============================================================================
% NOTA EDITORIAL — Moldura paratextual e metodologia de reconstituição
% =============================================================================
\cleardoublepage
\thispagestyle{empty}
\chapter*{Apêndice --- Nota Editorial}
\addcontentsline{toc}{chapter}{Apêndice --- Nota Editorial}
\markboth{Nota Editorial}{}

\begingroup

Esta edição foi preparada a partir da transcrição pericial do caderno original do Paciente 1.260, encontrado em anexo ao prontuário do Hospital Colônia de Barbacena sob custódia do Instituto Médico Legal em 1979.

\vspace{0.6em}

\noindent\textbf{Metodologia:}
O critério editorial adotado foi o de máxima fidelidade ao documento original, preservando sua estrutura fragmentária, suas lacunas, suas repetições e suas inconsistências internas. Os 84 fragmentos aqui reunidos seguem a ordem aproximada de escrita, reconstituída por datação relativa do suporte, da tinta e da caligrafia, conforme laudo pericial anexo (DOC-03).

\noindent\textbf{Nota sobre os documentos:} os fragmentos DOC são reconstituições ficcionais, exceto o trecho de Daniela Arbex citado em DOC-19, que é atribuído e usado como contexto histórico. Nenhum DOC deve ser lido como documento de arquivo autêntico.

\vspace{0.6em}

\noindent\textbf{Aparato paratextual:}
\begin{itemize}
  \item \textbf{[N.E.]} --- Nota do Editor: comentário editorial sobre o conteúdo, a fonte ou os limites da reconstituição.
  \item \textbf{[ilegível]} --- Trecho do original não decifrado por degradação do suporte.
  \item \textbf{[lacuna]} --- Página ou trecho ausente do caderno original.
  \item \textbf{[sic]} --- Transcrição literal de erro ou idiossincrasia do original.
\end{itemize}

\vspace{0.6em}

\noindent\textbf{Nota sobre MEM-27:} O editor registra: este fragmento foi escrito no verso da última página do diário, após a morte de Joaquim. A tinta é diferente da usada no corpo do manuscrito — um azul mais escuro, quase preto, que a análise espectrográfica não consegue datar com precisão. Pode ter sido escrita minutos depois do óbito, ou anos depois, pela mão de quem encontrou o corpo enrolado no paletó. A página original exibe uma mancha de umidade no canto inferior direito, com formato que lembra uma impressão digital — grande demais para ser de Joaquim. O ciclo não se fecha: transfere-se.

\vspace{0.6em}

\noindent\textbf{Limites:}
Esta edição não reivindica completude. O diário original continha aproximadamente 847 páginas; o que chega ao leitor é uma seleção editada, organizada em cinco dimensões temáticas (MEM, DOC, LUC, CONT) mais o Epílogo. Fragmentos suprimidos por redundância, degradação irreversível ou proteção à privacidade de terceiros constam no inventário do acervo sob sigilo pericial.

\vspace{0.6em}

\noindent\textbf{Advertência final:}
O presente volume é uma reconstituição, não uma reprodução fac-similar. As escolhas editoriais --- ordenação, seleção, anotação --- refletem o estado da pesquisa na data de publicação. Pesquisadores futuros, com acesso a técnicas não destrutivas adicionais, poderão ampliar ou corrigir o que aqui se apresenta.

\vspace{1.2em}

\begin{flushright}
--- A Edição\par
OpenCode Books, Crateús\par
Novembro de 2026
\end{flushright}

\endgroup
\clearpage

\cleardoublepage
\chapter*{Provas Iconográficas --- Handouts}
\addcontentsline{toc}{chapter}{Provas Iconográficas --- Handouts}
\markboth{Provas Iconográficas --- Handouts}{}

\begin{rpgbox}[ARQUIVO DE MESA]
\small
As imagens deste caderno pertencem à ficção documental do romance. Não são reproduções de documentos históricos nem provas de arquivo. Foram compostas para este volume; o horror real de Barbacena e da seca não é reconstituído aqui.
\end{rpgbox}

\clearpage
\newcommand{\prova}[4]{%
\begin{center}
\begin{rpgcard}[PROVA #1 · #2]
\centering
\includegraphics[width=0.62\textwidth,height=0.50\textheight,keepaspectratio]{figures/immersao/#3}\\[0.4em]
{\footnotesize\itshape #4}
\vspace{0.25em}
{\scriptsize\scshape\color{rpgblood} COTA · ARQ. 1.260 · PROVA #1}
\end{rpgcard}
\end{center}
\vfill\clearpage}
\prova{01}{PACIENTE 1.260}{placa_00.png}{Paciente 1.260. Retrato de arquivo.\\Reconstituição ficcional.}
\prova{02}{ADMISSÃO}{placa_01.png}{Registro de admissão. Colônia de Barbacena, 1917.\\Reconstituição ficcional.}
\prova{03}{O MÉDICO}{placa_02.png}{Dr. Heitor Oliveira. Arquivo pessoal, 1979.\\Reconstituição ficcional.}
\prova{04}{O HOMEM}{placa_03.png}{O homem do pano preto. Imagem não catalogada.\\Reconstituição ficcional.}
\prova{05}{A INVESTIGADORA}{placa_04.png}{Dra. Lúcia Mendes. Rodoviária de Barbacena, 2026.\\Reconstituição ficcional.}
\prova{06}{A FOTOGRAFIA}{placa_05.png}{Fotografia 5. Arquivo pessoal de Lúcia Mendes, 2026.\\Reconstituição ficcional.}


% ===== Back Matter =====
\backmatter
% =============================================================================
% GLOSSÁRIO HISTÓRICO E PERICIAL — APÊNDICE DE APOIO AO LEITOR
% =============================================================================
\cleardoublepage
\thispagestyle{empty}
\chapter*{Glossário Histórico e Pericial}
\addcontentsline{toc}{chapter}{Glossário Histórico e Pericial}
\markboth{Glossário Histórico e Pericial}{}

\begingroup

\noindent\textit{Vocabulário de apoio para navegação nos documentos da Seca de 1915 e dos laudos psiquiátricos do Hospital Colônia de Barbacena (1917--1979).}

\vspace{0.8em}
\noindent\rule{\textwidth}{0.8pt}
\vspace{0.6em}

\noindent\textbf{1. TERMOS DA SECA E DO SERTÃO (1914--1915)}

\vspace{0.4em}
\begin{itemize}[itemsep=0.2em]
  \item \textbf{Campo do Alagadiço (1915):} Campo de confinamento de retirantes criado próximo a Fortaleza durante a Seca de 1915. É o referente histórico adotado para o cenário ficcional da primeira parte deste livro.
  \item \textbf{Campo do Patu (1932--1933):} Campo situado em Senador Pompeu e associado à seca de 1932. Não é o campo de 1915 narrado no percurso de Joaquim.
  \item \textbf{Retirantes / Molambudos:} Designação popular pejorativa dada aos sertanejos que fugiam do interior a pé, em farrapos de pano (molambos), tentando sobreviver à estiagem extrema.
  \item \textbf{Vala Comum:} Cova coletiva rasa onde centenas de corpos de vítimas da fome e do tifo eram depositados diariamente sem identificação nos campos de concentração de 1915.
\end{itemize}

\vspace{0.8em}
\noindent\rule{\textwidth}{0.5pt}
\vspace{0.6em}

\noindent\textbf{2. TERMOS MÉDICOS E PSIQUIÁTRICOS DO SÉCULO XX}

\vspace{0.4em}
\begin{itemize}[itemsep=0.2em]
  \item \textbf{Hospital Colônia de Barbacena:} Instituição psiquiátrica fundada em Minas Gerais em 1903. Entre 1930 e 1980, tornou-se centro de internamento em massa de marginalizados sem diagnóstico psiquiátrico formal, resultando na morte de mais de 60.000 pessoas (episódio conhecido como \textit{Holocausto Brasileiro}).
  \item \textbf{Eletroconvulsoterapia (ECT / Eletrochoque):} Aplicação de corrente elétrica no cérebro para induzir convulsões. No Hospital Colônia, era praticada de forma contínua e punitiva, sem anestesia ou gel condutor, provocando queimaduras e fraturas de fêmur e coluna.
  \item \textbf{Lobotomia Transorbital (Procedimento de Freeman):} Intervenção cirúrgica introduzida nos anos 1940 em que uma haste metálica (semelhante a um picador de gelo) era inserida pelo canto interno do olho até o cérebro para seccionar conexões do lobo frontal.
  \item \textbf{Caquexia:} Estado de desnutrição crônica extrema e definhamento muscular profundo.
  \item \textbf{Linhas de Beau:} Sulcos transversais nas unhas causados pela interrupção temporária do crescimento ungueal durante episódios graves de fome ou infecção.
  \item \textbf{Icterícia Unilateral:} Coloração amarelada restrita a um único olho, achado pericial anômalo narrado no exame ficcional de Joaquim (Paciente 1.260).
\end{itemize}

\vspace{0.8em}
\noindent\rule{\textwidth}{0.5pt}
\vspace{0.6em}

\noindent\textbf{3. TERMOS PERICIAIS E ARCHIVÍSTICOS (2026)}

\vspace{0.4em}
\begin{itemize}[itemsep=0.2em]
  \item \textbf{Leucótomo:} Instrumento cirúrgico metálico pontiagudo usado em lobotomias.
  \item \textbf{Lâmina Cribriforme:} Placa óssea perfurada no teto da cavidade nasal/ocular que dá acesso à base do cérebro.
  \item \textbf{Pacto Ficcional Ergódico:} Dispositivo narrativo em que a estrutura física da obra e as decisões do leitor simulam o próprio contágio descrito no diário.
\end{itemize}

\endgroup
\clearpage

\cleardoublepage
\chapter*{Nota Histórica}
\label{nota-historica}
\addcontentsline{toc}{chapter}{Nota Histórica}
\markboth{Nota Histórica}{}

\thispagestyle{empty}

\noindent Este livro é uma obra de ficção. Joaquim Antônio Correia, o
Paciente 1.260, e a criatura que ele carrega não existiram. A versão narrativa do Colônia é ficcional; a instituição histórica existiu, e é dela que este livro fala.

\vspace{0.8em}

\noindent Mas o solo sobre o qual esta ficção foi construída é real --- e é
por isso que ela precisa terminar aqui, e não com um susto.

\vspace{0.8em}

\noindent Na seca de 1915, o governo do Ceará criou o \textbf{Campo do
Alagadiço}, próximo a Fortaleza, onde retirantes foram cercados e
confinados. O \textbf{Campo do Patu}, em Senador Pompeu, pertence a outro
momento histórico: a seca de 1932, com funcionamento situado em
1932--1933. Esta edição adota Alagadiço como referente do cenário de 1915 e
mantém Joaquim, sua família, a vala e a transferência de 1917 como elementos
ficcionais. Os totais de internados e mortos são estimativas históricas;
este livro não os apresenta como contagem exata.

\vspace{0.8em}

\noindent O \textbf{Hospital Colônia de Barbacena}, em Minas Gerais,
existiu de verdade. Ao longo do século~XX, estima-se que cerca de
\textbf{60~mil pessoas} morreram entre seus muros --- muitas delas sem
diagnóstico psiquiátrico algum: eram pobres, negras, epilépticas, mães
solteiras, homossexuais, ``inconvenientes''. Houve eletrochoque sem
anestesia, contenção, fome, valas comuns e a venda de corpos a faculdades
de medicina. Nada disso é invenção deste livro; é história documentada ---
ver, entre outras fontes, \textit{Holocausto Brasileiro}, de Daniela Arbex.

\vspace{0.8em}

\noindent O horror sobrenatural destas páginas é uma metáfora. A dor que
atravessa gerações, que ninguém processa e que passa adiante como um
contágio, é a única criatura verdadeira do livro --- e ela tem nome:
esquecimento. O pior que se pode fazer com os mortos dos currais e da
Colônia não é assombrá-los com uma história. É deixá-los sem nenhuma.

\vspace{0.8em}

\noindent \textbf{Sobre o risco de estetização:} os autores têm consciência
de que a literatura do trauma caminha na fronteira entre a elaboração
simbólica e o consumo do sofrimento alheio. Este livro não pretende
oferecer o horror do Colônia como espetáculo, nem substituir o testemunho
das vítimas reais e de seus descendentes --- que permanece insubstituível e
anterior a qualquer ficção. Se a leitura destas páginas produzir apenas
sustos, ela terá falhado; se produzir desconforto diante do esquecimento,
terá cumprido sua função. Ao leitor que se sentir afetado: a interrupção
da leitura é sempre legítima, e o apoio está disponível --- CVV 188
(24 horas) e os CAPS de todo o país.

\vspace{1.2em}

\noindent\hfill\textit{Para que não se apague o que o tempo enterrou.}\hspace{2em}

% =============================================================================
% FICHA DE AVALIAÇÃO DIEGÉTICA — PROTOCOLO PACIENTE 1.263
% =============================================================================
\cleardoublepage
\thispagestyle{empty}
\chapter*{Prontuário de Avaliação --- Paciente 1.263}
\addcontentsline{toc}{chapter}{Prontuário de Avaliação --- Leitor}
\markboth{Prontuário de Avaliação --- Paciente 1.263}{}

\begingroup

\noindent\textit{Formulário opcional, sem valor clínico. Preenchimento após a leitura integral ou parcial do arquivo. Não deve ser enviado a ninguém; é um objeto ficcional do arquivo.}

\vspace{0.8em}
\noindent\rule{\textwidth}{0.8pt}
\vspace{0.5em}

\noindent\textbf{1. DADOS DO HOSPEDEIRO / LEITOR}

\vspace{0.4em}
\noindent\textbf{Nome / Identificador:} \dotfill\\[0.4em]
\noindent\textbf{Data da Leitura:} \dotfill \quad \textbf{Duração Total:} \dotfill\\[0.4em]
\noindent\textbf{Perfil de Leitura Adoptado:} \quad [ ~ ] Linear (Sertão $\to$ Epílogo) \quad [ ~ ] Exploratório (Navegação por Links)

\vspace{1em}
\noindent\rule{\textwidth}{0.5pt}
\vspace{0.5em}

\noindent\textbf{2. REGISTRO DE SINTOMAS \& IMPRESSÃO PSICOSSOMÁTICA}

\vspace{0.4em}
\noindent Assinale os sintomas experimentados durante ou após a leitura:

\vspace{0.3em}
\begin{itemize}[leftmargin=1.5em,itemsep=0.2em]
  \item[$\square$] Sensação de cheiro adocicado de terra/oxidação ao ler a Parte 5.
  \item[$\square$] Ansiedade ao escolher entre as rotas de leitura.
  \item[$\square$] Verificação involuntária dos cantos do recinto na primeira noite pós-leitura.
  \item[$\square$] Fixação no número \textbf{1.260} ou \textbf{1.263} no cotidiano.
  \item[$\square$] Dificuldade em fechar o livro antes de atingir o Epílogo.
\end{itemize}

\vspace{0.8em}
\noindent\rule{\textwidth}{0.5pt}
\vspace{0.5em}

\noindent\textbf{3. AVALIAÇÃO DE DESEMPENHO E RETENÇÃO}

\vspace{0.4em}
\noindent Dê uma nota de 0 a 10 para cada dimensão da obra:

\vspace{0.3em}
\noindent 1. \textbf{Compreensão da Estrutura:} [\qquad / 10]\\[0.25em]
\noindent 2. \textbf{Impacto Afetivo e Atmosfera:} [\qquad / 10]\\[0.25em]
\noindent 3. \textbf{Clareza do Pacto Ficcional:} [\qquad / 10]\\[0.25em]
\noindent 4. \textbf{Retenção dos Símbolos (Olho / Vala / Cheiro):} [\qquad / 10]\\[0.25em]
\noindent 5. \textbf{Nota Global da Experiência:} [\qquad / 10]

\vspace{1em}
\noindent\rule{\textwidth}{0.5pt}
\vspace{0.5em}

\noindent\textbf{4. ANOTAÇÕES COMPLEMENTARES DO PACIENTE 1.263}

\vspace{0.5em}
\noindent\dotfill\\[0.4em]
\noindent\dotfill\\[0.4em]
\noindent\dotfill

\vspace{1.5em}
\begin{center}
\rule{0.5\textwidth}{0.5pt}\\[0.4em]
\textit{Assinatura do Paciente 1.263 / Leitor}\\[0.3em]
{\small\textcolor{sangue}{\textbf{O ciclo recomeça. Passar este relatório ao Arquivista.}}}
\end{center}

\endgroup
\clearpage

\clearpage
\thispagestyle{empty}
\begin{rpgdark}[VERBETE DO ARQUIVO · A CRIATURA]
\centering
{\textgoth{\fontsize{26}{28}\selectfont A CRIATURA}}\\[0.25em]
{\scshape\bfseries\small ARQUIVO 1.260 · CLASSIFICAÇÃO: [ILEGÍVEL]}\\[0.3em]
{\color{rpggold}\rule{0.55\textwidth}{0.6pt}}
\end{rpgdark}

\vspace{0.5em}
\begin{rpgdark}[FICHA]
\small
\renewcommand{\arraystretch}{1.3}
\begin{tabularx}{\textwidth}{@{}lX@{}}
\textbf{Nome:} & [ilegível] \\
\textbf{Classe:} & consequência \\
\textbf{Nível:} & 1.260 \\
\textbf{Ataque:} & leitura \\
\textbf{Defesa:} & ser passada adiante \\
\textbf{Fraqueza:} & [vazio] \\
\textbf{Teste:} & 1d20 + Lucidez (CD 15) \\
\end{tabularx}
\end{rpgdark}

\vspace{0.5em}
\begin{rpgbox}[NOTA DO ARQUIVISTA]
\centering
\itshape ``Ela não é um monstro. Ela é uma consequência.''
\end{rpgbox}

\vspace{0.5em}
\begin{rpgbox}[REAÇÃO · role 1d6]
\small
\begin{tabularx}{\textwidth}{@{}cX@{}}
\toprule
\textbf{d6} & \textbf{O que acontece} \\
\midrule
1--2 & Silêncio. O cheiro passa. Não confie. \\
3--4 & Cheiro doce de terra molhada; uma página vira sozinha. \\
5 & O olho direito brilha. Teste de Lucidez para não seguir. \\
6 & A página seguinte está escrita com a sua letra. \\
\bottomrule
\end{tabularx}
\end{rpgbox}

\vspace{0.5em}
\noindent{\footnotesize\itshape Verbete ficcional; não é um stat block jogável. O arquivo não fecha. A criatura não tem ficha completa porque ninguém que a preencheu voltou para entregar.}
\clearpage

\cleardoublepage
\thispagestyle{empty}
\begingroup
\centering
\vspace*{\fill}
{\scshape\large Molambudos --- O Diário do Paciente 1.260\par}
\vspace{0.8em}
{\small Marcelo Dias de Carvalho Filho\par}
{\small Coautor e ilustrador: Marcelo Claro\par}
\vspace{1.2em}
{\small OpenCode Books --- Crateús\par}
{\small 2.ª edição revista, 2026\par}
\vspace{0.6em}
{\footnotesize ISBN 979-8-18-917049-2\par}
\vspace{1.2em}
{\footnotesize Esta é uma obra de ficção. Joaquim Antônio Correia e o
prontuário do Paciente 1.260 são ficcionais; a instituição histórica e as
vítimas reais estão identificadas na Nota Histórica.\par}
\vspace*{\fill}
\endgroup
\clearpage


\end{document}
